% Options for packages loaded elsewhere
\PassOptionsToPackage{unicode}{hyperref}
\PassOptionsToPackage{hyphens}{url}
\PassOptionsToPackage{dvipsnames,svgnames,x11names}{xcolor}
%
\documentclass[
  10pt,
]{article}
\usepackage{amsmath,amssymb}
\usepackage{iftex}
\ifPDFTeX
  \usepackage[T1]{fontenc}
  \usepackage[utf8]{inputenc}
  \usepackage{textcomp} % provide euro and other symbols
\else % if luatex or xetex
  \usepackage{unicode-math} % this also loads fontspec
  \defaultfontfeatures{Scale=MatchLowercase}
  \defaultfontfeatures[\rmfamily]{Ligatures=TeX,Scale=1}
\fi
\usepackage{lmodern}
\ifPDFTeX\else
  % xetex/luatex font selection
\fi
% Use upquote if available, for straight quotes in verbatim environments
\IfFileExists{upquote.sty}{\usepackage{upquote}}{}
\IfFileExists{microtype.sty}{% use microtype if available
  \usepackage[]{microtype}
  \UseMicrotypeSet[protrusion]{basicmath} % disable protrusion for tt fonts
}{}
\makeatletter
\@ifundefined{KOMAClassName}{% if non-KOMA class
  \IfFileExists{parskip.sty}{%
    \usepackage{parskip}
  }{% else
    \setlength{\parindent}{0pt}
    \setlength{\parskip}{6pt plus 2pt minus 1pt}}
}{% if KOMA class
  \KOMAoptions{parskip=half}}
\makeatother
\usepackage{xcolor}
\usepackage{longtable,booktabs,array}
\usepackage{calc} % for calculating minipage widths
% Correct order of tables after \paragraph or \subparagraph
\usepackage{etoolbox}
\makeatletter
\patchcmd\longtable{\par}{\if@noskipsec\mbox{}\fi\par}{}{}
\makeatother
% Allow footnotes in longtable head/foot
\IfFileExists{footnotehyper.sty}{\usepackage{footnotehyper}}{\usepackage{footnote}}
\makesavenoteenv{longtable}
\setlength{\emergencystretch}{3em} % prevent overfull lines
\providecommand{\tightlist}{%
  \setlength{\itemsep}{0pt}\setlength{\parskip}{0pt}}
\setcounter{secnumdepth}{5}

\usepackage{amsmath,amssymb}
\usepackage{graphicx}
\usepackage{longtable,booktabs,array}
\usepackage{xcolor}
\usepackage[margin=2.0cm]{geometry}
\usepackage{microtype}
\setlength{\emergencystretch}{3em}
\setcounter{MaxMatrixCols}{20}
\pdfinfoomitdate=1
\pdftrailerid{}
\pdfsuppressptexinfo=-1
\renewcommand{\arraystretch}{1.12}
\usepackage{fvextra}
\RecustomVerbatimEnvironment{verbatim}{Verbatim}{breaklines,breakanywhere,fontsize=\small}
\fvset{breaklines,breakanywhere}
\allowdisplaybreaks
\sloppy
\ifLuaTeX
  \usepackage{selnolig}  % disable illegal ligatures
\fi
\IfFileExists{bookmark.sty}{\usepackage{bookmark}}{\usepackage{hyperref}}
\IfFileExists{xurl.sty}{\usepackage{xurl}}{} % add URL line breaks if available
\urlstyle{same}
\hypersetup{
  pdftitle={UM-4 analysis},
  pdfauthor={Underlying theory: Patrick L. Nash; Analysed and verified by Claude (Anthropic); Analysis of UM-4, 2026-10-25; Repository once-ere/Dirac\_claude, folder universe\_model\_4},
  colorlinks=true,
  linkcolor={blue},
  filecolor={Maroon},
  citecolor={Blue},
  urlcolor={blue},
  pdfcreator={LaTeX via pandoc}}

\title{UM-4 analysis}
\usepackage{etoolbox}
\makeatletter
\providecommand{\subtitle}[1]{% add subtitle to \maketitle
  \apptocmd{\@title}{\par {\large #1 \par}}{}{}
}
\makeatother
\subtitle{The x0 mass mechanism, dirac16complex as dark matter and as
dark energy, the open questions, what the author may have left out, and
UM-4 compared with UM-3}
\author{Underlying theory: Patrick L. Nash \and Analysed and verified by
Claude (Anthropic) \and Analysis of UM-4, 2026-10-25 \and Repository
once-ere/Dirac\_claude, folder universe\_model\_4}
\date{}

\begin{document}
\maketitle

{
\hypersetup{linkcolor=}
\setcounter{tocdepth}{3}
\tableofcontents
}
\hypertarget{the-fold-definition-and-examples}{%
\section{The fold: definition and
examples}\label{the-fold-definition-and-examples}}

\textbf{Definition.} The \emph{fold} is the seven-dimensional surface of
the author's eight-dimensional spacetime on which the cosine of the
angle \(z = 6Hx_0\) of the hidden direction vanishes:

\[
\mathcal{F} = \big\{X = (x_0, x_1, \dots, x_7):\ \cos(6Hx_0) = 0\big\},\qquad
\text{in one period } 0 < z < \pi\text{ of the metric: } z = \frac{\pi}{2},\quad x_0 = \frac{\pi}{12H} \tag{FD-1}
\] \textbf{Symbols in (FD-1).}

\begin{itemize}
\tightlist
\item
  \(\mathcal{F}\): The fold: the seven-dimensional surface
  \(\cos(6Hx_0) = 0\) of the eight-dimensional spacetime (FD-1), in one
  period \(0 < z < \pi\) the surface \(z = \pi/2\), \(x_0 = \pi/(12H)\),
  where \(\sigma = 0\); \(\vert_{\mathcal{F}}\) means restricted to the
  fold. Status: DERIVED.
\item
  \(X = \{x_0, \dots, x_7\}\): The author's coordinates of the
  eight-dimensional spacetime, \(X = \{x_0, x_1, \dots, x_7\}\) with the
  hidden direction first (Dirac\_Einstein-Lovelock\_4+4\_Nash.nb,
  section DEFCHART); \(X^M\) is the coordinate with the index \(M\) and
  \(\mathrm{d}^8X = \mathrm{d}x_0\,\mathrm{d}x_1\cdots\mathrm{d}x_7\)
  the coordinate volume element. Unit: length (each coordinate). Status:
  AUTHOR.
\item
  \(x_0\): The hidden space-like coordinate (listed first); the metric
  depends on it through \(z = 6Hx_0\), periodically with the period
  \(\pi/(6H)\). Unit: length. Status: AUTHOR.
\item
  \(x_1\): The first coordinate of ordinary 3-space. Unit: length.
  Status: AUTHOR.
\item
  \(x_7\): The third of the three extra time-like coordinates. Unit:
  length. Status: AUTHOR.
\item
  \(H\): The author's constant (an inverse length) that fixes the period
  \(\pi/(6H)\) of \(x_0\), the warp factor
  \(\mathrm{e}^{-2H\lvert\sigma\rvert}\) of the metric, the cosmological
  constant \(\Lambda = -21H^2\) and the tension
  \(\kappa\lambda_{\mathrm{F}} = 12H\) of the fold; in UM-4 a free
  parameter, bounded from below by the propagation of the gauge fields
  in the bulk. Unit: GeV (1/length). Value: \(H \gtrsim 2.5\) TeV (from
  the invisible width of the Z boson). Status: AUTHOR.
\item
  \(z\): The angle variable of the hidden direction, \(z = 6Hx_0\); one
  period of \(x_0\) is \(0 < z < \pi\), the fold is \(z = \pi/2\). Unit:
  dimensionless. Status: AUTHOR.
\item
  \(\pi\): The number pi, the ratio of the circumference of a circle to
  its diameter. Unit: dimensionless. Value: 3.141592654. Status:
  STANDARD.
\item
  \(\Lambda\): The eight-dimensional cosmological constant of UM-4,
  \(\Lambda = -21H^2\): the value for which the static block is a vacuum
  of Einstein gravity (its Einstein tensor is \(21H^2\) in every
  direction). The author writes the pure number \(\Lambda_{\mathrm{N}}\)
  in
  \(0 = -\Lambda_{\mathrm{N}} + \sum_p H^{-2p}w_{(p)}\mathcal{E}_{(p)}\),
  so that \(\Lambda = -\Lambda_{\mathrm{N}}H^2\). Unit: GeV\(^2\).
  Value: \(-21H^2\); the author's original \(-10H^2\); UM-3: \(-7H^2\).
  Status: AUTHOR (the term), DERIVED (the value).
\item
  \(\Lambda_{\mathrm{N}}\): The author's pure number in the vacuum
  equations of his notebook, \(\Lambda = -\Lambda_{\mathrm{N}}H^2\): his
  original value 10; 21 in UM-4 (Einstein gravity), 7 in UM-3. Unit:
  dimensionless. Value: 21 (the author's original 10). Status: AUTHOR.
\item
  \(w_{(p)}\): The Lovelock weights of the author's Einstein-Lovelock
  action, \(w_{(1)} = 1\) and \(w_{(2)}\), \(w_{(3)}\) of the
  Gauss-Bonnet and the cubic densities. UM-4 sets
  \(w_{(2)} = w_{(3)} = 0\) (Einstein gravity): the only weights for
  which gravitons propagate on the light cone on every background, not
  only around the static block. The author's original weights
  \(w_{(2)} = 1/32\), \(w_{(3)} = 1/1152\)
  (Dirac\_Einstein-Lovelock\_4+4\_Nash.nb, In{[}134{]}) keep the
  gravitons on the light cone only around the static block (with
  \(\Lambda = -161H^2/16\)); UM-3 used \((1, 1/20, 1/360)\). Unit:
  dimensionless. Value: \((1, 0, 0)\) in UM-4; the author's original
  \((1, 1/32, 1/1152)\). Status: AUTHOR (the form), DERIVED (the values
  that the order of 2026-10-25 selects).
\end{itemize}

\(X = \{x_0, \dots, x_7\}\) are the author's coordinates: \(x_0\) the
hidden space-like direction, \(x_1, x_2, x_3\) ordinary space, \(x_4\)
the time, \(x_5, x_6, x_7\) the three extra times; \(H\) is the author's
constant (an inverse length). The metric of UM-4 depends on \(x_0\) only
through \(\cot^2z\) and \((\sin^2z)^{1/6}\), so it repeats with the
period \(\pi/(6H)\) of \(x_0\), and every period \(0 < z < \pi\)
contains exactly one surface of \(\mathcal{F}\), at \(z = \pi/2\):
\(\mathcal{F}\) is the union of these copies, one per period (example
2), and ``the fold'' without a qualifier is the copy in the period that
contains our world. A point of the fold is labelled by the seven
coordinates \(x_1, \dots, x_7\), and the metric restricted to the fold
(\(\mathrm{d}x_0 = 0\), \(z = \pi/2\)) is

\[
g_{MN}\,\mathrm{d}X^M\mathrm{d}X^N\big\vert_{\mathcal{F}} = \mathrm{e}^{2a_4}\big(\mathrm{d}x_1^2 + \mathrm{d}x_2^2 + \mathrm{d}x_3^2\big)
 - \mathrm{d}x_4^2 - \mathrm{e}^{-2a_4}b_4^2\big(\mathrm{d}x_5^2 + \mathrm{d}x_6^2 + \mathrm{d}x_7^2\big) \tag{FD-2}
\] \textbf{Symbols in (FD-2).}

\begin{itemize}
\tightlist
\item
  \(g_{MN}\): The metric tensor of the eight-dimensional spacetime
  (signature \(+,+,+,+,-,-,-,-\) in the order \(x_0, \dots, x_7\)):
  \(g_{MN}\,\mathrm{d}X^M\mathrm{d}X^N\) is the squared interval,
  \(g^{MN}\) its inverse, \(\lvert g\rvert\) the absolute value of its
  determinant (\(\sqrt{\lvert g\rvert}\,\mathrm{d}^8X\) is the proper
  volume element). Unit: dimensionless. Status: AUTHOR.
\item
  \(M\): A coordinate index of eight dimensions (also with a numbered
  copy \(M_1, M_2, \dots\)): it runs over \(x_0, x_1, \dots, x_7\); a
  repeated index is summed. Status: STANDARD.
\item
  \(N\): A coordinate index of eight dimensions (also
  \(N_1, N_2, \dots\)), running over \(x_0, \dots, x_7\); a repeated
  index is summed. Status: STANDARD.
\item
  \(\mathrm{d}\): The differential: \(\mathrm{d}\) in front of a
  quantity is an infinitesimal change of it (as \(\mathrm{d}\sigma\), a
  small step of \(\sigma\)); the quotient of two such changes is a
  derivative; \(\mathrm{d}\) with an exponent, in front of several
  variables, is their volume element, the product of their
  differentials. Status: STANDARD.
\item
  \(X = \{x_0, \dots, x_7\}\): The author's coordinates of the
  eight-dimensional spacetime, \(X = \{x_0, x_1, \dots, x_7\}\) with the
  hidden direction first (Dirac\_Einstein-Lovelock\_4+4\_Nash.nb,
  section DEFCHART); \(X^M\) is the coordinate with the index \(M\) and
  \(\mathrm{d}^8X = \mathrm{d}x_0\,\mathrm{d}x_1\cdots\mathrm{d}x_7\)
  the coordinate volume element. Unit: length (each coordinate). Status:
  AUTHOR.
\item
  \(\mathcal{F}\): The fold: the seven-dimensional surface
  \(\cos(6Hx_0) = 0\) of the eight-dimensional spacetime (FD-1), in one
  period \(0 < z < \pi\) the surface \(z = \pi/2\), \(x_0 = \pi/(12H)\),
  where \(\sigma = 0\); \(\vert_{\mathcal{F}}\) means restricted to the
  fold. Status: DERIVED.
\item
  \(\mathrm{e}\): Euler's number, the base of the natural logarithm;
  \(\mathrm{e}^{X}\) is the exponential of \(X\). Unit: dimensionless.
  Value: 2.718281828. Status: STANDARD.
\item
  \(a_4\): The author's metric function of the time \(x_4\) that scales
  ordinary 3-space: on the fold \(\mathrm{e}^{a_4}\) is the scale factor
  of cosmology and \(a_4'\) the Hubble rate. Unit: dimensionless.
  Status: AUTHOR.
\item
  \(b_4\): The author's metric function of the time \(x_4\) in the
  extra-time components, \(b_4 > 0\), never identically 1 (the author's
  order of 2026-10-24); the scale factor (the factor of the lengths) of
  the extra times is \(\mathrm{e}^{-a_4}b_4\) on the fold and
  \(\mathrm{e}^{-a_4}b_4\,\mathrm{e}^{-H\lvert\sigma\rvert}\) at the
  distance \(\sigma\) from it (the square root of the warp factor times
  \(\mathrm{e}^{-a_4}b_4\)). Unit: dimensionless. Status: AUTHOR.
\item
  \(x_1\): The first coordinate of ordinary 3-space. Unit: length.
  Status: AUTHOR.
\item
  \(x_2\): The second coordinate of ordinary 3-space. Unit: length.
  Status: AUTHOR.
\item
  \(x_3\): The third coordinate of ordinary 3-space. Unit: length.
  Status: AUTHOR.
\item
  \(x_4\): The time; a prime on a function of \(x_4\) alone is
  \(\mathrm{d}/\mathrm{d}x_4\). Unit: length (= time, \(c = 1\)).
  Status: AUTHOR.
\item
  \(x_5\): The first of the three extra time-like coordinates. Unit:
  length. Status: AUTHOR.
\item
  \(x_6\): The second of the three extra time-like coordinates. Unit:
  length. Status: AUTHOR.
\item
  \(x_7\): The third of the three extra time-like coordinates. Unit:
  length. Status: AUTHOR.
\item
  \({}'\): The derivative of a function of the time \(x_4\) alone with
  respect to \(x_4\) (\(a_4' = \mathrm{d}a_4/\mathrm{d}x_4\); two primes
  the second derivative); on a function of \(\sigma\) alone (a profile)
  the derivative with respect to \(\sigma\). Status: STANDARD.
\end{itemize}

with three space-like and four time-like directions; \(a_4\) and \(b_4\)
are the author's two functions of the time \(x_4\). Notation: \(\ln\) is
the natural logarithm, \(\mathrm{e}\) Euler's number,
\(\exp(y) = \mathrm{e}^y\), \(\sin\), \(\cos\) and \(\cot = \cos/\sin\)
the trigonometric functions, \(\sinh\) and \(\tanh\) the hyperbolic
functions, \(\lvert y\rvert\) the absolute value; \(f(0^+)\) and
\(f(0^-)\) are the limits of \(f(\sigma)\) as \(\sigma\) approaches 0
from above and from below.

\textbf{Why this surface is called a fold.} Measured by the proper
distance along the hidden direction,

\[
\sigma = \mathrm{sgn}\Big(\frac{\pi}{2} - z\Big)\frac{\ln\sin z}{6H},\qquad \mathrm{d}\sigma = \lvert\cot z\rvert\,\mathrm{d}x_0,\qquad
\sin z = \mathrm{e}^{-6H\lvert\sigma\rvert} \tag{FD-3}
\] \textbf{Symbols in (FD-3).}

\begin{itemize}
\tightlist
\item
  \(\sigma\): The proper coordinate of the hidden direction:
  \(\mathrm{d}\sigma = \lvert\cot z\rvert\,\mathrm{d}x_0\),
  \(\sigma = \mathrm{sgn}(\pi/2 - z)\ln(\sin z)/(6H)\) on one period
  \(0 < z < \pi\); \(\sigma = 0\) is the fold \(z = \pi/2\),
  \(\sigma \to \mp\infty\) the tips \(z \to 0, \pi\); one period of
  \(x_0\) is the whole line \(-\infty < \sigma < \infty\). Unit: length.
  Status: DERIVED.
\item
  \(\mathrm{sgn}\): The sign function: \(+1\) for a positive argument
  and \(-1\) for a negative one. Status: STANDARD.
\item
  \(\pi\): The number pi, the ratio of the circumference of a circle to
  its diameter. Unit: dimensionless. Value: 3.141592654. Status:
  STANDARD.
\item
  \(z\): The angle variable of the hidden direction, \(z = 6Hx_0\); one
  period of \(x_0\) is \(0 < z < \pi\), the fold is \(z = \pi/2\). Unit:
  dimensionless. Status: AUTHOR.
\item
  \(H\): The author's constant (an inverse length) that fixes the period
  \(\pi/(6H)\) of \(x_0\), the warp factor
  \(\mathrm{e}^{-2H\lvert\sigma\rvert}\) of the metric, the cosmological
  constant \(\Lambda = -21H^2\) and the tension
  \(\kappa\lambda_{\mathrm{F}} = 12H\) of the fold; in UM-4 a free
  parameter, bounded from below by the propagation of the gauge fields
  in the bulk. Unit: GeV (1/length). Value: \(H \gtrsim 2.5\) TeV (from
  the invisible width of the Z boson). Status: AUTHOR.
\item
  \(\mathrm{d}\): The differential: \(\mathrm{d}\) in front of a
  quantity is an infinitesimal change of it (as \(\mathrm{d}\sigma\), a
  small step of \(\sigma\)); the quotient of two such changes is a
  derivative; \(\mathrm{d}\) with an exponent, in front of several
  variables, is their volume element, the product of their
  differentials. Status: STANDARD.
\item
  \(x_0\): The hidden space-like coordinate (listed first); the metric
  depends on it through \(z = 6Hx_0\), periodically with the period
  \(\pi/(6H)\). Unit: length. Status: AUTHOR.
\item
  \(\mathrm{e}\): Euler's number, the base of the natural logarithm;
  \(\mathrm{e}^{X}\) is the exponential of \(X\). Unit: dimensionless.
  Value: 2.718281828. Status: STANDARD.
\item
  \(\infty\): Infinity: as a limit of an integral or a sum it means that
  the range has no bound. Status: STANDARD.
\end{itemize}

(the fold is \(\sigma = 0\); the half-period \(0 < z < \pi/2\) is
\(\sigma < 0\), the half-period \(\pi/2 < z < \pi\) is \(\sigma > 0\);
the two ends of the period, the tips \(z \to 0\) and \(z \to \pi\), lie
at \(\sigma \to -\infty\) and \(\sigma \to +\infty\)), the metric of
UM-4 is

\[
g_{MN}\,\mathrm{d}X^M\mathrm{d}X^N = \mathrm{d}\sigma^2 + \mathrm{e}^{-2H\lvert\sigma\rvert}\Big[\mathrm{e}^{2a_4}\big(\mathrm{d}x_1^2
 + \mathrm{d}x_2^2 + \mathrm{d}x_3^2\big) - \mathrm{d}x_4^2 - \mathrm{e}^{-2a_4}b_4^2\big(\mathrm{d}x_5^2 + \mathrm{d}x_6^2 + \mathrm{d}x_7^2\big)\Big] \tag{FD-4}
\] \textbf{Symbols in (FD-4).}

\begin{itemize}
\tightlist
\item
  \(g_{MN}\): The metric tensor of the eight-dimensional spacetime
  (signature \(+,+,+,+,-,-,-,-\) in the order \(x_0, \dots, x_7\)):
  \(g_{MN}\,\mathrm{d}X^M\mathrm{d}X^N\) is the squared interval,
  \(g^{MN}\) its inverse, \(\lvert g\rvert\) the absolute value of its
  determinant (\(\sqrt{\lvert g\rvert}\,\mathrm{d}^8X\) is the proper
  volume element). Unit: dimensionless. Status: AUTHOR.
\item
  \(M\): A coordinate index of eight dimensions (also with a numbered
  copy \(M_1, M_2, \dots\)): it runs over \(x_0, x_1, \dots, x_7\); a
  repeated index is summed. Status: STANDARD.
\item
  \(N\): A coordinate index of eight dimensions (also
  \(N_1, N_2, \dots\)), running over \(x_0, \dots, x_7\); a repeated
  index is summed. Status: STANDARD.
\item
  \(\mathrm{d}\): The differential: \(\mathrm{d}\) in front of a
  quantity is an infinitesimal change of it (as \(\mathrm{d}\sigma\), a
  small step of \(\sigma\)); the quotient of two such changes is a
  derivative; \(\mathrm{d}\) with an exponent, in front of several
  variables, is their volume element, the product of their
  differentials. Status: STANDARD.
\item
  \(X = \{x_0, \dots, x_7\}\): The author's coordinates of the
  eight-dimensional spacetime, \(X = \{x_0, x_1, \dots, x_7\}\) with the
  hidden direction first (Dirac\_Einstein-Lovelock\_4+4\_Nash.nb,
  section DEFCHART); \(X^M\) is the coordinate with the index \(M\) and
  \(\mathrm{d}^8X = \mathrm{d}x_0\,\mathrm{d}x_1\cdots\mathrm{d}x_7\)
  the coordinate volume element. Unit: length (each coordinate). Status:
  AUTHOR.
\item
  \(\sigma\): The proper coordinate of the hidden direction:
  \(\mathrm{d}\sigma = \lvert\cot z\rvert\,\mathrm{d}x_0\),
  \(\sigma = \mathrm{sgn}(\pi/2 - z)\ln(\sin z)/(6H)\) on one period
  \(0 < z < \pi\); \(\sigma = 0\) is the fold \(z = \pi/2\),
  \(\sigma \to \mp\infty\) the tips \(z \to 0, \pi\); one period of
  \(x_0\) is the whole line \(-\infty < \sigma < \infty\). Unit: length.
  Status: DERIVED.
\item
  \(\mathrm{e}\): Euler's number, the base of the natural logarithm;
  \(\mathrm{e}^{X}\) is the exponential of \(X\). Unit: dimensionless.
  Value: 2.718281828. Status: STANDARD.
\item
  \(H\): The author's constant (an inverse length) that fixes the period
  \(\pi/(6H)\) of \(x_0\), the warp factor
  \(\mathrm{e}^{-2H\lvert\sigma\rvert}\) of the metric, the cosmological
  constant \(\Lambda = -21H^2\) and the tension
  \(\kappa\lambda_{\mathrm{F}} = 12H\) of the fold; in UM-4 a free
  parameter, bounded from below by the propagation of the gauge fields
  in the bulk. Unit: GeV (1/length). Value: \(H \gtrsim 2.5\) TeV (from
  the invisible width of the Z boson). Status: AUTHOR.
\item
  \(a_4\): The author's metric function of the time \(x_4\) that scales
  ordinary 3-space: on the fold \(\mathrm{e}^{a_4}\) is the scale factor
  of cosmology and \(a_4'\) the Hubble rate. Unit: dimensionless.
  Status: AUTHOR.
\item
  \(b_4\): The author's metric function of the time \(x_4\) in the
  extra-time components, \(b_4 > 0\), never identically 1 (the author's
  order of 2026-10-24); the scale factor (the factor of the lengths) of
  the extra times is \(\mathrm{e}^{-a_4}b_4\) on the fold and
  \(\mathrm{e}^{-a_4}b_4\,\mathrm{e}^{-H\lvert\sigma\rvert}\) at the
  distance \(\sigma\) from it (the square root of the warp factor times
  \(\mathrm{e}^{-a_4}b_4\)). Unit: dimensionless. Status: AUTHOR.
\item
  \(x_1\): The first coordinate of ordinary 3-space. Unit: length.
  Status: AUTHOR.
\item
  \(x_2\): The second coordinate of ordinary 3-space. Unit: length.
  Status: AUTHOR.
\item
  \(x_3\): The third coordinate of ordinary 3-space. Unit: length.
  Status: AUTHOR.
\item
  \(x_4\): The time; a prime on a function of \(x_4\) alone is
  \(\mathrm{d}/\mathrm{d}x_4\). Unit: length (= time, \(c = 1\)).
  Status: AUTHOR.
\item
  \(x_5\): The first of the three extra time-like coordinates. Unit:
  length. Status: AUTHOR.
\item
  \(x_6\): The second of the three extra time-like coordinates. Unit:
  length. Status: AUTHOR.
\item
  \(x_7\): The third of the three extra time-like coordinates. Unit:
  length. Status: AUTHOR.
\end{itemize}

It depends on \(\sigma\) only through \(\lvert\sigma\rvert\). The
function \(\mathcal{W}(\sigma) = \mathrm{e}^{-2H\lvert\sigma\rvert}\),
which multiplies the squared steps of all seven directions of the fold
(ordinary space, the time and the extra times), is the warp factor of
the metric (its square root \(\mathrm{e}^{-H\lvert\sigma\rvert}\)
multiplies the lengths). The time \(x_4\) is warped like the other
directions (the author's order of 2026-10-25, 13:15 PDT; in UM-3 it was
not), so that every surface \(\sigma = \text{const}\) has the light cone
of the fold. Four properties make the surface \(\sigma = 0\) a fold, and
together they single it out within the period:

\begin{enumerate}
\def\labelenumi{\arabic{enumi}.}
\tightlist
\item
  \emph{It is a mirror.} The reflection \(\sigma \to -\sigma\) (that is
  \(z \to \pi - z\), or \(x_0 \to \pi/(6H) - x_0\)) maps the metric to
  itself, and the fold is the only surface of the period that the
  reflection leaves in place: the two halves of the period are mirror
  images of each other, joined along the fold, as the two halves of a
  folded sheet of paper are joined along the crease.
\item
  \emph{The warp factor is largest on it.} The warp factor
  \((\sin^2z)^{1/6} = \mathrm{e}^{-2H\lvert\sigma\rvert}\) equals 1 on
  the fold and decreases on both sides.
\item
  \emph{The metric has a kink on it.} The metric is continuous across
  the fold, but its first derivative along \(\sigma\) jumps. The
  extrinsic curvature of the surfaces \(\sigma = \text{const}\) changes
  sign across the fold,
\end{enumerate}

\[
K^a{}_b = \frac12\,g^{ac}\,\partial_\sigma g_{cb} = -H\,\mathrm{sgn}(\sigma)\,\delta^a{}_b\quad(\text{all seven directions of the fold}) \tag{FD-5}
\] \textbf{Symbols in (FD-5).}

\begin{itemize}
\tightlist
\item
  \(K^a{}_b\): The extrinsic curvature of a surface of constant
  \(\sigma\) of the hidden direction,
  \(K^a{}_b = \frac12 g^{ac}\partial_\sigma g_{cb}\) for \(a, b\) along
  the surface (FD-5): \(-H\,\mathrm{sgn}(\sigma)\) along all seven
  directions of the fold for the author's form of the metric of UM-4;
  \(K = K^c{}_c\) its trace. Unit: GeV (1/length). Status: DERIVED.
\item
  \(a\): A frame (orthonormal) index running over \(0, 1, \dots, 7\)
  (frame direction \((0)\) along \(\sigma\), \((1), (2), (3)\) along
  ordinary space, \((4)\) along the time, \((5), (6), (7)\) along the
  extra times); raised and lowered with the frame metric
  \(\eta_{ab} = \mathrm{diag}(1, 1, 1, 1, -1, -1, -1, -1)\); on the fold
  (\(a, b \in \{1, \dots, 7\}\)) the frame of the fold; in the tensors
  of the fold written with its coordinates \(x^a\) (its metric
  \(\hat g_{ab}\) and the tensors built from it) a coordinate index of
  the fold, \(a \in \{1, \dots, 7\}\); a repeated index is summed.
  Status: STANDARD.
\item
  \(b\): A frame index (as \(a\)). Status: STANDARD.
\item
  \(c\): A frame index (as \(a\)). Status: STANDARD.
\item
  \(g_{MN}\): The metric tensor of the eight-dimensional spacetime
  (signature \(+,+,+,+,-,-,-,-\) in the order \(x_0, \dots, x_7\)):
  \(g_{MN}\,\mathrm{d}X^M\mathrm{d}X^N\) is the squared interval,
  \(g^{MN}\) its inverse, \(\lvert g\rvert\) the absolute value of its
  determinant (\(\sqrt{\lvert g\rvert}\,\mathrm{d}^8X\) is the proper
  volume element). Unit: dimensionless. Status: AUTHOR.
\item
  \(\partial\): A partial derivative: the derivative with respect to one
  variable at fixed values of the others (\(\partial_\sigma\) with
  respect to \(\sigma\); \(\partial_M\) and \(\partial_\mu\) with
  respect to the coordinate that carries the same index, of the eight
  dimensions or of the fold). Status: STANDARD.
\item
  \(\sigma\): The label \(\sigma\) in \(\partial_\sigma\): the
  derivative with respect to the proper hidden coordinate \(\sigma\).
  Status: STANDARD.
\item
  \(\sigma\): The proper coordinate of the hidden direction:
  \(\mathrm{d}\sigma = \lvert\cot z\rvert\,\mathrm{d}x_0\),
  \(\sigma = \mathrm{sgn}(\pi/2 - z)\ln(\sin z)/(6H)\) on one period
  \(0 < z < \pi\); \(\sigma = 0\) is the fold \(z = \pi/2\),
  \(\sigma \to \mp\infty\) the tips \(z \to 0, \pi\); one period of
  \(x_0\) is the whole line \(-\infty < \sigma < \infty\). Unit: length.
  Status: DERIVED.
\item
  \(H\): The author's constant (an inverse length) that fixes the period
  \(\pi/(6H)\) of \(x_0\), the warp factor
  \(\mathrm{e}^{-2H\lvert\sigma\rvert}\) of the metric, the cosmological
  constant \(\Lambda = -21H^2\) and the tension
  \(\kappa\lambda_{\mathrm{F}} = 12H\) of the fold; in UM-4 a free
  parameter, bounded from below by the propagation of the gauge fields
  in the bulk. Unit: GeV (1/length). Value: \(H \gtrsim 2.5\) TeV (from
  the invisible width of the Z boson). Status: AUTHOR.
\item
  \(\mathrm{sgn}\): The sign function: \(+1\) for a positive argument
  and \(-1\) for a negative one. Status: STANDARD.
\item
  \(\delta\): The Kronecker delta: with two indices \(\delta^M_N = 1\)
  if \(M = N\) and 0 otherwise (likewise \(\delta_{X\mathrm{L}}\) is 1
  if the label \(X\) of a half of dirac16complex is L and 0 otherwise);
  with two index strings of equal length the generalised Kronecker delta
  \(\delta^{M_1M_2\cdots}_{N_1N_2\cdots}\), the determinant of the
  square matrix of the Kronecker deltas \(\delta^{M_a}_{N_b}\) (\(+1\)
  or \(-1\) if the upper indices are an even or odd permutation of the
  lower ones, 0 otherwise). Status: STANDARD.
\item
  \(\eta_{ab}\): The flat metric of the frame,
  \(\eta = \mathrm{diag}(1, 1, 1, 1, -1, -1, -1, -1)\) (eta4488 of
  4+4.nb); \(\eta^{ab}\) is its inverse (the same matrix). Unit:
  dimensionless. Status: AUTHOR.
\item
  \(\hat g_{ab}\): The metric of the seven-dimensional fold (FD-2) in
  its coordinates \(x_1, \dots, x_7\) (or in its frame), \(\hat g^{ab}\)
  its inverse; \(\hat h_{\mu\nu}\) is its part along the four observable
  directions. Unit: dimensionless. Status: AUTHOR.
\item
  \(\mathrm{diag}\): A diagonal matrix with the listed entries on its
  diagonal (in the order of the coordinates \(x_0, x_1, \dots, x_7\) for
  a tensor of eight dimensions, \(x_1, \dots, x_7\) for a tensor of the
  fold, \(x_1, \dots, x_4\) for a four-dimensional tensor). Status:
  STANDARD.
\item
  \(x\): The coordinates of the fold: \(x_1, x_2, x_3\) of ordinary
  space, \(x_4\) the time, \(x_5, x_6, x_7\) the extra times; a field of
  the fold is a function of the four observable coordinates
  \(x = (x_1, x_2, x_3, x_4)\), and
  \(\mathrm{d}^4x = \mathrm{d}x_1\,\mathrm{d}x_2\,\mathrm{d}x_3\,\mathrm{d}x_4\).
  Unit: length (each coordinate). Status: AUTHOR.
\end{itemize}

(\(a, b\) frame indices along the fold); it jumps by \(-2H\) in every
direction of the fold. The field equations of gravity of UM-4,
\(\mathcal{E}_{(1)}{}^M{}_N + \Lambda\delta^M_N = \kappa\mathcal{T}^M{}_N\)
(Einstein gravity: the author's Einstein-Lovelock equations with the
weights of the Gauss-Bonnet and cubic terms set to zero;
\(\mathcal{E}_{(1)}\) the Einstein tensor, \(\Lambda = -21H^2\) the
cosmological constant, \(\kappa = 8\pi G_8\) the gravitational coupling
of eight dimensions, \(\mathcal{T}^M{}_N\) the energy-momentum tensor,
\(\delta^M_N\) the Kronecker delta), are linear in the second
\(\sigma\)-derivatives of the metric, so the jump puts a Dirac delta
function \(\delta_{\mathrm{D}}(\sigma)\) into them: the metric solves
them only if the fold carries a thin source \(\mathcal{S}^a{}_b\), an
energy-momentum tensor concentrated on \(\sigma = 0\), given by the
integral of the equations across the fold (the junction condition of
Israel),

\[
\kappa\,\mathcal{S}^a{}_b = \int_{0^-}^{0^+}\mathrm{d}\sigma\,\mathcal{E}_{(1)}{}^a{}_b
 = -\Big(\big(K^a{}_b\big\vert_{0^+} - K^a{}_b\big\vert_{0^-}\big) - \delta^a{}_b\big(K\big\vert_{0^+} - K\big\vert_{0^-}\big)\Big) \tag{FD-11}
\] \textbf{Symbols in (FD-11).}

\begin{itemize}
\tightlist
\item
  \(\kappa\): The eight-dimensional gravitational coupling,
  \(\kappa = 8\pi G_8\). Unit: GeV\(^{-6}\). Status: AUTHOR.
\item
  \(\mathcal{S}^a{}_b\): The thin source of the fold: the
  energy-momentum tensor concentrated on the fold, integrated across it,
  \(\kappa\mathcal{S}^a{}_b = \int_{0^-}^{0^+}\mathrm{d}\sigma\,\mathcal{E}_{(1)}{}^a{}_b\)
  (the junction condition of Israel); \(\mathcal{S}^4{}_4\) is minus its
  energy density, \(\mathcal{S}^1{}_1\) and \(\mathcal{S}^5{}_5\) its
  pressures along ordinary space and along the extra times (frame
  components along the fold, eight-dimensional normalisation: an energy
  per unit volume of the six other directions of the fold). Unit:
  GeV\(^7\). Value: \(\kappa\mathcal{S}^a{}_b = -12H\delta^a{}_b\) for
  every \(a_4\), \(b_4\) of the author's form. Status: DERIVED.
\item
  \(a\): A frame (orthonormal) index running over \(0, 1, \dots, 7\)
  (frame direction \((0)\) along \(\sigma\), \((1), (2), (3)\) along
  ordinary space, \((4)\) along the time, \((5), (6), (7)\) along the
  extra times); raised and lowered with the frame metric
  \(\eta_{ab} = \mathrm{diag}(1, 1, 1, 1, -1, -1, -1, -1)\); on the fold
  (\(a, b \in \{1, \dots, 7\}\)) the frame of the fold; in the tensors
  of the fold written with its coordinates \(x^a\) (its metric
  \(\hat g_{ab}\) and the tensors built from it) a coordinate index of
  the fold, \(a \in \{1, \dots, 7\}\); a repeated index is summed.
  Status: STANDARD.
\item
  \(b\): A frame index (as \(a\)). Status: STANDARD.
\item
  \(\mathrm{d}\): The differential: \(\mathrm{d}\) in front of a
  quantity is an infinitesimal change of it (as \(\mathrm{d}\sigma\), a
  small step of \(\sigma\)); the quotient of two such changes is a
  derivative; \(\mathrm{d}\) with an exponent, in front of several
  variables, is their volume element, the product of their
  differentials. Status: STANDARD.
\item
  \(\sigma\): The proper coordinate of the hidden direction:
  \(\mathrm{d}\sigma = \lvert\cot z\rvert\,\mathrm{d}x_0\),
  \(\sigma = \mathrm{sgn}(\pi/2 - z)\ln(\sin z)/(6H)\) on one period
  \(0 < z < \pi\); \(\sigma = 0\) is the fold \(z = \pi/2\),
  \(\sigma \to \mp\infty\) the tips \(z \to 0, \pi\); one period of
  \(x_0\) is the whole line \(-\infty < \sigma < \infty\). Unit: length.
  Status: DERIVED.
\item
  \(\mathcal{E}_{(p)}{}^M{}_N\): The Lovelock tensor of order \(p\),
  \(\mathcal{E}_{(p)}{}^M{}_N = -2^{-(p+1)}\,\delta^{MP_1\cdots P_{2p}}_{NQ_1\cdots Q_{2p}}\,R^{Q_1Q_2}{}_{P_1P_2}\cdots R^{Q_{2p-1}Q_{2p}}{}_{P_{2p-1}P_{2p}}\)
  (with both indices down it is \(1/\sqrt{\lvert g\rvert}\) times the
  variation of \(\sqrt{\lvert g\rvert}\,\mathcal{L}_{(p)}\) with respect
  to \(g^{MN}\); \(\mathcal{E}_{(1)}\) the Einstein tensor,
  \(\mathcal{E}_{(2)}\) the Gauss-Bonnet tensor, \(\mathcal{E}_{(3)}\)
  the cubic Lovelock tensor); with frame indices its frame components.
  Unit: 1/length\(^{2p}\). Status: STANDARD.
\item
  \(p\): The Lovelock order \(p = 1, 2, 3\) (Einstein, Gauss-Bonnet,
  cubic), as an index and as a number (UM-4 keeps only \(p = 1\)); in
  the Kohn-Sham equations also a second label of a Kohn-Sham orbital,
  with the Kronecker delta of the two labels. Status: STANDARD.
\item
  \(K^a{}_b\): The extrinsic curvature of a surface of constant
  \(\sigma\) of the hidden direction,
  \(K^a{}_b = \frac12 g^{ac}\partial_\sigma g_{cb}\) for \(a, b\) along
  the surface (FD-5): \(-H\,\mathrm{sgn}(\sigma)\) along all seven
  directions of the fold for the author's form of the metric of UM-4;
  \(K = K^c{}_c\) its trace. Unit: GeV (1/length). Status: DERIVED.
\item
  \(\delta\): The Kronecker delta: with two indices \(\delta^M_N = 1\)
  if \(M = N\) and 0 otherwise (likewise \(\delta_{X\mathrm{L}}\) is 1
  if the label \(X\) of a half of dirac16complex is L and 0 otherwise);
  with two index strings of equal length the generalised Kronecker delta
  \(\delta^{M_1M_2\cdots}_{N_1N_2\cdots}\), the determinant of the
  square matrix of the Kronecker deltas \(\delta^{M_a}_{N_b}\) (\(+1\)
  or \(-1\) if the upper indices are an even or odd permutation of the
  lower ones, 0 otherwise). Status: STANDARD.
\item
  \(G_8\): The eight-dimensional Newton constant
  (\(\kappa = 8\pi G_8\)). Unit: GeV\(^{-6}\). Status: AUTHOR.
\item
  \(R^{MN}{}_{PQ},\ R_{MN},\ R\): The curvature of the metric
  \(g_{MN}\): with four indices the Riemann tensor \(R^M{}_{NPQ}\) (in
  the convention in which the commutator of two covariant derivatives,
  acting on a vector, is the Riemann tensor, with its last two indices
  those of the two derivatives, contracted with the vector), with two
  indices the Ricci tensor \(R_{MN} = R^P{}_{MPN}\), without indices the
  scalar curvature \(R = g^{MN}R_{MN}\); with frame indices
  \(R^{(a)(b)}{}_{(c)(d)}\) its frame components,
  \(R^{(a)(b)}{}_{(a)(b)}\) (no sum) the sectional curvature of the
  frame plane \((a, b)\). Unit: 1/length\(^2\). Status: STANDARD.
\item
  \(\mathcal{L}_{(p)}\): The Lovelock density of order \(p\),
  \(\mathcal{L}_{(p)} = 2^{-p}\,\delta^{P_1\cdots P_{2p}}_{Q_1\cdots Q_{2p}}\,R^{Q_1Q_2}{}_{P_1P_2}\cdots R^{Q_{2p-1}Q_{2p}}{}_{P_{2p-1}P_{2p}}\),
  a scalar built from \(p\) curvature tensors: \(\mathcal{L}_{(1)} = R\)
  (Einstein-Hilbert), \(\mathcal{L}_{(2)}\) the Gauss-Bonnet density,
  \(\mathcal{L}_{(3)}\) the cubic Lovelock density. Unit:
  1/length\(^{2p}\). Status: STANDARD.
\end{itemize}

(\(K = K^c{}_c\) the trace of the extrinsic curvature), with
\(\mathcal{S}^4{}_4 = -(\text{its energy density})\),
\(\mathcal{S}^1{}_1 = \mathcal{S}^2{}_2 = \mathcal{S}^3{}_3\) its
pressure along ordinary space and
\(\mathcal{S}^5{}_5 = \mathcal{S}^6{}_6 = \mathcal{S}^7{}_7\) its
pressure along the extra times; \(\mathcal{S}\) is \(\mathcal{T}\)
integrated across the fold, an energy per unit volume of the six other
directions of the fold (unit GeV\(^7\)). For the static block (\(a_4\)
and \(b_4\) constant) and, to the first order in the curvature of the
fold divided by \(H^2\), for every fold

\[
\begin{aligned}
&\text{static block:}\quad \kappa\mathcal{S}^a{}_b = -12H\,\delta^a{}_b\quad\text{(a pure tension)},\qquad \kappa\lambda_{\mathrm{F}} = 12H,\\
&\text{any fold:}\quad \kappa\mathcal{S}^a{}_b = -12H\,\delta^a{}_b + \frac{2}{5H}\,\hat G^a{}_b + \dots
\end{aligned} \tag{FD-12}
\] \textbf{Symbols in (FD-12).}

\begin{itemize}
\tightlist
\item
  \(\kappa\): The eight-dimensional gravitational coupling,
  \(\kappa = 8\pi G_8\). Unit: GeV\(^{-6}\). Status: AUTHOR.
\item
  \(\mathcal{S}^a{}_b\): The thin source of the fold: the
  energy-momentum tensor concentrated on the fold, integrated across it,
  \(\kappa\mathcal{S}^a{}_b = \int_{0^-}^{0^+}\mathrm{d}\sigma\,\mathcal{E}_{(1)}{}^a{}_b\)
  (the junction condition of Israel); \(\mathcal{S}^4{}_4\) is minus its
  energy density, \(\mathcal{S}^1{}_1\) and \(\mathcal{S}^5{}_5\) its
  pressures along ordinary space and along the extra times (frame
  components along the fold, eight-dimensional normalisation: an energy
  per unit volume of the six other directions of the fold). Unit:
  GeV\(^7\). Value: \(\kappa\mathcal{S}^a{}_b = -12H\delta^a{}_b\) for
  every \(a_4\), \(b_4\) of the author's form. Status: DERIVED.
\item
  \(a\): A frame (orthonormal) index running over \(0, 1, \dots, 7\)
  (frame direction \((0)\) along \(\sigma\), \((1), (2), (3)\) along
  ordinary space, \((4)\) along the time, \((5), (6), (7)\) along the
  extra times); raised and lowered with the frame metric
  \(\eta_{ab} = \mathrm{diag}(1, 1, 1, 1, -1, -1, -1, -1)\); on the fold
  (\(a, b \in \{1, \dots, 7\}\)) the frame of the fold; in the tensors
  of the fold written with its coordinates \(x^a\) (its metric
  \(\hat g_{ab}\) and the tensors built from it) a coordinate index of
  the fold, \(a \in \{1, \dots, 7\}\); a repeated index is summed.
  Status: STANDARD.
\item
  \(b\): A frame index (as \(a\)). Status: STANDARD.
\item
  \(H\): The author's constant (an inverse length) that fixes the period
  \(\pi/(6H)\) of \(x_0\), the warp factor
  \(\mathrm{e}^{-2H\lvert\sigma\rvert}\) of the metric, the cosmological
  constant \(\Lambda = -21H^2\) and the tension
  \(\kappa\lambda_{\mathrm{F}} = 12H\) of the fold; in UM-4 a free
  parameter, bounded from below by the propagation of the gauge fields
  in the bulk. Unit: GeV (1/length). Value: \(H \gtrsim 2.5\) TeV (from
  the invisible width of the Z boson). Status: AUTHOR.
\item
  \(\delta\): The Kronecker delta: with two indices \(\delta^M_N = 1\)
  if \(M = N\) and 0 otherwise (likewise \(\delta_{X\mathrm{L}}\) is 1
  if the label \(X\) of a half of dirac16complex is L and 0 otherwise);
  with two index strings of equal length the generalised Kronecker delta
  \(\delta^{M_1M_2\cdots}_{N_1N_2\cdots}\), the determinant of the
  square matrix of the Kronecker deltas \(\delta^{M_a}_{N_b}\) (\(+1\)
  or \(-1\) if the upper indices are an even or odd permutation of the
  lower ones, 0 otherwise). Status: STANDARD.
\item
  \(\lambda_{\mathrm{F}}\): The tension that the static block requires
  of the fold, \(\kappa\lambda_{\mathrm{F}} = 12H\): the total vacuum
  energy of the seven-dimensional fold, made of the tension of the fold
  fluid and the vacuum energy of the fields of the fold (both have the
  pressure minus their energy density along ordinary space and along the
  extra times), per unit volume of the six other directions of the fold
  (eight-dimensional normalisation); in the four-dimensional
  normalisation it is \(\lambda_{\mathrm{F}}V_*\), and the excess of the
  total over \(\lambda_{\mathrm{F}}V_*\) is the dark energy
  \(\rho_\Lambda\). Unit: GeV\(^7\). Value: \(12H/\kappa\);
  \(\lambda_{\mathrm{F}}V_* = 30H^2M_{\mathrm{red}}^2 = 1.1\times10^{45}\)
  GeV\(^4\) for \(H = 2.5\) TeV. Status: DERIVED.
\item
  \(\hat G^\mu{}_\nu\): The Einstein tensor of the metric
  \(\hat g_{ab}\) of the fold (FD-2), with frame indices along its seven
  directions; with four-dimensional indices, \(\hat G^\mu{}_\nu\), its
  components along the four observable directions, which for frozen
  extra times are those of the Einstein tensor of the four-dimensional
  metric \(\hat h_{\mu\nu}\). Unit: 1/length\(^2\). Status: STANDARD.
\item
  \(M_{\mathrm{red}}\): The reduced Planck mass,
  \(M_{\mathrm{red}} = (8\pi G_{\mathrm{N}})^{-1/2}\). Unit: GeV. Value:
  \(2.435\times10^{18}\) GeV. Status: STANDARD.
\item
  \(V_*\): The physical volume of the extra times \(x_5, x_6, x_7\) on
  the fold in the late universe,
  \(V_* = V_{\mathrm{t}}(\mathrm{e}^{-a_4}b_4)^3\) (with
  \(V_{\mathrm{t}} = \int\mathrm{d}x_5\,\mathrm{d}x_6\,\mathrm{d}x_7\)
  the coordinate volume of the extra times) at a time when
  \(\mathrm{e}^{-a_4}b_4\) is constant; it converts the
  eight-dimensional normalisation into the four-dimensional one. Unit:
  length\(^3\) (GeV\(^{-3}\)). Status: ASSUMED.
\item
  \(\hat h_{\mu\nu}\): The four-dimensional metric of the fold,
  \(\hat h_{\mu\nu} = \mathrm{diag}(\mathrm{e}^{2a_4}, \mathrm{e}^{2a_4}, \mathrm{e}^{2a_4}, -1)\)
  in the coordinates \(x_1, \dots, x_4\) (the induced metric at
  \(\sigma = 0\) without the extra times); \(\hat h^{\mu\nu}\) its
  inverse, \(\hat h\) its determinant,
  \(\sqrt{-\hat h} = \mathrm{e}^{3a_4}\); \(\hat h^{(1)}_{\mu\nu}\) its
  first-order cosmological perturbation. Unit: dimensionless. Status:
  DERIVED.
\item
  \(\rho_\Lambda\): The vacuum energy density of the fold above its
  static tension, the dark energy of UM-4:
  \(\rho_\Lambda = \rho_{\mathrm{ff}} + \rho_{\mathrm{vac}} - \lambda_{\mathrm{F}}V_*\)
  (four-dimensional normalisation). Unit: GeV\(^4\). Value:
  \(0.685\,\rho_{\mathrm{c}}^{(0)} = 2.5\times10^{-47}\) GeV\(^4\)
  (Planck 2018). Status: DERIVED.
\item
  \(G_{\mathrm{N}}\): Newton's constant of four-dimensional gravity (the
  gravity of our world, on the fold),
  \(G_{\mathrm{N}} = 1/M_{\mathrm{Pl}}^2\) with the Planck mass
  \(M_{\mathrm{Pl}} = 1.221\times10^{19}\) GeV. Unit: GeV\(^{-2}\).
  Value: \(6.708\times10^{-39}\) GeV\(^{-2}\). Status: STANDARD.
\item
  \(V_{\mathrm{t}}\): The coordinate volume of the (compact) extra
  times,
  \(V_{\mathrm{t}} = \int\mathrm{d}x_5\,\mathrm{d}x_6\,\mathrm{d}x_7\).
  Unit: length\(^3\). Status: ASSUMED.
\item
  \(\rho_{\mathrm{c}}\): The critical density, the energy density of a
  spatially flat fold at the expansion rate \(H_4\);
  \(\rho_{\mathrm{c}}^{(0)}\) its value today. Unit: GeV\(^4\). Value:
  \(\rho_{\mathrm{c}}^{(0)} = 1.0537\times10^{-5}\,h_{100}^2\)
  GeV/cm\(^3\). Status: STANDARD.
\item
  \(\rho_{\mathrm{ff}}\): The energy density of the fold fluid. Unit:
  GeV\(^4\). Value:
  \(\lambda_{\mathrm{F}}V_* + \rho_\Lambda - \rho_{\mathrm{vac}}\):
  together with the vacuum energy \(\rho_{\mathrm{vac}}\) of the fields
  of the fold it makes up the tension \(\lambda_{\mathrm{F}}V_*\) that
  the static block requires plus the dark energy \(\rho_\Lambda\).
  Status: ASSUMED (the fluid), DERIVED (the value).
\item
  \(\rho_{\mathrm{vac}}\): The vacuum energy of the fields of the fold:
  the minimum of their potential (for the Standard Model the minimum of
  the Higgs potential with its loop corrections), a tension of the
  seven-dimensional fold. Unit: GeV\(^4\). Status: STANDARD.
\item
  \(H_4\): The expansion rate (Hubble rate) of the fold, \(H_4 = a_4'\),
  as a function of the time or of the redshift; \(H_4^{(0)}\) its value
  today, the Hubble constant. Unit: GeV (1/time). Value:
  \(H_4^{(0)} = 100\,h_{100}\) km/s/Mpc with \(h_{100} = 0.674\) (Planck
  2018): 67.4 km/s/Mpc \(= 1.44\times10^{-42}\) GeV. Status: STANDARD.
\item
  \(M_{\mathrm{Pl}}\): The Planck mass,
  \(M_{\mathrm{Pl}} = G_{\mathrm{N}}^{-1/2}\). Unit: GeV. Value:
  \(1.221\times10^{19}\) GeV. Status: STANDARD.
\item
  \(h_{100}\): The present expansion rate of the fold in units of 100
  km/s/Mpc. Unit: dimensionless. Value: 0.674 (Planck 2018). Status:
  STANDARD.
\end{itemize}

(\(\hat G^a{}_b\) the Einstein tensor of the metric FD-2 of the fold):
the tension plus matter that obeys the Einstein equations of seven
dimensions, \(\hat G^a{}_b = \kappa_7\mathsf{t}^a{}_b\) with
\(\kappa_7 = 5\kappa H/2\) (\(\mathsf{t}^a{}_b\) the matter of the
fold), which for frozen extra times are the Friedmann equations of
general relativity. The fold is therefore not only a place but a
physical object: a thin layer that carries energy and momentum (a
brane). 4. \emph{The author's coordinate \(x_0\) degenerates on it.} The
metric component \(g_{x_0x_0} = \cot^2(6Hx_0)\) vanishes on the fold and
nowhere else in the period; near the fold the proper distance grows with
the square of the coordinate distance,

\[
\lvert\sigma\rvert = 3H\Big(x_0 - \frac{\pi}{12H}\Big)^2\Big[1 + 6H^2\Big(x_0 - \frac{\pi}{12H}\Big)^2 + \dots\Big] \tag{FD-6}
\] \textbf{Symbols in (FD-6).}

\begin{itemize}
\tightlist
\item
  \(\sigma\): The proper coordinate of the hidden direction:
  \(\mathrm{d}\sigma = \lvert\cot z\rvert\,\mathrm{d}x_0\),
  \(\sigma = \mathrm{sgn}(\pi/2 - z)\ln(\sin z)/(6H)\) on one period
  \(0 < z < \pi\); \(\sigma = 0\) is the fold \(z = \pi/2\),
  \(\sigma \to \mp\infty\) the tips \(z \to 0, \pi\); one period of
  \(x_0\) is the whole line \(-\infty < \sigma < \infty\). Unit: length.
  Status: DERIVED.
\item
  \(H\): The author's constant (an inverse length) that fixes the period
  \(\pi/(6H)\) of \(x_0\), the warp factor
  \(\mathrm{e}^{-2H\lvert\sigma\rvert}\) of the metric, the cosmological
  constant \(\Lambda = -21H^2\) and the tension
  \(\kappa\lambda_{\mathrm{F}} = 12H\) of the fold; in UM-4 a free
  parameter, bounded from below by the propagation of the gauge fields
  in the bulk. Unit: GeV (1/length). Value: \(H \gtrsim 2.5\) TeV (from
  the invisible width of the Z boson). Status: AUTHOR.
\item
  \(x_0\): The hidden space-like coordinate (listed first); the metric
  depends on it through \(z = 6Hx_0\), periodically with the period
  \(\pi/(6H)\). Unit: length. Status: AUTHOR.
\item
  \(\pi\): The number pi, the ratio of the circumference of a circle to
  its diameter. Unit: dimensionless. Value: 3.141592654. Status:
  STANDARD.
\end{itemize}

so the two coordinate values on either side of \(x_0 = \pi/(12H)\) at
the same coordinate distance lie at the same proper distance from the
fold: the coordinate line of \(x_0\) is folded onto the line of proper
distances (a fold map, example 5).

\textbf{What lives on the fold.} In UM-4 the Higgs doublet (with the
author's ad hoc vacuum, a time-dependent choice of its vacuum value),
three generations of quarks and leptons, three right-handed neutrinos
and the fold fluid (its tension and the pressure along the extra times
that keeps their volume fixed) live on the fold: they are fields at
\(\sigma = 0\), functions of \(x_1, \dots, x_4\) (no field depends on
the extra times). The gauge fields of
\(SU(3)_C \times SU(2)_L \times U(1)_Y\) live in the bulk; their
massless modes are bound to the fold, so that the photon and the gluons
of the four-dimensional world are waves of the bulk that the fold holds,
and photons that leave the fold travel through the bulk at the speed of
light. dirac16complex (the author's spinor field of eight dimensions)
and gravity live in the bulk \(\sigma \ne 0\) as well as on the fold;
the states \(\chi\) of dirac16complex that the fold holds (in UM-4
long-lived resonances) are the part of it near the fold, and the portals
(couplings between dirac16complex and the fields of the fold) act at
\(\sigma = 0\). Integrated over the extra times, the fold is the
four-dimensional universe that we observe, with the metric
\(\mathrm{diag}(\mathrm{e}^{2a_4}, \mathrm{e}^{2a_4}, \mathrm{e}^{2a_4}, -1)\)
in \(x_1, \dots, x_4\). In UM-4 files ``the fold'' means this surface;
the examples below use the word for other surfaces with the same
properties. The word is not used in its other senses: a set of equal
energy levels is called degenerate, and a growth by the factor
\(\mathrm{e}\) is called an e-folding.

\textbf{Examples of folds that are meaningful for UM-4.}

\begin{enumerate}
\def\labelenumi{\arabic{enumi}.}
\tightlist
\item
  \emph{Our fold.} The fold \(z = \pi/2\) of the period that contains
  our world (FD-1): the seven-dimensional world-volume of the quarks,
  leptons and the Higgs field, and after the integration over the extra
  times the four-dimensional universe that we observe.
\item
  \emph{The copies of the fold.} Because the metric repeats with the
  period \(\pi/(6H)\), every surface
  \(x_0 = \pi/(12H) + (\text{an integer multiple of } \pi/(6H))\) is a
  copy of the fold, one per period. Neighbouring copies are separated by
  a tip (\(z\) a multiple of \(\pi\)), where the warp factor
  \((\sin^2z)^{1/6}\) vanishes and \(g_{x_0x_0} = \cot^2z\) diverges,
  and which lies at an infinite proper distance from both: in proper
  distance every copy has its own complete line
  \(-\infty < \sigma < \infty\). A compact direction of length \(L\)
  would give every field a tower of particles with the masses
  \(2\pi n/L\), \(n = 1, 2, \dots\) (a Kaluza-Klein tower); because the
  copies are infinitely far apart, the period \(\pi/(6H)\) of \(x_0\)
  produces no such tower with the spacing \(12H\): every field of the
  bulk has instead a continuum of modes that starts at zero mass, and
  the fold holds its massless mode.
\item
  \emph{A folded sheet of paper.} The crease of a folded sheet is a fold
  in this sense: the two halves are mirror images (property 1) joined
  along the crease, the sheet is continuous there and its direction
  changes abruptly (property 3). In UM-4 the direction across the crease
  is the hidden direction \(\sigma\), the seven directions along the
  crease are \(x_1, \dots, x_7\), and the crease is the fold.
\item
  \emph{A folded function.} The warp factor
  \(\mathrm{e}^{-2H\lvert\sigma\rvert}\) is the simplest folded
  function: continuous, mirror symmetric, largest at \(\sigma = 0\),
  with a jump of its derivative,
\end{enumerate}

\[
\frac{\mathrm{d}}{\mathrm{d}\sigma}\mathrm{e}^{-2H\lvert\sigma\rvert} = -2H\,\mathrm{sgn}(\sigma)\,\mathrm{e}^{-2H\lvert\sigma\rvert},\qquad
\frac{\mathrm{d}^2}{\mathrm{d}\sigma^2}\mathrm{e}^{-2H\lvert\sigma\rvert} = 4H^2\,\mathrm{e}^{-2H\lvert\sigma\rvert} - 4H\,\delta_{\mathrm{D}}(\sigma) \tag{FD-7}
\] \textbf{Symbols in (FD-7).}

\begin{itemize}
\tightlist
\item
  \(\mathrm{d}\): The differential: \(\mathrm{d}\) in front of a
  quantity is an infinitesimal change of it (as \(\mathrm{d}\sigma\), a
  small step of \(\sigma\)); the quotient of two such changes is a
  derivative; \(\mathrm{d}\) with an exponent, in front of several
  variables, is their volume element, the product of their
  differentials. Status: STANDARD.
\item
  \(\sigma\): The proper coordinate of the hidden direction:
  \(\mathrm{d}\sigma = \lvert\cot z\rvert\,\mathrm{d}x_0\),
  \(\sigma = \mathrm{sgn}(\pi/2 - z)\ln(\sin z)/(6H)\) on one period
  \(0 < z < \pi\); \(\sigma = 0\) is the fold \(z = \pi/2\),
  \(\sigma \to \mp\infty\) the tips \(z \to 0, \pi\); one period of
  \(x_0\) is the whole line \(-\infty < \sigma < \infty\). Unit: length.
  Status: DERIVED.
\item
  \(\mathrm{e}\): Euler's number, the base of the natural logarithm;
  \(\mathrm{e}^{X}\) is the exponential of \(X\). Unit: dimensionless.
  Value: 2.718281828. Status: STANDARD.
\item
  \(H\): The author's constant (an inverse length) that fixes the period
  \(\pi/(6H)\) of \(x_0\), the warp factor
  \(\mathrm{e}^{-2H\lvert\sigma\rvert}\) of the metric, the cosmological
  constant \(\Lambda = -21H^2\) and the tension
  \(\kappa\lambda_{\mathrm{F}} = 12H\) of the fold; in UM-4 a free
  parameter, bounded from below by the propagation of the gauge fields
  in the bulk. Unit: GeV (1/length). Value: \(H \gtrsim 2.5\) TeV (from
  the invisible width of the Z boson). Status: AUTHOR.
\item
  \(\mathrm{sgn}\): The sign function: \(+1\) for a positive argument
  and \(-1\) for a negative one. Status: STANDARD.
\item
  \(\delta_{\mathrm{D}}\): The Dirac delta function of \(\sigma\): zero
  for \(\sigma \ne 0\), with the integral 1 over any range that contains
  \(\sigma = 0\); the integral of \(\delta_{\mathrm{D}}(\sigma)\) times
  a smooth function of \(\sigma\) is the value of that function at
  \(\sigma = 0\). Unit: 1/length (for \(\delta_{\mathrm{D}}(\sigma)\)).
  Status: STANDARD.
\end{itemize}

whose delta function is the origin of the thin source of property 3.
Every metric function of FD-4 depends on \(\sigma\) through
\(\lvert\sigma\rvert\) and is folded at \(\sigma = 0\) in this sense. 5.
\emph{A fold map.} In the mathematics of smooth maps the fold (Whitney's
fold singularity) is the map \(u \mapsto u^2\) of a line onto a
half-line: it folds the line at \(u = 0\), two points \(\pm u\) land on
the same point, and the derivative of the map vanishes on the fold. Near
the fold the map from the author's coordinate \(x_0\) to the distance
\(\lvert\sigma\rvert\) from the fold is of this kind (FD-6), and
\(g_{x_0x_0} = \cot^2z\), the square of the derivative
\(\mathrm{d}\sigma/\mathrm{d}x_0\), vanishes there. 6. \emph{The brane
of Randall and Sundrum.} The five-dimensional metric

\[
\mathrm{d}y_{\mathrm{RS}}^2 + \mathrm{e}^{-2k_{\mathrm{RS}}\lvert y_{\mathrm{RS}}\rvert}\big(\mathrm{d}x_1^2 + \mathrm{d}x_2^2 + \mathrm{d}x_3^2
 - \mathrm{d}x_4^2\big),\qquad \lambda_{\mathrm{RS}} = \frac{6k_{\mathrm{RS}}}{\kappa_5},\qquad \Lambda_{\mathrm{RS}} = -\frac{6k_{\mathrm{RS}}^2}{\kappa_5} \tag{FD-8}
\] \textbf{Symbols in (FD-8).}

\begin{itemize}
\tightlist
\item
  \(\mathrm{d}\): The differential: \(\mathrm{d}\) in front of a
  quantity is an infinitesimal change of it (as \(\mathrm{d}\sigma\), a
  small step of \(\sigma\)); the quotient of two such changes is a
  derivative; \(\mathrm{d}\) with an exponent, in front of several
  variables, is their volume element, the product of their
  differentials. Status: STANDARD.
\item
  \(y_{\mathrm{RS}}\): The proper coordinate across the brane of the
  five-dimensional model of Randall and Sundrum (FD-8, an example, not
  UM-4). Unit: length. Status: STANDARD.
\item
  \(\mathrm{e}\): Euler's number, the base of the natural logarithm;
  \(\mathrm{e}^{X}\) is the exponential of \(X\). Unit: dimensionless.
  Value: 2.718281828. Status: STANDARD.
\item
  \(k_{\mathrm{RS}}\): The inverse curvature length of the bulk of the
  model of Randall and Sundrum: the warp factor of the brane is
  \(\mathrm{e}^{-2k_{\mathrm{RS}}\lvert y_{\mathrm{RS}}\rvert}\) (FD-8,
  an example, not UM-4). Unit: GeV (1/length). Status: STANDARD.
\item
  \(x_1\): The first coordinate of ordinary 3-space. Unit: length.
  Status: AUTHOR.
\item
  \(x_2\): The second coordinate of ordinary 3-space. Unit: length.
  Status: AUTHOR.
\item
  \(x_3\): The third coordinate of ordinary 3-space. Unit: length.
  Status: AUTHOR.
\item
  \(x_4\): The time; a prime on a function of \(x_4\) alone is
  \(\mathrm{d}/\mathrm{d}x_4\). Unit: length (= time, \(c = 1\)).
  Status: AUTHOR.
\item
  \(\lambda_{\mathrm{RS}}\): The tension (energy per unit volume of
  ordinary space) of the brane of the model of Randall and Sundrum,
  \(\lambda_{\mathrm{RS}} = 6k_{\mathrm{RS}}/\kappa_5\) (FD-8, an
  example, not UM-4). Unit: GeV\(^4\). Status: STANDARD.
\item
  \(\kappa_5\): The gravitational coupling of five dimensions in the
  model of Randall and Sundrum, \(8\pi\) times the Newton constant of
  five dimensions (FD-8, an example, not UM-4). Unit: GeV\(^{-3}\).
  Status: STANDARD.
\item
  \(\Lambda_{\mathrm{RS}}\): The vacuum energy density of the bulk of
  the model of Randall and Sundrum,
  \(\Lambda_{\mathrm{RS}} = -6k_{\mathrm{RS}}^2/\kappa_5\) (FD-8, an
  example, not UM-4). Unit: GeV\(^5\). Status: STANDARD.
\end{itemize}

has a fold at \(y_{\mathrm{RS}} = 0\) with the properties 1 to 3 (its
coordinate \(y_{\mathrm{RS}}\) is already the proper distance, so
property 4 does not arise); Einstein's equations require there a brane
of the tension \(\lambda_{\mathrm{RS}}\) (Israel's junction condition)
and in the bulk the vacuum energy \(\Lambda_{\mathrm{RS}}\), and the
bulk is otherwise empty. The fold of UM-4 is of the same kind, with
every direction of the fold warped by the same factor as in the model of
Randall and Sundrum: its static block is a space of constant curvature
\(-H^2\) in all eight directions, a vacuum of Einstein gravity with
\(\Lambda = -21H^2\), with a fold of pure tension (FD-12), and the
graviton is bound to the fold as there. It differs in that the fold has
seven dimensions, three of them extra times, in that the space is not
cut at the brane (Randall and Sundrum identify the two sides of their
brane, one physical side; the fold of UM-4 has two distinct sides,
example 7), and in that the gauge fields live in the bulk as well. 7.
\emph{A particle held by a thin layer.} A particle held by an attractive
thin layer has a wave function that falls off exponentially with the
distance on both sides, and the layer makes its derivative jump (for the
first-order Dirac equation, the wave function itself). In UM-4 the
states \(\chi\) of dirac16complex near the fold are of this kind only
near the fold: at rest and without their self-interaction (the field is
\(\Psi = \mathrm{e}^{(7/2)H\lvert\sigma\rvert}\chi(\sigma)\mathrm{e}^{-\mathrm{i}\omega x_4}\))

\[
\begin{aligned}
&\chi(\sigma) = \mathrm{e}^{-q_{\mathrm{b}}\lvert\sigma\rvert}\chi(0^+)\ \ (0 < \sigma \ll \sigma_{\mathrm{t}}),\qquad
\chi(\sigma) = \mathrm{e}^{-q_{\mathrm{b}}\lvert\sigma\rvert}\chi(0^-)\ \ (0 < -\sigma \ll \sigma_{\mathrm{t}}),\\
&\chi(0^+) = \exp\Big(\gamma^{(0)}\big(\mu_{\mathrm{L}}P_{\mathrm{L}} + \mu_{\mathrm{R}}P_{\mathrm{R}}\big)\Big)\chi(0^-),\qquad
\omega_0 = \omega_{\mathrm{r}} - \frac{\mathrm{i}}{2}\,\Gamma_\chi,\qquad \omega_{\mathrm{r}} \to \omega_{\mathrm{b}} = \sqrt{m_\Psi^2 - q_{\mathrm{b}}^2}
\end{aligned} \tag{FD-9}
\] \textbf{Symbols in (FD-9).}

\begin{itemize}
\tightlist
\item
  \(\chi(\sigma)\): The profile of a mode of dirac16complex along the
  hidden direction,
  \(\Psi = \mathrm{e}^{\frac72H\lvert\sigma\rvert}\chi(\sigma)\mathrm{e}^{\mathrm{i}(\vec k\cdot\vec x - \omega x_4)}\)
  (a 16-component spinor function of \(\sigma\)); as a four-dimensional
  field, \(\chi\) denotes the particle of the level of dirac16complex at
  the fold (a long-lived resonance) and \(\bar\chi\) its antiparticle.
  Unit: GeV\(^{2}\) (with the \(V_*\) normalisation). Status: DERIVED.
\item
  \(\sigma\): The proper coordinate of the hidden direction:
  \(\mathrm{d}\sigma = \lvert\cot z\rvert\,\mathrm{d}x_0\),
  \(\sigma = \mathrm{sgn}(\pi/2 - z)\ln(\sin z)/(6H)\) on one period
  \(0 < z < \pi\); \(\sigma = 0\) is the fold \(z = \pi/2\),
  \(\sigma \to \mp\infty\) the tips \(z \to 0, \pi\); one period of
  \(x_0\) is the whole line \(-\infty < \sigma < \infty\). Unit: length.
  Status: DERIVED.
\item
  \(\mathrm{e}\): Euler's number, the base of the natural logarithm;
  \(\mathrm{e}^{X}\) is the exponential of \(X\). Unit: dimensionless.
  Value: 2.718281828. Status: STANDARD.
\item
  \(q_{\mathrm{b}}\): The inverse width of a state of dirac16complex at
  the fold: inside the turning point its profile falls as
  \(\mathrm{e}^{-q_{\mathrm{b}}\lvert\sigma\rvert}\). Unit: GeV. Status:
  DERIVED.
\item
  \(\sigma_{\mathrm{t}}\): The turning point of a state of
  dirac16complex of the frequency \(\omega\): the depth at which the
  warped frequency equals the bulk mass,
  \(\lvert\omega\rvert\mathrm{e}^{H\sigma_{\mathrm{t}}} = m_\Psi\);
  beyond it every solution oscillates. Unit: length. Status: DERIVED.
\item
  \(\gamma^M,\ \gamma^{(a)}\): The Dirac matrices of dirac16complex
  along the coordinates, \(\gamma^M = \mathbf{e}_a{}^MT^a\),
  \(\gamma_M = g_{MN}\gamma^N\); \(\gamma^{(a)} = T^a\) along the frame
  direction \((a)\) (\(\gamma^{(0)}\) along \(\sigma\), \(\gamma^{(4)}\)
  along the time). Unit: dimensionless. Status: AUTHOR.
\item
  \(\mu_X,\ \mu_{\mathrm{L}},\ \mu_{\mathrm{R}}\): The dimensionless
  strengths of the fold for the left-handed and right-handed halves of
  dirac16complex, set by the mean fields of the Standard Model at the
  fold: the jump of the profile across the fold is the exponential of
  \(\gamma^{(0)}(\mu_{\mathrm{L}}P_{\mathrm{L}} + \mu_{\mathrm{R}}P_{\mathrm{R}})\);
  negative values attract. Unit: dimensionless. Status: DERIVED.
\item
  \(P_{\mathrm{L}}\): The author's left-handed chiral projector,
  \(P_{\mathrm{L}} = (\mathbf{1} - \Gamma_8)/2\). Unit: dimensionless.
  Status: AUTHOR.
\item
  \(P_{\mathrm{R}}\): The author's right-handed chiral projector,
  \(P_{\mathrm{R}} = (\mathbf{1} + \Gamma_8)/2\). Unit: dimensionless.
  Status: AUTHOR.
\item
  \(\omega_0\): The energy at rest (zero momentum along ordinary space
  and along the extra times) of a state of dirac16complex at the fold;
  complex for a resonance,
  \(\omega_0 = \omega_{\mathrm{r}} - \frac{\mathrm{i}}{2}\Gamma_\chi\).
  Unit: GeV. Status: DERIVED.
\item
  \(\omega_{\mathrm{r}}\): The mass of the particle \(\chi\): the real
  part of the energy at rest of the resonance of dirac16complex at the
  fold,
  \(\omega_{\mathrm{r}} = \omega_{\mathrm{b}}(1 - H/(2q_{\mathrm{b}})) + \dots\).
  Unit: GeV. Status: DERIVED.
\item
  \(\mathrm{i}\): The imaginary unit, \(\mathrm{i}^2 = -1\). Unit:
  dimensionless. Status: STANDARD.
\item
  \(\Gamma_\chi\): The width of the particle \(\chi\) at rest: the rate
  of its decay into the waves of dirac16complex in the bulk. Unit: GeV.
  Value: below \(1/t_0 = 1.5\times10^{-42}\) GeV for dark matter.
  Status: DERIVED.
\item
  \(\omega_{\mathrm{b}}\): The energy at zero momentum of the state of
  dirac16complex at the fold in the limit \(H \to 0\) (the bound state
  of UM-3): the particle \(\chi\) of UM-4 has the mass
  \(\omega_{\mathrm{r}} = \omega_{\mathrm{b}}(1 - H/(2q_{\mathrm{b}})) + \dots\),
  equal to \(\omega_{\mathrm{b}}\) up to the relative correction
  \(H/(2q_{\mathrm{b}})\). Unit: GeV. Status: DERIVED.
\item
  \(m_\Psi\): The mass parameter of dirac16complex in the bulk. Unit:
  GeV. Status: RECORD.
\item
  \(\Gamma_8\): The chirality matrix of eight dimensions,
  \(\Gamma_8 = T_0T_1\cdots T_7\) (the author's T16\^{}A{[}8{]})
  \(= \mathrm{diag}(-\mathbf{1}_8, \mathbf{1}_8)\);
  \(\Gamma_8^2 = \mathbf{1}\); it anticommutes with every \(T_a\). Unit:
  dimensionless. Status: AUTHOR.
\item
  \(\Psi,\ \bar\Psi\): dirac16complex: the author's 16-component complex
  spinor field of eight dimensions;
  \(\bar\Psi = \Psi^\dagger\sigma_{16}\) is its Dirac adjoint; in a term
  of the fold, \(\Psi\) is its value at \(\sigma = 0\). Unit:
  GeV\(^{7/2}\). Status: AUTHOR.
\item
  \(\mathbf{1},\ \mathbf{1}_8\): The identity matrix of the size fixed
  by the context (\(16 \times 16\) on dirac16complex, \(4 \times 4\) on
  a four-dimensional spinor, \(3 \times 3\) in generation or flavour
  space); \(\mathbf{1}_8\) is the \(8 \times 8\) identity matrix. Unit:
  dimensionless. Status: STANDARD.
\item
  \(\mathbf{e}_a{}^M,\ \mathbf{e}^a{}_M\): The vielbein (frame field) of
  the metric: \(g_{MN} = \eta_{ab}\,\mathbf{e}^a{}_M\mathbf{e}^b{}_N\);
  \(\mathbf{e}_a{}^M\) is its inverse;
  \(\mathbf{e}_{bN} = \eta_{bc}\mathbf{e}^c{}_N\); \(\mathbf{e}^{(a)}\)
  is the frame one-form \(\mathbf{e}^a{}_M\,\mathrm{d}X^M\). Unit:
  dimensionless. Status: STANDARD.
\item
  \(\omega\): A frequency: the energy of a mode
  \(\propto \mathrm{e}^{-\mathrm{i}\omega x_4}\) (\(\omega_o\) of the
  orbital \(o\)). Unit: GeV. Status: STANDARD.
\item
  \(\vec k\): The (comoving) momentum along ordinary space, \(\vec k^2\)
  its square. Unit: GeV. Status: STANDARD.
\item
  \(\vec x\): The position in ordinary space,
  \(\vec x = (x_1, x_2, x_3)\). Unit: length. Status: AUTHOR.
\item
  \(t_0\): The age of the universe. Unit: time. Value: 13.8 Gyr. Status:
  STANDARD.
\item
  \(\sigma_{16}\): The matrix of the Dirac adjoint,
  \(\sigma_{16} = T_0T_1T_2T_3\), real and symmetric; \(\sigma_{16}T_a\)
  is real and antisymmetric. Unit: dimensionless. Status: AUTHOR.
\item
  \({}^\dagger\): Hermitian conjugation: the complex conjugate of the
  transpose of a matrix or spinor; for an operator on the space of
  states, its adjoint with respect to the scalar product of that space
  (for the quantised dirac16complex the indefinite Krein scalar
  product). Status: STANDARD.
\item
  \(o\): The label of a single-particle state of dirac16complex: a
  Kohn-Sham orbital, or a mode of the quantised field; a sum over \(o\)
  runs over the orbitals or the modes. Status: STANDARD.
\end{itemize}

where \(\mu_{\mathrm{L}}\), \(\mu_{\mathrm{R}}\) are the strengths of
the fold for the left- and right-handed halves of dirac16complex (set by
the fields of the fold),
\(q_{\mathrm{b}} = -m_\Psi(\mu_{\mathrm{L}} + \mu_{\mathrm{R}})\tanh\varsigma/(2\varsigma)\),
\(\varsigma = \sqrt{\mu_{\mathrm{L}}\mu_{\mathrm{R}}}\)
(\(q_{\mathrm{b}} = -m_\Psi\mu_{\mathrm{L}}/2\) for
\(\mu_{\mathrm{R}} = 0\)), and \(\omega_{\mathrm{r}}\) tends to the
energy \(\omega_{\mathrm{b}}\) of the bound state of UM-3 when
\(H/q_{\mathrm{b}} \to 0\). Because the time is warped, the energy
\(\omega\) of the field enters its equation multiplied by
\(\mathrm{e}^{H\lvert\sigma\rvert}\): beyond the turning point
\(\sigma_{\mathrm{t}}\)
(\(\omega_{\mathrm{r}}\mathrm{e}^{H\sigma_{\mathrm{t}}} = m_\Psi\)) the
profile oscillates as a wave that leaves the fold, so the state is a
resonance of the width \(\Gamma_\chi\), which falls like
\(\mathrm{e}^{-2F(\upsilon_{\mathrm{b}})m_\Psi/H}\) with
\(F(\upsilon) = \ln((1 + \sqrt{1 - \upsilon^2})/\upsilon) - \sqrt{1 - \upsilon^2}\)
and \(\upsilon_{\mathrm{b}} = \omega_{\mathrm{b}}/m_\Psi\) (the
tunnelling through the region between the fold and
\(\sigma_{\mathrm{t}}\)). The profile is not a function of
\(\lvert\sigma\rvert\): it jumps at the fold,
\(\chi(0^+) \ne \chi(0^-)\). 8. \emph{The fold as a mirror for
dirac16complex.} The reflection across the fold acts on the profiles
\(\chi(\sigma)\) in two ways,

\[
(\mathsf{R}_1\chi)(\sigma) = T_0T_1T_2T_3\,\chi(-\sigma),\qquad (\mathsf{R}_2\chi)(\sigma) = T_0T_5T_6T_7\,\chi(-\sigma) \tag{FD-10}
\] \textbf{Symbols in (FD-10).}

\begin{itemize}
\tightlist
\item
  \(\mathsf{R}_1\): The reflection across the fold acting on a profile
  of dirac16complex together with the reflection of ordinary space,
  \((\mathsf{R}_1\chi)(\sigma) = T_0T_1T_2T_3\,\chi(-\sigma)\) (FD-10):
  a symmetry of dirac16complex, of gravity and of the mean field of the
  gas of the particles \(\chi\); the states at the fold of Krein charge
  \(-1\) are even, those of Krein charge \(+1\) odd. Unit:
  dimensionless. Status: DERIVED.
\item
  \(\mathsf{R}_2\): The reflection across the fold acting on a profile
  of dirac16complex together with the reflection of the extra times,
  \((\mathsf{R}_2\chi)(\sigma) = T_0T_5T_6T_7\,\chi(-\sigma)\) (FD-10):
  it maps the states at the fold of one Krein charge onto those of the
  other. Unit: dimensionless. Status: DERIVED.
\item
  \(\chi(\sigma)\): The profile of a mode of dirac16complex along the
  hidden direction,
  \(\Psi = \mathrm{e}^{\frac72H\lvert\sigma\rvert}\chi(\sigma)\mathrm{e}^{\mathrm{i}(\vec k\cdot\vec x - \omega x_4)}\)
  (a 16-component spinor function of \(\sigma\)); as a four-dimensional
  field, \(\chi\) denotes the particle of the level of dirac16complex at
  the fold (a long-lived resonance) and \(\bar\chi\) its antiparticle.
  Unit: GeV\(^{2}\) (with the \(V_*\) normalisation). Status: DERIVED.
\item
  \(\sigma\): The proper coordinate of the hidden direction:
  \(\mathrm{d}\sigma = \lvert\cot z\rvert\,\mathrm{d}x_0\),
  \(\sigma = \mathrm{sgn}(\pi/2 - z)\ln(\sin z)/(6H)\) on one period
  \(0 < z < \pi\); \(\sigma = 0\) is the fold \(z = \pi/2\),
  \(\sigma \to \mp\infty\) the tips \(z \to 0, \pi\); one period of
  \(x_0\) is the whole line \(-\infty < \sigma < \infty\). Unit: length.
  Status: DERIVED.
\item
  \(T_a,\ T^a\): The author's eight real \(16 \times 16\) Dirac matrices
  T16\^{}A{[}a{]} of 4+4.nb (\(a = 0, \dots, 7\); \(T_0, \dots, T_7\)),
  with \(\{T_a, T_b\} = 2\eta_{ab}\mathbf{1}\); \(T^a = \eta^{ab}T_b\).
  Unit: dimensionless. Status: AUTHOR.
\end{itemize}

(\(T_a\) the author's Dirac matrices; for states at rest
\(\mathsf{R}_1\) is the reflection across the fold combined with the
reflection of ordinary space, \(\mathsf{R}_2\) the reflection across the
fold combined with the reflection of the extra times). The Krein charge
of a state is its share of the conserved charge of dirac16complex, which
is indefinite: a state can carry the charge \(-1\) or \(+1\).
\(\mathsf{R}_1\) is a symmetry of dirac16complex, of gravity and of a
mean field of a gas of the particles \(\chi\) that is the same on both
sides of the fold: the four resonances of Krein charge \(-1\) (the
particles \(\chi\)) are even under it and the four of Krein charge
\(+1\) are odd, so such a mean field cannot mix the two signs of the
indefinite charge. \(\mathsf{R}_2\) maps the states of one Krein charge
onto those of the other and makes the two sets degenerate without the
self-interaction. 9. \emph{The planes of Horava and Witten.} In M-theory
(the eleven-dimensional theory whose limits include the superstring
theories) on an interval (a circle whose points are identified with
their mirror images), the two ten-dimensional planes left in place by
the reflection are mirror surfaces (property 1); each carries one
\(E_8\) gauge group while gravity fills the bulk. As in UM-4, gravity
fills the bulk; unlike UM-4, the gauge fields live on the mirror planes
(in UM-4 they live in the bulk), and the two sides of each plane are
identified (one physical side).

\textbf{What is not a fold.}

\begin{itemize}
\tightlist
\item
  \emph{The tips} \(z = 0\) and \(z = \pi\) (mod \(\pi\)): the mirror
  exchanges them within a period, but there the warp factor vanishes,
  \(g_{x_0x_0}\) diverges, and they lie at an infinite proper distance;
  they carry no thin source.
\item
  \emph{A smooth maximum of a warp factor}, such as
  \(1/\cosh^2(H\sigma)\): mirror symmetric and largest at
  \(\sigma = 0\), but without a kink and without a thin source.
\end{itemize}

\hypertarget{how-to-read-this-document}{%
\section{How to read this document}\label{how-to-read-this-document}}

\hypertarget{what-this-document-contains}{%
\subsection{What this document
contains}\label{what-this-document-contains}}

The author asked for three analyses of his universe model: the mass
mechanism along the hidden coordinate \(x_0\) (the period \(\pi/(6H)\)
of the metric and the uncertainty relation), whether dirac16complex can
be the dark matter, and whether it can be the dark energy. He also asked
to be told what he may have left out, and (on 2026-10-25) for answers to
the questions that the model had left open. On 2026-10-25 at 13:15 PDT
he ordered a new metric, in which the time \(x_4\) is warped like the
other directions of the fold and photons and gravitons propagate at the
speed of light in the bulk; the model built on it is the universe model
UM-4. This document gives the analyses and the answers for UM-4, and
compares UM-4 with its predecessor UM-3. All of them use the metric of
UM-4 (FE-27) and Einstein gravity in eight dimensions with the
cosmological constant \(\Lambda = -21H^2\) (the author's
Einstein-Lovelock theory with the weights of its Gauss-Bonnet and cubic
terms set to zero, the only choice of the weights for which gravitons
travel on the light cone on every background). It is self-contained:
every equation it uses is written in it, and every symbol of every
display equation is defined right after it; a symbol that one of these
definitions uses is defined in this document as well. The equations keep
the tag they carry in every UM-4 document: FE-n for a field equation or
a result of the field equations, IN-n for a term of the action, S-n for
an equation of the solution of UM-4; the equations that only this
document needs carry the tags A-1, A-2, and so on. The sections are:

\begin{itemize}
\tightlist
\item
  \textbf{The answers in brief.}
\item
  \textbf{The \(x_0\) mass mechanism}: the period of \(x_0\), the proper
  length of the hidden direction, the spectra that replace the tower of
  a compact dimension, and what the observer on the fold sees of them.
\item
  \textbf{dirac16complex as dark matter}: the stability and the lifetime
  of its particles, their gravitation on the fold, the equation of state
  of their Kohn-Sham states, the relic density, direct detection, and
  what is left of the window of parameters.
\item
  \textbf{dirac16complex and the dark energy}: why the gas of its
  particles cannot be dark energy, what the dark energy of UM-4 is, and
  how it compares with the baryon acoustic oscillations measured by DESI
  DR2 (the second data release of the Dark Energy Spectroscopic
  Instrument).
\item
  \textbf{The open questions answered}: whether the indefinite Krein
  norm of dirac16complex is stable, what fields that depend on the extra
  times do, whether anything moves faster than light, the primordial
  era, and the lightest neutrino mass and the Dirac phase.
\item
  \textbf{What the author may have left out}: the additions that UM-4
  needed, and the open problems.
\item
  \textbf{UM-4 compared with UM-3}: what the warped time changed.
\end{itemize}

\hypertarget{the-status-labels}{%
\subsection{The status labels}\label{the-status-labels}}

Every result carries one of the labels \textbf{AUTHOR} (from the
author's notebooks \texttt{4+4.nb} and
\texttt{Dirac\_Einstein-Lovelock\_4+4\_Nash.nb} and his orders of
2026-10-24 and 2026-10-25), \textbf{RECORD} (from the record of this
repository: the Lagrangian of dirac16complex with its canonical spin
connection and contact self-interaction, the Kohn-Sham treatment),
\textbf{STANDARD} (textbook physics or a published measurement, named),
\textbf{ASSUMED} (a choice of UM-4 without derivation),
\textbf{INVENTED} (a new interaction of UM-4: the portals of
dirac16complex), \textbf{DERIVED} (computed and verified in UM-4: the
program and the check that verify it are named), \textbf{ESTIMATE} (an
order-of-magnitude evaluation) and \textbf{OPEN} (not settled).

\hypertarget{every-equation-is-followed-by-its-symbols}{%
\subsection{Every equation is followed by its
symbols}\label{every-equation-is-followed-by-its-symbols}}

Every display equation is followed by the list of EVERY symbol that
appears in it, with its meaning, unit (natural units, in which the
reduced Planck constant \(\hbar\) and the speed of light \(c\) are 1),
value and status. The author's coordinates are
\(X = \{x_0, x_1, \dots, x_7\}\), always with the hidden direction
first: \(x_0\) the hidden space-like coordinate (the metric depends on
it through \(z = 6Hx_0\)), \(x_1, x_2, x_3\) ordinary space, \(x_4\) the
time, \(x_5, x_6, x_7\) three extra time-like coordinates; the signature
is \((+,+,+,+,-,-,-,-)\) in this order. \(H\) is the author's constant
(an inverse length), a free parameter of UM-4 bounded from below by
about 2.5 TeV. A prime on a function of the time \(x_4\) alone is its
derivative with respect to \(x_4\), on a profile (a function of the
proper coordinate \(\sigma\) of the hidden direction alone) its
derivative with respect to \(\sigma\). The Higgs doublet, the quarks and
the leptons live on the fold of the metric (the surface \(z = \pi/2\),
where \(\sigma = 0\); ASSUMED); gravity, dirac16complex and the gauge
fields of the Standard Model live in the bulk (the photon there by the
author's order; ASSUMED for the other gauge fields).

\hypertarget{the-answers-in-brief}{%
\subsection{The answers in brief}\label{the-answers-in-brief}}

\begin{enumerate}
\def\labelenumi{\arabic{enumi}.}
\tightlist
\item
  \textbf{The \(x_0\) mass mechanism.} The coordinate \(x_0\) has the
  period \(\pi/(6H)\), but one period is a hidden direction of infinite
  proper length, warped towards the fold: there is no compact dimension,
  no tower of Kaluza-Klein masses with the spacing \(12H\), and no
  uncertainty bound \(p_{\mathrm{h}} \gtrsim 6H\) on the momentum along
  the hidden direction. Because the time is warped like every other
  direction of the fold, the hidden direction gives no gap either (UM-3
  had the gaps \(3H\) for scalars and \(2H\) for gauge fields): every
  massless field of the bulk has one massless mode bound to the fold and
  a continuum of massive modes that starts at zero mass. The continuum
  reaches the fold with a weight that vanishes like the fourth power of
  the mass over \(H\) for the graviton and the scalars and like its
  square for the gauge fields; the observer on the fold sees it only as
  corrections to Newton's and Coulomb's laws at distances of the order
  of \(1/H\) and as decays of the W and Z bosons into the bulk, which
  bound \(H\) from below by 2.5 TeV. For dirac16complex the continuum
  starts at zero frequency, and the fold holds long-lived resonances
  (the particles \(\chi\)) with the mass
  \(\omega_{\mathrm{r}} \simeq \omega_{\mathrm{b}}(1 - H/(2q_{\mathrm{b}}))\),
  \(\omega_{\mathrm{b}} = \sqrt{m_\Psi^2 - q_{\mathrm{b}}^2}\), and a
  width that is exponentially small in \(m_\Psi/H\). Their masses come
  from the binding to the fold, not from the period of \(x_0\)
  (DERIVED).
\item
  \textbf{Dark matter.} The particles \(\chi\) cannot decay into the
  Standard Model alone, but they tunnel into the bulk: they live longer
  than the universe only if \(m_\Psi/H\) is above about
  \(1.1\times10^3\) (for the binding \(q_{\mathrm{b}} = m_\Psi/2\)), so
  with \(H > 2.5\) TeV they are heavier than \(2.4\times10^6\) GeV. They
  gravitate on the fold like matter of the fold, and a dilute gas of
  them is pressureless: cold dark matter. But they cannot be a thermal
  relic (unitarity allows a thermal relic only up to \(8.9\times10^4\)
  GeV), and bound to the fold by the author's Higgs portal alone they
  would have to be bound so weakly (direct detection) that they would
  live long enough only for \(H < 2.1\) GeV. dirac16complex can be the
  dark matter of UM-4 only with a binding to the fold that does not
  couple \(\chi\) to nucleons and with a non-thermal origin of its
  abundance, neither of which the action of UM-4 specifies (DERIVED,
  ESTIMATE, OPEN).
\item
  \textbf{Dark energy.} No: the gas of the particles \(\chi\) has a
  positive pressure that vanishes in the dilute limit, never a negative
  one. The dark energy of UM-4 is the excess of the tension of the fold
  over the value \(\lambda_{\mathrm{F}}\) that the static bulk requires,
  a vacuum energy with the equation of state \(w = -1\); the observed
  value needs the tension tuned to one part in \(10^{92}\) for
  \(H = 2.5\) TeV. It fits the cosmic microwave background (CMB) but
  keeps the tension of the cosmological-constant model with the baryon
  acoustic oscillations (BAO) of DESI DR2, 2.3 standard deviations in
  DESI's own analysis (DERIVED: \(w = -1\); STANDARD: the tension; OPEN:
  the tuning).
\item
  \textbf{The open questions.} The Krein norm: the vacuum is stable;
  whether a gas of \(\chi\) is classically stable is not decided by the
  computations; the chiral couplings of the fold (the author's
  left-handed Higgs portal among them) produce ghosts, so UM-4 with them
  is not unitary. The extra times: every field with a momentum along
  them is tachyonic, and compact extra times do not help; UM-4 needs a
  gauge principle that removes that momentum. Faster than light: in UM-4
  nothing is: every massless field of the bulk travels at the speed of
  light, every state of dirac16complex obeys
  \(\omega^2 = \omega_0^2 + k^2\) exactly, no signal through the static
  bulk arrives earlier than light along the fold, and around the
  expanding fold the largest advance is \(4.5\times10^{-94}\) of the
  travel time over the Hubble distance; the four time-like directions
  still allow closed time-like curves. The primordial era: the author's
  vacuum is not a solution; the exact de Sitter fold inflates the extra
  times as well; an inflation with frozen extra times needs a Weyl
  tensor of the bulk and, in the regime of general relativity, a rate
  below \(H/10\); the inflaton is not specified. The neutrinos: the
  Dirac phase drops out of every result; the lightest mass is below 23
  meV and changes only the capture rate of the cosmic neutrinos
  (DERIVED, ESTIMATE; OPEN: the stability of the gas, a unitary
  coupling, the gauge principle, the inflaton).
\item
  \textbf{What was left out.} UM-4 had to add right-handed quarks and
  Yukawa couplings, three right-handed neutrinos (accepted by the
  author), a thin fluid on the fold (a tension and a pressure along the
  extra times that holds their volume fixed), the gauge field that the
  author's ad hoc Higgs vacuum needs, the choice of the gravity
  (Einstein gravity) and the gauge fields in the bulk; the dark matter
  of UM-4 needs a binding of \(\chi\) to the fold and an origin of its
  abundance that the action does not contain. The last sections list the
  additions and the open problems and compare UM-4 with UM-3.
\end{enumerate}

\hypertarget{the-x_0-mass-mechanism}{%
\section{\texorpdfstring{The \(x_0\) mass
mechanism}{The x\_0 mass mechanism}}\label{the-x_0-mass-mechanism}}

\hypertarget{the-period-of-x_0-and-the-proper-length-of-the-hidden-direction}{%
\subsection{\texorpdfstring{The period of \(x_0\) and the proper length
of the hidden
direction}{The period of x\_0 and the proper length of the hidden direction}}\label{the-period-of-x_0-and-the-proper-length-of-the-hidden-direction}}

\[
\begin{aligned}
&g_{MN} = \mathrm{diag}\Big(\cot^2z,\ \mathrm{e}^{2a_4}(\sin^2z)^{1/6},\ \mathrm{e}^{2a_4}(\sin^2z)^{1/6},
 \ \mathrm{e}^{2a_4}(\sin^2z)^{1/6},\ -(\sin^2z)^{1/6},\\
&\qquad\qquad -\mathrm{e}^{-2a_4}b_4^2(\sin^2z)^{1/6},\ -\mathrm{e}^{-2a_4}b_4^2(\sin^2z)^{1/6},
 \ -\mathrm{e}^{-2a_4}b_4^2(\sin^2z)^{1/6}\Big),\qquad z = 6Hx_0,\\
&\sigma = \mathrm{sgn}\Big(\frac{\pi}{2} - z\Big)\frac{\ln\sin z}{6H},\qquad \mathrm{d}\sigma = \lvert\cot z\rvert\,\mathrm{d}x_0,\\
&g_{MN}\,\mathrm{d}X^M\mathrm{d}X^N = \mathrm{d}\sigma^2 + \mathrm{e}^{-2H\lvert\sigma\rvert}\Big[\mathrm{e}^{2a_4}\big(\mathrm{d}x_1^2
 + \mathrm{d}x_2^2 + \mathrm{d}x_3^2\big) - \mathrm{d}x_4^2 - \mathrm{e}^{-2a_4}b_4^2\big(\mathrm{d}x_5^2 + \mathrm{d}x_6^2 + \mathrm{d}x_7^2\big)\Big]
\end{aligned} \tag{FE-27}
\] \textbf{Symbols in (FE-27).}

\begin{itemize}
\tightlist
\item
  \(g_{MN}\): The metric tensor of the eight-dimensional spacetime
  (signature \(+,+,+,+,-,-,-,-\) in the order \(x_0, \dots, x_7\)):
  \(g_{MN}\,\mathrm{d}X^M\mathrm{d}X^N\) is the squared interval,
  \(g^{MN}\) its inverse, \(\lvert g\rvert\) the absolute value of its
  determinant (\(\sqrt{\lvert g\rvert}\,\mathrm{d}^8X\) is the proper
  volume element). Unit: dimensionless. Status: AUTHOR.
\item
  \(M\): A coordinate index of eight dimensions (also with a numbered
  copy \(M_1, M_2, \dots\)): it runs over \(x_0, x_1, \dots, x_7\); a
  repeated index is summed. Status: STANDARD.
\item
  \(N\): A coordinate index of eight dimensions (also
  \(N_1, N_2, \dots\)), running over \(x_0, \dots, x_7\); a repeated
  index is summed. Status: STANDARD.
\item
  \(\mathrm{diag}\): A diagonal matrix with the listed entries on its
  diagonal (in the order of the coordinates \(x_0, x_1, \dots, x_7\) for
  a tensor of eight dimensions, \(x_1, \dots, x_7\) for a tensor of the
  fold, \(x_1, \dots, x_4\) for a four-dimensional tensor). Status:
  STANDARD.
\item
  \(z\): The angle variable of the hidden direction, \(z = 6Hx_0\); one
  period of \(x_0\) is \(0 < z < \pi\), the fold is \(z = \pi/2\). Unit:
  dimensionless. Status: AUTHOR.
\item
  \(\mathrm{e}\): Euler's number, the base of the natural logarithm;
  \(\mathrm{e}^{X}\) is the exponential of \(X\). Unit: dimensionless.
  Value: 2.718281828. Status: STANDARD.
\item
  \(a_4\): The author's metric function of the time \(x_4\) that scales
  ordinary 3-space: on the fold \(\mathrm{e}^{a_4}\) is the scale factor
  of cosmology and \(a_4'\) the Hubble rate. Unit: dimensionless.
  Status: AUTHOR.
\item
  \(b_4\): The author's metric function of the time \(x_4\) in the
  extra-time components, \(b_4 > 0\), never identically 1 (the author's
  order of 2026-10-24); the scale factor (the factor of the lengths) of
  the extra times is \(\mathrm{e}^{-a_4}b_4\) on the fold and
  \(\mathrm{e}^{-a_4}b_4\,\mathrm{e}^{-H\lvert\sigma\rvert}\) at the
  distance \(\sigma\) from it (the square root of the warp factor times
  \(\mathrm{e}^{-a_4}b_4\)). Unit: dimensionless. Status: AUTHOR.
\item
  \(H\): The author's constant (an inverse length) that fixes the period
  \(\pi/(6H)\) of \(x_0\), the warp factor
  \(\mathrm{e}^{-2H\lvert\sigma\rvert}\) of the metric, the cosmological
  constant \(\Lambda = -21H^2\) and the tension
  \(\kappa\lambda_{\mathrm{F}} = 12H\) of the fold; in UM-4 a free
  parameter, bounded from below by the propagation of the gauge fields
  in the bulk. Unit: GeV (1/length). Value: \(H \gtrsim 2.5\) TeV (from
  the invisible width of the Z boson). Status: AUTHOR.
\item
  \(x_0\): The hidden space-like coordinate (listed first); the metric
  depends on it through \(z = 6Hx_0\), periodically with the period
  \(\pi/(6H)\). Unit: length. Status: AUTHOR.
\item
  \(\sigma\): The proper coordinate of the hidden direction:
  \(\mathrm{d}\sigma = \lvert\cot z\rvert\,\mathrm{d}x_0\),
  \(\sigma = \mathrm{sgn}(\pi/2 - z)\ln(\sin z)/(6H)\) on one period
  \(0 < z < \pi\); \(\sigma = 0\) is the fold \(z = \pi/2\),
  \(\sigma \to \mp\infty\) the tips \(z \to 0, \pi\); one period of
  \(x_0\) is the whole line \(-\infty < \sigma < \infty\). Unit: length.
  Status: DERIVED.
\item
  \(\mathrm{sgn}\): The sign function: \(+1\) for a positive argument
  and \(-1\) for a negative one. Status: STANDARD.
\item
  \(\pi\): The number pi, the ratio of the circumference of a circle to
  its diameter. Unit: dimensionless. Value: 3.141592654. Status:
  STANDARD.
\item
  \(\mathrm{d}\): The differential: \(\mathrm{d}\) in front of a
  quantity is an infinitesimal change of it (as \(\mathrm{d}\sigma\), a
  small step of \(\sigma\)); the quotient of two such changes is a
  derivative; \(\mathrm{d}\) with an exponent, in front of several
  variables, is their volume element, the product of their
  differentials. Status: STANDARD.
\item
  \(X = \{x_0, \dots, x_7\}\): The author's coordinates of the
  eight-dimensional spacetime, \(X = \{x_0, x_1, \dots, x_7\}\) with the
  hidden direction first (Dirac\_Einstein-Lovelock\_4+4\_Nash.nb,
  section DEFCHART); \(X^M\) is the coordinate with the index \(M\) and
  \(\mathrm{d}^8X = \mathrm{d}x_0\,\mathrm{d}x_1\cdots\mathrm{d}x_7\)
  the coordinate volume element. Unit: length (each coordinate). Status:
  AUTHOR.
\item
  \(x_1\): The first coordinate of ordinary 3-space. Unit: length.
  Status: AUTHOR.
\item
  \(x_2\): The second coordinate of ordinary 3-space. Unit: length.
  Status: AUTHOR.
\item
  \(x_3\): The third coordinate of ordinary 3-space. Unit: length.
  Status: AUTHOR.
\item
  \(x_4\): The time; a prime on a function of \(x_4\) alone is
  \(\mathrm{d}/\mathrm{d}x_4\). Unit: length (= time, \(c = 1\)).
  Status: AUTHOR.
\item
  \(x_5\): The first of the three extra time-like coordinates. Unit:
  length. Status: AUTHOR.
\item
  \(x_6\): The second of the three extra time-like coordinates. Unit:
  length. Status: AUTHOR.
\item
  \(x_7\): The third of the three extra time-like coordinates. Unit:
  length. Status: AUTHOR.
\end{itemize}

The metric of UM-4 and its proper chart. The metric depends on \(x_0\)
only through \(z = 6Hx_0\), so its coefficients repeat with the
coordinate period \(\pi/(6H)\): one period is \(0 < z < \pi\), with the
fold at \(z = \pi/2\) and the tips at \(z = 0\) and \(z = \pi\), where
\(g_{x_0x_0} = \cot^2z\) diverges and the warp factor
\((\sin^2z)^{1/6}\) vanishes. In the proper coordinate \(\sigma\) the
warp factor \(\mathrm{e}^{-2H\lvert\sigma\rvert}\) multiplies the whole
metric of the fold, the time included (UM-3: the time was not warped).
Status: AUTHOR (the metric and the order of 2026-10-25, 13:15 PDT),
DERIVED (the chart, \texttt{um4\_geometry.py}, G1).

\[
\int_0^{\pi/(6H)}\sqrt{\cot^2(6Hx_0)}\;\mathrm{d}x_0 = \frac{1}{6H}\int_0^{\pi}\lvert\cot z\rvert\,\mathrm{d}z = \infty,\qquad
\sigma(z) = \mathrm{sgn}\Big(\frac{\pi}{2} - z\Big)\frac{\ln\sin z}{6H} \to \mp\infty\ \ (z \to 0,\ \pi) \tag{A-1}
\] \textbf{Symbols in (A-1).}

\begin{itemize}
\tightlist
\item
  \(x_0\): The hidden space-like coordinate (listed first); the metric
  depends on it through \(z = 6Hx_0\), periodically with the period
  \(\pi/(6H)\). Unit: length. Status: AUTHOR.
\item
  \(H\): The author's constant (an inverse length) that fixes the period
  \(\pi/(6H)\) of \(x_0\), the warp factor
  \(\mathrm{e}^{-2H\lvert\sigma\rvert}\) of the metric, the cosmological
  constant \(\Lambda = -21H^2\) and the tension
  \(\kappa\lambda_{\mathrm{F}} = 12H\) of the fold; in UM-4 a free
  parameter, bounded from below by the propagation of the gauge fields
  in the bulk. Unit: GeV (1/length). Value: \(H \gtrsim 2.5\) TeV (from
  the invisible width of the Z boson). Status: AUTHOR.
\item
  \(\mathrm{d}\): The differential: \(\mathrm{d}\) in front of a
  quantity is an infinitesimal change of it (as \(\mathrm{d}\sigma\), a
  small step of \(\sigma\)); the quotient of two such changes is a
  derivative; \(\mathrm{d}\) with an exponent, in front of several
  variables, is their volume element, the product of their
  differentials. Status: STANDARD.
\item
  \(z\): The angle variable of the hidden direction, \(z = 6Hx_0\); one
  period of \(x_0\) is \(0 < z < \pi\), the fold is \(z = \pi/2\). Unit:
  dimensionless. Status: AUTHOR.
\item
  \(\pi\): The number pi, the ratio of the circumference of a circle to
  its diameter. Unit: dimensionless. Value: 3.141592654. Status:
  STANDARD.
\item
  \(\infty\): Infinity: as a limit of an integral or a sum it means that
  the range has no bound. Status: STANDARD.
\item
  \(\sigma\): The proper coordinate of the hidden direction:
  \(\mathrm{d}\sigma = \lvert\cot z\rvert\,\mathrm{d}x_0\),
  \(\sigma = \mathrm{sgn}(\pi/2 - z)\ln(\sin z)/(6H)\) on one period
  \(0 < z < \pi\); \(\sigma = 0\) is the fold \(z = \pi/2\),
  \(\sigma \to \mp\infty\) the tips \(z \to 0, \pi\); one period of
  \(x_0\) is the whole line \(-\infty < \sigma < \infty\). Unit: length.
  Status: DERIVED.
\item
  \(\mathrm{sgn}\): The sign function: \(+1\) for a positive argument
  and \(-1\) for a negative one. Status: STANDARD.
\end{itemize}

The proper length of one period of \(x_0\) is infinite:
\(\lvert\cot z\rvert\) diverges at the tips like \(1/z\), so its
integral diverges logarithmically, and the proper coordinate \(\sigma\)
runs over the whole line. The hidden direction is not a circle of
circumference \(\pi/(6H)\) but an infinite line, warped by
\(\mathrm{e}^{-H\lvert\sigma\rvert}\) away from the fold; the tips are
infinitely far away. The component \(g_{x_0x_0}\) is the author's and is
the same in UM-3 and UM-4. Status: DERIVED (\texttt{um4\_geometry.py},
G1: the proper chart on both halves of the period).

\[
\begin{aligned}
&p_{\mathrm{h}} \gtrsim \frac{\pi}{\pi/(6H)} = 6H,\qquad m_{\mathrm{KK}}^2 = m_\Psi^2 + \big(12H\,\mathsf{k}\big)^2\ \ (\mathsf{k} = 0, \pm1, \pm2, \dots)\\
&\text{(for a compact direction of length } \pi/(6H)\text{, not UM-4)}
\end{aligned} \tag{A-2}
\] \textbf{Symbols in (A-2).}

\begin{itemize}
\tightlist
\item
  \(p_{\mathrm{h}}\): The momentum along the hidden direction of a
  particle confined to it, if that direction were compact (not UM-4).
  Unit: GeV. Status: STANDARD.
\item
  \(\pi\): The number pi, the ratio of the circumference of a circle to
  its diameter. Unit: dimensionless. Value: 3.141592654. Status:
  STANDARD.
\item
  \(H\): The author's constant (an inverse length) that fixes the period
  \(\pi/(6H)\) of \(x_0\), the warp factor
  \(\mathrm{e}^{-2H\lvert\sigma\rvert}\) of the metric, the cosmological
  constant \(\Lambda = -21H^2\) and the tension
  \(\kappa\lambda_{\mathrm{F}} = 12H\) of the fold; in UM-4 a free
  parameter, bounded from below by the propagation of the gauge fields
  in the bulk. Unit: GeV (1/length). Value: \(H \gtrsim 2.5\) TeV (from
  the invisible width of the Z boson). Status: AUTHOR.
\item
  \(m_{\mathrm{KK}}\): The masses of the tower of modes of a field on a
  compact hidden direction (Kaluza-Klein masses; not UM-4). Unit: GeV.
  Status: STANDARD.
\item
  \(m_\Psi\): The mass parameter of dirac16complex in the bulk. Unit:
  GeV. Status: RECORD.
\item
  \(\mathsf{k}\): An integer that numbers the modes of a compact
  direction. Unit: dimensionless. Status: STANDARD.
\end{itemize}

The mass mechanism one would expect if \(x_0\) were a compact direction
of proper length \(\pi/(6H)\): a particle confined to it has, by the
uncertainty relation, a momentum \(p_{\mathrm{h}}\) along \(x_0\) of at
least about \(\pi\) divided by the length, \(6H\), and a field periodic
in \(x_0\) has a tower of masses \(m_{\mathrm{KK}}\) with the spacing
\(12H\) (the Kaluza-Klein masses, numbered by the integer
\(\mathsf{k}\)). Status: STANDARD (for a compact direction). Neither
applies to UM-4, because the proper length is infinite (A-1).

\hypertarget{what-the-hidden-direction-produces-instead}{%
\subsection{What the hidden direction produces
instead}\label{what-the-hidden-direction-produces-instead}}

\[
\begin{aligned}
&\phi_{\mathrm{bulk}} = u(\sigma)\,\mathrm{e}^{\mathrm{i}(kx_3 - \omega x_4)}:\qquad
u'' - 7H\,\mathrm{sgn}(\sigma)\,u' + \mathrm{e}^{2H\lvert\sigma\rvert}\,m^2\,u = 0,\qquad m^2 = \omega^2 - k^2,\\
&u_0 = 1,\quad \int_{-\infty}^{\infty}\mathrm{d}\sigma\,\mathrm{e}^{-5H\lvert\sigma\rvert}\,u_0^2 = \frac{2}{5H}
\end{aligned} \tag{FE-73}
\] \textbf{Symbols in (FE-73).}

\begin{itemize}
\tightlist
\item
  \(\phi_{\mathrm{bulk}}\): A massless scalar field of the bulk,
  \(\nabla_M\nabla^M\phi_{\mathrm{bulk}} = 0\) (an illustration: each
  polarisation of the graviton obeys its equation). Unit: GeV\(^3\).
  Status: STANDARD.
\item
  \(u(\sigma),\ u_0\): The profile along the hidden direction of a mode
  of a field of the bulk (a scalar, a gauge field or a polarisation of
  the graviton): the field is \(u(\sigma)\) times a field of the fold;
  \(u_0 = 1\) is the profile of the massless mode (the zero mode); in a
  propagator the response of the field to a unit source on the fold.
  Unit: dimensionless (a relative amplitude); length in a propagator.
  Status: DERIVED.
\item
  \(\sigma\): The proper coordinate of the hidden direction:
  \(\mathrm{d}\sigma = \lvert\cot z\rvert\,\mathrm{d}x_0\),
  \(\sigma = \mathrm{sgn}(\pi/2 - z)\ln(\sin z)/(6H)\) on one period
  \(0 < z < \pi\); \(\sigma = 0\) is the fold \(z = \pi/2\),
  \(\sigma \to \mp\infty\) the tips \(z \to 0, \pi\); one period of
  \(x_0\) is the whole line \(-\infty < \sigma < \infty\). Unit: length.
  Status: DERIVED.
\item
  \(\mathrm{e}\): Euler's number, the base of the natural logarithm;
  \(\mathrm{e}^{X}\) is the exponential of \(X\). Unit: dimensionless.
  Value: 2.718281828. Status: STANDARD.
\item
  \(\mathrm{i}\): The imaginary unit, \(\mathrm{i}^2 = -1\). Unit:
  dimensionless. Status: STANDARD.
\item
  \(k\): A wave number: the momentum of a plane wave along \(x_3\),
  \(\mathrm{e}^{\mathrm{i}(kx_3 - \omega x_4)}\); in the cosmological
  perturbations the comoving wave number of a Fourier mode,
  \(k = \lvert\vec k\rvert\). Unit: GeV (1/length). Status: STANDARD.
\item
  \(x_3\): The third coordinate of ordinary 3-space. Unit: length.
  Status: AUTHOR.
\item
  \(\omega\): A frequency: the energy of a mode
  \(\propto \mathrm{e}^{-\mathrm{i}\omega x_4}\) (\(\omega_o\) of the
  orbital \(o\)). Unit: GeV. Status: STANDARD.
\item
  \(x_4\): The time; a prime on a function of \(x_4\) alone is
  \(\mathrm{d}/\mathrm{d}x_4\). Unit: length (= time, \(c = 1\)).
  Status: AUTHOR.
\item
  \({}'\): The derivative of a function of the time \(x_4\) alone with
  respect to \(x_4\) (\(a_4' = \mathrm{d}a_4/\mathrm{d}x_4\); two primes
  the second derivative); on a function of \(\sigma\) alone (a profile)
  the derivative with respect to \(\sigma\). Status: STANDARD.
\item
  \(H\): The author's constant (an inverse length) that fixes the period
  \(\pi/(6H)\) of \(x_0\), the warp factor
  \(\mathrm{e}^{-2H\lvert\sigma\rvert}\) of the metric, the cosmological
  constant \(\Lambda = -21H^2\) and the tension
  \(\kappa\lambda_{\mathrm{F}} = 12H\) of the fold; in UM-4 a free
  parameter, bounded from below by the propagation of the gauge fields
  in the bulk. Unit: GeV (1/length). Value: \(H \gtrsim 2.5\) TeV (from
  the invisible width of the Z boson). Status: AUTHOR.
\item
  \(\mathrm{sgn}\): The sign function: \(+1\) for a positive argument
  and \(-1\) for a negative one. Status: STANDARD.
\item
  \(m\): A mass: \(m_i\) the mass of the neutrino mass state \(i\);
  \(\Delta m_{ij}^2 = m_i^2 - m_j^2\) the squared-mass difference of the
  states \(i\) and \(j\) (\(\Delta m_{21}^2\), \(\Delta m_{31}^2\),
  \(\Delta m_{32}^2\)); \(\Delta m^2\) without labels the difference of
  the two states of a resonance, \(m_2^2 - m_1^2\), multiplied by
  \(v_{\mathrm{loc}}^2/v^2\) where this differs from 1; \(m\) without a
  label the mass of a mode of a field of the bulk as the fold sees it,
  \(m^2 = \omega^2 - k^2\); for the W and Z bosons their complex pole
  mass, \(m^2 = m_0^2 - \mathrm{i}m_0\Gamma_{\mathrm{bulk}}\) to the
  leading order; in the relic abundance of the particles \(\chi\) their
  mass. Unit: GeV. Status: STANDARD.
\item
  \(\infty\): Infinity: as a limit of an integral or a sum it means that
  the range has no bound. Status: STANDARD.
\item
  \(\mathrm{d}\): The differential: \(\mathrm{d}\) in front of a
  quantity is an infinitesimal change of it (as \(\mathrm{d}\sigma\), a
  small step of \(\sigma\)); the quotient of two such changes is a
  derivative; \(\mathrm{d}\) with an exponent, in front of several
  variables, is their volume element, the product of their
  differentials. Status: STANDARD.
\end{itemize}

\[
\begin{aligned}
&B_\sigma = 0,\quad B_a = u(\sigma)\,\hat B_a(x):\qquad
u'' - 5H\,\mathrm{sgn}(\sigma)\,u' + \mathrm{e}^{2H\lvert\sigma\rvert}\,m^2\,u = 0,\qquad
\eta^{ac}\,\partial_a\hat B_{cb} = m^2\,\hat B_b,\\
&u_0 = 1,\quad \int_{-\infty}^{\infty}\mathrm{d}\sigma\,\mathrm{e}^{-3H\lvert\sigma\rvert}\,u_0^2 = \frac{2}{3H}
\end{aligned} \tag{FE-74}
\] \textbf{Symbols in (FE-74).}

\begin{itemize}
\tightlist
\item
  \(B_M,\ B_\mu,\ B_{MN}\): The hypercharge gauge field of \(U(1)_Y\):
  in the bulk \(B_M\) (geometric normalisation), on the fold the
  canonically normalised four-dimensional field \(B_\mu(x)\) of its
  massless mode, \(g_1B_\mu(x) = B_\mu(\sigma = 0, x)\); \(B_{MN}\),
  \(B_{\mu\nu}\) their field strengths. Unit: GeV (field), GeV\(^2\)
  (field strength). Status: AUTHOR.
\item
  \(\sigma\): The label \(\sigma\) in \(\partial_\sigma\): the
  derivative with respect to the proper hidden coordinate \(\sigma\).
  Status: STANDARD.
\item
  \(a\): A frame (orthonormal) index running over \(0, 1, \dots, 7\)
  (frame direction \((0)\) along \(\sigma\), \((1), (2), (3)\) along
  ordinary space, \((4)\) along the time, \((5), (6), (7)\) along the
  extra times); raised and lowered with the frame metric
  \(\eta_{ab} = \mathrm{diag}(1, 1, 1, 1, -1, -1, -1, -1)\); on the fold
  (\(a, b \in \{1, \dots, 7\}\)) the frame of the fold; in the tensors
  of the fold written with its coordinates \(x^a\) (its metric
  \(\hat g_{ab}\) and the tensors built from it) a coordinate index of
  the fold, \(a \in \{1, \dots, 7\}\); a repeated index is summed.
  Status: STANDARD.
\item
  \(u(\sigma),\ u_0\): The profile along the hidden direction of a mode
  of a field of the bulk (a scalar, a gauge field or a polarisation of
  the graviton): the field is \(u(\sigma)\) times a field of the fold;
  \(u_0 = 1\) is the profile of the massless mode (the zero mode); in a
  propagator the response of the field to a unit source on the fold.
  Unit: dimensionless (a relative amplitude); length in a propagator.
  Status: DERIVED.
\item
  \(\sigma\): The proper coordinate of the hidden direction:
  \(\mathrm{d}\sigma = \lvert\cot z\rvert\,\mathrm{d}x_0\),
  \(\sigma = \mathrm{sgn}(\pi/2 - z)\ln(\sin z)/(6H)\) on one period
  \(0 < z < \pi\); \(\sigma = 0\) is the fold \(z = \pi/2\),
  \(\sigma \to \mp\infty\) the tips \(z \to 0, \pi\); one period of
  \(x_0\) is the whole line \(-\infty < \sigma < \infty\). Unit: length.
  Status: DERIVED.
\item
  \(\hat B_a,\ \hat B_{ab}\): The field of the fold of a mode of the
  hypercharge field of the bulk (of any gauge field of the bulk to the
  first order in the field), \(B_a = u(\sigma)\hat B_a(x)\);
  \(\hat B_{ab} = \partial_a\hat B_b - \partial_b\hat B_a\) its field
  strength. Unit: GeV (field), GeV\(^2\) (field strength). Status:
  STANDARD.
\item
  \(x\): The coordinates of the fold: \(x_1, x_2, x_3\) of ordinary
  space, \(x_4\) the time, \(x_5, x_6, x_7\) the extra times; a field of
  the fold is a function of the four observable coordinates
  \(x = (x_1, x_2, x_3, x_4)\), and
  \(\mathrm{d}^4x = \mathrm{d}x_1\,\mathrm{d}x_2\,\mathrm{d}x_3\,\mathrm{d}x_4\).
  Unit: length (each coordinate). Status: AUTHOR.
\item
  \({}'\): The derivative of a function of the time \(x_4\) alone with
  respect to \(x_4\) (\(a_4' = \mathrm{d}a_4/\mathrm{d}x_4\); two primes
  the second derivative); on a function of \(\sigma\) alone (a profile)
  the derivative with respect to \(\sigma\). Status: STANDARD.
\item
  \(H\): The author's constant (an inverse length) that fixes the period
  \(\pi/(6H)\) of \(x_0\), the warp factor
  \(\mathrm{e}^{-2H\lvert\sigma\rvert}\) of the metric, the cosmological
  constant \(\Lambda = -21H^2\) and the tension
  \(\kappa\lambda_{\mathrm{F}} = 12H\) of the fold; in UM-4 a free
  parameter, bounded from below by the propagation of the gauge fields
  in the bulk. Unit: GeV (1/length). Value: \(H \gtrsim 2.5\) TeV (from
  the invisible width of the Z boson). Status: AUTHOR.
\item
  \(\mathrm{sgn}\): The sign function: \(+1\) for a positive argument
  and \(-1\) for a negative one. Status: STANDARD.
\item
  \(\mathrm{e}\): Euler's number, the base of the natural logarithm;
  \(\mathrm{e}^{X}\) is the exponential of \(X\). Unit: dimensionless.
  Value: 2.718281828. Status: STANDARD.
\item
  \(m\): A mass: \(m_i\) the mass of the neutrino mass state \(i\);
  \(\Delta m_{ij}^2 = m_i^2 - m_j^2\) the squared-mass difference of the
  states \(i\) and \(j\) (\(\Delta m_{21}^2\), \(\Delta m_{31}^2\),
  \(\Delta m_{32}^2\)); \(\Delta m^2\) without labels the difference of
  the two states of a resonance, \(m_2^2 - m_1^2\), multiplied by
  \(v_{\mathrm{loc}}^2/v^2\) where this differs from 1; \(m\) without a
  label the mass of a mode of a field of the bulk as the fold sees it,
  \(m^2 = \omega^2 - k^2\); for the W and Z bosons their complex pole
  mass, \(m^2 = m_0^2 - \mathrm{i}m_0\Gamma_{\mathrm{bulk}}\) to the
  leading order; in the relic abundance of the particles \(\chi\) their
  mass. Unit: GeV. Status: STANDARD.
\item
  \(\eta_{ab}\): The flat metric of the frame,
  \(\eta = \mathrm{diag}(1, 1, 1, 1, -1, -1, -1, -1)\) (eta4488 of
  4+4.nb); \(\eta^{ab}\) is its inverse (the same matrix). Unit:
  dimensionless. Status: AUTHOR.
\item
  \(c\): A frame index (as \(a\)). Status: STANDARD.
\item
  \(\partial\): A partial derivative: the derivative with respect to one
  variable at fixed values of the others (\(\partial_\sigma\) with
  respect to \(\sigma\); \(\partial_M\) and \(\partial_\mu\) with
  respect to the coordinate that carries the same index, of the eight
  dimensions or of the fold). Status: STANDARD.
\item
  \(b\): A frame index (as \(a\)). Status: STANDARD.
\item
  \(\infty\): Infinity: as a limit of an integral or a sum it means that
  the range has no bound. Status: STANDARD.
\item
  \(\mathrm{d}\): The differential: \(\mathrm{d}\) in front of a
  quantity is an infinitesimal change of it (as \(\mathrm{d}\sigma\), a
  small step of \(\sigma\)); the quotient of two such changes is a
  derivative; \(\mathrm{d}\) with an exponent, in front of several
  variables, is their volume element, the product of their
  differentials. Status: STANDARD.
\end{itemize}

The massless fields of the bulk: a scalar (FE-73; each polarisation of
the graviton obeys the same equation) and a gauge field (FE-74, written
for the hypercharge field), each a profile \(u(\sigma)\) along the
hidden direction times a field of the fold of the mass \(m\). Because
the warp factor multiplies the frequency as well as the momentum, they
enter only through \(m^2 = \omega^2 - k^2\): every mode is a field of
the fold with the Lorentz-invariant dispersion \(\omega^2 = k^2 + m^2\).
The massless mode has the profile 1 and a finite kinetic weight, so it
is bound to the fold and travels at the speed of light at every
wavelength. Status: DERIVED (\texttt{um4\_bulk.py}, B1 and B2).

\[
-\mathrm{e}^{nH\lvert\sigma\rvert}\,\partial_\sigma\big(\mathrm{e}^{-nH\lvert\sigma\rvert}\,\partial_\sigma u\big)
 + \mathrm{e}^{2H\lvert\sigma\rvert}\,\mathsf{p}^2\,u = \delta_{\mathrm{D}}(\sigma),\qquad
\mathbb{G}_n(\mathsf{p}) = u(0) = \frac{\mathrm{K}_{n/2}(\mathsf{p}/H)}{2\mathsf{p}\,\mathrm{K}_{n/2-1}(\mathsf{p}/H)} \tag{FE-79}
\] \textbf{Symbols in (FE-79).}

\begin{itemize}
\tightlist
\item
  \(\mathrm{e}\): Euler's number, the base of the natural logarithm;
  \(\mathrm{e}^{X}\) is the exponential of \(X\). Unit: dimensionless.
  Value: 2.718281828. Status: STANDARD.
\item
  \(n\): The label of a field of the bulk in its profile equation: the
  exponent of the weight \(\mathrm{e}^{-nH\lvert\sigma\rvert}\) of its
  \(\sigma\)-derivatives, \(n = 7\) for a scalar and for each
  polarisation of the graviton, \(n = 5\) for a gauge field. Status:
  STANDARD.
\item
  \(H\): The author's constant (an inverse length) that fixes the period
  \(\pi/(6H)\) of \(x_0\), the warp factor
  \(\mathrm{e}^{-2H\lvert\sigma\rvert}\) of the metric, the cosmological
  constant \(\Lambda = -21H^2\) and the tension
  \(\kappa\lambda_{\mathrm{F}} = 12H\) of the fold; in UM-4 a free
  parameter, bounded from below by the propagation of the gauge fields
  in the bulk. Unit: GeV (1/length). Value: \(H \gtrsim 2.5\) TeV (from
  the invisible width of the Z boson). Status: AUTHOR.
\item
  \(\sigma\): The proper coordinate of the hidden direction:
  \(\mathrm{d}\sigma = \lvert\cot z\rvert\,\mathrm{d}x_0\),
  \(\sigma = \mathrm{sgn}(\pi/2 - z)\ln(\sin z)/(6H)\) on one period
  \(0 < z < \pi\); \(\sigma = 0\) is the fold \(z = \pi/2\),
  \(\sigma \to \mp\infty\) the tips \(z \to 0, \pi\); one period of
  \(x_0\) is the whole line \(-\infty < \sigma < \infty\). Unit: length.
  Status: DERIVED.
\item
  \(\partial\): A partial derivative: the derivative with respect to one
  variable at fixed values of the others (\(\partial_\sigma\) with
  respect to \(\sigma\); \(\partial_M\) and \(\partial_\mu\) with
  respect to the coordinate that carries the same index, of the eight
  dimensions or of the fold). Status: STANDARD.
\item
  \(\sigma\): The label \(\sigma\) in \(\partial_\sigma\): the
  derivative with respect to the proper hidden coordinate \(\sigma\).
  Status: STANDARD.
\item
  \(u(\sigma),\ u_0\): The profile along the hidden direction of a mode
  of a field of the bulk (a scalar, a gauge field or a polarisation of
  the graviton): the field is \(u(\sigma)\) times a field of the fold;
  \(u_0 = 1\) is the profile of the massless mode (the zero mode); in a
  propagator the response of the field to a unit source on the fold.
  Unit: dimensionless (a relative amplitude); length in a propagator.
  Status: DERIVED.
\item
  \(\mathsf{p}\): The Euclidean momentum along the fold of a source on
  the fold: \(\mathsf{p}^2\) the square of its seven-dimensional
  momentum (for a static source spread uniformly over the extra times
  the square of its momentum along ordinary space);
  \(\mathsf{p} = -\mathrm{i}m\) for a time-like momentum of the mass
  \(m\) with outgoing waves into the bulk. Unit: GeV. Status: STANDARD.
\item
  \(\delta_{\mathrm{D}}\): The Dirac delta function of \(\sigma\): zero
  for \(\sigma \ne 0\), with the integral 1 over any range that contains
  \(\sigma = 0\); the integral of \(\delta_{\mathrm{D}}(\sigma)\) times
  a smooth function of \(\sigma\) is the value of that function at
  \(\sigma = 0\). Unit: 1/length (for \(\delta_{\mathrm{D}}(\sigma)\)).
  Status: STANDARD.
\item
  \(\mathbb{G}_n,\ \mathbb{G}_7,\ \mathbb{G}_5\): The propagator of a
  massless field of the bulk from one point of the fold to another, in
  the momentum space of the fold: the value at the fold of its response
  to a unit source on the fold; \(\mathbb{G}_7\) for each polarisation
  of the graviton and for a scalar, \(\mathbb{G}_5\) for a gauge field.
  Unit: GeV\(^{-1}\). Status: DERIVED.
\item
  \(\mathrm{K}\): The modified Bessel function of the second kind (of
  the order written as its subscript): the solution of the modified
  Bessel equation that decays exponentially for a large argument; for a
  half-integer order it is an exponential times a polynomial in the
  inverse of the argument, which makes the propagators of the bulk
  rational functions. Unit: dimensionless. Status: STANDARD.
\end{itemize}

The propagator of a massless field of the bulk from one point of the
fold to another: the response at the fold to a unit source on the fold,
for the Euclidean momentum \(\mathsf{p}\) along the fold. Status:
DERIVED (\texttt{um4\_propagators.py}, P1).

\[
\begin{aligned}
&u \propto \zeta_{\mathrm{w}}^{n/2}\Big(\mathrm{Y}_{n/2-1}\big(m/H\big)\,\mathrm{J}_{n/2}\big(m\zeta_{\mathrm{w}}\big) - \mathrm{J}_{n/2-1}\big(m/H\big)\,\mathrm{Y}_{n/2}\big(m\zeta_{\mathrm{w}}\big)\Big),\qquad
 \zeta_{\mathrm{w}} = \frac{\mathrm{e}^{H\lvert\sigma\rvert}}{H},\\
&\mathbb{G}_n(\mathsf{p}) = \frac{(n - 2)H}{2\mathsf{p}^2} + \int_0^\infty\mathrm{d}m\,\frac{\mathfrak{s}_n(m)}{\mathsf{p}^2 + m^2},\qquad
 \mathfrak{s}_n(m) = \frac{2H}{\pi^2m\,\big[\mathrm{J}_{n/2-1}(m/H)^2 + \mathrm{Y}_{n/2-1}(m/H)^2\big]},\\
&\mathfrak{s}_7(m) \simeq \frac{(m/H)^4}{9\pi},\quad \mathfrak{s}_5(m) \simeq \frac{(m/H)^2}{\pi}\quad(m \ll H),\qquad \mathfrak{s}_n(m) \to \frac{1}{\pi}\quad(m \gg H)
\end{aligned} \tag{A-3}
\] \textbf{Symbols in (A-3).}

\begin{itemize}
\tightlist
\item
  \(u(\sigma),\ u_0\): The profile along the hidden direction of a mode
  of a field of the bulk (a scalar, a gauge field or a polarisation of
  the graviton): the field is \(u(\sigma)\) times a field of the fold;
  \(u_0 = 1\) is the profile of the massless mode (the zero mode); in a
  propagator the response of the field to a unit source on the fold.
  Unit: dimensionless (a relative amplitude); length in a propagator.
  Status: DERIVED.
\item
  \(\zeta_{\mathrm{w}}\): The warped distance from the fold,
  \(\zeta_{\mathrm{w}} = \mathrm{e}^{H\lvert\sigma\rvert}/H\) (\(1/H\)
  at the fold): \(\omega\zeta_{\mathrm{w}}\) is the argument of the
  exact solutions of dirac16complex, and \(m\zeta_{\mathrm{w}}\) that of
  the massive modes of the mass \(m\) of the massless fields of the
  bulk. Unit: length. Status: DERIVED.
\item
  \(n\): The label of a field of the bulk in its profile equation: the
  exponent of the weight \(\mathrm{e}^{-nH\lvert\sigma\rvert}\) of its
  \(\sigma\)-derivatives, \(n = 7\) for a scalar and for each
  polarisation of the graviton, \(n = 5\) for a gauge field. Status:
  STANDARD.
\item
  \(\mathrm{Y}\): The Bessel function of the second kind (of the order
  written as its subscript): the second solution of Bessel's equation,
  singular at zero argument; for a half-integer order an elementary
  function. Unit: dimensionless. Status: STANDARD.
\item
  \(m\): A mass: \(m_i\) the mass of the neutrino mass state \(i\);
  \(\Delta m_{ij}^2 = m_i^2 - m_j^2\) the squared-mass difference of the
  states \(i\) and \(j\) (\(\Delta m_{21}^2\), \(\Delta m_{31}^2\),
  \(\Delta m_{32}^2\)); \(\Delta m^2\) without labels the difference of
  the two states of a resonance, \(m_2^2 - m_1^2\), multiplied by
  \(v_{\mathrm{loc}}^2/v^2\) where this differs from 1; \(m\) without a
  label the mass of a mode of a field of the bulk as the fold sees it,
  \(m^2 = \omega^2 - k^2\); for the W and Z bosons their complex pole
  mass, \(m^2 = m_0^2 - \mathrm{i}m_0\Gamma_{\mathrm{bulk}}\) to the
  leading order; in the relic abundance of the particles \(\chi\) their
  mass. Unit: GeV. Status: STANDARD.
\item
  \(H\): The author's constant (an inverse length) that fixes the period
  \(\pi/(6H)\) of \(x_0\), the warp factor
  \(\mathrm{e}^{-2H\lvert\sigma\rvert}\) of the metric, the cosmological
  constant \(\Lambda = -21H^2\) and the tension
  \(\kappa\lambda_{\mathrm{F}} = 12H\) of the fold; in UM-4 a free
  parameter, bounded from below by the propagation of the gauge fields
  in the bulk. Unit: GeV (1/length). Value: \(H \gtrsim 2.5\) TeV (from
  the invisible width of the Z boson). Status: AUTHOR.
\item
  \(\mathrm{J}\): The Bessel function of the first kind (of the order
  written as its subscript): the solution of Bessel's equation that is
  finite at zero argument; for a half-integer order it is an elementary
  function (sines, cosines and powers of its argument). Unit:
  dimensionless. Status: STANDARD.
\item
  \(\mathrm{e}\): Euler's number, the base of the natural logarithm;
  \(\mathrm{e}^{X}\) is the exponential of \(X\). Unit: dimensionless.
  Value: 2.718281828. Status: STANDARD.
\item
  \(\sigma\): The proper coordinate of the hidden direction:
  \(\mathrm{d}\sigma = \lvert\cot z\rvert\,\mathrm{d}x_0\),
  \(\sigma = \mathrm{sgn}(\pi/2 - z)\ln(\sin z)/(6H)\) on one period
  \(0 < z < \pi\); \(\sigma = 0\) is the fold \(z = \pi/2\),
  \(\sigma \to \mp\infty\) the tips \(z \to 0, \pi\); one period of
  \(x_0\) is the whole line \(-\infty < \sigma < \infty\). Unit: length.
  Status: DERIVED.
\item
  \(\mathbb{G}_n,\ \mathbb{G}_7,\ \mathbb{G}_5\): The propagator of a
  massless field of the bulk from one point of the fold to another, in
  the momentum space of the fold: the value at the fold of its response
  to a unit source on the fold; \(\mathbb{G}_7\) for each polarisation
  of the graviton and for a scalar, \(\mathbb{G}_5\) for a gauge field.
  Unit: GeV\(^{-1}\). Status: DERIVED.
\item
  \(\mathsf{p}\): The Euclidean momentum along the fold of a source on
  the fold: \(\mathsf{p}^2\) the square of its seven-dimensional
  momentum (for a static source spread uniformly over the extra times
  the square of its momentum along ordinary space);
  \(\mathsf{p} = -\mathrm{i}m\) for a time-like momentum of the mass
  \(m\) with outgoing waves into the bulk. Unit: GeV. Status: STANDARD.
\item
  \(\mathrm{d}\): The differential: \(\mathrm{d}\) in front of a
  quantity is an infinitesimal change of it (as \(\mathrm{d}\sigma\), a
  small step of \(\sigma\)); the quotient of two such changes is a
  derivative; \(\mathrm{d}\) with an exponent, in front of several
  variables, is their volume element, the product of their
  differentials. Status: STANDARD.
\item
  \(\infty\): Infinity: as a limit of an integral or a sum it means that
  the range has no bound. Status: STANDARD.
\item
  \(\mathfrak{s}_n(m),\ \mathfrak{s}_7,\ \mathfrak{s}_5\): The spectral
  density at the fold of the massive modes of a massless field of the
  bulk (\(n = 7\) for a scalar and each polarisation of the graviton,
  \(n = 5\) for a gauge field): the square of the value at the fold of
  the normalised even profile of the mode of the mass \(m\), per unit of
  \(m\); the weight with which the mode of the mass \(m\) propagates
  from one point of the fold to another. Unit: dimensionless. Value:
  \(\mathfrak{s}_7 \simeq (m/H)^4/(9\pi)\),
  \(\mathfrak{s}_5 \simeq (m/H)^2/\pi\) for \(m \ll H\); \(1/\pi\) for
  \(m \gg H\). Status: DERIVED.
\item
  \(\pi\): The number pi, the ratio of the circumference of a circle to
  its diameter. Unit: dimensionless. Value: 3.141592654. Status:
  STANDARD.
\item
  \(v_{\mathrm{loc}}\): The local value of the electroweak vacuum,
  \(v_{\mathrm{loc}}^2 = 2\langle\Phi^\dagger\Phi\rangle\): it sets the
  local masses of the quarks and leptons (each mass times
  \(v_{\mathrm{loc}}/v\)) and the local Fermi constant;
  \(v_{\mathrm{loc}} = v\) in the Standard-Model vacuum. Unit: GeV.
  Value: \(v\) in the late universe outside dense dirac16complex.
  Status: DERIVED.
\end{itemize}

The massive modes, the masses that the hidden direction produces. With
the warped distance \(\zeta_{\mathrm{w}}\) from the fold the profile
equations FE-73 (\(n = 7\)) and FE-74 (\(n = 5\)) become Bessel's
equation: their solutions are \(\zeta_{\mathrm{w}}^{n/2}\) times the
Bessel functions of the order \(n/2\) of \(m\zeta_{\mathrm{w}}\). The
combination of the first line is the even mode, which has no kink at the
fold (\(u' = 0\) on both sides of it, as for a field without a source at
the fold); the odd modes vanish at the fold and do not reach it, so the
even modes are the ones that the fold sees. Far from the fold every one
of them oscillates, for every \(m > 0\), and the kinetic weight
\(\mathrm{e}^{-(n-2)H\lvert\sigma\rvert}\mathrm{d}\sigma\) is the
measure of the Hankel transform in \(\zeta_{\mathrm{w}}\): the massive
modes form a continuum that starts at zero mass, with no gap and no
tower, beside the one massless mode. The second line is the spectral
representation of the propagator FE-79: the massless mode with the
weight \((n - 2)H/2\) plus the continuum, each mass with the weight
\(\mathfrak{s}_n(m)\), the square of the value at the fold of its
normalised profile. For masses far below \(H\) the weight vanishes like
\((m/H)^4\) for the graviton and the scalars and like \((m/H)^2\) for
the gauge fields: the light modes of the continuum live far from the
fold, and only modes heavier than \(H\) reach the fold with full
strength. (A record of UM-3, whose time was not warped: there the
massless fields of the bulk had a gap above their massless mode, \(3H\)
for a scalar and \(2H\) for a gauge field; at zero momentum the profile
equation of the scalar,
\(u'' - 6H\,\mathrm{sgn}(\sigma)\,u' + \omega^2u = 0\), has no
oscillating solutions below \(\omega = 3H\).) Status: DERIVED
(\texttt{um4\_bulk.py}, B10: the exact solutions with sympy on both
halves, the measure, and the spectral representation, which reproduces
FE-79 to \(6\times10^{-16}\) at \(\mathsf{p}/H = 0.3\), 1 and 3).

\[
\mathcal{V}_7(r) = \frac{5H}{8\pi r}\Big[1 + \frac{16}{15\pi}\,\frac{1}{(Hr)^5} + \dots\Big],\qquad
\mathcal{V}_5(r) = \frac{3H}{8\pi r}\Big[1 + \frac{4}{3\pi}\,\frac{1}{(Hr)^3} + \dots\Big] \tag{FE-81}
\] \textbf{Symbols in (FE-81).}

\begin{itemize}
\tightlist
\item
  \(\mathcal{V}_7(r),\ \mathcal{V}_5(r)\): The static potentials between
  two unit point sources on the fold at the distance \(r\) (spread
  uniformly over the extra times) from the exchange of the graviton
  (\(\mathcal{V}_7\)) and of a gauge field (\(\mathcal{V}_5\)) of the
  bulk: the three-dimensional Fourier transforms of \(\mathbb{G}_7\) and
  \(\mathbb{G}_5\). Unit: GeV\(^2\). Status: DERIVED.
\item
  \(r\): A distance from a centre: the distance from the centre of the
  star along the path of a neutrino; in Newton's and Coulomb's laws the
  distance between two point sources on the fold. Unit: length. Status:
  STANDARD.
\item
  \(H\): The author's constant (an inverse length) that fixes the period
  \(\pi/(6H)\) of \(x_0\), the warp factor
  \(\mathrm{e}^{-2H\lvert\sigma\rvert}\) of the metric, the cosmological
  constant \(\Lambda = -21H^2\) and the tension
  \(\kappa\lambda_{\mathrm{F}} = 12H\) of the fold; in UM-4 a free
  parameter, bounded from below by the propagation of the gauge fields
  in the bulk. Unit: GeV (1/length). Value: \(H \gtrsim 2.5\) TeV (from
  the invisible width of the Z boson). Status: AUTHOR.
\item
  \(\pi\): The number pi, the ratio of the circumference of a circle to
  its diameter. Unit: dimensionless. Value: 3.141592654. Status:
  STANDARD.
\end{itemize}

\[
\Big(1 + M_{\mathrm{loc}}^2\,\mathbb{G}_5(\mathsf{p})\Big)\Big\vert_{\mathsf{p} = -\mathrm{i}m} = 0:\qquad
m^2 = m_0^2\Big(1 - \frac{m_0^2}{3H^2}\Big) - \mathrm{i}\,\frac{m_0^5}{3H^3} + \dots,\qquad
m_0^2 = \frac32\,H M_{\mathrm{loc}}^2,\qquad \Gamma_{\mathrm{bulk}} = \frac{m_0^4}{3H^3} \tag{FE-82}
\] \textbf{Symbols in (FE-82).}

\begin{itemize}
\tightlist
\item
  \(M_{\mathrm{loc}}\): The mass term that the Higgs vacuum on the fold
  gives a gauge field of the bulk:
  \(M_{\mathrm{loc}}^2\delta_{\mathrm{D}}(\sigma)\) in its profile
  equation; for the W,
  \(M_{\mathrm{loc}}^2 = \mathsf{g}_2^2(v^2 + \mathcal{A}^2)/(4V_*)\).
  Unit: GeV\(^{1/2}\) (\(M_{\mathrm{loc}}^2\) in GeV). Status: DERIVED.
\item
  \(\mathbb{G}_n,\ \mathbb{G}_7,\ \mathbb{G}_5\): The propagator of a
  massless field of the bulk from one point of the fold to another, in
  the momentum space of the fold: the value at the fold of its response
  to a unit source on the fold; \(\mathbb{G}_7\) for each polarisation
  of the graviton and for a scalar, \(\mathbb{G}_5\) for a gauge field.
  Unit: GeV\(^{-1}\). Status: DERIVED.
\item
  \(\mathsf{p}\): The Euclidean momentum along the fold of a source on
  the fold: \(\mathsf{p}^2\) the square of its seven-dimensional
  momentum (for a static source spread uniformly over the extra times
  the square of its momentum along ordinary space);
  \(\mathsf{p} = -\mathrm{i}m\) for a time-like momentum of the mass
  \(m\) with outgoing waves into the bulk. Unit: GeV. Status: STANDARD.
\item
  \(\mathrm{i}\): The imaginary unit, \(\mathrm{i}^2 = -1\). Unit:
  dimensionless. Status: STANDARD.
\item
  \(m\): A mass: \(m_i\) the mass of the neutrino mass state \(i\);
  \(\Delta m_{ij}^2 = m_i^2 - m_j^2\) the squared-mass difference of the
  states \(i\) and \(j\) (\(\Delta m_{21}^2\), \(\Delta m_{31}^2\),
  \(\Delta m_{32}^2\)); \(\Delta m^2\) without labels the difference of
  the two states of a resonance, \(m_2^2 - m_1^2\), multiplied by
  \(v_{\mathrm{loc}}^2/v^2\) where this differs from 1; \(m\) without a
  label the mass of a mode of a field of the bulk as the fold sees it,
  \(m^2 = \omega^2 - k^2\); for the W and Z bosons their complex pole
  mass, \(m^2 = m_0^2 - \mathrm{i}m_0\Gamma_{\mathrm{bulk}}\) to the
  leading order; in the relic abundance of the particles \(\chi\) their
  mass. Unit: GeV. Status: STANDARD.
\item
  \(m_0\): The mass of the W or the Z boson without the massive modes of
  the bulk (the mass of the massless mode with the Higgs vacuum),
  \(m_0^2 = \frac32HM_{\mathrm{loc}}^2\);
  \(m_0^2 = g_2^2(v^2 + \mathcal{A}^2)/4\) for the W. Unit: GeV. Status:
  DERIVED.
\item
  \(H\): The author's constant (an inverse length) that fixes the period
  \(\pi/(6H)\) of \(x_0\), the warp factor
  \(\mathrm{e}^{-2H\lvert\sigma\rvert}\) of the metric, the cosmological
  constant \(\Lambda = -21H^2\) and the tension
  \(\kappa\lambda_{\mathrm{F}} = 12H\) of the fold; in UM-4 a free
  parameter, bounded from below by the propagation of the gauge fields
  in the bulk. Unit: GeV (1/length). Value: \(H \gtrsim 2.5\) TeV (from
  the invisible width of the Z boson). Status: AUTHOR.
\item
  \(\Gamma_{\mathrm{bulk}}\): The width of the W or the Z boson into the
  bulk (its decay into the massive modes of the gauge fields of the
  bulk), \(\Gamma_{\mathrm{bulk}} = m_0(m_0/H)^3/3\). Unit: GeV. Value:
  1.4 MeV for the Z and 0.9 MeV for the W at \(H = 2.5\) TeV. Status:
  DERIVED.
\end{itemize}

What the observer on the fold sees of the continuum. Between two point
sources on the fold at the distance \(r\) the massless mode gives
Newton's law (the first potential, multiplied by the gravitational
coupling) and Coulomb's law (the second, multiplied by the square of the
gauge coupling), and the continuum adds the relative corrections
\(16/(15\pi(Hr)^5) = 0.3395/(Hr)^5\) and
\(4/(3\pi(Hr)^3) = 0.4244/(Hr)^3\): they are the integrals of the
weights \(\mathfrak{s}_7\) and \(\mathfrak{s}_5\) of A-3 against the
Yukawa factor \(\mathrm{e}^{-mr}\) of each mass, divided by the weight
of the massless mode, and the computed integrals approach them (0.33945
and 0.42429 at \(Hr = 200\)). For \(H \ge 2.5\) TeV they are
\(3\times10^{-74}\) at 52 micrometres, the shortest distance at which
torsion balances test Newton's law, and \(1.4\times10^{-27}\) at the
Bohr radius. The mass that the Higgs vacuum gives the gauge fields at
the fold turns the W and Z bosons into resonances of the bulk (FE-82):
they can decay into the continuum, with the width
\(\Gamma_{\mathrm{bulk}} = m_0^4/(3H^3)\), and the measured invisible
width of the Z boson (the number of light neutrinos
\(2.9963 \pm 0.0074\), which allows at most 1.4 MeV of extra invisible
width) requires \(H > 2.5\) TeV (2534 GeV), the bound used in every UM-4
document. Status: DERIVED (\texttt{um4\_propagators.py}, P3 and P4;
\texttt{um4\_bulk.py}, B10), ESTIMATE (the bounds, P5).

\[
\begin{aligned}
&\chi' = \mathbb{K}\,\chi,\qquad \mathbb{K} = \gamma^{(0)}\big(m_\Psi + \mathrm{i}\omega\,\mathrm{e}^{H\lvert\sigma\rvert}\gamma^{(4)}\big),\qquad
\mathbb{K}^2 = \big(m_\Psi^2 - \omega^2\mathrm{e}^{2H\lvert\sigma\rvert}\big)\mathbf{1},\\
&q = \sqrt{m_\Psi^2 - \omega^2\mathrm{e}^{2H\lvert\sigma\rvert}},\qquad \lvert\omega\rvert\,\mathrm{e}^{H\sigma_{\mathrm{t}}} = m_\Psi
\end{aligned} \tag{FE-37}
\] \textbf{Symbols in (FE-37).}

\begin{itemize}
\tightlist
\item
  \(\chi(\sigma)\): The profile of a mode of dirac16complex along the
  hidden direction,
  \(\Psi = \mathrm{e}^{\frac72H\lvert\sigma\rvert}\chi(\sigma)\mathrm{e}^{\mathrm{i}(\vec k\cdot\vec x - \omega x_4)}\)
  (a 16-component spinor function of \(\sigma\)); as a four-dimensional
  field, \(\chi\) denotes the particle of the level of dirac16complex at
  the fold (a long-lived resonance) and \(\bar\chi\) its antiparticle.
  Unit: GeV\(^{2}\) (with the \(V_*\) normalisation). Status: DERIVED.
\item
  \({}'\): The derivative of a function of the time \(x_4\) alone with
  respect to \(x_4\) (\(a_4' = \mathrm{d}a_4/\mathrm{d}x_4\); two primes
  the second derivative); on a function of \(\sigma\) alone (a profile)
  the derivative with respect to \(\sigma\). Status: STANDARD.
\item
  \(\mathbb{K}\): The transfer matrix of the profile equation at zero
  momentum, \(\chi' = \mathbb{K}\chi\). Unit: GeV. Status: DERIVED.
\item
  \(\gamma^M,\ \gamma^{(a)}\): The Dirac matrices of dirac16complex
  along the coordinates, \(\gamma^M = \mathbf{e}_a{}^MT^a\),
  \(\gamma_M = g_{MN}\gamma^N\); \(\gamma^{(a)} = T^a\) along the frame
  direction \((a)\) (\(\gamma^{(0)}\) along \(\sigma\), \(\gamma^{(4)}\)
  along the time). Unit: dimensionless. Status: AUTHOR.
\item
  \(m_\Psi\): The mass parameter of dirac16complex in the bulk. Unit:
  GeV. Status: RECORD.
\item
  \(\mathrm{i}\): The imaginary unit, \(\mathrm{i}^2 = -1\). Unit:
  dimensionless. Status: STANDARD.
\item
  \(\omega\): A frequency: the energy of a mode
  \(\propto \mathrm{e}^{-\mathrm{i}\omega x_4}\) (\(\omega_o\) of the
  orbital \(o\)). Unit: GeV. Status: STANDARD.
\item
  \(\mathrm{e}\): Euler's number, the base of the natural logarithm;
  \(\mathrm{e}^{X}\) is the exponential of \(X\). Unit: dimensionless.
  Value: 2.718281828. Status: STANDARD.
\item
  \(H\): The author's constant (an inverse length) that fixes the period
  \(\pi/(6H)\) of \(x_0\), the warp factor
  \(\mathrm{e}^{-2H\lvert\sigma\rvert}\) of the metric, the cosmological
  constant \(\Lambda = -21H^2\) and the tension
  \(\kappa\lambda_{\mathrm{F}} = 12H\) of the fold; in UM-4 a free
  parameter, bounded from below by the propagation of the gauge fields
  in the bulk. Unit: GeV (1/length). Value: \(H \gtrsim 2.5\) TeV (from
  the invisible width of the Z boson). Status: AUTHOR.
\item
  \(\sigma\): The proper coordinate of the hidden direction:
  \(\mathrm{d}\sigma = \lvert\cot z\rvert\,\mathrm{d}x_0\),
  \(\sigma = \mathrm{sgn}(\pi/2 - z)\ln(\sin z)/(6H)\) on one period
  \(0 < z < \pi\); \(\sigma = 0\) is the fold \(z = \pi/2\),
  \(\sigma \to \mp\infty\) the tips \(z \to 0, \pi\); one period of
  \(x_0\) is the whole line \(-\infty < \sigma < \infty\). Unit: length.
  Status: DERIVED.
\item
  \(\mathbf{1},\ \mathbf{1}_8\): The identity matrix of the size fixed
  by the context (\(16 \times 16\) on dirac16complex, \(4 \times 4\) on
  a four-dimensional spinor, \(3 \times 3\) in generation or flavour
  space); \(\mathbf{1}_8\) is the \(8 \times 8\) identity matrix. Unit:
  dimensionless. Status: STANDARD.
\item
  \(q\): The local decay rate of a profile of dirac16complex along
  \(\sigma\) at the frequency \(\omega\),
  \(q = \sqrt{m_\Psi^2 - \omega^2\mathrm{e}^{2H\lvert\sigma\rvert}}\)
  (imaginary beyond the turning point, where the profile oscillates).
  Unit: GeV. Status: DERIVED.
\item
  \(\sigma_{\mathrm{t}}\): The turning point of a state of
  dirac16complex of the frequency \(\omega\): the depth at which the
  warped frequency equals the bulk mass,
  \(\lvert\omega\rvert\mathrm{e}^{H\sigma_{\mathrm{t}}} = m_\Psi\);
  beyond it every solution oscillates. Unit: length. Status: DERIVED.
\end{itemize}

\[
\begin{aligned}
&\chi(0^+) = \exp\Big(\gamma^{(0)}\big(\mu_{\mathrm{L}}P_{\mathrm{L}} + \mu_{\mathrm{R}}P_{\mathrm{R}}\big)\Big)\,\chi(0^-),\\
&\mu_X = \frac{1}{V_*}\Big[c_\Phi\,\langle\Phi^\dagger\Phi\rangle\,\delta_{X\mathrm{L}} + \frac14\big(c^{G}_X\langle G^2\rangle
 + c^{W}_X\langle W^2\rangle + c^{B}_X\langle B^2\rangle\big) + 2\omega\sum_f c^{f}_X\,n_f\Big],\qquad
 \langle\Phi^\dagger\Phi\rangle = \frac{v^2 + \mathcal{A}^2}{2}
\end{aligned} \tag{FE-39}
\] \textbf{Symbols in (FE-39).}

\begin{itemize}
\tightlist
\item
  \(\chi(\sigma)\): The profile of a mode of dirac16complex along the
  hidden direction,
  \(\Psi = \mathrm{e}^{\frac72H\lvert\sigma\rvert}\chi(\sigma)\mathrm{e}^{\mathrm{i}(\vec k\cdot\vec x - \omega x_4)}\)
  (a 16-component spinor function of \(\sigma\)); as a four-dimensional
  field, \(\chi\) denotes the particle of the level of dirac16complex at
  the fold (a long-lived resonance) and \(\bar\chi\) its antiparticle.
  Unit: GeV\(^{2}\) (with the \(V_*\) normalisation). Status: DERIVED.
\item
  \(\gamma^M,\ \gamma^{(a)}\): The Dirac matrices of dirac16complex
  along the coordinates, \(\gamma^M = \mathbf{e}_a{}^MT^a\),
  \(\gamma_M = g_{MN}\gamma^N\); \(\gamma^{(a)} = T^a\) along the frame
  direction \((a)\) (\(\gamma^{(0)}\) along \(\sigma\), \(\gamma^{(4)}\)
  along the time). Unit: dimensionless. Status: AUTHOR.
\item
  \(\mu_X,\ \mu_{\mathrm{L}},\ \mu_{\mathrm{R}}\): The dimensionless
  strengths of the fold for the left-handed and right-handed halves of
  dirac16complex, set by the mean fields of the Standard Model at the
  fold: the jump of the profile across the fold is the exponential of
  \(\gamma^{(0)}(\mu_{\mathrm{L}}P_{\mathrm{L}} + \mu_{\mathrm{R}}P_{\mathrm{R}})\);
  negative values attract. Unit: dimensionless. Status: DERIVED.
\item
  \(P_{\mathrm{L}}\): The author's left-handed chiral projector,
  \(P_{\mathrm{L}} = (\mathbf{1} - \Gamma_8)/2\). Unit: dimensionless.
  Status: AUTHOR.
\item
  \(P_{\mathrm{R}}\): The author's right-handed chiral projector,
  \(P_{\mathrm{R}} = (\mathbf{1} + \Gamma_8)/2\). Unit: dimensionless.
  Status: AUTHOR.
\item
  \(X\): The label of a half of dirac16complex:
  \(X \in \{\mathrm{L}, \mathrm{R}\}\) (left-handed or right-handed
  half); a sum over \(X\) runs over both. Status: STANDARD.
\item
  \(V_*\): The physical volume of the extra times \(x_5, x_6, x_7\) on
  the fold in the late universe,
  \(V_* = V_{\mathrm{t}}(\mathrm{e}^{-a_4}b_4)^3\) (with
  \(V_{\mathrm{t}} = \int\mathrm{d}x_5\,\mathrm{d}x_6\,\mathrm{d}x_7\)
  the coordinate volume of the extra times) at a time when
  \(\mathrm{e}^{-a_4}b_4\) is constant; it converts the
  eight-dimensional normalisation into the four-dimensional one. Unit:
  length\(^3\) (GeV\(^{-3}\)). Status: ASSUMED.
\item
  \(c_\Phi\): The Higgs-portal coupling of the left-handed half of
  dirac16complex (the right-handed half has none). Unit: GeV\(^{-5}\).
  Status: INVENTED.
\item
  \(\Phi\): The Higgs doublet (two complex components, hypercharge
  \(1/2\)), a field of the fold. Unit: GeV. Status: AUTHOR.
\item
  \({}^\dagger\): Hermitian conjugation: the complex conjugate of the
  transpose of a matrix or spinor; for an operator on the space of
  states, its adjoint with respect to the scalar product of that space
  (for the quantised dirac16complex the indefinite Krein scalar
  product). Status: STANDARD.
\item
  \(\langle\cdots\rangle\): The mean (expectation) value of an operator
  in the state of the medium (the vacuum, a plasma, the Kohn-Sham
  state). Status: STANDARD.
\item
  \(\delta\): The Kronecker delta: with two indices \(\delta^M_N = 1\)
  if \(M = N\) and 0 otherwise (likewise \(\delta_{X\mathrm{L}}\) is 1
  if the label \(X\) of a half of dirac16complex is L and 0 otherwise);
  with two index strings of equal length the generalised Kronecker delta
  \(\delta^{M_1M_2\cdots}_{N_1N_2\cdots}\), the determinant of the
  square matrix of the Kronecker deltas \(\delta^{M_a}_{N_b}\) (\(+1\)
  or \(-1\) if the upper indices are an even or odd permutation of the
  lower ones, 0 otherwise). Status: STANDARD.
\item
  \(c^{G}_X\): The gluon-kinetic portal coupling of the half \(X\) of
  dirac16complex. Unit: GeV\(^{-7}\). Status: INVENTED.
\item
  \(\langle G^2\rangle\): The mean square of the gluon field strength,
  \(\langle G^A_{\mu\nu}G^{A\mu\nu}\rangle\) (the gluon condensate in
  the QCD vacuum, a thermal value in the plasma). Unit: GeV\(^4\).
  Status: STANDARD.
\item
  \(G^A_M,\ G^A_\mu,\ G^A_{MN}\): The gluon field of \(SU(3)_C\)
  (\(A = 1, \dots, 8\)): in the bulk \(G^A_M\), a field of eight
  dimensions in its geometric normalisation (the covariant derivative
  \(\partial - \mathrm{i}G\)), and on the fold the canonically
  normalised four-dimensional field \(G^A_\mu(x)\) of its massless mode,
  \(g_3G^A_\mu(x) = G^A_\mu(\sigma = 0, x)\); \(G^A_{MN}\) and
  \(G^A_{\mu\nu}\) their field strengths. Unit: GeV (field), GeV\(^2\)
  (field strength). Status: AUTHOR.
\item
  \(c^{W}_X\): The \(W\)-kinetic portal coupling of the half \(X\).
  Unit: GeV\(^{-7}\). Status: INVENTED.
\item
  \(\langle W^2\rangle\): The mean square of the \(SU(2)\) field
  strength, \(\langle W^I_{\mu\nu}W^{I\mu\nu}\rangle\). Unit: GeV\(^4\).
  Status: STANDARD.
\item
  \(W^I_M,\ W^I_\mu,\ W^I_{MN}\): The gauge field of \(SU(2)_L\)
  (\(I = 1, 2, 3\)): in the bulk \(W^I_M\) (geometric normalisation), on
  the fold the canonically normalised four-dimensional field
  \(W^I_\mu(x)\) of its massless mode,
  \(g_2W^I_\mu(x) = W^I_\mu(\sigma = 0, x)\); \(W^I_{MN}\),
  \(W^I_{\mu\nu}\) their field strengths;
  \(W^\pm_\mu = (W^1_\mu \mp \mathrm{i}W^2_\mu)/\sqrt2\). Unit: GeV
  (field), GeV\(^2\) (field strength). Status: AUTHOR.
\item
  \(c^{B}_X\): The \(B\)-kinetic portal coupling of the half \(X\).
  Unit: GeV\(^{-7}\). Status: INVENTED.
\item
  \(\langle B^2\rangle\): The mean square of the hypercharge field
  strength, \(\langle B_{\mu\nu}B^{\mu\nu}\rangle\). Unit: GeV\(^4\).
  Status: STANDARD.
\item
  \(B_M,\ B_\mu,\ B_{MN}\): The hypercharge gauge field of \(U(1)_Y\):
  in the bulk \(B_M\) (geometric normalisation), on the fold the
  canonically normalised four-dimensional field \(B_\mu(x)\) of its
  massless mode, \(g_1B_\mu(x) = B_\mu(\sigma = 0, x)\); \(B_{MN}\),
  \(B_{\mu\nu}\) their field strengths. Unit: GeV (field), GeV\(^2\)
  (field strength). Status: AUTHOR.
\item
  \(\omega\): A frequency: the energy of a mode
  \(\propto \mathrm{e}^{-\mathrm{i}\omega x_4}\) (\(\omega_o\) of the
  orbital \(o\)). Unit: GeV. Status: STANDARD.
\item
  \(f\): The label of a kind of Standard-Model multiplet:
  \(f \in \{Q, u, d, \ell, e, \nu\}\) (quark doublets, up-quark
  singlets, down-quark singlets, lepton doublets, charged-lepton
  singlets, right-handed neutrinos); a sum over \(f\) runs over the six.
  Status: STANDARD.
\item
  \(c^{f}_X\): The current-portal coupling of the half \(X\) of
  dirac16complex with the fermion-number current of the multiplets of
  kind \(f\) (flavour blind); \(c^{\ell}_X\) and \(c^{\nu}_X\) are those
  of the lepton doublets and of the right-handed neutrinos. Unit:
  GeV\(^{-7}\). Status: INVENTED.
\item
  \(n_f\): The number density of the Standard-Model fermions of kind
  \(f\) (particles minus antiparticles, summed over generations and
  gauge components), the time component of the mean current
  \(\langle J_f^\mu\rangle\) in the rest frame of the medium. Unit:
  GeV\(^3\). Status: STANDARD.
\item
  \(v\): The vacuum expectation value of the lower component of the
  Higgs doublet, \(v^2 = -\mu^2/\lambda\). Unit: GeV. Value: 246.22 GeV
  \(= (\sqrt2\,G_{\mathrm{F}})^{-1/2}\) with the Fermi constant
  \(G_{\mathrm{F}} = 1.1663788\times10^{-5}\) GeV\(^{-2}\) of the muon
  lifetime (Particle Data Group). Status: STANDARD.
\item
  \(\mathcal{A}\): The author's ad hoc upper component of the Higgs
  vacuum, almightyAdHoc,
  \(\mathcal{A}(x_4) = v_{\mathcal{A}}\mathrm{e}^{-2a_{40}x_4}\). Unit:
  GeV. Status: AUTHOR.
\end{itemize}

\[
\begin{aligned}
&q_{\mathrm{b}} = -m_\Psi\,\big(\mu_{\mathrm{L}} + \mu_{\mathrm{R}}\big)\,\frac{\tanh\varsigma}{2\varsigma},\qquad \varsigma = \sqrt{\mu_{\mathrm{L}}\mu_{\mathrm{R}}},\qquad
\omega_{\mathrm{b}} = \sqrt{m_\Psi^2 - q_{\mathrm{b}}^2},\\
&\omega_{\mathrm{r}} = \omega_{\mathrm{b}}\Big(1 - \frac{H}{2q_{\mathrm{b}}}\Big) + \dots,\qquad
\lvert\chi_{\mathrm{b}}(\sigma)\rvert \propto \mathrm{e}^{-q_{\mathrm{b}}\lvert\sigma\rvert}\quad\big(\lvert\sigma\rvert < \sigma_{\mathrm{t}}\big)
\end{aligned} \tag{FE-40}
\] \textbf{Symbols in (FE-40).}

\begin{itemize}
\tightlist
\item
  \(q_{\mathrm{b}}\): The inverse width of a state of dirac16complex at
  the fold: inside the turning point its profile falls as
  \(\mathrm{e}^{-q_{\mathrm{b}}\lvert\sigma\rvert}\). Unit: GeV. Status:
  DERIVED.
\item
  \(m_\Psi\): The mass parameter of dirac16complex in the bulk. Unit:
  GeV. Status: RECORD.
\item
  \(\mu_X,\ \mu_{\mathrm{L}},\ \mu_{\mathrm{R}}\): The dimensionless
  strengths of the fold for the left-handed and right-handed halves of
  dirac16complex, set by the mean fields of the Standard Model at the
  fold: the jump of the profile across the fold is the exponential of
  \(\gamma^{(0)}(\mu_{\mathrm{L}}P_{\mathrm{L}} + \mu_{\mathrm{R}}P_{\mathrm{R}})\);
  negative values attract. Unit: dimensionless. Status: DERIVED.
\item
  \(\varsigma\): The auxiliary variable of the bound-state condition,
  \(\varsigma = \sqrt{\mu_{\mathrm{L}}\mu_{\mathrm{R}}}\)
  (\(\varsigma_- = \sqrt{-\mu_{\mathrm{L}}\mu_{\mathrm{R}}}\) when
  \(\mu_{\mathrm{L}}\mu_{\mathrm{R}} < 0\)). Unit: dimensionless.
  Status: DERIVED.
\item
  \(\omega_{\mathrm{b}}\): The energy at zero momentum of the state of
  dirac16complex at the fold in the limit \(H \to 0\) (the bound state
  of UM-3): the particle \(\chi\) of UM-4 has the mass
  \(\omega_{\mathrm{r}} = \omega_{\mathrm{b}}(1 - H/(2q_{\mathrm{b}})) + \dots\),
  equal to \(\omega_{\mathrm{b}}\) up to the relative correction
  \(H/(2q_{\mathrm{b}})\). Unit: GeV. Status: DERIVED.
\item
  \(\omega_{\mathrm{r}}\): The mass of the particle \(\chi\): the real
  part of the energy at rest of the resonance of dirac16complex at the
  fold,
  \(\omega_{\mathrm{r}} = \omega_{\mathrm{b}}(1 - H/(2q_{\mathrm{b}})) + \dots\).
  Unit: GeV. Status: DERIVED.
\item
  \(H\): The author's constant (an inverse length) that fixes the period
  \(\pi/(6H)\) of \(x_0\), the warp factor
  \(\mathrm{e}^{-2H\lvert\sigma\rvert}\) of the metric, the cosmological
  constant \(\Lambda = -21H^2\) and the tension
  \(\kappa\lambda_{\mathrm{F}} = 12H\) of the fold; in UM-4 a free
  parameter, bounded from below by the propagation of the gauge fields
  in the bulk. Unit: GeV (1/length). Value: \(H \gtrsim 2.5\) TeV (from
  the invisible width of the Z boson). Status: AUTHOR.
\item
  \(\chi_{\mathrm{b}}\): The profile of a state of dirac16complex at the
  fold (a resonance): it decays as
  \(\mathrm{e}^{-q_{\mathrm{b}}\lvert\sigma\rvert}\) out to the turning
  point and continues beyond it as an outgoing wave of exponentially
  small amplitude. Unit: GeV\(^{2}\). Status: DERIVED.
\item
  \(\mathrm{e}\): Euler's number, the base of the natural logarithm;
  \(\mathrm{e}^{X}\) is the exponential of \(X\). Unit: dimensionless.
  Value: 2.718281828. Status: STANDARD.
\item
  \(\sigma\): The proper coordinate of the hidden direction:
  \(\mathrm{d}\sigma = \lvert\cot z\rvert\,\mathrm{d}x_0\),
  \(\sigma = \mathrm{sgn}(\pi/2 - z)\ln(\sin z)/(6H)\) on one period
  \(0 < z < \pi\); \(\sigma = 0\) is the fold \(z = \pi/2\),
  \(\sigma \to \mp\infty\) the tips \(z \to 0, \pi\); one period of
  \(x_0\) is the whole line \(-\infty < \sigma < \infty\). Unit: length.
  Status: DERIVED.
\item
  \(\sigma_{\mathrm{t}}\): The turning point of a state of
  dirac16complex of the frequency \(\omega\): the depth at which the
  warped frequency equals the bulk mass,
  \(\lvert\omega\rvert\mathrm{e}^{H\sigma_{\mathrm{t}}} = m_\Psi\);
  beyond it every solution oscillates. Unit: length. Status: DERIVED.
\end{itemize}

\[
\omega_0 = \omega_{\mathrm{r}} - \frac{\mathrm{i}}{2}\,\Gamma_\chi,\qquad
\Gamma_\chi \simeq p_{\mathrm{w}}\,m_\Psi\exp\Big(-2F(\upsilon_{\mathrm{b}})\,\frac{m_\Psi}{H}\Big),\qquad
F(\upsilon) = \ln\frac{1 + \sqrt{1 - \upsilon^2}}{\upsilon} - \sqrt{1 - \upsilon^2},\qquad \upsilon_{\mathrm{b}} = \frac{\omega_{\mathrm{b}}}{m_\Psi} \tag{FE-86}
\] \textbf{Symbols in (FE-86).}

\begin{itemize}
\tightlist
\item
  \(\omega_0\): The energy at rest (zero momentum along ordinary space
  and along the extra times) of a state of dirac16complex at the fold;
  complex for a resonance,
  \(\omega_0 = \omega_{\mathrm{r}} - \frac{\mathrm{i}}{2}\Gamma_\chi\).
  Unit: GeV. Status: DERIVED.
\item
  \(\omega_{\mathrm{r}}\): The mass of the particle \(\chi\): the real
  part of the energy at rest of the resonance of dirac16complex at the
  fold,
  \(\omega_{\mathrm{r}} = \omega_{\mathrm{b}}(1 - H/(2q_{\mathrm{b}})) + \dots\).
  Unit: GeV. Status: DERIVED.
\item
  \(\mathrm{i}\): The imaginary unit, \(\mathrm{i}^2 = -1\). Unit:
  dimensionless. Status: STANDARD.
\item
  \(\Gamma_\chi\): The width of the particle \(\chi\) at rest: the rate
  of its decay into the waves of dirac16complex in the bulk. Unit: GeV.
  Value: below \(1/t_0 = 1.5\times10^{-42}\) GeV for dark matter.
  Status: DERIVED.
\item
  \(p_{\mathrm{w}}\): The prefactor of the width of the particles
  \(\chi\) in its tunnelling form
  \(\Gamma_\chi \approx p_{\mathrm{w}}m_\Psi\mathrm{e}^{-2F(\upsilon_{\mathrm{b}})m_\Psi/H}\),
  taken from the computed widths at \(m_\Psi/H = 40\). Unit:
  dimensionless. Value: 0.14. Status: ESTIMATE.
\item
  \(m_\Psi\): The mass parameter of dirac16complex in the bulk. Unit:
  GeV. Status: RECORD.
\item
  \(\mathrm{e}\): Euler's number, the base of the natural logarithm;
  \(\mathrm{e}^{X}\) is the exponential of \(X\). Unit: dimensionless.
  Value: 2.718281828. Status: STANDARD.
\item
  \(F(\upsilon)\): The tunnelling exponent of the width of the particles
  \(\chi\) per unit of \(m_\Psi/H\): the integral of the local decay
  rate \(q\) over the region inside the turning point, divided by
  \(m_\Psi/H\), as a function of the ratio \(\upsilon\) of the energy to
  the bulk mass. Unit: dimensionless. Value: 0.049 for
  \(\upsilon = 0.866\); \((q_{\mathrm{b}}/\omega_{\mathrm{b}})^3/3\) for
  weak binding. Status: DERIVED.
\item
  \(\upsilon,\ \upsilon_{\mathrm{b}}\): The argument of the tunnelling
  exponent \(F\), a ratio of an energy to the bulk mass \(m_\Psi\);
  \(\upsilon_{\mathrm{b}} = \omega_{\mathrm{b}}/m_\Psi\) that of the
  particle \(\chi\). Unit: dimensionless. Status: DERIVED.
\item
  \(\omega_{\mathrm{b}}\): The energy at zero momentum of the state of
  dirac16complex at the fold in the limit \(H \to 0\) (the bound state
  of UM-3): the particle \(\chi\) of UM-4 has the mass
  \(\omega_{\mathrm{r}} = \omega_{\mathrm{b}}(1 - H/(2q_{\mathrm{b}})) + \dots\),
  equal to \(\omega_{\mathrm{b}}\) up to the relative correction
  \(H/(2q_{\mathrm{b}})\). Unit: GeV. Status: DERIVED.
\item
  \(H\): The author's constant (an inverse length) that fixes the period
  \(\pi/(6H)\) of \(x_0\), the warp factor
  \(\mathrm{e}^{-2H\lvert\sigma\rvert}\) of the metric, the cosmological
  constant \(\Lambda = -21H^2\) and the tension
  \(\kappa\lambda_{\mathrm{F}} = 12H\) of the fold; in UM-4 a free
  parameter, bounded from below by the propagation of the gauge fields
  in the bulk. Unit: GeV (1/length). Value: \(H \gtrsim 2.5\) TeV (from
  the invisible width of the Z boson). Status: AUTHOR.
\end{itemize}

dirac16complex, the author's field, on the static block (the late-time
state of the metric, in which \(a_4\) and \(b_4\) are constant). At zero
momentum (FE-37) the local gap
\(m_\Psi^2 - \omega^2\mathrm{e}^{2H\lvert\sigma\rvert}\) of its profile
closes at the turning point \(\sigma_{\mathrm{t}}\), because the warped
frequency grows away from the fold; beyond it every solution oscillates.
Its continuum therefore starts at zero frequency, and the fold, which
acts on it as a chirality-dependent mass concentrated at \(\sigma = 0\)
with the strengths \(\mu_{\mathrm{L}}\) and \(\mu_{\mathrm{R}}\) set by
the mean fields of the Standard Model (FE-39), holds no bound states but
resonances (FE-40, FE-86): eight states, four of each Krein sign, whose
profile decays as \(\mathrm{e}^{-q_{\mathrm{b}}\lvert\sigma\rvert}\) out
to the turning point and continues beyond it as an outgoing wave of
exponentially small amplitude. Their mass \(\omega_{\mathrm{r}}\) is the
energy \(\omega_{\mathrm{b}}\) of the bound states of UM-3 lowered by
the warped time, and their width \(\Gamma_\chi\) is the rate of the
tunnelling through the region inside the turning point. Computed with
the author's exact \(16 \times 16\) matrices for a fold bound by the
Higgs portal alone (\(\mu_{\mathrm{L}} = -1\),
\(\mu_{\mathrm{R}} = 0\)): \(\omega_{\mathrm{r}}/m_\Psi = 0.7785\),
0.8175, 0.8333 and 0.8418 and \(\Gamma_\chi/m_\Psi = 0.0587\), 0.0180,
0.00670 and 0.00263 for \(m_\Psi/H = 10\), 20, 30 and 40, approaching
the bound state of UM-3, 0.8660, as \(H/m_\Psi \to 0\). Status: DERIVED
(\texttt{um4\_spectra.py}, U1 to U4, U7), ESTIMATE (the prefactor
\(p_{\mathrm{w}}\)).

\[
\ell_{\mathrm{b}} = \sqrt{\big\langle\sigma^2\big\rangle} = \frac{1}{\sqrt2\,q_{\mathrm{b}}},\qquad
m_\Psi - \omega_{\mathrm{r}} = \big(m_\Psi - \omega_{\mathrm{b}}\big) + \frac{\omega_{\mathrm{b}}H}{2q_{\mathrm{b}}} + \dots,\qquad
m_\Psi - \omega_{\mathrm{b}} \approx \frac{q_{\mathrm{b}}^2}{2m_\Psi} = \frac{1}{4m_\Psi\,\ell_{\mathrm{b}}^2}\quad\big(q_{\mathrm{b}} \ll m_\Psi\big) \tag{A-4}
\] \textbf{Symbols in (A-4).}

\begin{itemize}
\tightlist
\item
  \(\ell_{\mathrm{b}}\): The width of a state of dirac16complex at the
  fold along the hidden direction: the root mean square of \(\sigma\) in
  its density \(\mathrm{e}^{-2q_{\mathrm{b}}\lvert\sigma\rvert}\) inside
  the turning point. Unit: length. Value:
  \(1/(\sqrt2\,q_{\mathrm{b}})\). Status: DERIVED.
\item
  \(\langle\cdots\rangle\): The mean (expectation) value of an operator
  in the state of the medium (the vacuum, a plasma, the Kohn-Sham
  state). Status: STANDARD.
\item
  \(\sigma\): The proper coordinate of the hidden direction:
  \(\mathrm{d}\sigma = \lvert\cot z\rvert\,\mathrm{d}x_0\),
  \(\sigma = \mathrm{sgn}(\pi/2 - z)\ln(\sin z)/(6H)\) on one period
  \(0 < z < \pi\); \(\sigma = 0\) is the fold \(z = \pi/2\),
  \(\sigma \to \mp\infty\) the tips \(z \to 0, \pi\); one period of
  \(x_0\) is the whole line \(-\infty < \sigma < \infty\). Unit: length.
  Status: DERIVED.
\item
  \(q_{\mathrm{b}}\): The inverse width of a state of dirac16complex at
  the fold: inside the turning point its profile falls as
  \(\mathrm{e}^{-q_{\mathrm{b}}\lvert\sigma\rvert}\). Unit: GeV. Status:
  DERIVED.
\item
  \(m_\Psi\): The mass parameter of dirac16complex in the bulk. Unit:
  GeV. Status: RECORD.
\item
  \(\omega_{\mathrm{r}}\): The mass of the particle \(\chi\): the real
  part of the energy at rest of the resonance of dirac16complex at the
  fold,
  \(\omega_{\mathrm{r}} = \omega_{\mathrm{b}}(1 - H/(2q_{\mathrm{b}})) + \dots\).
  Unit: GeV. Status: DERIVED.
\item
  \(\omega_{\mathrm{b}}\): The energy at zero momentum of the state of
  dirac16complex at the fold in the limit \(H \to 0\) (the bound state
  of UM-3): the particle \(\chi\) of UM-4 has the mass
  \(\omega_{\mathrm{r}} = \omega_{\mathrm{b}}(1 - H/(2q_{\mathrm{b}})) + \dots\),
  equal to \(\omega_{\mathrm{b}}\) up to the relative correction
  \(H/(2q_{\mathrm{b}})\). Unit: GeV. Status: DERIVED.
\item
  \(H\): The author's constant (an inverse length) that fixes the period
  \(\pi/(6H)\) of \(x_0\), the warp factor
  \(\mathrm{e}^{-2H\lvert\sigma\rvert}\) of the metric, the cosmological
  constant \(\Lambda = -21H^2\) and the tension
  \(\kappa\lambda_{\mathrm{F}} = 12H\) of the fold; in UM-4 a free
  parameter, bounded from below by the propagation of the gauge fields
  in the bulk. Unit: GeV (1/length). Value: \(H \gtrsim 2.5\) TeV (from
  the invisible width of the Z boson). Status: AUTHOR.
\end{itemize}

The width and the mass of the particles \(\chi\): their width
\(\ell_{\mathrm{b}}\) along the hidden direction (the root mean square
of \(\sigma\) in the density
\(\mathrm{e}^{-2q_{\mathrm{b}}\lvert\sigma\rvert}\) inside the turning
point) and their mass deficit below the bulk mass. The states obey the
ordinary uncertainty relation: the mean square of the momentum along
\(\sigma\) in the profile
\(\mathrm{e}^{-q_{\mathrm{b}}\lvert\sigma\rvert}\) is
\(q_{\mathrm{b}}^2\), so that
\(\langle p_\sigma^2\rangle\langle\sigma^2\rangle = 1/2\). The narrower
the state, the larger the deficit, because a narrower state is held by a
stronger attraction of the fold: the energy falls not because the
particle is confined (in the compact case confinement raises it) but
because the fold binds it more strongly. The warped time lowers the mass
further by \(\omega_{\mathrm{b}}H\) times the mean distance
\(\langle\lvert\sigma\rvert\rangle = 1/(2q_{\mathrm{b}})\) of the state
from the fold (the redshift of its energy, averaged over its profile).
The \(x_0\) mass of dirac16complex is therefore a binding mass, set by
the bulk mass and by the strengths of the fold terms, not by the period
of \(x_0\). Status: DERIVED (\texttt{um4\_kohnsham.py}, K1: the level
equals \(\omega_{\mathrm{b}}(1 - H/(2q_{\mathrm{b}}))\) up to a term of
the second order in \(H\)).

\hypertarget{the-answer}{%
\subsection{The answer}\label{the-answer}}

The period \(\pi/(6H)\) of the coordinate \(x_0\) does not produce a
mass by the uncertainty relation, because one period is an infinite
proper distance. In UM-4 the hidden direction produces neither a tower
nor a gap: every massless field of the bulk has one massless mode bound
to the fold and a continuum of massive modes from zero mass upwards,
which reaches the fold with the weight \(\mathfrak{s}_n(m)\), vanishing
like \((m/H)^4\) (graviton, scalars) or \((m/H)^2\) (gauge fields) for
light modes; the fold sees it as the corrections \(0.34/(Hr)^5\) to
Newton's law and \(0.42/(Hr)^3\) to Coulomb's law and as the widths of
the W and Z bosons into the bulk. For dirac16complex the fold holds
eight long-lived resonances of mass
\(\omega_{\mathrm{r}} \simeq \omega_{\mathrm{b}}(1 - H/(2q_{\mathrm{b}}))\),
which are lighter the narrower they are, because a narrower state is
bound more strongly, and which decay into the bulk at a rate
exponentially small in \(m_\Psi/H\). Status: DERIVED.

\hypertarget{dirac16complex-as-dark-matter}{%
\section{dirac16complex as dark
matter}\label{dirac16complex-as-dark-matter}}

\hypertarget{stability-and-lifetime}{%
\subsection{Stability and lifetime}\label{stability-and-lifetime}}

\[
\begin{aligned}
\mathcal{L}_{\mathrm{P}} = {}&-c_\Phi\,\big(\Phi^\dagger\Phi\big)\,\bar\Psi P_{\mathrm{L}}\Psi
 - \sum_{X}\Big[\frac14\big(c^{G}_X\,G^A_{\mu\nu}G^{A\mu\nu} + c^{W}_X\,W^I_{\mu\nu}W^{I\mu\nu}
 + c^{B}_X\,B_{\mu\nu}B^{\mu\nu}\big)\,\bar\Psi P_X\Psi\\
&\qquad + \frac{d_X}{2}\,B_{\mu\nu}\,\mathcal{D}_X^{\mu\nu} + \sum_{f} c^{f}_X\,J_f^\mu\,j_{X\mu}\Big],\\
\mathcal{D}_X^{\mu\nu} = {}&\mathrm{i}\,\bar\Psi\,\Sigma^{\mu\nu}P_X\Psi,\qquad
j_X^\mu = \mathrm{i}\big(\bar\Psi P_X D^\mu\Psi - (D^\mu\bar\Psi)P_X\Psi\big),\qquad
J_f^\mu = \bar\psi_f\,\hat\gamma^\mu\psi_f
\end{aligned} \tag{IN-16}
\] \textbf{Symbols in (IN-16).}

\begin{itemize}
\tightlist
\item
  \(\mathcal{L}_{\mathrm{P}}\): The Lagrangian of the portals of
  dirac16complex at the fold. Unit: GeV\(^4\). Status: INVENTED.
\item
  \(c_\Phi\): The Higgs-portal coupling of the left-handed half of
  dirac16complex (the right-handed half has none). Unit: GeV\(^{-5}\).
  Status: INVENTED.
\item
  \(\Phi\): The Higgs doublet (two complex components, hypercharge
  \(1/2\)), a field of the fold. Unit: GeV. Status: AUTHOR.
\item
  \({}^\dagger\): Hermitian conjugation: the complex conjugate of the
  transpose of a matrix or spinor; for an operator on the space of
  states, its adjoint with respect to the scalar product of that space
  (for the quantised dirac16complex the indefinite Krein scalar
  product). Status: STANDARD.
\item
  \(\Psi,\ \bar\Psi\): dirac16complex: the author's 16-component complex
  spinor field of eight dimensions;
  \(\bar\Psi = \Psi^\dagger\sigma_{16}\) is its Dirac adjoint; in a term
  of the fold, \(\Psi\) is its value at \(\sigma = 0\). Unit:
  GeV\(^{7/2}\). Status: AUTHOR.
\item
  \(P_{\mathrm{L}}\): The author's left-handed chiral projector,
  \(P_{\mathrm{L}} = (\mathbf{1} - \Gamma_8)/2\). Unit: dimensionless.
  Status: AUTHOR.
\item
  \(X\): The label of a half of dirac16complex:
  \(X \in \{\mathrm{L}, \mathrm{R}\}\) (left-handed or right-handed
  half); a sum over \(X\) runs over both. Status: STANDARD.
\item
  \(c^{G}_X\): The gluon-kinetic portal coupling of the half \(X\) of
  dirac16complex. Unit: GeV\(^{-7}\). Status: INVENTED.
\item
  \(G^A_M,\ G^A_\mu,\ G^A_{MN}\): The gluon field of \(SU(3)_C\)
  (\(A = 1, \dots, 8\)): in the bulk \(G^A_M\), a field of eight
  dimensions in its geometric normalisation (the covariant derivative
  \(\partial - \mathrm{i}G\)), and on the fold the canonically
  normalised four-dimensional field \(G^A_\mu(x)\) of its massless mode,
  \(g_3G^A_\mu(x) = G^A_\mu(\sigma = 0, x)\); \(G^A_{MN}\) and
  \(G^A_{\mu\nu}\) their field strengths. Unit: GeV (field), GeV\(^2\)
  (field strength). Status: AUTHOR.
\item
  \(A\): A colour (adjoint SU(3)) index running over \(1, \dots, 8\); a
  repeated index is summed. Status: STANDARD.
\item
  \(\mu\): A four-dimensional coordinate index running over the
  observable coordinates \(x_1, x_2, x_3, x_4\) of the fold (signature
  \(+,+,+,-\)); a repeated index is summed. Status: STANDARD.
\item
  \(\nu\): A four-dimensional coordinate index (as \(\mu\)). Status:
  STANDARD.
\item
  \(c^{W}_X\): The \(W\)-kinetic portal coupling of the half \(X\).
  Unit: GeV\(^{-7}\). Status: INVENTED.
\item
  \(W^I_M,\ W^I_\mu,\ W^I_{MN}\): The gauge field of \(SU(2)_L\)
  (\(I = 1, 2, 3\)): in the bulk \(W^I_M\) (geometric normalisation), on
  the fold the canonically normalised four-dimensional field
  \(W^I_\mu(x)\) of its massless mode,
  \(g_2W^I_\mu(x) = W^I_\mu(\sigma = 0, x)\); \(W^I_{MN}\),
  \(W^I_{\mu\nu}\) their field strengths;
  \(W^\pm_\mu = (W^1_\mu \mp \mathrm{i}W^2_\mu)/\sqrt2\). Unit: GeV
  (field), GeV\(^2\) (field strength). Status: AUTHOR.
\item
  \(I\): A weak-isospin (adjoint SU(2)) index running over \(1, 2, 3\);
  a repeated index is summed. Status: STANDARD.
\item
  \(c^{B}_X\): The \(B\)-kinetic portal coupling of the half \(X\).
  Unit: GeV\(^{-7}\). Status: INVENTED.
\item
  \(B_M,\ B_\mu,\ B_{MN}\): The hypercharge gauge field of \(U(1)_Y\):
  in the bulk \(B_M\) (geometric normalisation), on the fold the
  canonically normalised four-dimensional field \(B_\mu(x)\) of its
  massless mode, \(g_1B_\mu(x) = B_\mu(\sigma = 0, x)\); \(B_{MN}\),
  \(B_{\mu\nu}\) their field strengths. Unit: GeV (field), GeV\(^2\)
  (field strength). Status: AUTHOR.
\item
  \(P_X\): The chiral projector of the half \(X\):
  \(P_X = P_{\mathrm{L}}\) for \(X = \mathrm{L}\), \(P_{\mathrm{R}}\)
  for \(X = \mathrm{R}\). Unit: dimensionless. Status: AUTHOR.
\item
  \(d_X\): The dipole-portal coupling of the half \(X\) of
  dirac16complex with the hypercharge field strength. Unit:
  GeV\(^{-5}\). Status: INVENTED.
\item
  \(\mathcal{D}_X^{\mu\nu}\): The dipole (antisymmetric tensor) density
  of the half \(X\) of dirac16complex at the fold,
  \(\mathcal{D}_X^{\mu\nu} = \mathrm{i}\bar\Psi\Sigma^{\mu\nu}P_X\Psi\)
  (real, DERIVED). Unit: GeV\(^7\). Status: DERIVED.
\item
  \(f\): The label of a kind of Standard-Model multiplet:
  \(f \in \{Q, u, d, \ell, e, \nu\}\) (quark doublets, up-quark
  singlets, down-quark singlets, lepton doublets, charged-lepton
  singlets, right-handed neutrinos); a sum over \(f\) runs over the six.
  Status: STANDARD.
\item
  \(c^{f}_X\): The current-portal coupling of the half \(X\) of
  dirac16complex with the fermion-number current of the multiplets of
  kind \(f\) (flavour blind); \(c^{\ell}_X\) and \(c^{\nu}_X\) are those
  of the lepton doublets and of the right-handed neutrinos. Unit:
  GeV\(^{-7}\). Status: INVENTED.
\item
  \(J_f^\mu\): The fermion-number current of the Standard-Model
  multiplets of kind \(f\), \(J_f^\mu = \bar\psi_f\hat\gamma^\mu\psi_f\)
  summed over the three generations and the gauge components; with other
  labels the currents that are the sources of the gauge fields. Unit:
  GeV\(^3\). Status: STANDARD.
\item
  \(j_X^\mu\): The convective current of the half \(X\) of
  dirac16complex,
  \(j_X^\mu = \mathrm{i}(\bar\Psi P_XD^\mu\Psi - (D^\mu\bar\Psi)P_X\Psi)\)
  (real and joining the half with itself, DERIVED). Unit: GeV\(^8\).
  Status: DERIVED.
\item
  \(\mathrm{i}\): The imaginary unit, \(\mathrm{i}^2 = -1\). Unit:
  dimensionless. Status: STANDARD.
\item
  \(\Sigma^{ab},\ \Sigma^{\mu\nu}\): The generators of Spin(4,4) on
  dirac16complex, \(\Sigma^{ab} = \frac14[T^a, T^b]\) (the scaled
  commutators of the author's matrices);
  \(\Sigma^{\mu\nu} = \mathbf{e}_a{}^\mu\mathbf{e}_b{}^\nu\Sigma^{ab}\)
  along the fold. Unit: dimensionless. Status: AUTHOR.
\item
  \(D_M,\ D_\mu\): The covariant derivative: on dirac16complex
  \(D_M = \partial_M + \Omega_M\) (with
  \(D_M\bar\Psi = \partial_M\bar\Psi - \bar\Psi\Omega_M\)); on the
  fields of the fold the gauge and spin covariant derivative \(D_\mu\)
  (the partial derivative plus the spin connection of the fold and the
  gauge fields times their couplings and generators);
  \(D^\mu = \hat h^{\mu\nu}D_\nu\); on an adjoint field
  \((D_\mu X)^A = \partial_\mu X^A + g_3f^{ABC}G^B_\mu X^C\) (likewise
  for \(SU(2)\) with \(g_2\epsilon^{IJK}\)). Unit: 1/length. Status:
  STANDARD.
\item
  \(\psi,\ \bar\psi\): A Standard-Model fermion field (a
  four-dimensional Weyl field of the fold with three generations,
  written as a four-component spinor of fixed chirality);
  \(\bar\psi = \psi^\dagger\beta\); \(\psi_f\) the fields of the kind
  \(f\). Unit: GeV\(^{3/2}\). Status: STANDARD.
\item
  \(\hat\gamma^\mu\): The four-dimensional complex \(4 \times 4\) Dirac
  matrices of the fold: \(\hat\gamma^{(a)}\) along the frame directions
  of the fold and \(\hat\gamma^\mu\) along its coordinates (the frame
  matrices contracted with the inverse frame field of the fold),
  \(\{\hat\gamma^\mu, \hat\gamma^\nu\} = -2\hat h^{\mu\nu}\mathbf{1}\);
  \(\hat\gamma_\mu = \hat h_{\mu\nu}\hat\gamma^\nu\). Unit:
  dimensionless. Status: STANDARD.
\end{itemize}

\[
\begin{aligned}
&Q_\Psi = \int\mathrm{d}\sigma\,\mathrm{d}^3x\,\sqrt{\lvert g\rvert}\;\mathrm{e}^{H\lvert\sigma\rvert}\,\Psi^\dagger\mathfrak{B}\Psi
 = b_4^3\int\mathrm{d}\sigma\,\mathrm{d}^3x\;\mathrm{e}^{H\lvert\sigma\rvert}\,\chi^\dagger\mathfrak{B}\chi,\\
&\mathfrak{B} = \mathrm{i}\,\sigma_{16}\,\gamma^{(4)},\qquad \mathfrak{B}^\dagger = \mathfrak{B},\qquad \mathfrak{B}^2 = \mathbf{1},\qquad
E_o = -\omega_o\,\mathcal{Q}_o
\end{aligned} \tag{FE-43}
\] \textbf{Symbols in (FE-43).}

\begin{itemize}
\tightlist
\item
  \(Q_\Psi\): The conserved charge of dirac16complex (the Noether charge
  of a change of the phase of \(\Psi\) by a constant), a Krein form
  (indefinite); as written it is the charge per unit coordinate volume
  of the extra times, on which no field depends. Unit: GeV\(^3\).
  Status: DERIVED.
\item
  \(\mathrm{d}\): The differential: \(\mathrm{d}\) in front of a
  quantity is an infinitesimal change of it (as \(\mathrm{d}\sigma\), a
  small step of \(\sigma\)); the quotient of two such changes is a
  derivative; \(\mathrm{d}\) with an exponent, in front of several
  variables, is their volume element, the product of their
  differentials. Status: STANDARD.
\item
  \(\sigma\): The proper coordinate of the hidden direction:
  \(\mathrm{d}\sigma = \lvert\cot z\rvert\,\mathrm{d}x_0\),
  \(\sigma = \mathrm{sgn}(\pi/2 - z)\ln(\sin z)/(6H)\) on one period
  \(0 < z < \pi\); \(\sigma = 0\) is the fold \(z = \pi/2\),
  \(\sigma \to \mp\infty\) the tips \(z \to 0, \pi\); one period of
  \(x_0\) is the whole line \(-\infty < \sigma < \infty\). Unit: length.
  Status: DERIVED.
\item
  \(x\): The coordinates of the fold: \(x_1, x_2, x_3\) of ordinary
  space, \(x_4\) the time, \(x_5, x_6, x_7\) the extra times; a field of
  the fold is a function of the four observable coordinates
  \(x = (x_1, x_2, x_3, x_4)\), and
  \(\mathrm{d}^4x = \mathrm{d}x_1\,\mathrm{d}x_2\,\mathrm{d}x_3\,\mathrm{d}x_4\).
  Unit: length (each coordinate). Status: AUTHOR.
\item
  \(g_{MN}\): The metric tensor of the eight-dimensional spacetime
  (signature \(+,+,+,+,-,-,-,-\) in the order \(x_0, \dots, x_7\)):
  \(g_{MN}\,\mathrm{d}X^M\mathrm{d}X^N\) is the squared interval,
  \(g^{MN}\) its inverse, \(\lvert g\rvert\) the absolute value of its
  determinant (\(\sqrt{\lvert g\rvert}\,\mathrm{d}^8X\) is the proper
  volume element). Unit: dimensionless. Status: AUTHOR.
\item
  \(\mathrm{e}\): Euler's number, the base of the natural logarithm;
  \(\mathrm{e}^{X}\) is the exponential of \(X\). Unit: dimensionless.
  Value: 2.718281828. Status: STANDARD.
\item
  \(H\): The author's constant (an inverse length) that fixes the period
  \(\pi/(6H)\) of \(x_0\), the warp factor
  \(\mathrm{e}^{-2H\lvert\sigma\rvert}\) of the metric, the cosmological
  constant \(\Lambda = -21H^2\) and the tension
  \(\kappa\lambda_{\mathrm{F}} = 12H\) of the fold; in UM-4 a free
  parameter, bounded from below by the propagation of the gauge fields
  in the bulk. Unit: GeV (1/length). Value: \(H \gtrsim 2.5\) TeV (from
  the invisible width of the Z boson). Status: AUTHOR.
\item
  \(\Psi,\ \bar\Psi\): dirac16complex: the author's 16-component complex
  spinor field of eight dimensions;
  \(\bar\Psi = \Psi^\dagger\sigma_{16}\) is its Dirac adjoint; in a term
  of the fold, \(\Psi\) is its value at \(\sigma = 0\). Unit:
  GeV\(^{7/2}\). Status: AUTHOR.
\item
  \({}^\dagger\): Hermitian conjugation: the complex conjugate of the
  transpose of a matrix or spinor; for an operator on the space of
  states, its adjoint with respect to the scalar product of that space
  (for the quantised dirac16complex the indefinite Krein scalar
  product). Status: STANDARD.
\item
  \(\mathfrak{B}\): The Krein operator,
  \(\mathfrak{B} = \mathrm{i}\sigma_{16}\gamma^{(4)}\): a Hermitian
  involution with the signature (8, 8); \(\Psi^\dagger\mathfrak{B}\Psi\)
  is the charge density. Unit: dimensionless. Status: DERIVED.
\item
  \(b_4\): The author's metric function of the time \(x_4\) in the
  extra-time components, \(b_4 > 0\), never identically 1 (the author's
  order of 2026-10-24); the scale factor (the factor of the lengths) of
  the extra times is \(\mathrm{e}^{-a_4}b_4\) on the fold and
  \(\mathrm{e}^{-a_4}b_4\,\mathrm{e}^{-H\lvert\sigma\rvert}\) at the
  distance \(\sigma\) from it (the square root of the warp factor times
  \(\mathrm{e}^{-a_4}b_4\)). Unit: dimensionless. Status: AUTHOR.
\item
  \(\chi(\sigma)\): The profile of a mode of dirac16complex along the
  hidden direction,
  \(\Psi = \mathrm{e}^{\frac72H\lvert\sigma\rvert}\chi(\sigma)\mathrm{e}^{\mathrm{i}(\vec k\cdot\vec x - \omega x_4)}\)
  (a 16-component spinor function of \(\sigma\)); as a four-dimensional
  field, \(\chi\) denotes the particle of the level of dirac16complex at
  the fold (a long-lived resonance) and \(\bar\chi\) its antiparticle.
  Unit: GeV\(^{2}\) (with the \(V_*\) normalisation). Status: DERIVED.
\item
  \(\mathrm{i}\): The imaginary unit, \(\mathrm{i}^2 = -1\). Unit:
  dimensionless. Status: STANDARD.
\item
  \(\sigma_{16}\): The matrix of the Dirac adjoint,
  \(\sigma_{16} = T_0T_1T_2T_3\), real and symmetric; \(\sigma_{16}T_a\)
  is real and antisymmetric. Unit: dimensionless. Status: AUTHOR.
\item
  \(\gamma^M,\ \gamma^{(a)}\): The Dirac matrices of dirac16complex
  along the coordinates, \(\gamma^M = \mathbf{e}_a{}^MT^a\),
  \(\gamma_M = g_{MN}\gamma^N\); \(\gamma^{(a)} = T^a\) along the frame
  direction \((a)\) (\(\gamma^{(0)}\) along \(\sigma\), \(\gamma^{(4)}\)
  along the time). Unit: dimensionless. Status: AUTHOR.
\item
  \(\mathbf{1},\ \mathbf{1}_8\): The identity matrix of the size fixed
  by the context (\(16 \times 16\) on dirac16complex, \(4 \times 4\) on
  a four-dimensional spinor, \(3 \times 3\) in generation or flavour
  space); \(\mathbf{1}_8\) is the \(8 \times 8\) identity matrix. Unit:
  dimensionless. Status: STANDARD.
\item
  \(E\): An energy: \(E_o\) the energy of the Kohn-Sham orbital \(o\);
  in the neutrino sector the energy of a neutrino. Unit: GeV. Status:
  STANDARD.
\item
  \(o\): The label of a single-particle state of dirac16complex: a
  Kohn-Sham orbital, or a mode of the quantised field; a sum over \(o\)
  runs over the orbitals or the modes. Status: STANDARD.
\item
  \(\omega\): A frequency: the energy of a mode
  \(\propto \mathrm{e}^{-\mathrm{i}\omega x_4}\) (\(\omega_o\) of the
  orbital \(o\)). Unit: GeV. Status: STANDARD.
\item
  \(\mathcal{Q}_o\): The charge \(Q_\Psi\) (the Krein charge) of the
  single-particle state \(o\) of dirac16complex, a Kohn-Sham orbital or
  a mode: \(+1\) or \(-1\) for a normalised state. Unit: dimensionless.
  Status: DERIVED.
\end{itemize}

The portals of dirac16complex to the Standard Model (IN-16) and its
conserved charge (FE-43). Every portal term, like the bulk Lagrangian
and the self-interaction, contains dirac16complex only in the bilinears
\(\bar\Psi\cdots\Psi\): the action is invariant under a constant change
of the phase of \(\Psi\), and the charge \(Q_\Psi\) is conserved. A
particle \(\chi\) can therefore annihilate only with an antiparticle,
and it cannot decay into the Standard Model alone. Conservation of
\(Q_\Psi\) does not by itself make it stable, for two reasons. The first
is the indefinite Krein form of the charge: with
\(E_o = -\omega_o\mathcal{Q}_o\), a particle of positive energy (charge
\(-1\) at \(+\omega_{\mathrm{r}}\)) could turn into a state of negative
energy with the same charge and Standard-Model quanta that carry away
the difference; the quantisation of the charge settles this (the section
on the open questions, FE-67): in the positive-energy quantisation every
quantum carries a positive energy, so no such decay exists and the
vacuum is stable. The second is new in UM-4: the particle \(\chi\) is a
resonance of the fold and decays into the waves of dirac16complex in the
bulk, which carry its charge away from the fold (FE-86). Status: DERIVED
(no decay into the Standard Model alone, from the form of IN-16; the
Krein signature, \texttt{um4\_fields.py}, F3, and
\texttt{um4\_spectra.py}, U7; the quantisation, \texttt{um4\_krein.py},
KR1; the width, \texttt{um4\_spectra.py}, U3 and U4).

\[
\Gamma_\chi < \frac{1}{t_0}\quad\Leftrightarrow\quad
\frac{m_\Psi}{H} > \frac{\ln\big(p_{\mathrm{w}}\,m_\Psi\,t_0\big)}{2F(\upsilon_{\mathrm{b}})},\qquad
F(\upsilon_{\mathrm{b}}) \simeq \frac13\Big(\frac{q_{\mathrm{b}}}{\omega_{\mathrm{b}}}\Big)^3 - \frac{3}{10}\Big(\frac{q_{\mathrm{b}}}{\omega_{\mathrm{b}}}\Big)^5\quad
\Big(\frac{q_{\mathrm{b}}}{\omega_{\mathrm{b}}} \ll 1\Big) \tag{S-6}
\] \textbf{Symbols in (S-6).}

\begin{itemize}
\tightlist
\item
  \(\Gamma_\chi\): The width of the particle \(\chi\) at rest: the rate
  of its decay into the waves of dirac16complex in the bulk. Unit: GeV.
  Value: below \(1/t_0 = 1.5\times10^{-42}\) GeV for dark matter.
  Status: DERIVED.
\item
  \(t_0\): The age of the universe. Unit: time. Value: 13.8 Gyr. Status:
  STANDARD.
\item
  \(m_\Psi\): The mass parameter of dirac16complex in the bulk. Unit:
  GeV. Status: RECORD.
\item
  \(H\): The author's constant (an inverse length) that fixes the period
  \(\pi/(6H)\) of \(x_0\), the warp factor
  \(\mathrm{e}^{-2H\lvert\sigma\rvert}\) of the metric, the cosmological
  constant \(\Lambda = -21H^2\) and the tension
  \(\kappa\lambda_{\mathrm{F}} = 12H\) of the fold; in UM-4 a free
  parameter, bounded from below by the propagation of the gauge fields
  in the bulk. Unit: GeV (1/length). Value: \(H \gtrsim 2.5\) TeV (from
  the invisible width of the Z boson). Status: AUTHOR.
\item
  \(p_{\mathrm{w}}\): The prefactor of the width of the particles
  \(\chi\) in its tunnelling form
  \(\Gamma_\chi \approx p_{\mathrm{w}}m_\Psi\mathrm{e}^{-2F(\upsilon_{\mathrm{b}})m_\Psi/H}\),
  taken from the computed widths at \(m_\Psi/H = 40\). Unit:
  dimensionless. Value: 0.14. Status: ESTIMATE.
\item
  \(F(\upsilon)\): The tunnelling exponent of the width of the particles
  \(\chi\) per unit of \(m_\Psi/H\): the integral of the local decay
  rate \(q\) over the region inside the turning point, divided by
  \(m_\Psi/H\), as a function of the ratio \(\upsilon\) of the energy to
  the bulk mass. Unit: dimensionless. Value: 0.049 for
  \(\upsilon = 0.866\); \((q_{\mathrm{b}}/\omega_{\mathrm{b}})^3/3\) for
  weak binding. Status: DERIVED.
\item
  \(\upsilon,\ \upsilon_{\mathrm{b}}\): The argument of the tunnelling
  exponent \(F\), a ratio of an energy to the bulk mass \(m_\Psi\);
  \(\upsilon_{\mathrm{b}} = \omega_{\mathrm{b}}/m_\Psi\) that of the
  particle \(\chi\). Unit: dimensionless. Status: DERIVED.
\item
  \(q_{\mathrm{b}}\): The inverse width of a state of dirac16complex at
  the fold: inside the turning point its profile falls as
  \(\mathrm{e}^{-q_{\mathrm{b}}\lvert\sigma\rvert}\). Unit: GeV. Status:
  DERIVED.
\item
  \(\omega_{\mathrm{b}}\): The energy at zero momentum of the state of
  dirac16complex at the fold in the limit \(H \to 0\) (the bound state
  of UM-3): the particle \(\chi\) of UM-4 has the mass
  \(\omega_{\mathrm{r}} = \omega_{\mathrm{b}}(1 - H/(2q_{\mathrm{b}})) + \dots\),
  equal to \(\omega_{\mathrm{b}}\) up to the relative correction
  \(H/(2q_{\mathrm{b}})\). Unit: GeV. Status: DERIVED.
\end{itemize}

The lifetime that the dark matter requires. The particles \(\chi\) that
make up the dark matter today must live longer than the age
\(t_0 = 13.8\) Gyr of the universe,
\(\Gamma_\chi < 1/t_0 = 1.5\times10^{-42}\) GeV. With the width of FE-86
this bounds \(m_\Psi/H\) from below; for weak binding the tunnelling
exponent is small,
\(F \simeq (q_{\mathrm{b}}/\omega_{\mathrm{b}})^3/3\), and the bound
grows like \((\omega_{\mathrm{b}}/q_{\mathrm{b}})^3\). With \(H\) at its
lower bound 2534 GeV:

\begin{longtable}[]{@{}
  >{\raggedright\arraybackslash}p{(\columnwidth - 8\tabcolsep) * \real{0.2000}}
  >{\raggedright\arraybackslash}p{(\columnwidth - 8\tabcolsep) * \real{0.2000}}
  >{\raggedright\arraybackslash}p{(\columnwidth - 8\tabcolsep) * \real{0.2000}}
  >{\raggedright\arraybackslash}p{(\columnwidth - 8\tabcolsep) * \real{0.2000}}
  >{\raggedright\arraybackslash}p{(\columnwidth - 8\tabcolsep) * \real{0.2000}}@{}}
\toprule\noalign{}
\begin{minipage}[b]{\linewidth}\raggedright
binding
\end{minipage} & \begin{minipage}[b]{\linewidth}\raggedright
lifetime
\end{minipage} & \begin{minipage}[b]{\linewidth}\raggedright
\(F(\upsilon_{\mathrm{b}})\)
\end{minipage} & \begin{minipage}[b]{\linewidth}\raggedright
\(m_\Psi/H\) above
\end{minipage} & \begin{minipage}[b]{\linewidth}\raggedright
\(\omega_{\mathrm{b}}\) above (GeV)
\end{minipage} \\
\midrule\noalign{}
\endhead
\bottomrule\noalign{}
\endlastfoot
\(q_{\mathrm{b}} = m_\Psi/2\)
(\(q_{\mathrm{b}}/\omega_{\mathrm{b}} = 0.577\)) & \(t_0\) & 0.0493 &
\(1.11\times10^3\) & \(2.43\times10^6\) \\
\(q_{\mathrm{b}} = m_\Psi/2\) & \(100\,t_0\) & 0.0493 &
\(1.15\times10^3\) & \(2.53\times10^6\) \\
\(q_{\mathrm{b}}/\omega_{\mathrm{b}} = 0.1\) & \(t_0\) &
\(3.30\times10^{-4}\) & \(1.73\times10^5\) & \(4.36\times10^8\) \\
\(q_{\mathrm{b}}/\omega_{\mathrm{b}} = 0.1\) & \(100\,t_0\) &
\(3.30\times10^{-4}\) & \(1.80\times10^5\) & \(4.53\times10^8\) \\
\end{longtable}

The particles \(\chi\) of UM-4 are therefore heavier than about
\(2.4\times10^6\) GeV, more for weaker binding. A particle moving with
the momentum \(\vec k\) has the energy \(\sqrt{\omega_0^2 + \vec k^2}\)
and its width is dilated by the factor
\(\omega_{\mathrm{r}}/\sqrt{\omega_{\mathrm{r}}^2 + \vec k^2}\), as for
any unstable particle. Status: ESTIMATE (\texttt{um4\_spectra.py}, U6).

\hypertarget{gravitation-on-the-fold}{%
\subsection{Gravitation on the fold}\label{gravitation-on-the-fold}}

\[
\kappa\int_{\text{layer}}\mathcal{T}_\Psi{}^a{}_b\,\mathrm{d}\sigma = \frac{\kappa}{\mathcal{Z}\,V_*}\,\mathcal{T}_\chi{}^a{}_b\qquad\big(H \ll q_{\mathrm{b}}\big) \tag{A-5}
\] \textbf{Symbols in (A-5).}

\begin{itemize}
\tightlist
\item
  \(\kappa\): The eight-dimensional gravitational coupling,
  \(\kappa = 8\pi G_8\). Unit: GeV\(^{-6}\). Status: AUTHOR.
\item
  \(\mathcal{T}_\Psi\): The energy-momentum tensor of dirac16complex:
  \(-2/\sqrt{\lvert g\rvert}\) times the variation of the action
  \(S_\Psi\) of dirac16complex with respect to \(g^{MN}\). Unit:
  GeV\(^8\). Status: STANDARD.
\item
  \(a\): A frame (orthonormal) index running over \(0, 1, \dots, 7\)
  (frame direction \((0)\) along \(\sigma\), \((1), (2), (3)\) along
  ordinary space, \((4)\) along the time, \((5), (6), (7)\) along the
  extra times); raised and lowered with the frame metric
  \(\eta_{ab} = \mathrm{diag}(1, 1, 1, 1, -1, -1, -1, -1)\); on the fold
  (\(a, b \in \{1, \dots, 7\}\)) the frame of the fold; in the tensors
  of the fold written with its coordinates \(x^a\) (its metric
  \(\hat g_{ab}\) and the tensors built from it) a coordinate index of
  the fold, \(a \in \{1, \dots, 7\}\); a repeated index is summed.
  Status: STANDARD.
\item
  \(b\): A frame index (as \(a\)). Status: STANDARD.
\item
  \(\mathrm{d}\): The differential: \(\mathrm{d}\) in front of a
  quantity is an infinitesimal change of it (as \(\mathrm{d}\sigma\), a
  small step of \(\sigma\)); the quotient of two such changes is a
  derivative; \(\mathrm{d}\) with an exponent, in front of several
  variables, is their volume element, the product of their
  differentials. Status: STANDARD.
\item
  \(\sigma\): The proper coordinate of the hidden direction:
  \(\mathrm{d}\sigma = \lvert\cot z\rvert\,\mathrm{d}x_0\),
  \(\sigma = \mathrm{sgn}(\pi/2 - z)\ln(\sin z)/(6H)\) on one period
  \(0 < z < \pi\); \(\sigma = 0\) is the fold \(z = \pi/2\),
  \(\sigma \to \mp\infty\) the tips \(z \to 0, \pi\); one period of
  \(x_0\) is the whole line \(-\infty < \sigma < \infty\). Unit: length.
  Status: DERIVED.
\item
  \(\mathcal{Z}\): The fold factor,
  \(\mathcal{Z}(x_4) = V_{\mathrm{t}}(\mathrm{e}^{-a_4}b_4)^3/V_*\): the
  factor by which the size of the extra times multiplies every term of
  the fold and the kinetic terms of the massless modes of the fields of
  the bulk; \(\mathcal{Z} = 1\) when the extra times are frozen at their
  late value. Unit: dimensionless. Value: 1 in the late universe.
  Status: DERIVED.
\item
  \(V_*\): The physical volume of the extra times \(x_5, x_6, x_7\) on
  the fold in the late universe,
  \(V_* = V_{\mathrm{t}}(\mathrm{e}^{-a_4}b_4)^3\) (with
  \(V_{\mathrm{t}} = \int\mathrm{d}x_5\,\mathrm{d}x_6\,\mathrm{d}x_7\)
  the coordinate volume of the extra times) at a time when
  \(\mathrm{e}^{-a_4}b_4\) is constant; it converts the
  eight-dimensional normalisation into the four-dimensional one. Unit:
  length\(^3\) (GeV\(^{-3}\)). Status: ASSUMED.
\item
  \(\mathcal{T}_\chi{}^a{}_b\): The four-dimensional energy-momentum
  tensor of the gas of the particles \(\chi\) of dirac16complex (its
  long-lived states at the fold), per proper volume of ordinary space at
  the fold (energy density \(\rho_\chi\)). Unit: GeV\(^4\). Status:
  DERIVED.
\item
  \(H\): The author's constant (an inverse length) that fixes the period
  \(\pi/(6H)\) of \(x_0\), the warp factor
  \(\mathrm{e}^{-2H\lvert\sigma\rvert}\) of the metric, the cosmological
  constant \(\Lambda = -21H^2\) and the tension
  \(\kappa\lambda_{\mathrm{F}} = 12H\) of the fold; in UM-4 a free
  parameter, bounded from below by the propagation of the gauge fields
  in the bulk. Unit: GeV (1/length). Value: \(H \gtrsim 2.5\) TeV (from
  the invisible width of the Z boson). Status: AUTHOR.
\item
  \(q_{\mathrm{b}}\): The inverse width of a state of dirac16complex at
  the fold: inside the turning point its profile falls as
  \(\mathrm{e}^{-q_{\mathrm{b}}\lvert\sigma\rvert}\). Unit: GeV. Status:
  DERIVED.
\item
  \(S_\Psi\): The action of dirac16complex in the bulk (kinetic term
  with the canonical spin connection, mass term, contact
  self-interaction). Unit: dimensionless. Status: RECORD.
\end{itemize}

The particles \(\chi\) live in the bulk, in a layer of the width
\(1/q_{\mathrm{b}}\) around the fold. The junction condition of the fold
integrates the bulk field equations across the fold: the jump of the
extrinsic curvature across the fold equals \(\kappa\) times the
energy-momentum tensor of everything that lives on the fold. If the
layer is thin compared with the scale \(1/H\) on which the metric
varies, integrating across the whole layer gives that jump the
additional source A-5, where \(\mathcal{T}_\chi\) is the
four-dimensional energy-momentum tensor of the gas of \(\chi\) (per
proper volume of ordinary space) and \(\mathcal{Z}V_*\) the proper
volume of the extra times at the fold (\(\mathcal{Z} = 1\) with frozen
extra times). This is exactly the form in which the matter of the fold
enters (the energy-momentum tensor of the Standard Model divided by the
volume of the extra times): the particles \(\chi\) gravitate on the fold
like matter of the fold, with the same Newton constant, and enter the
Friedmann equations FE-59 as \(\rho_\chi\). In UM-4 the bulk is a vacuum
of Einstein gravity with linear gravitons, so the response of the bulk
to the part of the source away from \(\sigma = 0\) is linear in it, and
the corrections to A-5 are of the relative order \(H/q_{\mathrm{b}}\),
at most \(2\times10^{-3}\) for the long-lived particles
(\(m_\Psi/H > 1.1\times10^3\) with \(q_{\mathrm{b}} = m_\Psi/2\), less
for weaker binding) (in UM-3, whose static bulk was a degenerate vacuum
without a linear part, their size was not known). Status: DERIVED (the
integration across a thin layer), ESTIMATE (the size of the
corrections).

\hypertarget{the-equation-of-state-of-the-kohn-sham-states}{%
\subsection{The equation of state of the Kohn-Sham
states}\label{the-equation-of-state-of-the-kohn-sham-states}}

\[
\begin{aligned}
&\Psi = \mathrm{e}^{\frac72H\lvert\sigma\rvert}\,\chi_o(\sigma)\,\mathrm{e}^{-\mathrm{i}\omega x_4},\qquad
 \chi_o' = \gamma^{(0)}\Big(\mathrm{i}\omega\,\mathrm{e}^{H\lvert\sigma\rvert}\gamma^{(4)} + m_\Psi + \mathcal{M}(\sigma)\Big)\chi_o\qquad(\sigma \ne 0),\\
&\mathcal{M}(\sigma) = G_{\mathrm{KS}}\,\mathrm{e}^{7H\lvert\sigma\rvert}\Big(\hat\rho(\sigma)\,\mathbf{1} - \hat{\mathcal{N}}(\sigma)\Big),\qquad
 \hat\rho = \frac14\sum_{o=1}^{4}\bar\chi_o\chi_o,\qquad \hat{\mathcal{N}} = \frac14\sum_{o=1}^{4}\chi_o\bar\chi_o,\\
&\chi_o(0^+) = \exp\Big(\gamma^{(0)}\big(\mu_{\mathrm{L}}P_{\mathrm{L}} + \mu_{\mathrm{R}}P_{\mathrm{R}}\big)\Big)\chi_o(0^-),\qquad
 \int\mathrm{d}\sigma\,\mathrm{e}^{H\lvert\sigma\rvert}\,\chi_o^\dagger\mathfrak{B}\chi_p = -\delta_{op},\qquad G_{\mathrm{KS}} = \frac{\lambda_\Psi\,n_\chi}{V_*}
\end{aligned} \tag{S-7}
\] \textbf{Symbols in (S-7).}

\begin{itemize}
\tightlist
\item
  \(\Psi,\ \bar\Psi\): dirac16complex: the author's 16-component complex
  spinor field of eight dimensions;
  \(\bar\Psi = \Psi^\dagger\sigma_{16}\) is its Dirac adjoint; in a term
  of the fold, \(\Psi\) is its value at \(\sigma = 0\). Unit:
  GeV\(^{7/2}\). Status: AUTHOR.
\item
  \(\mathrm{e}\): Euler's number, the base of the natural logarithm;
  \(\mathrm{e}^{X}\) is the exponential of \(X\). Unit: dimensionless.
  Value: 2.718281828. Status: STANDARD.
\item
  \(H\): The author's constant (an inverse length) that fixes the period
  \(\pi/(6H)\) of \(x_0\), the warp factor
  \(\mathrm{e}^{-2H\lvert\sigma\rvert}\) of the metric, the cosmological
  constant \(\Lambda = -21H^2\) and the tension
  \(\kappa\lambda_{\mathrm{F}} = 12H\) of the fold; in UM-4 a free
  parameter, bounded from below by the propagation of the gauge fields
  in the bulk. Unit: GeV (1/length). Value: \(H \gtrsim 2.5\) TeV (from
  the invisible width of the Z boson). Status: AUTHOR.
\item
  \(\sigma\): The proper coordinate of the hidden direction:
  \(\mathrm{d}\sigma = \lvert\cot z\rvert\,\mathrm{d}x_0\),
  \(\sigma = \mathrm{sgn}(\pi/2 - z)\ln(\sin z)/(6H)\) on one period
  \(0 < z < \pi\); \(\sigma = 0\) is the fold \(z = \pi/2\),
  \(\sigma \to \mp\infty\) the tips \(z \to 0, \pi\); one period of
  \(x_0\) is the whole line \(-\infty < \sigma < \infty\). Unit: length.
  Status: DERIVED.
\item
  \(\chi_o(\sigma),\ \bar\chi_o\): The profile of the occupied Kohn-Sham
  orbital \(o\) (\(o = 1, \dots, 4\), the four positive-energy states of
  the level at the fold) of a homogeneous gas of the particles \(\chi\)
  at rest on the static block: a 16-component spinor function of
  \(\sigma\); the orbital of each particle is
  \(\mathrm{e}^{\frac72H\lvert\sigma\rvert}\chi_o(\sigma)\mathrm{e}^{-\mathrm{i}\omega x_4}\)
  divided by the square root of the volume of ordinary space and of
  \(V_*\); normalised to the Krein charge \(-1\) with the measure of
  UM-4,
  \(\int\mathrm{d}\sigma\,\mathrm{e}^{H\lvert\sigma\rvert}\chi_o^\dagger\mathfrak{B}\chi_o = -1\);
  \(\bar\chi_o = \chi_o^\dagger\sigma_{16}\) is its Dirac adjoint. Unit:
  GeV\(^{1/2}\). Status: DERIVED.
\item
  \(o\): The label of a single-particle state of dirac16complex: a
  Kohn-Sham orbital, or a mode of the quantised field; a sum over \(o\)
  runs over the orbitals or the modes. Status: STANDARD.
\item
  \(p\): The Lovelock order \(p = 1, 2, 3\) (Einstein, Gauss-Bonnet,
  cubic), as an index and as a number (UM-4 keeps only \(p = 1\)); in
  the Kohn-Sham equations also a second label of a Kohn-Sham orbital,
  with the Kronecker delta of the two labels. Status: STANDARD.
\item
  \(\mathrm{i}\): The imaginary unit, \(\mathrm{i}^2 = -1\). Unit:
  dimensionless. Status: STANDARD.
\item
  \(\omega\): A frequency: the energy of a mode
  \(\propto \mathrm{e}^{-\mathrm{i}\omega x_4}\) (\(\omega_o\) of the
  orbital \(o\)). Unit: GeV. Status: STANDARD.
\item
  \(x_4\): The time; a prime on a function of \(x_4\) alone is
  \(\mathrm{d}/\mathrm{d}x_4\). Unit: length (= time, \(c = 1\)).
  Status: AUTHOR.
\item
  \({}'\): The derivative of a function of the time \(x_4\) alone with
  respect to \(x_4\) (\(a_4' = \mathrm{d}a_4/\mathrm{d}x_4\); two primes
  the second derivative); on a function of \(\sigma\) alone (a profile)
  the derivative with respect to \(\sigma\). Status: STANDARD.
\item
  \(\gamma^M,\ \gamma^{(a)}\): The Dirac matrices of dirac16complex
  along the coordinates, \(\gamma^M = \mathbf{e}_a{}^MT^a\),
  \(\gamma_M = g_{MN}\gamma^N\); \(\gamma^{(a)} = T^a\) along the frame
  direction \((a)\) (\(\gamma^{(0)}\) along \(\sigma\), \(\gamma^{(4)}\)
  along the time). Unit: dimensionless. Status: AUTHOR.
\item
  \(m_\Psi\): The mass parameter of dirac16complex in the bulk. Unit:
  GeV. Status: RECORD.
\item
  \(\mathcal{M}(\sigma)\): The mean field of the self-interaction in the
  Kohn-Sham equation of the gas: the \(16 \times 16\) matrix
  \(\lambda_\Psi(n_{\mathrm{s}}\mathbf{1} - \mathcal{N})\) at each
  \(\sigma\), which adds to the mass \(m_\Psi\); positive (repulsive)
  for \(\lambda_\Psi > 0\). Unit: GeV. Status: DERIVED.
\item
  \(G_{\mathrm{KS}}\): The dimensionless strength of the
  self-interaction of the gas of particles \(\chi\) in the Kohn-Sham
  states: the self-coupling \(\lambda_\Psi\) times the number density
  \(n_\chi\) of the particles on the fold divided by the volume \(V_*\)
  of the extra times, \(G_{\mathrm{KS}} = \lambda_\Psi n_\chi/V_*\) (in
  units of \(m_\Psi\) when numbers are given). Unit: dimensionless.
  Value: \(10^{-6}\) to 0.5 in the solutions. Status: DERIVED.
\item
  \(\hat\rho(\sigma)\): The scalar density of the occupied orbitals,
  \(\hat\rho = \frac14\sum_o\bar\chi_o\chi_o\) (per particle); the
  scalar density of the gas is
  \(n_{\mathrm{s}} = (n_\chi/V_*)\,\mathrm{e}^{7H\lvert\sigma\rvert}\hat\rho\).
  Unit: GeV. Status: DERIVED.
\item
  \(\mathbf{1},\ \mathbf{1}_8\): The identity matrix of the size fixed
  by the context (\(16 \times 16\) on dirac16complex, \(4 \times 4\) on
  a four-dimensional spinor, \(3 \times 3\) in generation or flavour
  space); \(\mathbf{1}_8\) is the \(8 \times 8\) identity matrix. Unit:
  dimensionless. Status: STANDARD.
\item
  \(\hat{\mathcal{N}}(\sigma)\): The one-body density matrix of the
  occupied orbitals,
  \(\hat{\mathcal{N}} = \frac14\sum_o\chi_o\bar\chi_o\) (a
  \(16 \times 16\) matrix at each \(\sigma\), per particle); the density
  matrix of the gas is
  \(\mathcal{N} = (n_\chi/V_*)\,\mathrm{e}^{7H\lvert\sigma\rvert}\hat{\mathcal{N}}\).
  Unit: GeV. Status: DERIVED.
\item
  \(\mu_X,\ \mu_{\mathrm{L}},\ \mu_{\mathrm{R}}\): The dimensionless
  strengths of the fold for the left-handed and right-handed halves of
  dirac16complex, set by the mean fields of the Standard Model at the
  fold: the jump of the profile across the fold is the exponential of
  \(\gamma^{(0)}(\mu_{\mathrm{L}}P_{\mathrm{L}} + \mu_{\mathrm{R}}P_{\mathrm{R}})\);
  negative values attract. Unit: dimensionless. Status: DERIVED.
\item
  \(P_{\mathrm{L}}\): The author's left-handed chiral projector,
  \(P_{\mathrm{L}} = (\mathbf{1} - \Gamma_8)/2\). Unit: dimensionless.
  Status: AUTHOR.
\item
  \(P_{\mathrm{R}}\): The author's right-handed chiral projector,
  \(P_{\mathrm{R}} = (\mathbf{1} + \Gamma_8)/2\). Unit: dimensionless.
  Status: AUTHOR.
\item
  \(\mathrm{d}\): The differential: \(\mathrm{d}\) in front of a
  quantity is an infinitesimal change of it (as \(\mathrm{d}\sigma\), a
  small step of \(\sigma\)); the quotient of two such changes is a
  derivative; \(\mathrm{d}\) with an exponent, in front of several
  variables, is their volume element, the product of their
  differentials. Status: STANDARD.
\item
  \({}^\dagger\): Hermitian conjugation: the complex conjugate of the
  transpose of a matrix or spinor; for an operator on the space of
  states, its adjoint with respect to the scalar product of that space
  (for the quantised dirac16complex the indefinite Krein scalar
  product). Status: STANDARD.
\item
  \(\mathfrak{B}\): The Krein operator,
  \(\mathfrak{B} = \mathrm{i}\sigma_{16}\gamma^{(4)}\): a Hermitian
  involution with the signature (8, 8); \(\Psi^\dagger\mathfrak{B}\Psi\)
  is the charge density. Unit: dimensionless. Status: DERIVED.
\item
  \(\delta\): The Kronecker delta: with two indices \(\delta^M_N = 1\)
  if \(M = N\) and 0 otherwise (likewise \(\delta_{X\mathrm{L}}\) is 1
  if the label \(X\) of a half of dirac16complex is L and 0 otherwise);
  with two index strings of equal length the generalised Kronecker delta
  \(\delta^{M_1M_2\cdots}_{N_1N_2\cdots}\), the determinant of the
  square matrix of the Kronecker deltas \(\delta^{M_a}_{N_b}\) (\(+1\)
  or \(-1\) if the upper indices are an even or odd permutation of the
  lower ones, 0 otherwise). Status: STANDARD.
\item
  \(\lambda_\Psi\): The coupling of the contact self-interaction of
  dirac16complex. Unit: GeV\(^{-6}\). Status: RECORD.
\item
  \(n_\chi\): The number density of the dark-matter particles \(\chi\)
  (particles plus antiparticles). Unit: GeV\(^3\). Status: DERIVED.
\item
  \(V_*\): The physical volume of the extra times \(x_5, x_6, x_7\) on
  the fold in the late universe,
  \(V_* = V_{\mathrm{t}}(\mathrm{e}^{-a_4}b_4)^3\) (with
  \(V_{\mathrm{t}} = \int\mathrm{d}x_5\,\mathrm{d}x_6\,\mathrm{d}x_7\)
  the coordinate volume of the extra times) at a time when
  \(\mathrm{e}^{-a_4}b_4\) is constant; it converts the
  eight-dimensional normalisation into the four-dimensional one. Unit:
  length\(^3\) (GeV\(^{-3}\)). Status: ASSUMED.
\item
  \(\mathcal{N}\): The one-body density matrix of dirac16complex,
  \(\mathcal{N} = \sum_o\mathsf{n}_o\varphi_o\bar\varphi_o\) (a
  \(16 \times 16\) matrix at each point). Unit: GeV\(^7\). Status:
  DERIVED.
\item
  \(n_{\mathrm{s}}\): The scalar density of dirac16complex in the
  Kohn-Sham state,
  \(n_{\mathrm{s}} = \langle\bar\Psi\Psi\rangle = \mathrm{tr}\,\mathcal{N}\).
  Unit: GeV\(^7\). Status: DERIVED.
\item
  \(\varphi_o,\ \bar\varphi_o\): A Kohn-Sham orbital of dirac16complex
  (a single-particle solution of the Kohn-Sham equation: the Dirac
  equation of dirac16complex with the mean field of the gas added to its
  mass) and its Dirac adjoint. Unit: GeV\(^{7/2}\). Status: RECORD.
\item
  \(\mathsf{n}_o\): The occupation of the Kohn-Sham orbital \(o\)
  (between 0 and 1). Unit: dimensionless. Status: RECORD.
\end{itemize}

\[
\begin{aligned}
&\mathsf{E}_{\mathrm{int}} = \frac{G_{\mathrm{KS}}}{2}\int\mathrm{d}\sigma\,\mathrm{e}^{7H\lvert\sigma\rvert}\Big(\hat\rho^2 - \mathrm{tr}\,\hat{\mathcal{N}}^2\Big),\qquad
 \epsilon_{\mathrm{p}} = \omega - \mathsf{E}_{\mathrm{int}},\qquad P_{\mathrm{KS}} = q_\chi = n_\chi\,\mathsf{E}_{\mathrm{int}},\\
&\Delta\omega = G_{\mathrm{KS}}\int\mathrm{d}\sigma\,\mathrm{e}^{7H\lvert\sigma\rvert}\Big(\hat\rho^2 - \mathrm{tr}\,\hat{\mathcal{N}}^2\Big)\Big\vert_{G_{\mathrm{KS}} = 0}\qquad\Big(q_{\mathrm{b}} > \frac74H\Big)
\end{aligned} \tag{S-9}
\] \textbf{Symbols in (S-9).}

\begin{itemize}
\tightlist
\item
  \(\mathsf{E}_{\mathrm{int}}\): The interaction energy per particle of
  the gas of \(\chi\): the mean self-interaction energy
  \(\frac{\lambda_\Psi}{2}(n_{\mathrm{s}}^2 - \mathrm{tr}\,\mathcal{N}^2)\)
  integrated across the hidden direction, per particle. Unit: GeV.
  Status: DERIVED.
\item
  \(G_{\mathrm{KS}}\): The dimensionless strength of the
  self-interaction of the gas of particles \(\chi\) in the Kohn-Sham
  states: the self-coupling \(\lambda_\Psi\) times the number density
  \(n_\chi\) of the particles on the fold divided by the volume \(V_*\)
  of the extra times, \(G_{\mathrm{KS}} = \lambda_\Psi n_\chi/V_*\) (in
  units of \(m_\Psi\) when numbers are given). Unit: dimensionless.
  Value: \(10^{-6}\) to 0.5 in the solutions. Status: DERIVED.
\item
  \(\mathrm{d}\): The differential: \(\mathrm{d}\) in front of a
  quantity is an infinitesimal change of it (as \(\mathrm{d}\sigma\), a
  small step of \(\sigma\)); the quotient of two such changes is a
  derivative; \(\mathrm{d}\) with an exponent, in front of several
  variables, is their volume element, the product of their
  differentials. Status: STANDARD.
\item
  \(\sigma\): The proper coordinate of the hidden direction:
  \(\mathrm{d}\sigma = \lvert\cot z\rvert\,\mathrm{d}x_0\),
  \(\sigma = \mathrm{sgn}(\pi/2 - z)\ln(\sin z)/(6H)\) on one period
  \(0 < z < \pi\); \(\sigma = 0\) is the fold \(z = \pi/2\),
  \(\sigma \to \mp\infty\) the tips \(z \to 0, \pi\); one period of
  \(x_0\) is the whole line \(-\infty < \sigma < \infty\). Unit: length.
  Status: DERIVED.
\item
  \(\mathrm{e}\): Euler's number, the base of the natural logarithm;
  \(\mathrm{e}^{X}\) is the exponential of \(X\). Unit: dimensionless.
  Value: 2.718281828. Status: STANDARD.
\item
  \(H\): The author's constant (an inverse length) that fixes the period
  \(\pi/(6H)\) of \(x_0\), the warp factor
  \(\mathrm{e}^{-2H\lvert\sigma\rvert}\) of the metric, the cosmological
  constant \(\Lambda = -21H^2\) and the tension
  \(\kappa\lambda_{\mathrm{F}} = 12H\) of the fold; in UM-4 a free
  parameter, bounded from below by the propagation of the gauge fields
  in the bulk. Unit: GeV (1/length). Value: \(H \gtrsim 2.5\) TeV (from
  the invisible width of the Z boson). Status: AUTHOR.
\item
  \(\hat\rho(\sigma)\): The scalar density of the occupied orbitals,
  \(\hat\rho = \frac14\sum_o\bar\chi_o\chi_o\) (per particle); the
  scalar density of the gas is
  \(n_{\mathrm{s}} = (n_\chi/V_*)\,\mathrm{e}^{7H\lvert\sigma\rvert}\hat\rho\).
  Unit: GeV. Status: DERIVED.
\item
  \(\mathrm{tr}\): The trace of a matrix: the sum of its diagonal
  entries. Status: STANDARD.
\item
  \(\hat{\mathcal{N}}(\sigma)\): The one-body density matrix of the
  occupied orbitals,
  \(\hat{\mathcal{N}} = \frac14\sum_o\chi_o\bar\chi_o\) (a
  \(16 \times 16\) matrix at each \(\sigma\), per particle); the density
  matrix of the gas is
  \(\mathcal{N} = (n_\chi/V_*)\,\mathrm{e}^{7H\lvert\sigma\rvert}\hat{\mathcal{N}}\).
  Unit: GeV. Status: DERIVED.
\item
  \(\epsilon_{\mathrm{p}}\): The energy per particle of the gas of
  \(\chi\): the orbital energy minus the interaction energy per
  particle, which the orbital energy counts twice. Unit: GeV. Status:
  DERIVED.
\item
  \(\omega\): A frequency: the energy of a mode
  \(\propto \mathrm{e}^{-\mathrm{i}\omega x_4}\) (\(\omega_o\) of the
  orbital \(o\)). Unit: GeV. Status: STANDARD.
\item
  \(P_{\mathrm{KS}}\): The pressure of the gas of \(\chi\) along
  ordinary space in the Kohn-Sham state: the particles are at rest, so
  it is the mean self-interaction,
  \(P_{\mathrm{KS}} = n_\chi\mathsf{E}_{\mathrm{int}}\). Unit:
  GeV\(^4\). Status: DERIVED.
\item
  \(q_\chi\): The pressure of the particles \(\chi\) of dirac16complex
  along the extra times, integrated across the fold (their on-shell
  Lagrangian, the self-interaction energy per volume; zero for a dilute
  gas). Unit: GeV\(^4\). Status: DERIVED.
\item
  \(n_\chi\): The number density of the dark-matter particles \(\chi\)
  (particles plus antiparticles). Unit: GeV\(^3\). Status: DERIVED.
\item
  \(\Delta\): A change, difference or small variation of the quantity
  that follows (for example an interval of the time, a perturbation of a
  density, a difference of two squared masses, or an addition to an
  effective number of species). Status: STANDARD.
\item
  \(q_{\mathrm{b}}\): The inverse width of a state of dirac16complex at
  the fold: inside the turning point its profile falls as
  \(\mathrm{e}^{-q_{\mathrm{b}}\lvert\sigma\rvert}\). Unit: GeV. Status:
  DERIVED.
\end{itemize}

The Kohn-Sham equations of a homogeneous gas of the particles \(\chi\)
at rest on the fold (S-7) and its energy and pressure (S-9). Each
particle occupies one of the four positive-energy states of the level at
the fold. The profiles \(\chi_o\) obey the Dirac equation along the
hidden direction with the warped frequency
\(\omega\mathrm{e}^{H\lvert\sigma\rvert}\) and the mean field
\(\mathcal{M}\) of the self-interaction added to the mass, where
\(\hat\rho\) and \(\hat{\mathcal{N}}\) are the scalar density and the
density matrix of the profiles and \(\mathrm{e}^{7H\lvert\sigma\rvert}\)
turns them into proper densities; the Standard Model acts through the
matching at the fold, and \(G_{\mathrm{KS}} = \lambda_\Psi n_\chi/V_*\)
is the dimensionless strength of the self-interaction (the self-coupling
\(\lambda_\Psi\) times the number density \(n_\chi\) of the particles on
the fold, divided by the volume \(V_*\) of the extra times).
\(\mathsf{E}_{\mathrm{int}}\) is the interaction energy per particle,
\(\epsilon_{\mathrm{p}}\) the energy per particle (the orbital energy
counts the interaction twice) and \(P_{\mathrm{KS}}\) the pressure: the
particles are at rest and do not depend on the extra times, so only the
self-interaction contributes, along ordinary space and along the extra
times alike. Status: DERIVED (from the field equation of dirac16complex
with its self-interaction), RECORD (the Kohn-Sham treatment).

The Kohn-Sham equations were solved (in units of \(m_\Psi\)) deep in the
regime of long-lived particles, \(H = 10^{-4}\) (\(m_\Psi/H = 10^4\)),
with the Higgs portal of the vacuum (\(\mu_{\mathrm{L}} = -0.6\),
\(\mu_{\mathrm{R}} = 0\), so \(q_{\mathrm{b}} = 0.3 = 3000H\)), for
\(G_{\mathrm{KS}}\) from \(10^{-6}\) to 0.5. The self-interaction is
repulsive for \(\lambda_\Psi > 0\), the level rises, and the pressure is
positive: the equation of state
\(w_\chi = P_{\mathrm{KS}}/(n_\chi\epsilon_{\mathrm{p}})\) is
\(5.36\times10^{-8}\), \(5.36\times10^{-6}\), \(5.32\times10^{-4}\) and
\(5.00\times10^{-3}\) at \(G_{\mathrm{KS}} = 10^{-6}\), \(10^{-4}\),
\(10^{-2}\) and \(10^{-1}\); in the dilute limit the shift of the level
is linear in \(G_{\mathrm{KS}}\) (the exponent 1.0000) and equals
first-order perturbation theory; far from the fold the proper density
decays, so the gas stays at the fold. The repulsion also raises the
level towards \(m_\Psi\) and lowers the barrier through which the
particles tunnel into the bulk: at \(m_\Psi/H = 10^4\) the gas lives
longer than the universe only for \(G_{\mathrm{KS}} < 0.147\). At a
weakly bound point (\(q_{\mathrm{b}} = 0.03\,m_\Psi\),
\(H = 10^{-7}m_\Psi\)) the self-consistent shift approaches first-order
perturbation theory as \(G_{\mathrm{KS}} \to 0\), and
\(w_\chi = 5.6\times10^{-7}\), \(5.6\times10^{-6}\) and
\(5.1\times10^{-5}\) at \(G_{\mathrm{KS}} = 10^{-4}\), \(10^{-3}\) and
\(10^{-2}\). Status: DERIVED (\texttt{um4\_kohnsham.py}, K2 to K5, K7 to
K9).

How dilute is the gas of the universe? With the self-coupling of the
natural size \(\lambda_\Psi = m_\Psi^{-6}\) and extra times of the size
\(1/m_\Psi\) (\(V_* = m_\Psi^{-3}\)),
\(G_{\mathrm{KS}} = n_\chi/m_\Psi^3\). For the lightest long-lived
particles (\(\omega_{\mathrm{b}} = 2.43\times10^6\) GeV with
\(q_{\mathrm{b}} = m_\Psi/2\), so \(m_\Psi = 2.81\times10^6\) GeV) and
the density of dark matter near the Sun (0.4 GeV/cm\(^3\)) this is
\(G_{\mathrm{KS}} = 5.7\times10^{-68}\), and \(1.8\times10^{-73}\) for
the mean density of the universe today; extra times of the Planck size
(\(V_* = M_{\mathrm{Pl}}^{-3}\)) give \(4.7\times10^{-30}\) and
\(1.5\times10^{-35}\). The cosmic gas of \(\chi\) is therefore deep in
the dilute limit, where \(w_\chi = 0.0536\,G_{\mathrm{KS}}\)
(\(3\times10^{-69}\) near the Sun for the natural sizes), and far below
the coupling at which the self-interaction would shorten its life: it is
cold dark matter. Status: ESTIMATE (\texttt{um4\_kohnsham.py}, K10).

\hypertarget{the-relic-density-direct-detection-and-the-window}{%
\subsection{The relic density, direct detection and the
window}\label{the-relic-density-direct-detection-and-the-window}}

\[
\frac{\mathrm{d}n_\chi}{\mathrm{d}x_4} + 3H_4\,n_\chi = -\langle\sigma_{\mathrm{ann}}\mathsf{v}\rangle\big(n_\chi^2 - n_{\chi,\mathrm{eq}}^2\big),\qquad
n_{\chi,\mathrm{eq}} = g_\chi\Big(\frac{\omega_{\mathrm{b}}\mathsf{T}}{2\pi}\Big)^{3/2}\mathrm{e}^{-\omega_{\mathrm{b}}/\mathsf{T}},\qquad
\Omega_\chi h_{100}^2 \simeq \frac{1.07\times10^9\ \text{GeV}^{-1}\,x_{\mathrm{f}}}{\sqrt{g_*}\,M_{\mathrm{Pl}}\,\langle\sigma_{\mathrm{ann}}\mathsf{v}\rangle} \tag{FE-61}
\] \textbf{Symbols in (FE-61).}

\begin{itemize}
\tightlist
\item
  \(\mathrm{d}\): The differential: \(\mathrm{d}\) in front of a
  quantity is an infinitesimal change of it (as \(\mathrm{d}\sigma\), a
  small step of \(\sigma\)); the quotient of two such changes is a
  derivative; \(\mathrm{d}\) with an exponent, in front of several
  variables, is their volume element, the product of their
  differentials. Status: STANDARD.
\item
  \(n_\chi\): The number density of the dark-matter particles \(\chi\)
  (particles plus antiparticles). Unit: GeV\(^3\). Status: DERIVED.
\item
  \(x_4\): The time; a prime on a function of \(x_4\) alone is
  \(\mathrm{d}/\mathrm{d}x_4\). Unit: length (= time, \(c = 1\)).
  Status: AUTHOR.
\item
  \(H_4\): The expansion rate (Hubble rate) of the fold, \(H_4 = a_4'\),
  as a function of the time or of the redshift; \(H_4^{(0)}\) its value
  today, the Hubble constant. Unit: GeV (1/time). Value:
  \(H_4^{(0)} = 100\,h_{100}\) km/s/Mpc with \(h_{100} = 0.674\) (Planck
  2018): 67.4 km/s/Mpc \(= 1.44\times10^{-42}\) GeV. Status: STANDARD.
\item
  \(\sigma_{\mathrm{ann}}\): The annihilation cross-section of the
  particles \(\chi\) into Standard-Model particles (through the portals)
  in the normalisation of the relic abundance, where \(n_\chi\) counts
  the four particle and the four antiparticle states of the level at the
  fold: \(\langle\sigma_{\mathrm{ann}}\mathsf{v}\rangle\) is a quarter
  of the annihilation rate \((\sigma\mathsf{v})_{\chi\bar\chi}\) of one
  particle \(\chi_a\) with one antiparticle \(\bar\chi_a\). Unit:
  GeV\(^{-2}\). Status: DERIVED.
\item
  \(\mathsf{v}\): The relative velocity of two annihilating particles;
  \(\langle\sigma_{\mathrm{ann}}\mathsf{v}\rangle\) is the thermal
  average of the product. Unit: dimensionless. Status: STANDARD.
\item
  \(\langle\cdots\rangle\): The mean (expectation) value of an operator
  in the state of the medium (the vacuum, a plasma, the Kohn-Sham
  state). Status: STANDARD.
\item
  \(n_{\chi,\mathrm{eq}}\): The number density of \(\chi\) in thermal
  equilibrium with the plasma (non-relativistic). Unit: GeV\(^3\).
  Status: STANDARD.
\item
  \(g_\chi\): The number of internal states of the level of \(\chi\) at
  the fold that are in thermal equilibrium. Unit: dimensionless. Value:
  8 (four particles and four antiparticles). Status: DERIVED.
\item
  \(\omega_{\mathrm{b}}\): The energy at zero momentum of the state of
  dirac16complex at the fold in the limit \(H \to 0\) (the bound state
  of UM-3): the particle \(\chi\) of UM-4 has the mass
  \(\omega_{\mathrm{r}} = \omega_{\mathrm{b}}(1 - H/(2q_{\mathrm{b}})) + \dots\),
  equal to \(\omega_{\mathrm{b}}\) up to the relative correction
  \(H/(2q_{\mathrm{b}})\). Unit: GeV. Status: DERIVED.
\item
  \(\mathsf{T}\): The temperature of the Standard-Model plasma (of the
  photons after the annihilation of electrons and positrons). Unit: GeV.
  Status: STANDARD.
\item
  \(\pi\): The number pi, the ratio of the circumference of a circle to
  its diameter. Unit: dimensionless. Value: 3.141592654. Status:
  STANDARD.
\item
  \(\mathrm{e}\): Euler's number, the base of the natural logarithm;
  \(\mathrm{e}^{X}\) is the exponential of \(X\). Unit: dimensionless.
  Value: 2.718281828. Status: STANDARD.
\item
  \(\Omega_\chi\): The density parameter of the dark matter today. Unit:
  dimensionless. Value: \(\Omega_\chi h_{100}^2 = 0.120\) (Planck 2018).
  Status: DERIVED.
\item
  \(h_{100}\): The present expansion rate of the fold in units of 100
  km/s/Mpc. Unit: dimensionless. Value: 0.674 (Planck 2018). Status:
  STANDARD.
\item
  \(x_{\mathrm{f}}\): The ratio of the mass of \(\chi\) to the
  temperature at which its annihilations stop (freeze-out). Unit:
  dimensionless. Value: 23 to 30 (for \(\omega_{\mathrm{b}}\) = 10 GeV
  to 10 TeV). Status: STANDARD.
\item
  \(g_*\): The effective number of relativistic degrees of freedom at
  the freeze-out (\(g_\rho\) there). Unit: dimensionless. Status:
  STANDARD.
\item
  \(M_{\mathrm{Pl}}\): The Planck mass,
  \(M_{\mathrm{Pl}} = G_{\mathrm{N}}^{-1/2}\). Unit: GeV. Value:
  \(1.221\times10^{19}\) GeV. Status: STANDARD.
\end{itemize}

\[
\sigma_{\mathrm{pair}}\mathsf{v} \le \frac{4\pi}{\omega_{\mathrm{b}}^2\,\mathsf{v}},\qquad \Big\langle\frac{1}{\mathsf{v}}\Big\rangle = \sqrt{\frac{\mathsf{x}}{\pi}}:\qquad
\Omega_\chi h_{100}^2 = 0.120\ \ \text{at}\ \ \omega_{\mathrm{b}} = 8.87\times10^4\ \text{GeV} \tag{S-20}
\] \textbf{Symbols in (S-20).}

\begin{itemize}
\tightlist
\item
  \(\sigma_{\mathrm{pair}}\): The cross-section of the annihilation of
  one particle \(\chi_a\) with one antiparticle \(\bar\chi_a\) of the
  same species, summed over the final states;
  \(\sigma_{\mathrm{pair}}\mathsf{v}\) its product with the relative
  velocity (a constant for s-wave annihilation). Unit: GeV\(^{-2}\).
  Status: DERIVED.
\item
  \(\mathsf{v}\): The relative velocity of two annihilating particles;
  \(\langle\sigma_{\mathrm{ann}}\mathsf{v}\rangle\) is the thermal
  average of the product. Unit: dimensionless. Status: STANDARD.
\item
  \(\omega_{\mathrm{b}}\): The energy at zero momentum of the state of
  dirac16complex at the fold in the limit \(H \to 0\) (the bound state
  of UM-3): the particle \(\chi\) of UM-4 has the mass
  \(\omega_{\mathrm{r}} = \omega_{\mathrm{b}}(1 - H/(2q_{\mathrm{b}})) + \dots\),
  equal to \(\omega_{\mathrm{b}}\) up to the relative correction
  \(H/(2q_{\mathrm{b}})\). Unit: GeV. Status: DERIVED.
\item
  \(\pi\): The number pi, the ratio of the circumference of a circle to
  its diameter. Unit: dimensionless. Value: 3.141592654. Status:
  STANDARD.
\item
  \(\langle\cdots\rangle\): The mean (expectation) value of an operator
  in the state of the medium (the vacuum, a plasma, the Kohn-Sham
  state). Status: STANDARD.
\item
  \(\mathsf{x}\): The ratio of the mass of \(\chi\) to the temperature,
  \(\mathsf{x} = \omega_{\mathrm{b}}/\mathsf{T}\) (the time variable of
  the freeze-out). Unit: dimensionless. Status: STANDARD.
\item
  \(\Omega_\chi\): The density parameter of the dark matter today. Unit:
  dimensionless. Value: \(\Omega_\chi h_{100}^2 = 0.120\) (Planck 2018).
  Status: DERIVED.
\item
  \(h_{100}\): The present expansion rate of the fold in units of 100
  km/s/Mpc. Unit: dimensionless. Value: 0.674 (Planck 2018). Status:
  STANDARD.
\end{itemize}

The relic density of a thermal relic (FE-61, solved numerically for
\(g_\chi = 8\): the four particle and the four antiparticle states of
the level at the fold) needs
\(\langle\sigma_{\mathrm{ann}}\mathsf{v}\rangle = 2.2\times10^{-26}\)
cm\(^3\)/s for every mass from 10 GeV to 10 TeV, that is
\((\sigma\mathsf{v})_{\chi\bar\chi} \approx 9\times10^{-26}\) cm\(^3\)/s
for the annihilation of one \(\chi_a\) with one \(\bar\chi_a\).
Partial-wave unitarity bounds the s-wave annihilation of one \(\chi\)
with one \(\bar\chi\) (S-20, with the thermal average of
\(1/\mathsf{v}\) over the Maxwell-Boltzmann distribution): with the
largest annihilation that unitarity allows, the relic density equals the
measured one at \(\omega_{\mathrm{b}} = 8.87\times10^4\) GeV, and
heavier particles would be overproduced by freeze-out (the classic bound
of K. Griest and M. Kamionkowski, Phys. Rev.~Lett. 64 (1990) 615). The
particles \(\chi\) of UM-4 live longer than the universe only above
\(2.43\times10^6\) GeV (S-6), 27 times more: they cannot be a thermal
relic, and their abundance needs a non-thermal origin. Status: DERIVED
(\texttt{um4\_thermal.py}, T2 and T7), OPEN (the origin of the
abundance).

\[
\begin{aligned}
&\mathfrak{g}_{\Phi} = \frac{\partial\omega_{\mathrm{b}}}{\partial\mu_{\mathrm{L}}}\,\frac{c_\Phi}{V_*},\qquad
 \mathfrak{g}_{G} = \sum_X\frac{\partial\omega_{\mathrm{b}}}{\partial\mu_X}\,\frac{c^{G}_X}{4V_*},\qquad
 \mathfrak{g}_{f} = \sum_X\frac{\partial\omega_{\mathrm{b}}}{\partial\mu_X}\,\frac{2\omega_{\mathrm{b}}\,c^{f}_X}{V_*},\qquad
 \frac{\partial\omega_{\mathrm{b}}}{\partial\mu_X} = -\frac{q_{\mathrm{b}}}{\omega_{\mathrm{b}}}\,\frac{\partial q_{\mathrm{b}}}{\partial\mu_X},\\
&\Delta\omega_{\mathrm{b}} = \mathfrak{g}_{\Phi}\,\Delta\langle\Phi^\dagger\Phi\rangle + \mathfrak{g}_{G}\,\Delta\langle G^2\rangle + \sum_f\mathfrak{g}_{f}\,\Delta n_f,\qquad
 \mathfrak{g}_{h\chi} = \mathfrak{g}_{\Phi}\,v = -\frac{m_\Psi^2\,\mu_{\mathrm{L}}^2}{2\,\omega_{\mathrm{b}}\,v}\quad(\mu_{\mathrm{R}} = 0)
\end{aligned} \tag{FE-47}
\] \textbf{Symbols in (FE-47).}

\begin{itemize}
\tightlist
\item
  \(\mathfrak{g}_{\Phi},\ \mathfrak{g}_{G},\ \mathfrak{g}_{f},\ \mathfrak{g}_{h\chi}\):
  The four-dimensional couplings of the particles \(\chi\) of
  dirac16complex to the Standard Model: to \(\Phi^\dagger\Phi\), to the
  gluon kinetic density, to the number densities of the fermions of kind
  \(f\), and the Yukawa coupling of the Higgs boson to \(\chi\). Unit:
  \(\mathfrak{g}_{\Phi}\): GeV\(^{-1}\); \(\mathfrak{g}_{G}\):
  GeV\(^{-3}\); \(\mathfrak{g}_{f}\): GeV\(^{-2}\);
  \(\mathfrak{g}_{h\chi}\): dimensionless. Status: DERIVED.
\item
  \(\partial\): A partial derivative: the derivative with respect to one
  variable at fixed values of the others (\(\partial_\sigma\) with
  respect to \(\sigma\); \(\partial_M\) and \(\partial_\mu\) with
  respect to the coordinate that carries the same index, of the eight
  dimensions or of the fold). Status: STANDARD.
\item
  \(\omega_{\mathrm{b}}\): The energy at zero momentum of the state of
  dirac16complex at the fold in the limit \(H \to 0\) (the bound state
  of UM-3): the particle \(\chi\) of UM-4 has the mass
  \(\omega_{\mathrm{r}} = \omega_{\mathrm{b}}(1 - H/(2q_{\mathrm{b}})) + \dots\),
  equal to \(\omega_{\mathrm{b}}\) up to the relative correction
  \(H/(2q_{\mathrm{b}})\). Unit: GeV. Status: DERIVED.
\item
  \(\mu_X,\ \mu_{\mathrm{L}},\ \mu_{\mathrm{R}}\): The dimensionless
  strengths of the fold for the left-handed and right-handed halves of
  dirac16complex, set by the mean fields of the Standard Model at the
  fold: the jump of the profile across the fold is the exponential of
  \(\gamma^{(0)}(\mu_{\mathrm{L}}P_{\mathrm{L}} + \mu_{\mathrm{R}}P_{\mathrm{R}})\);
  negative values attract. Unit: dimensionless. Status: DERIVED.
\item
  \(c_\Phi\): The Higgs-portal coupling of the left-handed half of
  dirac16complex (the right-handed half has none). Unit: GeV\(^{-5}\).
  Status: INVENTED.
\item
  \(V_*\): The physical volume of the extra times \(x_5, x_6, x_7\) on
  the fold in the late universe,
  \(V_* = V_{\mathrm{t}}(\mathrm{e}^{-a_4}b_4)^3\) (with
  \(V_{\mathrm{t}} = \int\mathrm{d}x_5\,\mathrm{d}x_6\,\mathrm{d}x_7\)
  the coordinate volume of the extra times) at a time when
  \(\mathrm{e}^{-a_4}b_4\) is constant; it converts the
  eight-dimensional normalisation into the four-dimensional one. Unit:
  length\(^3\) (GeV\(^{-3}\)). Status: ASSUMED.
\item
  \(c^{G}_X\): The gluon-kinetic portal coupling of the half \(X\) of
  dirac16complex. Unit: GeV\(^{-7}\). Status: INVENTED.
\item
  \(X\): The label of a half of dirac16complex:
  \(X \in \{\mathrm{L}, \mathrm{R}\}\) (left-handed or right-handed
  half); a sum over \(X\) runs over both. Status: STANDARD.
\item
  \(c^{f}_X\): The current-portal coupling of the half \(X\) of
  dirac16complex with the fermion-number current of the multiplets of
  kind \(f\) (flavour blind); \(c^{\ell}_X\) and \(c^{\nu}_X\) are those
  of the lepton doublets and of the right-handed neutrinos. Unit:
  GeV\(^{-7}\). Status: INVENTED.
\item
  \(f\): The label of a kind of Standard-Model multiplet:
  \(f \in \{Q, u, d, \ell, e, \nu\}\) (quark doublets, up-quark
  singlets, down-quark singlets, lepton doublets, charged-lepton
  singlets, right-handed neutrinos); a sum over \(f\) runs over the six.
  Status: STANDARD.
\item
  \(q_{\mathrm{b}}\): The inverse width of a state of dirac16complex at
  the fold: inside the turning point its profile falls as
  \(\mathrm{e}^{-q_{\mathrm{b}}\lvert\sigma\rvert}\). Unit: GeV. Status:
  DERIVED.
\item
  \(\Delta\): A change, difference or small variation of the quantity
  that follows (for example an interval of the time, a perturbation of a
  density, a difference of two squared masses, or an addition to an
  effective number of species). Status: STANDARD.
\item
  \(\Phi\): The Higgs doublet (two complex components, hypercharge
  \(1/2\)), a field of the fold. Unit: GeV. Status: AUTHOR.
\item
  \({}^\dagger\): Hermitian conjugation: the complex conjugate of the
  transpose of a matrix or spinor; for an operator on the space of
  states, its adjoint with respect to the scalar product of that space
  (for the quantised dirac16complex the indefinite Krein scalar
  product). Status: STANDARD.
\item
  \(\langle\cdots\rangle\): The mean (expectation) value of an operator
  in the state of the medium (the vacuum, a plasma, the Kohn-Sham
  state). Status: STANDARD.
\item
  \(\langle G^2\rangle\): The mean square of the gluon field strength,
  \(\langle G^A_{\mu\nu}G^{A\mu\nu}\rangle\) (the gluon condensate in
  the QCD vacuum, a thermal value in the plasma). Unit: GeV\(^4\).
  Status: STANDARD.
\item
  \(G^A_M,\ G^A_\mu,\ G^A_{MN}\): The gluon field of \(SU(3)_C\)
  (\(A = 1, \dots, 8\)): in the bulk \(G^A_M\), a field of eight
  dimensions in its geometric normalisation (the covariant derivative
  \(\partial - \mathrm{i}G\)), and on the fold the canonically
  normalised four-dimensional field \(G^A_\mu(x)\) of its massless mode,
  \(g_3G^A_\mu(x) = G^A_\mu(\sigma = 0, x)\); \(G^A_{MN}\) and
  \(G^A_{\mu\nu}\) their field strengths. Unit: GeV (field), GeV\(^2\)
  (field strength). Status: AUTHOR.
\item
  \(n_f\): The number density of the Standard-Model fermions of kind
  \(f\) (particles minus antiparticles, summed over generations and
  gauge components), the time component of the mean current
  \(\langle J_f^\mu\rangle\) in the rest frame of the medium. Unit:
  GeV\(^3\). Status: STANDARD.
\item
  \(v\): The vacuum expectation value of the lower component of the
  Higgs doublet, \(v^2 = -\mu^2/\lambda\). Unit: GeV. Value: 246.22 GeV
  \(= (\sqrt2\,G_{\mathrm{F}})^{-1/2}\) with the Fermi constant
  \(G_{\mathrm{F}} = 1.1663788\times10^{-5}\) GeV\(^{-2}\) of the muon
  lifetime (Particle Data Group). Status: STANDARD.
\item
  \(m_\Psi\): The mass parameter of dirac16complex in the bulk. Unit:
  GeV. Status: RECORD.
\end{itemize}

\[
\begin{aligned}
&\sigma_{\mathrm{pair}}\mathsf{v} = \sum_f\frac{\mathfrak{g}_f^2\,\omega_{\mathrm{b}}^2}{2\pi},\qquad
 \langle\sigma_{\mathrm{ann}}\mathsf{v}\rangle = \frac{\sigma_{\mathrm{pair}}\mathsf{v}}{4},\qquad
 \sigma_{\mathrm{SI}} = \frac{\mu_{\mathrm{red}}^2\,f_{\mathrm{N}}^2}{\pi},\qquad f_{\mathrm{N}} = 3\,\mathfrak{g}_{\mathrm{q}}\ \ (\text{quark portals}),\\
&f_{\mathrm{N}} = \mathfrak{g}_{\mathrm{p}} = \frac{\alpha_{\mathrm{em}}}{3\pi}\,\mathfrak{g}_{\mathrm{l}}\sum_{l}\ln\frac{\omega_{\mathrm{b}}^2}{\max\big(m_l^2,\ q_{\mathrm{t}}^2\big)}\ \ (\text{lepton portals}),\qquad
 f_{\mathrm{N}} = \frac{\mathfrak{g}_{h\chi}\,f_{\mathrm{H}}\,m_{\mathrm{N}}}{v\,m_h^2}\ \ (\text{Higgs portal})
\end{aligned} \tag{S-21}
\] \textbf{Symbols in (S-21).}

\begin{itemize}
\tightlist
\item
  \(\sigma_{\mathrm{pair}}\): The cross-section of the annihilation of
  one particle \(\chi_a\) with one antiparticle \(\bar\chi_a\) of the
  same species, summed over the final states;
  \(\sigma_{\mathrm{pair}}\mathsf{v}\) its product with the relative
  velocity (a constant for s-wave annihilation). Unit: GeV\(^{-2}\).
  Status: DERIVED.
\item
  \(\mathsf{v}\): The relative velocity of two annihilating particles;
  \(\langle\sigma_{\mathrm{ann}}\mathsf{v}\rangle\) is the thermal
  average of the product. Unit: dimensionless. Status: STANDARD.
\item
  \(\mathfrak{g}_{\Phi},\ \mathfrak{g}_{G},\ \mathfrak{g}_{f},\ \mathfrak{g}_{h\chi}\):
  The four-dimensional couplings of the particles \(\chi\) of
  dirac16complex to the Standard Model: to \(\Phi^\dagger\Phi\), to the
  gluon kinetic density, to the number densities of the fermions of kind
  \(f\), and the Yukawa coupling of the Higgs boson to \(\chi\). Unit:
  \(\mathfrak{g}_{\Phi}\): GeV\(^{-1}\); \(\mathfrak{g}_{G}\):
  GeV\(^{-3}\); \(\mathfrak{g}_{f}\): GeV\(^{-2}\);
  \(\mathfrak{g}_{h\chi}\): dimensionless. Status: DERIVED.
\item
  \(f\): The label of a kind of Standard-Model multiplet:
  \(f \in \{Q, u, d, \ell, e, \nu\}\) (quark doublets, up-quark
  singlets, down-quark singlets, lepton doublets, charged-lepton
  singlets, right-handed neutrinos); a sum over \(f\) runs over the six.
  Status: STANDARD.
\item
  \(\omega_{\mathrm{b}}\): The energy at zero momentum of the state of
  dirac16complex at the fold in the limit \(H \to 0\) (the bound state
  of UM-3): the particle \(\chi\) of UM-4 has the mass
  \(\omega_{\mathrm{r}} = \omega_{\mathrm{b}}(1 - H/(2q_{\mathrm{b}})) + \dots\),
  equal to \(\omega_{\mathrm{b}}\) up to the relative correction
  \(H/(2q_{\mathrm{b}})\). Unit: GeV. Status: DERIVED.
\item
  \(\pi\): The number pi, the ratio of the circumference of a circle to
  its diameter. Unit: dimensionless. Value: 3.141592654. Status:
  STANDARD.
\item
  \(\sigma_{\mathrm{ann}}\): The annihilation cross-section of the
  particles \(\chi\) into Standard-Model particles (through the portals)
  in the normalisation of the relic abundance, where \(n_\chi\) counts
  the four particle and the four antiparticle states of the level at the
  fold: \(\langle\sigma_{\mathrm{ann}}\mathsf{v}\rangle\) is a quarter
  of the annihilation rate \((\sigma\mathsf{v})_{\chi\bar\chi}\) of one
  particle \(\chi_a\) with one antiparticle \(\bar\chi_a\). Unit:
  GeV\(^{-2}\). Status: DERIVED.
\item
  \(\langle\cdots\rangle\): The mean (expectation) value of an operator
  in the state of the medium (the vacuum, a plasma, the Kohn-Sham
  state). Status: STANDARD.
\item
  \(\sigma_{\mathrm{SI}}\): The spin-independent cross-section of the
  scattering of \(\chi\) on a nucleon (normalised per nucleon, as the
  direct-detection limits are quoted). Unit: GeV\(^{-2}\) (cm\(^2\)).
  Status: ESTIMATE.
\item
  \(\mu_{\mathrm{red}}\): The reduced mass of \(\chi\) and a nucleon,
  \(\omega_{\mathrm{b}}m_{\mathrm{N}}/(\omega_{\mathrm{b}} + m_{\mathrm{N}})\).
  Unit: GeV. Status: STANDARD.
\item
  \(f_{\mathrm{N}}\): The effective coupling of \(\chi\) to a nucleon in
  the coherent scattering (a vector or scalar contact interaction).
  Unit: GeV\(^{-2}\). Status: ESTIMATE.
\item
  \(\mathfrak{g}_{\mathrm{q}},\ \mathfrak{g}_{\mathrm{l}},\ \mathfrak{g}_{\mathrm{p}}\):
  The current-portal couplings \(\mathfrak{g}_f\) taken equal for all
  chiral quark states (\(\mathfrak{g}_{\mathrm{q}}\)) or for the nine
  chiral lepton states that take part, the six of the lepton doublets
  and the three right-handed charged leptons
  (\(\mathfrak{g}_{\mathrm{l}}\)), and the coupling to the proton that
  the lepton loop induces (\(\mathfrak{g}_{\mathrm{p}}\)). Unit:
  GeV\(^{-2}\). Status: ESTIMATE.
\item
  \(\alpha_{\mathrm{em}}\): The fine-structure constant. Unit:
  dimensionless. Value: 1/137.036. Status: STANDARD.
\item
  \(l\): The multipole index and number (\(l = 0, 1, 2, \dots\)) of an
  expansion in Legendre polynomials of the direction of a photon or a
  neutrino relative to the wave vector, or of the CMB sky in spherical
  harmonics. Status: STANDARD.
\item
  \(m\): A mass: \(m_i\) the mass of the neutrino mass state \(i\);
  \(\Delta m_{ij}^2 = m_i^2 - m_j^2\) the squared-mass difference of the
  states \(i\) and \(j\) (\(\Delta m_{21}^2\), \(\Delta m_{31}^2\),
  \(\Delta m_{32}^2\)); \(\Delta m^2\) without labels the difference of
  the two states of a resonance, \(m_2^2 - m_1^2\), multiplied by
  \(v_{\mathrm{loc}}^2/v^2\) where this differs from 1; \(m\) without a
  label the mass of a mode of a field of the bulk as the fold sees it,
  \(m^2 = \omega^2 - k^2\); for the W and Z bosons their complex pole
  mass, \(m^2 = m_0^2 - \mathrm{i}m_0\Gamma_{\mathrm{bulk}}\) to the
  leading order; in the relic abundance of the particles \(\chi\) their
  mass. Unit: GeV. Status: STANDARD.
\item
  \(q_{\mathrm{t}}\): The momentum transfer of the scattering of
  \(\chi\) on a xenon nucleus in a detector. Unit: GeV. Value: 0.05 GeV.
  Status: STANDARD.
\item
  \(m_{\mathrm{N}}\): The mass of a nucleon. Unit: GeV. Value: 0.939
  GeV. Status: STANDARD.
\item
  \(v\): The vacuum expectation value of the lower component of the
  Higgs doublet, \(v^2 = -\mu^2/\lambda\). Unit: GeV. Value: 246.22 GeV
  \(= (\sqrt2\,G_{\mathrm{F}})^{-1/2}\) with the Fermi constant
  \(G_{\mathrm{F}} = 1.1663788\times10^{-5}\) GeV\(^{-2}\) of the muon
  lifetime (Particle Data Group). Status: STANDARD.
\item
  \(m_h\): The mass of the Higgs boson. Unit: GeV. Value: 125.25 GeV
  (Particle Data Group 2022). Status: STANDARD.
\item
  \(f_{\mathrm{H}}\): The nucleon matrix element of the Higgs coupling:
  the fraction of the nucleon mass carried by the quark masses and,
  through the trace anomaly, by the gluons. Unit: dimensionless. Value:
  0.30. Status: STANDARD.
\item
  \(G_{\mathrm{F}}\): The Fermi constant,
  \(G_{\mathrm{F}} = 1/(\sqrt2\,v^2)\) in the Standard-Model vacuum (its
  measured value); in the equations of the neutrino flavours it stands
  for the local Fermi constant of the ad hoc vacuum. Unit: GeV\(^{-2}\).
  Value: \(1.1664\times10^{-5}\) GeV\(^{-2}\). Status: STANDARD.
\end{itemize}

The couplings of the particles \(\chi\) to the Standard Model (FE-47)
and their annihilation and scattering on nucleons (S-21). The
Higgs-portal Yukawa coupling is fixed by the binding: with the Higgs
portal alone (\(\mu_{\mathrm{R}} = 0\), so
\(q_{\mathrm{b}} = -m_\Psi\mu_{\mathrm{L}}/2\) by FE-40),
\(\lvert\mathfrak{g}_{h\chi}\rvert = 2q_{\mathrm{b}}^2/(\omega_{\mathrm{b}}v)\).
The current portals are free (through the INVENTED couplings
\(c^{f}_X\)). Against the limit of the LUX-ZEPLIN experiment (4.2
tonne-years of xenon, Phys. Rev.~Lett. 135 (2025) 011802:
\(2.2\times10^{-48}\) cm\(^2\) at its minimum, 40 GeV, at 90 per cent
confidence, scaled in proportion to the mass above it), quark portals
strong enough to set a relic density are excluded from 100 GeV to 10
TeV, lepton portals give a loop-suppressed scattering that is allowed
above about 1.5 TeV, and the Higgs portal requires a weak binding,
\(q_{\mathrm{b}} \lesssim 0.06\,\omega_{\mathrm{b}}\) at 100 GeV to
\(0.02\,\omega_{\mathrm{b}}\) at 10 TeV. Status: STANDARD (the
formulas), ESTIMATE (\texttt{um4\_thermal.py}, T3 and T4).

\[
\frac{q_{\mathrm{b}}}{\omega_{\mathrm{b}}}\Big\vert_{\max} \propto \omega_{\mathrm{b}}^{-1/4},\qquad
H < H_{\max} = \frac{2F(\upsilon_{\mathrm{b}})\,m_\Psi}{\ln\big(p_{\mathrm{w}}\,m_\Psi\,t_0\big)} \tag{S-22}
\] \textbf{Symbols in (S-22).}

\begin{itemize}
\tightlist
\item
  \(q_{\mathrm{b}}\): The inverse width of a state of dirac16complex at
  the fold: inside the turning point its profile falls as
  \(\mathrm{e}^{-q_{\mathrm{b}}\lvert\sigma\rvert}\). Unit: GeV. Status:
  DERIVED.
\item
  \(\omega_{\mathrm{b}}\): The energy at zero momentum of the state of
  dirac16complex at the fold in the limit \(H \to 0\) (the bound state
  of UM-3): the particle \(\chi\) of UM-4 has the mass
  \(\omega_{\mathrm{r}} = \omega_{\mathrm{b}}(1 - H/(2q_{\mathrm{b}})) + \dots\),
  equal to \(\omega_{\mathrm{b}}\) up to the relative correction
  \(H/(2q_{\mathrm{b}})\). Unit: GeV. Status: DERIVED.
\item
  \(H\): The author's constant (an inverse length) that fixes the period
  \(\pi/(6H)\) of \(x_0\), the warp factor
  \(\mathrm{e}^{-2H\lvert\sigma\rvert}\) of the metric, the cosmological
  constant \(\Lambda = -21H^2\) and the tension
  \(\kappa\lambda_{\mathrm{F}} = 12H\) of the fold; in UM-4 a free
  parameter, bounded from below by the propagation of the gauge fields
  in the bulk. Unit: GeV (1/length). Value: \(H \gtrsim 2.5\) TeV (from
  the invisible width of the Z boson). Status: AUTHOR.
\item
  \(H_{\max}\): The largest value of the author's constant \(H\) for
  which particles \(\chi\) bound to the fold by the Higgs portal, with
  the weakest binding that direct detection requires, live longer than
  the age \(t_0\) of the universe. Unit: GeV. Value: from
  \(4.7\times10^{-4}\) GeV at \(\omega_{\mathrm{b}} = 10^4\) GeV to 2.1
  GeV at the Planck mass. Status: ESTIMATE.
\item
  \(F(\upsilon)\): The tunnelling exponent of the width of the particles
  \(\chi\) per unit of \(m_\Psi/H\): the integral of the local decay
  rate \(q\) over the region inside the turning point, divided by
  \(m_\Psi/H\), as a function of the ratio \(\upsilon\) of the energy to
  the bulk mass. Unit: dimensionless. Value: 0.049 for
  \(\upsilon = 0.866\); \((q_{\mathrm{b}}/\omega_{\mathrm{b}})^3/3\) for
  weak binding. Status: DERIVED.
\item
  \(\upsilon,\ \upsilon_{\mathrm{b}}\): The argument of the tunnelling
  exponent \(F\), a ratio of an energy to the bulk mass \(m_\Psi\);
  \(\upsilon_{\mathrm{b}} = \omega_{\mathrm{b}}/m_\Psi\) that of the
  particle \(\chi\). Unit: dimensionless. Status: DERIVED.
\item
  \(m_\Psi\): The mass parameter of dirac16complex in the bulk. Unit:
  GeV. Status: RECORD.
\item
  \(p_{\mathrm{w}}\): The prefactor of the width of the particles
  \(\chi\) in its tunnelling form
  \(\Gamma_\chi \approx p_{\mathrm{w}}m_\Psi\mathrm{e}^{-2F(\upsilon_{\mathrm{b}})m_\Psi/H}\),
  taken from the computed widths at \(m_\Psi/H = 40\). Unit:
  dimensionless. Value: 0.14. Status: ESTIMATE.
\item
  \(t_0\): The age of the universe. Unit: time. Value: 13.8 Gyr. Status:
  STANDARD.
\end{itemize}

The particles \(\chi\) bound by the Higgs portal cannot be the dark
matter of UM-4. With the binding of the author's specification (the
Higgs portal of the left-handed half) the largest
\(q_{\mathrm{b}}/\omega_{\mathrm{b}}\) that the LUX-ZEPLIN limit allows
falls like \(\omega_{\mathrm{b}}^{-1/4}\) (the coupling
\(\mathfrak{g}_{h\chi}\) grows like
\(q_{\mathrm{b}}^2/\omega_{\mathrm{b}}\), the limit on the cross-section
like \(\omega_{\mathrm{b}}\)), and with that binding the lifetime
condition S-6 holds only for \(H\) below \(H_{\max}\):
\(1.8\times10^{-3}\) GeV at \(\omega_{\mathrm{b}} = 2.4\times10^6\) GeV
(where \(q_{\mathrm{b}}/\omega_{\mathrm{b}}\) may be at most 0.0049),
0.12 GeV at \(10^{14}\) GeV and 2.1 GeV at the Planck mass
(\(1.2\times10^{19}\) GeV). \(H_{\max}\) grows only like
\(\omega_{\mathrm{b}}^{1/4}\) and stays far below the bound \(H > 2.5\)
TeV at every mass. Status: ESTIMATE (\texttt{um4\_thermal.py}, T8: the
limit extrapolated in proportion to the mass, the prefactor
\(p_{\mathrm{w}}\) from the computed widths).

\[
\omega_{\mathrm{r}} \simeq \omega_{\mathrm{b}} > 2.4\times10^6\ \text{GeV} > 8.87\times10^4\ \text{GeV},\qquad
H_{\max} \le 2.1\ \text{GeV} < 2534\ \text{GeV} \le H \tag{A-6}
\] \textbf{Symbols in (A-6).}

\begin{itemize}
\tightlist
\item
  \(\omega_{\mathrm{r}}\): The mass of the particle \(\chi\): the real
  part of the energy at rest of the resonance of dirac16complex at the
  fold,
  \(\omega_{\mathrm{r}} = \omega_{\mathrm{b}}(1 - H/(2q_{\mathrm{b}})) + \dots\).
  Unit: GeV. Status: DERIVED.
\item
  \(\omega_{\mathrm{b}}\): The energy at zero momentum of the state of
  dirac16complex at the fold in the limit \(H \to 0\) (the bound state
  of UM-3): the particle \(\chi\) of UM-4 has the mass
  \(\omega_{\mathrm{r}} = \omega_{\mathrm{b}}(1 - H/(2q_{\mathrm{b}})) + \dots\),
  equal to \(\omega_{\mathrm{b}}\) up to the relative correction
  \(H/(2q_{\mathrm{b}})\). Unit: GeV. Status: DERIVED.
\item
  \(H_{\max}\): The largest value of the author's constant \(H\) for
  which particles \(\chi\) bound to the fold by the Higgs portal, with
  the weakest binding that direct detection requires, live longer than
  the age \(t_0\) of the universe. Unit: GeV. Value: from
  \(4.7\times10^{-4}\) GeV at \(\omega_{\mathrm{b}} = 10^4\) GeV to 2.1
  GeV at the Planck mass. Status: ESTIMATE.
\item
  \(H\): The author's constant (an inverse length) that fixes the period
  \(\pi/(6H)\) of \(x_0\), the warp factor
  \(\mathrm{e}^{-2H\lvert\sigma\rvert}\) of the metric, the cosmological
  constant \(\Lambda = -21H^2\) and the tension
  \(\kappa\lambda_{\mathrm{F}} = 12H\) of the fold; in UM-4 a free
  parameter, bounded from below by the propagation of the gauge fields
  in the bulk. Unit: GeV (1/length). Value: \(H \gtrsim 2.5\) TeV (from
  the invisible width of the Z boson). Status: AUTHOR.
\end{itemize}

What is left of the window. The first chain: the lightest particles
\(\chi\) that live longer than the universe (S-6 with the lower bound of
\(H\)) are heavier than the heaviest thermal relic that unitarity allows
(S-20). The second: the largest \(H\) for which particles bound by the
Higgs portal both live long enough and escape direct detection (S-22, at
any mass up to the Planck mass) lies below the lower bound of \(H\) from
the width of the Z boson. dirac16complex can therefore be the dark
matter of UM-4 only if three conditions hold that the action of UM-4
does not supply:

\begin{itemize}
\tightlist
\item
  a binding of \(\chi\) to the fold that does not couple it to nucleons
  (the gauge-kinetic portals with a small nucleon coupling, or a
  strength \(\mu_{\mathrm{R}}\) of the fold for the right-handed half
  that the Higgs field does not supply), strong enough for
  \(q_{\mathrm{b}}\) of the order of \(m_\Psi\) (then
  \(m_\Psi/H > 1.1\times10^3\) suffices);
\item
  a non-thermal origin of the abundance
  \(\Omega_\chi h_{100}^2 = 0.120\) (for example production at the end
  of an inflation, which UM-4 does not compute);
\item
  right-handed neutrinos that never thermalise: their portal couplings
  must be much smaller than those of the lepton doublets, or they would
  add 0.14 to 0.24 to the effective number of neutrino species.
\end{itemize}

Within these conditions the particles \(\chi\) are heavier than
\(2.4\times10^6\) GeV, live longer than the universe, gravitate as
matter of the fold and form a cold gas. Status: ESTIMATE (the window),
OPEN (the binding, the origin, the couplings of the right-handed
neutrinos).

\hypertarget{the-answer-1}{%
\subsection{The answer}\label{the-answer-1}}

dirac16complex can be the dark matter of UM-4 only in part of what that
requires: its particles cannot decay into the Standard Model alone, live
longer than the universe if \(m_\Psi/H > 1.1\times10^3\) (so they are
heavier than \(2.4\times10^6\) GeV), gravitate as matter of the fold,
and form a cold, pressureless gas. But they cannot have the measured
density as a thermal relic, and the binding of the author's
specification (the Higgs portal) cannot hold them at the fold while they
escape direct detection. The dark matter of UM-4 needs a binding of
\(\chi\) to the fold without a nucleon coupling and a non-thermal origin
of its abundance, both outside the action of UM-4. Status: DERIVED (no
decay into the Standard Model alone, the resonances, the equation of
state), ESTIMATE (the lifetime condition, the window), OPEN (the
binding, the origin of the abundance, the classical stability of the
gas, which the indefinite charge leaves undecided).

\hypertarget{dirac16complex-and-the-dark-energy}{%
\section{dirac16complex and the dark
energy}\label{dirac16complex-and-the-dark-energy}}

\[
\begin{aligned}
&H_4^2 = \frac{8\pi G_{\mathrm{N}}}{3}\big(\rho_{\mathrm{SM}} + \rho_\chi + \rho_\Lambda\big),\qquad
 2H_4' + 3H_4^2 = -8\pi G_{\mathrm{N}}\,\mathcal{P},\qquad
 \frac{1}{8\pi G_{\mathrm{N}}} = \frac{V_*}{\kappa_7} + \frac{1}{\kappa_{\mathrm{ind}}},\qquad \kappa_7 = \frac52\,\kappa H,\\
&\rho_\Lambda = \rho_{\mathrm{ff}} + \rho_{\mathrm{vac}} - \lambda_{\mathrm{F}}V_*,\qquad
 q_{\mathrm{SM}} + q_\chi + q_{\mathrm{P}} + q_{\mathrm{ff}} + \lambda_{\mathrm{F}}V_* + \frac{\rho - 3\mathcal{P}}{2} = 0,\\
&q_{\mathrm{SM}} = \big\langle\mathcal{L}_{\mathrm{H}} + \mathcal{L}_{\mathrm{f}} + \mathcal{L}_{\mathrm{Y}}\big\rangle,\qquad
 q_{\mathrm{P}} = \langle\mathcal{L}_{\mathrm{P}}\rangle
\end{aligned} \tag{FE-59}
\] \textbf{Symbols in (FE-59).}

\begin{itemize}
\tightlist
\item
  \(H_4\): The expansion rate (Hubble rate) of the fold, \(H_4 = a_4'\),
  as a function of the time or of the redshift; \(H_4^{(0)}\) its value
  today, the Hubble constant. Unit: GeV (1/time). Value:
  \(H_4^{(0)} = 100\,h_{100}\) km/s/Mpc with \(h_{100} = 0.674\) (Planck
  2018): 67.4 km/s/Mpc \(= 1.44\times10^{-42}\) GeV. Status: STANDARD.
\item
  \({}'\): The derivative of a function of the time \(x_4\) alone with
  respect to \(x_4\) (\(a_4' = \mathrm{d}a_4/\mathrm{d}x_4\); two primes
  the second derivative); on a function of \(\sigma\) alone (a profile)
  the derivative with respect to \(\sigma\). Status: STANDARD.
\item
  \(\pi\): The number pi, the ratio of the circumference of a circle to
  its diameter. Unit: dimensionless. Value: 3.141592654. Status:
  STANDARD.
\item
  \(G_{\mathrm{N}}\): Newton's constant of four-dimensional gravity (the
  gravity of our world, on the fold),
  \(G_{\mathrm{N}} = 1/M_{\mathrm{Pl}}^2\) with the Planck mass
  \(M_{\mathrm{Pl}} = 1.221\times10^{19}\) GeV. Unit: GeV\(^{-2}\).
  Value: \(6.708\times10^{-39}\) GeV\(^{-2}\). Status: STANDARD.
\item
  \(\rho_{\mathrm{SM}}\): The energy density of the Standard-Model
  particles. Unit: GeV\(^4\). Status: STANDARD.
\item
  \(\rho_\chi\): The energy density of the dark matter, the particles
  \(\chi\) of dirac16complex; \(\rho_\chi^{(0)}\) its value today. Unit:
  GeV\(^4\). Status: DERIVED.
\item
  \(\rho_\Lambda\): The vacuum energy density of the fold above its
  static tension, the dark energy of UM-4:
  \(\rho_\Lambda = \rho_{\mathrm{ff}} + \rho_{\mathrm{vac}} - \lambda_{\mathrm{F}}V_*\)
  (four-dimensional normalisation). Unit: GeV\(^4\). Value:
  \(0.685\,\rho_{\mathrm{c}}^{(0)} = 2.5\times10^{-47}\) GeV\(^4\)
  (Planck 2018). Status: DERIVED.
\item
  \(\mathcal{P}\): The total pressure of the matter of the fold;
  \(\Delta\mathcal{P}\) its perturbation. Unit: GeV\(^4\). Status:
  STANDARD.
\item
  \(V_*\): The physical volume of the extra times \(x_5, x_6, x_7\) on
  the fold in the late universe,
  \(V_* = V_{\mathrm{t}}(\mathrm{e}^{-a_4}b_4)^3\) (with
  \(V_{\mathrm{t}} = \int\mathrm{d}x_5\,\mathrm{d}x_6\,\mathrm{d}x_7\)
  the coordinate volume of the extra times) at a time when
  \(\mathrm{e}^{-a_4}b_4\) is constant; it converts the
  eight-dimensional normalisation into the four-dimensional one. Unit:
  length\(^3\) (GeV\(^{-3}\)). Status: ASSUMED.
\item
  \(\kappa_7\): The gravitational coupling of the fold: to the first
  order in its curvature over \(H^2\) the fold obeys the Einstein
  equations of seven dimensions
  \(\hat G^a{}_b = \kappa_7\mathsf{t}^a{}_b\) with
  \(\kappa_7 = 5\kappa H/2\), the coupling of the massless graviton of
  the bulk bound to the fold; four-dimensional gravity on the fold has
  \(8\pi G_{\mathrm{N}} = \kappa_7/V_*\) (plus the induced term, which
  is negligible). Unit: GeV\(^{-5}\). Value: \(5\kappa H/2\). Status:
  DERIVED.
\item
  \(\kappa_{\mathrm{ind}}\): The coupling of the Einstein-Hilbert term
  \(\hat R/(2\kappa_{\mathrm{ind}})\) that the quantum fluctuations of
  the fields of the fold may induce,
  \(\kappa_{\mathrm{ind}} = 8\pi G_{\mathrm{ind}}\); it adds to the
  gravity of the fold,
  \(1/(8\pi G_{\mathrm{N}}) = V_*/\kappa_7 + 1/\kappa_{\mathrm{ind}}\).
  Unit: GeV\(^{-2}\). Value: not fixed by UM-4; negligible for a cutoff
  of the order of \(H\). Status: ESTIMATE.
\item
  \(\kappa\): The eight-dimensional gravitational coupling,
  \(\kappa = 8\pi G_8\). Unit: GeV\(^{-6}\). Status: AUTHOR.
\item
  \(H\): The author's constant (an inverse length) that fixes the period
  \(\pi/(6H)\) of \(x_0\), the warp factor
  \(\mathrm{e}^{-2H\lvert\sigma\rvert}\) of the metric, the cosmological
  constant \(\Lambda = -21H^2\) and the tension
  \(\kappa\lambda_{\mathrm{F}} = 12H\) of the fold; in UM-4 a free
  parameter, bounded from below by the propagation of the gauge fields
  in the bulk. Unit: GeV (1/length). Value: \(H \gtrsim 2.5\) TeV (from
  the invisible width of the Z boson). Status: AUTHOR.
\item
  \(\rho_{\mathrm{ff}}\): The energy density of the fold fluid. Unit:
  GeV\(^4\). Value:
  \(\lambda_{\mathrm{F}}V_* + \rho_\Lambda - \rho_{\mathrm{vac}}\):
  together with the vacuum energy \(\rho_{\mathrm{vac}}\) of the fields
  of the fold it makes up the tension \(\lambda_{\mathrm{F}}V_*\) that
  the static block requires plus the dark energy \(\rho_\Lambda\).
  Status: ASSUMED (the fluid), DERIVED (the value).
\item
  \(\rho_{\mathrm{vac}}\): The vacuum energy of the fields of the fold:
  the minimum of their potential (for the Standard Model the minimum of
  the Higgs potential with its loop corrections), a tension of the
  seven-dimensional fold. Unit: GeV\(^4\). Status: STANDARD.
\item
  \(\lambda_{\mathrm{F}}\): The tension that the static block requires
  of the fold, \(\kappa\lambda_{\mathrm{F}} = 12H\): the total vacuum
  energy of the seven-dimensional fold, made of the tension of the fold
  fluid and the vacuum energy of the fields of the fold (both have the
  pressure minus their energy density along ordinary space and along the
  extra times), per unit volume of the six other directions of the fold
  (eight-dimensional normalisation); in the four-dimensional
  normalisation it is \(\lambda_{\mathrm{F}}V_*\), and the excess of the
  total over \(\lambda_{\mathrm{F}}V_*\) is the dark energy
  \(\rho_\Lambda\). Unit: GeV\(^7\). Value: \(12H/\kappa\);
  \(\lambda_{\mathrm{F}}V_* = 30H^2M_{\mathrm{red}}^2 = 1.1\times10^{45}\)
  GeV\(^4\) for \(H = 2.5\) TeV. Status: DERIVED.
\item
  \(q_{\mathrm{SM}}\): The pressure of the Standard Model along the
  extra times at the fold: the mean on-shell Lagrangian density of the
  Standard-Model fields (the extra-time component of the energy-momentum
  tensor of every term of the fold is its Lagrangian, because the fields
  of the fold do not depend on the extra times). Unit: GeV\(^4\).
  Status: DERIVED (for radiation and free fields), ASSUMED (for the
  baryons).
\item
  \(q_\chi\): The pressure of the particles \(\chi\) of dirac16complex
  along the extra times, integrated across the fold (their on-shell
  Lagrangian, the self-interaction energy per volume; zero for a dilute
  gas). Unit: GeV\(^4\). Status: DERIVED.
\item
  \(q_{\mathrm{P}}\): The pressure of the portals along the extra times:
  their mean on-shell Lagrangian density. Unit: GeV\(^4\). Status:
  DERIVED.
\item
  \(q_{\mathrm{ff}}\): The pressure of the fold fluid along the extra
  times: minus its energy density (its tension) plus the pressure
  \(q_{\mathrm{hold}}\) that holds the volume of the extra times fixed.
  Unit: GeV\(^4\). Value: \(-\rho_{\mathrm{ff}} + q_{\mathrm{hold}}\).
  Status: ASSUMED (the fluid), DERIVED (the value).
\item
  \(\rho\): The total energy density of the matter of the fold (all
  species); \(\Delta\rho\) its perturbation. Unit: GeV\(^4\). Status:
  STANDARD.
\item
  \(\langle\cdots\rangle\): The mean (expectation) value of an operator
  in the state of the medium (the vacuum, a plasma, the Kohn-Sham
  state). Status: STANDARD.
\item
  \(\mathcal{L}_{\mathrm{H}}\): The Lagrangian of the Higgs doublet on
  the fold. Unit: GeV\(^4\). Status: AUTHOR.
\item
  \(\mathcal{L}_{\mathrm{f}}\): The kinetic Lagrangian of the quarks and
  leptons on the fold. Unit: GeV\(^4\). Status: STANDARD.
\item
  \(\mathcal{L}_{\mathrm{Y}}\): The Yukawa Lagrangian of the quarks and
  leptons with the Higgs doublet. Unit: GeV\(^4\). Status: STANDARD.
\item
  \(\mathcal{L}_{\mathrm{P}}\): The Lagrangian of the portals of
  dirac16complex at the fold. Unit: GeV\(^4\). Status: INVENTED.
\end{itemize}

\[
\mathcal{S}^a{}_b = -\big(\lambda_{\mathrm{F}} + \Delta\lambda_{\mathrm{F}}\big)\delta^a{}_b + \mathsf{t}^a{}_b:\qquad
\hat G^a{}_b = \kappa_7\big(\mathsf{t}^a{}_b - \Delta\lambda_{\mathrm{F}}\,\delta^a{}_b\big),\qquad
\rho_\Lambda = V_*\,\Delta\lambda_{\mathrm{F}},\qquad
\frac{\Delta\lambda_{\mathrm{F}}}{\lambda_{\mathrm{F}}} = \frac{\rho_\Lambda}{30H^2M_{\mathrm{red}}^2} = 2.2\times10^{-92} \tag{S-4}
\] \textbf{Symbols in (S-4).}

\begin{itemize}
\tightlist
\item
  \(\mathcal{S}^a{}_b\): The thin source of the fold: the
  energy-momentum tensor concentrated on the fold, integrated across it,
  \(\kappa\mathcal{S}^a{}_b = \int_{0^-}^{0^+}\mathrm{d}\sigma\,\mathcal{E}_{(1)}{}^a{}_b\)
  (the junction condition of Israel); \(\mathcal{S}^4{}_4\) is minus its
  energy density, \(\mathcal{S}^1{}_1\) and \(\mathcal{S}^5{}_5\) its
  pressures along ordinary space and along the extra times (frame
  components along the fold, eight-dimensional normalisation: an energy
  per unit volume of the six other directions of the fold). Unit:
  GeV\(^7\). Value: \(\kappa\mathcal{S}^a{}_b = -12H\delta^a{}_b\) for
  every \(a_4\), \(b_4\) of the author's form. Status: DERIVED.
\item
  \(a\): A frame (orthonormal) index running over \(0, 1, \dots, 7\)
  (frame direction \((0)\) along \(\sigma\), \((1), (2), (3)\) along
  ordinary space, \((4)\) along the time, \((5), (6), (7)\) along the
  extra times); raised and lowered with the frame metric
  \(\eta_{ab} = \mathrm{diag}(1, 1, 1, 1, -1, -1, -1, -1)\); on the fold
  (\(a, b \in \{1, \dots, 7\}\)) the frame of the fold; in the tensors
  of the fold written with its coordinates \(x^a\) (its metric
  \(\hat g_{ab}\) and the tensors built from it) a coordinate index of
  the fold, \(a \in \{1, \dots, 7\}\); a repeated index is summed.
  Status: STANDARD.
\item
  \(b\): A frame index (as \(a\)). Status: STANDARD.
\item
  \(\lambda_{\mathrm{F}}\): The tension that the static block requires
  of the fold, \(\kappa\lambda_{\mathrm{F}} = 12H\): the total vacuum
  energy of the seven-dimensional fold, made of the tension of the fold
  fluid and the vacuum energy of the fields of the fold (both have the
  pressure minus their energy density along ordinary space and along the
  extra times), per unit volume of the six other directions of the fold
  (eight-dimensional normalisation); in the four-dimensional
  normalisation it is \(\lambda_{\mathrm{F}}V_*\), and the excess of the
  total over \(\lambda_{\mathrm{F}}V_*\) is the dark energy
  \(\rho_\Lambda\). Unit: GeV\(^7\). Value: \(12H/\kappa\);
  \(\lambda_{\mathrm{F}}V_* = 30H^2M_{\mathrm{red}}^2 = 1.1\times10^{45}\)
  GeV\(^4\) for \(H = 2.5\) TeV. Status: DERIVED.
\item
  \(\Delta\): A change, difference or small variation of the quantity
  that follows (for example an interval of the time, a perturbation of a
  density, a difference of two squared masses, or an addition to an
  effective number of species). Status: STANDARD.
\item
  \(\delta\): The Kronecker delta: with two indices \(\delta^M_N = 1\)
  if \(M = N\) and 0 otherwise (likewise \(\delta_{X\mathrm{L}}\) is 1
  if the label \(X\) of a half of dirac16complex is L and 0 otherwise);
  with two index strings of equal length the generalised Kronecker delta
  \(\delta^{M_1M_2\cdots}_{N_1N_2\cdots}\), the determinant of the
  square matrix of the Kronecker deltas \(\delta^{M_a}_{N_b}\) (\(+1\)
  or \(-1\) if the upper indices are an even or odd permutation of the
  lower ones, 0 otherwise). Status: STANDARD.
\item
  \(\mathsf{t}^a{}_b,\ \mathsf{t}\): The matter of the fold above its
  tension: the thin source minus the tension,
  \(\mathcal{S}^a{}_b = -\lambda_{\mathrm{F}}\delta^a{}_b + \mathsf{t}^a{}_b\)
  (eight-dimensional normalisation); \(\mathsf{t} = \mathsf{t}^c{}_c\)
  its trace; with frozen extra times
  \(\mathsf{t}^a{}_b = \mathrm{diag}(P_7, P_7, P_7, -\rho_7, q_7, q_7, q_7)\).
  Unit: GeV\(^7\). Status: DERIVED.
\item
  \(\hat G^\mu{}_\nu\): The Einstein tensor of the metric
  \(\hat g_{ab}\) of the fold (FD-2), with frame indices along its seven
  directions; with four-dimensional indices, \(\hat G^\mu{}_\nu\), its
  components along the four observable directions, which for frozen
  extra times are those of the Einstein tensor of the four-dimensional
  metric \(\hat h_{\mu\nu}\). Unit: 1/length\(^2\). Status: STANDARD.
\item
  \(\kappa_7\): The gravitational coupling of the fold: to the first
  order in its curvature over \(H^2\) the fold obeys the Einstein
  equations of seven dimensions
  \(\hat G^a{}_b = \kappa_7\mathsf{t}^a{}_b\) with
  \(\kappa_7 = 5\kappa H/2\), the coupling of the massless graviton of
  the bulk bound to the fold; four-dimensional gravity on the fold has
  \(8\pi G_{\mathrm{N}} = \kappa_7/V_*\) (plus the induced term, which
  is negligible). Unit: GeV\(^{-5}\). Value: \(5\kappa H/2\). Status:
  DERIVED.
\item
  \(\rho_\Lambda\): The vacuum energy density of the fold above its
  static tension, the dark energy of UM-4:
  \(\rho_\Lambda = \rho_{\mathrm{ff}} + \rho_{\mathrm{vac}} - \lambda_{\mathrm{F}}V_*\)
  (four-dimensional normalisation). Unit: GeV\(^4\). Value:
  \(0.685\,\rho_{\mathrm{c}}^{(0)} = 2.5\times10^{-47}\) GeV\(^4\)
  (Planck 2018). Status: DERIVED.
\item
  \(V_*\): The physical volume of the extra times \(x_5, x_6, x_7\) on
  the fold in the late universe,
  \(V_* = V_{\mathrm{t}}(\mathrm{e}^{-a_4}b_4)^3\) (with
  \(V_{\mathrm{t}} = \int\mathrm{d}x_5\,\mathrm{d}x_6\,\mathrm{d}x_7\)
  the coordinate volume of the extra times) at a time when
  \(\mathrm{e}^{-a_4}b_4\) is constant; it converts the
  eight-dimensional normalisation into the four-dimensional one. Unit:
  length\(^3\) (GeV\(^{-3}\)). Status: ASSUMED.
\item
  \(H\): The author's constant (an inverse length) that fixes the period
  \(\pi/(6H)\) of \(x_0\), the warp factor
  \(\mathrm{e}^{-2H\lvert\sigma\rvert}\) of the metric, the cosmological
  constant \(\Lambda = -21H^2\) and the tension
  \(\kappa\lambda_{\mathrm{F}} = 12H\) of the fold; in UM-4 a free
  parameter, bounded from below by the propagation of the gauge fields
  in the bulk. Unit: GeV (1/length). Value: \(H \gtrsim 2.5\) TeV (from
  the invisible width of the Z boson). Status: AUTHOR.
\item
  \(M_{\mathrm{red}}\): The reduced Planck mass,
  \(M_{\mathrm{red}} = (8\pi G_{\mathrm{N}})^{-1/2}\). Unit: GeV. Value:
  \(2.435\times10^{18}\) GeV. Status: STANDARD.
\end{itemize}

The Einstein equations of the fold with frozen extra times (FE-59) and
the dark energy (S-4). The fold of the static block is a pure tension
\(\lambda_{\mathrm{F}}\) (\(\kappa\lambda_{\mathrm{F}} = 12H\), in the
four-dimensional normalisation
\(\lambda_{\mathrm{F}}V_* = 30H^2M_{\mathrm{red}}^2\)); a tension above
it by \(\Delta\lambda_{\mathrm{F}}\) gives the fold a cosmological
constant, the vacuum energy
\(\rho_\Lambda = V_*\Delta\lambda_{\mathrm{F}}\) with the equation of
state \(w = -1\) (its pressure minus its energy density, exactly). The
dark energy of UM-4 is this excess: the energy of the fold fluid and the
vacuum energy of the fields of the fold minus the critical tension.
Status: DERIVED (\texttt{um4\_foldsources.py}, FS4; with the ASSUMED
fold fluid).

\[
\rho_\Lambda = \Omega_\Lambda\,\rho_{\mathrm{c}}^{(0)} = 0.685\cdot\frac{3\big(H_4^{(0)}\big)^2}{8\pi G_{\mathrm{N}}} = 2.5\times10^{-47}\ \text{GeV}^4 = \big(2.2\ \text{meV}\big)^4 \tag{A-7}
\] \textbf{Symbols in (A-7).}

\begin{itemize}
\tightlist
\item
  \(\rho_\Lambda\): The vacuum energy density of the fold above its
  static tension, the dark energy of UM-4:
  \(\rho_\Lambda = \rho_{\mathrm{ff}} + \rho_{\mathrm{vac}} - \lambda_{\mathrm{F}}V_*\)
  (four-dimensional normalisation). Unit: GeV\(^4\). Value:
  \(0.685\,\rho_{\mathrm{c}}^{(0)} = 2.5\times10^{-47}\) GeV\(^4\)
  (Planck 2018). Status: DERIVED.
\item
  \(\Omega_\Lambda\): The density parameter of the dark energy (the
  vacuum energy of the fold) today. Unit: dimensionless. Value: 0.685
  (Planck 2018). Status: DERIVED.
\item
  \(\rho_{\mathrm{c}}\): The critical density, the energy density of a
  spatially flat fold at the expansion rate \(H_4\);
  \(\rho_{\mathrm{c}}^{(0)}\) its value today. Unit: GeV\(^4\). Value:
  \(\rho_{\mathrm{c}}^{(0)} = 1.0537\times10^{-5}\,h_{100}^2\)
  GeV/cm\(^3\). Status: STANDARD.
\item
  \(H_4\): The expansion rate (Hubble rate) of the fold, \(H_4 = a_4'\),
  as a function of the time or of the redshift; \(H_4^{(0)}\) its value
  today, the Hubble constant. Unit: GeV (1/time). Value:
  \(H_4^{(0)} = 100\,h_{100}\) km/s/Mpc with \(h_{100} = 0.674\) (Planck
  2018): 67.4 km/s/Mpc \(= 1.44\times10^{-42}\) GeV. Status: STANDARD.
\item
  \(\pi\): The number pi, the ratio of the circumference of a circle to
  its diameter. Unit: dimensionless. Value: 3.141592654. Status:
  STANDARD.
\item
  \(G_{\mathrm{N}}\): Newton's constant of four-dimensional gravity (the
  gravity of our world, on the fold),
  \(G_{\mathrm{N}} = 1/M_{\mathrm{Pl}}^2\) with the Planck mass
  \(M_{\mathrm{Pl}} = 1.221\times10^{19}\) GeV. Unit: GeV\(^{-2}\).
  Value: \(6.708\times10^{-39}\) GeV\(^{-2}\). Status: STANDARD.
\end{itemize}

The vacuum energy that the measurements require (\(H_4^{(0)} = 67.36\)
km/s/Mpc \(= 1.44\times10^{-42}\) GeV). Status: STANDARD.

Can dirac16complex be the dark energy? A dark energy needs a pressure
close to minus its energy density, an equation of state \(w\) (pressure
over energy density) close to \(-1\). The gas of the particles \(\chi\)
has the pressure \(P_{\mathrm{KS}} = n_\chi\mathsf{E}_{\mathrm{int}}\)
of its self-interaction (S-9), which is positive for a repulsive
self-interaction and vanishes in the dilute limit
(\(w_\chi = 5.4\times10^{-8}\) at \(G_{\mathrm{KS}} = 10^{-6}\)). An
attractive self-interaction (\(G_{\mathrm{KS}} < 0\), whose states
converge as well) would give by S-9 a negative pressure of the same
small size, since the integral of S-9 is positive (0.102 per unit of
\(G_{\mathrm{KS}}\) in the dilute limit), and the gas of the universe
has \(G_{\mathrm{KS}}\) below \(10^{-29}\): \(w_\chi\) is far from
\(-1\) in every case. The gas of \(\chi\) is dark matter, not dark
energy. The other contribution of dirac16complex to the vacuum energy,
the energy of its unoccupied and its negative-energy states (the sea,
which the Kohn-Sham states leave out), is a renormalisation of the
tension of the fold and of the cosmological constant of the bulk that
UM-4 does not compute; its sign and size are OPEN, and with the
indefinite Krein norm it is not even clear that it is bounded. The dark
energy of UM-4 is therefore the excess of the tension, a free parameter:
its value \((2.2\ \text{meV})^4\) is accommodated by tuning the tension
to
\(\Delta\lambda_{\mathrm{F}}/\lambda_{\mathrm{F}} = 2.2\times10^{-92}\)
for \(H = 2.5\) TeV (smaller still for a larger \(H\)), not explained:
the cosmological-constant problem of the brane of Randall and Sundrum.
With the Planck 2018 parameters UM-4 reproduces the CMB
(\(\chi^2 = 584\) for 613 high-multipole bins); against the BAO of DESI
DR2 it has \(\chi^2 = 32.0\) for 13 measurements, and DESI alone prefers
a lower matter density (0.297 instead of 0.315, \(\chi^2 = 10.5\) for 11
degrees of freedom): with the errors of the CMB DESI quotes a tension of
2.3 standard deviations between the BAO and the CMB in the
cosmological-constant model, which UM-4 keeps, and UM-4 cannot produce
the evolving dark energy that DESI DR2 with the CMB prefers at 3.1
standard deviations, an equation of state \(w(a) = w_0 + w_a(1 - a)\)
that changes with the scale factor \(a = \mathrm{e}^{a_4}\) of the fold
(with its present value \(w_0 > -1\) and its slope \(w_a < 0\)). Status:
DERIVED (\texttt{um4\_kohnsham.py}, K5, K7 and K10;
\texttt{um4\_cosmology.py}, C4 and C5: \(w = -1\)), STANDARD (DESI's
significances), OPEN (the sea, the tuning).

\hypertarget{the-answer-2}{%
\subsection{The answer}\label{the-answer-2}}

No: dirac16complex is not the dark energy of UM-4. Its particles form
cold dark matter; the dark energy is the excess of the tension of the
fold, a free parameter with \(w = -1\), tuned to one part in
\(10^{92}\).

\hypertarget{the-open-questions-answered}{%
\section{The open questions
answered}\label{the-open-questions-answered}}

On 2026-10-25 the author asked for an analysis and an answer to the
questions that the model had left open: whether the indefinite Krein
norm of the charge of dirac16complex is stable; what fields that depend
on the extra times do (they would be tachyonic); whether the particles
\(\chi\) move faster than light; the primordial era, for which the model
had used the values of Planck 2018 as inputs; and the lightest neutrino
mass and the Dirac phase, which the model had set to zero. The answers
follow for UM-4; where the warped time changed an answer, the change is
stated.

\hypertarget{is-the-indefinite-krein-norm-stable}{%
\subsection{Is the indefinite Krein norm
stable?}\label{is-the-indefinite-krein-norm-stable}}

\[
\big\{\mathsf{a}_o, \mathsf{a}_o^\dagger\big\} = -\mathcal{Q}_o,\qquad
\mathsf{H}_\Psi = \sum_o\big(-\mathcal{Q}_o\,\omega_o\big)\,\mathsf{a}_o^\dagger\mathsf{a}_o,\qquad
\big[\mathsf{H}_\Psi, \mathsf{a}_o^\dagger\big] = \omega_o\,\mathsf{a}_o^\dagger \tag{FE-67}
\] \textbf{Symbols in (FE-67).}

\begin{itemize}
\tightlist
\item
  \(\mathsf{a}_o,\ \mathsf{a}_o^\dagger\): The annihilation and the
  creation operator of the mode \(o\) of the quantised dirac16complex
  (frequency \(\omega_o\), Krein charge \(\mathcal{Q}_o\)); the dagger
  is the adjoint with respect to the indefinite (Krein) scalar product.
  Unit: dimensionless. Status: STANDARD.
\item
  \(o\): The label of a single-particle state of dirac16complex: a
  Kohn-Sham orbital, or a mode of the quantised field; a sum over \(o\)
  runs over the orbitals or the modes. Status: STANDARD.
\item
  \({}^\dagger\): Hermitian conjugation: the complex conjugate of the
  transpose of a matrix or spinor; for an operator on the space of
  states, its adjoint with respect to the scalar product of that space
  (for the quantised dirac16complex the indefinite Krein scalar
  product). Status: STANDARD.
\item
  \(\mathcal{Q}_o\): The charge \(Q_\Psi\) (the Krein charge) of the
  single-particle state \(o\) of dirac16complex, a Kohn-Sham orbital or
  a mode: \(+1\) or \(-1\) for a normalised state. Unit: dimensionless.
  Status: DERIVED.
\item
  \(\mathsf{H}_\Psi\): The Hamiltonian (the energy operator) of the
  quantised dirac16complex. Unit: GeV. Status: DERIVED.
\item
  \(\omega\): A frequency: the energy of a mode
  \(\propto \mathrm{e}^{-\mathrm{i}\omega x_4}\) (\(\omega_o\) of the
  orbital \(o\)). Unit: GeV. Status: STANDARD.
\end{itemize}

The quantised dirac16complex (with FE-43). Each mode \(o\) of the field,
of the frequency \(\omega_o\) and the Krein charge
\(\mathcal{Q}_o = \pm1\), has the anticommutator \(-\mathcal{Q}_o\) (the
canonical anticommutator of the field is minus the Krein operator, the
sign chosen so that the particles \(\chi\), the occupied states of
charge \(-1\), have a positive norm), and every quantum raises the
energy by its frequency. In the positive-energy quantisation, in which
the vacuum is annihilated by \(\mathsf{a}_o\) for \(\omega_o > 0\) and
by \(\mathsf{a}_o^\dagger\) for \(\omega_o < 0\), every particle and
every hole carries a positive energy and the norm \(-\mathcal{Q}_o\) of
its mode. The particles and the holes of the modes of charge \(-1\) are
physical; those of the modes of charge \(+1\) are ghosts, states of
negative norm. A fermion cannot trade its norm for its energy, because
an anticommutator \(-1\) has no representation with a positive norm: no
choice of the vacuum removes the ghosts. Status: DERIVED
(\texttt{um4\_krein.py}, KR1: the Krein operator has eight eigenvalues
\(+1\) and eight \(-1\)).

Four results decide the question; they were computed for the Kohn-Sham
gas of UM-4 (\(m_\Psi/H = 10^4\), \(q_{\mathrm{b}} = 0.3\,m_\Psi\)).

\begin{enumerate}
\def\labelenumi{\arabic{enumi}.}
\tightlist
\item
  \emph{The vacuum is stable.} Every quantum, physical or ghost, carries
  a positive energy, so the vacuum cannot decay into pairs of quanta of
  opposite energy, and a particle \(\chi\) cannot fall into a state of
  negative energy. Status: DERIVED (KR1, KR7).
\item
  \emph{Whether the gas is classically stable is not decided.} At rest
  the transformations that keep the states at the fold form a Lie
  algebra of dimension 32: 16 compact generators keep the two charges
  apart, 16 boosts mix them (among them the nine boosts \(\Sigma^{ir}\)
  between ordinary space and the extra times). The boosts map the
  self-consistent gas onto other self-consistent gases (at
  \(G_{\mathrm{KS}} = 0.1\) the residual of the Kohn-Sham equation stays
  \(3.5\times10^{-5}\), as without the boost, and the charges stay
  \(-1\)): they are exact flat directions, whose frequency in the linear
  response is zero. The linear response truncated to the level at the
  fold (the admixtures of the four empty orbitals of charge \(+1\),
  split from the occupied level by \(3.491\times10^{-3}\)) has
  frequencies of the order \(10^{-4}\,m_\Psi\), two of them imaginary
  (up to \(2.05\times10^{-4}\)), but of the size of its truncation
  error: it does not decide whether the dilute gas is stable. Status:
  DERIVED (KR2, KR4).
\item
  \emph{The partner level.} In the self-consistent mean field the level
  of charge \(+1\) next to the occupied one is real and bound up to
  \(G_{\mathrm{KS}} = 0.3\), the largest coupling computed, and it moves
  away from the occupied level as the self-interaction grows (the gap is
  0.0004, 0.0035, 0.0064 and 0.0087 at \(G_{\mathrm{KS}} = 0.01\), 0.1,
  0.2 and 0.3): the two cannot mix, because they have opposite parity
  under the mirror of the fold (FD-10) in a mirror-even mean field. (In
  UM-3, computed at \(H = 0.05\), 0.15 and \(0.3\,m_\Psi\), this level
  turned complex from \(G_{\mathrm{KS}} = 0.3\), 0.7 and 1 on, with
  complex levels that were not converged in the size of the box; states
  with \(q_{\mathrm{b}}\) of the order of \(H\) are broad resonances in
  UM-4, not long-lived particles.) Status: DERIVED (KR3, KR5).
\item
  \emph{The couplings of the fold decide unitarity.} At the fold, at
  rest, the 16 chirality-blind operators (products of
  \(\gamma^{(1)}, \dots, \gamma^{(4)}\)) never connect a physical state
  with a ghost; the 16 chiral ones (the same times \(\Gamma_8\)) always
  do. The author's Higgs portal acts on the left-handed half only
  (\(P_{\mathrm{L}}\)): its amplitude to create a physical \(\chi\)
  together with a ghost antiparticle (an s-wave, 0.148) is 0.94 times
  its elastic amplitude (0.157), while a physical pair needs a p-wave.
  For moving particles the scalar, the density and the vector currents
  still do not mix, the magnetic dipole does, in proportion to the
  momentum. Ghosts are therefore produced wherever the chiral portals
  act, and chirality-blind couplings are necessary for unitarity but not
  shown to be sufficient. Status: DERIVED (KR6, KR7).
\end{enumerate}

The answer. The indefinite Krein norm does not make the vacuum decay.
Whether it destabilises a gas of \(\chi\) is not decided by these
computations: the directions that mix the two charges are exactly flat,
and the truncated response is inconclusive. The real problem is
unitarity: with couplings at the fold that distinguish the two
chiralities, the author's left-handed Higgs portal among them, UM-4
produces ghosts and is not a unitary quantum theory. A unitary UM-4
would need couplings of dirac16complex to the fold that do not
distinguish the chiralities, and these are necessary but not shown to be
sufficient (OPEN). The warped time does not change this answer.

\hypertarget{fields-that-depend-on-the-extra-times}{%
\subsection{Fields that depend on the extra
times}\label{fields-that-depend-on-the-extra-times}}

\[
\mathrm{e}^{2H\lvert\sigma\rvert}\Big(\omega^2 - \mathrm{e}^{-2a_4}\,\vec k^2 + \mathrm{e}^{2a_4}\,b_4^{-2}\,\vec\kappa^2\Big) = m_\Psi^2 + p_\sigma^2,\qquad
\omega^2 = \omega_0^2 + \vec k^2 - \vec\kappa^2\quad\big(a_4 = 0,\ b_4 = 1\big) \tag{FE-38}
\] \textbf{Symbols in (FE-38).}

\begin{itemize}
\tightlist
\item
  \(\mathrm{e}\): Euler's number, the base of the natural logarithm;
  \(\mathrm{e}^{X}\) is the exponential of \(X\). Unit: dimensionless.
  Value: 2.718281828. Status: STANDARD.
\item
  \(H\): The author's constant (an inverse length) that fixes the period
  \(\pi/(6H)\) of \(x_0\), the warp factor
  \(\mathrm{e}^{-2H\lvert\sigma\rvert}\) of the metric, the cosmological
  constant \(\Lambda = -21H^2\) and the tension
  \(\kappa\lambda_{\mathrm{F}} = 12H\) of the fold; in UM-4 a free
  parameter, bounded from below by the propagation of the gauge fields
  in the bulk. Unit: GeV (1/length). Value: \(H \gtrsim 2.5\) TeV (from
  the invisible width of the Z boson). Status: AUTHOR.
\item
  \(\sigma\): The proper coordinate of the hidden direction:
  \(\mathrm{d}\sigma = \lvert\cot z\rvert\,\mathrm{d}x_0\),
  \(\sigma = \mathrm{sgn}(\pi/2 - z)\ln(\sin z)/(6H)\) on one period
  \(0 < z < \pi\); \(\sigma = 0\) is the fold \(z = \pi/2\),
  \(\sigma \to \mp\infty\) the tips \(z \to 0, \pi\); one period of
  \(x_0\) is the whole line \(-\infty < \sigma < \infty\). Unit: length.
  Status: DERIVED.
\item
  \(\omega\): A frequency: the energy of a mode
  \(\propto \mathrm{e}^{-\mathrm{i}\omega x_4}\) (\(\omega_o\) of the
  orbital \(o\)). Unit: GeV. Status: STANDARD.
\item
  \(a_4\): The author's metric function of the time \(x_4\) that scales
  ordinary 3-space: on the fold \(\mathrm{e}^{a_4}\) is the scale factor
  of cosmology and \(a_4'\) the Hubble rate. Unit: dimensionless.
  Status: AUTHOR.
\item
  \(\vec k\): The (comoving) momentum along ordinary space, \(\vec k^2\)
  its square. Unit: GeV. Status: STANDARD.
\item
  \(b_4\): The author's metric function of the time \(x_4\) in the
  extra-time components, \(b_4 > 0\), never identically 1 (the author's
  order of 2026-10-24); the scale factor (the factor of the lengths) of
  the extra times is \(\mathrm{e}^{-a_4}b_4\) on the fold and
  \(\mathrm{e}^{-a_4}b_4\,\mathrm{e}^{-H\lvert\sigma\rvert}\) at the
  distance \(\sigma\) from it (the square root of the warp factor times
  \(\mathrm{e}^{-a_4}b_4\)). Unit: dimensionless. Status: AUTHOR.
\item
  \(\vec\kappa\): A momentum along the extra times (three components);
  the fields of UM-4 have \(\vec\kappa = 0\). Unit: GeV. Status:
  DERIVED.
\item
  \(m_\Psi\): The mass parameter of dirac16complex in the bulk. Unit:
  GeV. Status: RECORD.
\item
  \(p_\sigma\): The momentum along the hidden direction (local wave
  number of a profile along \(\sigma\)). Unit: GeV. Status: DERIVED.
\item
  \(\omega_0\): The energy at rest (zero momentum along ordinary space
  and along the extra times) of a state of dirac16complex at the fold;
  complex for a resonance,
  \(\omega_0 = \omega_{\mathrm{r}} - \frac{\mathrm{i}}{2}\Gamma_\chi\).
  Unit: GeV. Status: DERIVED.
\end{itemize}

The local dispersion relation of dirac16complex with the momentum
\(p_\sigma\) along the hidden direction, \(\vec k\) along ordinary space
and \(\vec\kappa\) along the extra times, and the exact dispersion of
every state at the fold on the static block. In UM-4 the warp multiplies
the frequency and both momenta alike, so the boosts along ordinary space
and the rotations of the time into an extra time (generated by
\(\frac12\gamma^{(4)}\gamma^{(5)}\), which commutes with
\(\gamma^{(0)}\), with the chirality and with the mass term) are
symmetries: a state at the fold with the energy \(\omega_0\) at rest has
\(\omega^2 = \omega_0^2 + \vec k^2 - \vec\kappa^2\) exactly. A momentum
along a time-like direction enters with the sign of the frequency: a
state with \(\lvert\vec\kappa\rvert > \omega_0\) has \(\omega^2 < 0\)
and grows exponentially in \(x_4\), a tachyon. Computed with the
author's matrices (\(H = 10^{-4}m_\Psi\),
\(q_{\mathrm{b}} = 0.3\,m_\Psi\), \(\omega_0 = 0.95378\,m_\Psi\)):
\(\omega = 0.94852\), 0.90537 and 0.74141 \(m_\Psi\) for
\(\lvert\vec\kappa\rvert = 0.1\), 0.3 and 0.6 \(m_\Psi\), equal to the
prediction to ten digits. (In UM-3 the momentum along the extra times
was multiplied by the warp but the frequency was not, so the states with
such a momentum also leaked into the bulk; in UM-4 a state with
\(\lvert\vec\kappa\rvert < \omega_0\) is the resonance at rest, rotated
in the plane of the time and an extra time.) The same holds for every
field, in the bulk as on the fold: a field of mass \(m\) has
\(\omega^2 = m^2 + \vec k^2 - \vec\kappa^2\) on the static fold, and a
massless field has growing modes for every
\(\lvert\vec\kappa\rvert > \lvert\vec k\rvert\). The metric does not
depend on \(x_5, x_6, x_7\), so \(\vec\kappa\) is conserved. Status:
DERIVED (\texttt{um4\_extratimes.py}, EX1 and EX2).

\[
\Gamma_{\mathrm{f}} = \frac{-3H_{\mathrm{i}} + \sqrt{9H_{\mathrm{i}}^2 + 4l_{\mathrm{t}}^2}}{2},\qquad
\Gamma_{\mathrm{t}} = \sqrt{l_{\mathrm{t}}^2 - k_{\mathrm{p}}^2},\qquad L_{\mathrm{t}} \ge 2\pi t_0 = 26.6\ \text{Gpc} \tag{A-8}
\] \textbf{Symbols in (A-8).}

\begin{itemize}
\tightlist
\item
  \(\Gamma_{\mathrm{f}}\): The growth rate of a mode of a massless field
  of the fold with the physical momentum \(l_{\mathrm{t}}\) along the
  extra times and no momentum along ordinary space, on the branch of
  frozen extra times with the constant expansion rate \(H_{\mathrm{i}}\)
  of ordinary space. Unit: GeV (1/time). Status: DERIVED.
\item
  \(H_{\mathrm{i}}\): A constant expansion rate of ordinary space: of
  the de Sitter fold (whose extra times expand at the same rate) or of
  an inflation of the fold with frozen extra times. Unit: GeV. Status:
  DERIVED (FE-87), ASSUMED (FE-88).
\item
  \(l_{\mathrm{t}}\): The physical momentum of a mode along the extra
  times. Unit: GeV. Status: STANDARD.
\item
  \(\Gamma_{\mathrm{t}}\): The growth rate of a mode of a massless field
  of the fold that has a momentum along the extra times, in the late
  universe. Unit: GeV (1/time). Status: DERIVED.
\item
  \(k_{\mathrm{p}}\): The physical momentum of a mode along ordinary
  space. Unit: GeV. Status: STANDARD.
\item
  \(L_{\mathrm{t}}\): The physical period of the extra times, if they
  are compact. Unit: length. Status: ASSUMED.
\item
  \(\pi\): The number pi, the ratio of the circumference of a circle to
  its diameter. Unit: dimensionless. Value: 3.141592654. Status:
  STANDARD.
\item
  \(t_0\): The age of the universe. Unit: time. Value: 13.8 Gyr. Status:
  STANDARD.
\end{itemize}

How fast they grow. With frozen extra times, the background of the late
universe and of an inflation of the fold, the physical momentum
\(l_{\mathrm{t}}\) along the extra times is constant, and a mode without
a momentum along ordinary space grows exponentially, with the expansion
rate \(H_{\mathrm{i}}\) of ordinary space at the rate
\(\Gamma_{\mathrm{f}}\) (computed: 0.646586 for \(H_{\mathrm{i}} = 0.3\)
and \(l_{\mathrm{t}} = 1\), equal to the formula). In the late universe,
with \(a_4\) and \(b_4\) constant, a mode with the physical momenta
\(l_{\mathrm{t}}\) along the extra times and \(k_{\mathrm{p}}\) along
ordinary space grows at the rate \(\Gamma_{\mathrm{t}}\). If the extra
times were compact with the physical period \(L_{\mathrm{t}}\), the
slowest mode (\(k_{\mathrm{p}} = 0\),
\(l_{\mathrm{t}} = 2\pi/L_{\mathrm{t}}\)) would grow by less than one
e-folding over the age \(t_0 = 13.8\) Gyr of the universe only if
\(L_{\mathrm{t}} \ge 2\pi t_0\), 26.6 Gpc; but the \(n\)-th harmonic,
\(l_{\mathrm{t}} = 2\pi n/L_{\mathrm{t}}\), grows \(n\) times faster, so
for every \(L_{\mathrm{t}}\) infinitely many modes grow without bound,
and a compact time-like direction is itself a closed time-like curve.
Status: DERIVED (EX3, EX4).

The answer. Fields that depend on the extra times are tachyonic: every
field, in the bulk and on the fold, has growing modes once it carries a
momentum along the extra times larger than its energy at rest. UM-4
keeps only fields without that momentum (ASSUMED). This is a consistent
truncation (the momentum is conserved, so the sector without it is
invariant) but not a stable one: interactions of fields without that
momentum create pairs with the opposite momenta \(\vec\kappa\) and
\(-\vec\kappa\), which are tachyonic (EX5; an ESTIMATE, no rate is
computed). Making the extra times compact does not help. UM-4 needs a
gauge principle that removes the momentum along the extra times of every
field, as the gauge invariance of electromagnetism removes the
unphysical polarisations of the photon; UM-4 does not contain one
(OPEN).

\hypertarget{does-anything-move-faster-than-light}{%
\subsection{Does anything move faster than
light?}\label{does-anything-move-faster-than-light}}

In UM-3 the answer was yes: its time was not warped, so the light cone
of the bulk was wider than that of the fold, bulk scalars and photons
moved 1.22 and 1.41 times as fast as light on the fold, the bound
particles \(\chi\) had the speed coefficient
\(c_{\mathrm{b}}^2 = q_{\mathrm{b}}/(q_{\mathrm{b}} - H)\) along
ordinary space and a group velocity up to 1.0046 times the speed of
light, and their loops passed a Lorentz violation to the electron unless
\(q_{\mathrm{b}}/H\) was above about \(10^{13}\). This was the reason
for the author's order of 2026-10-25, 13:15 PDT, and UM-4 answers it.

\[
\mathsf{S} = \frac12\,\gamma^{(3)}\gamma^{(4)}:\qquad
\big[\mathsf{S}, \gamma^{(0)}\big] = \big[\mathsf{S}, \Gamma_8\big] = 0,\qquad
\big[\mathsf{S}, \gamma^{(3)}\big] = -\gamma^{(4)},\qquad \big[\mathsf{S}, \gamma^{(4)}\big] = -\gamma^{(3)}\qquad\Rightarrow\qquad
\omega^2 = \omega_0^2 + k^2 \tag{FE-78}
\] \textbf{Symbols in (FE-78).}

\begin{itemize}
\tightlist
\item
  \(\mathsf{S}\): The generator of the boosts along \(x_3\) on
  dirac16complex, \(\mathsf{S} = \frac12\gamma^{(3)}\gamma^{(4)}\).
  Unit: dimensionless. Status: DERIVED.
\item
  \(\gamma^M,\ \gamma^{(a)}\): The Dirac matrices of dirac16complex
  along the coordinates, \(\gamma^M = \mathbf{e}_a{}^MT^a\),
  \(\gamma_M = g_{MN}\gamma^N\); \(\gamma^{(a)} = T^a\) along the frame
  direction \((a)\) (\(\gamma^{(0)}\) along \(\sigma\), \(\gamma^{(4)}\)
  along the time). Unit: dimensionless. Status: AUTHOR.
\item
  \(\Gamma_8\): The chirality matrix of eight dimensions,
  \(\Gamma_8 = T_0T_1\cdots T_7\) (the author's T16\^{}A{[}8{]})
  \(= \mathrm{diag}(-\mathbf{1}_8, \mathbf{1}_8)\);
  \(\Gamma_8^2 = \mathbf{1}\); it anticommutes with every \(T_a\). Unit:
  dimensionless. Status: AUTHOR.
\item
  \(\omega\): A frequency: the energy of a mode
  \(\propto \mathrm{e}^{-\mathrm{i}\omega x_4}\) (\(\omega_o\) of the
  orbital \(o\)). Unit: GeV. Status: STANDARD.
\item
  \(\omega_0\): The energy at rest (zero momentum along ordinary space
  and along the extra times) of a state of dirac16complex at the fold;
  complex for a resonance,
  \(\omega_0 = \omega_{\mathrm{r}} - \frac{\mathrm{i}}{2}\Gamma_\chi\).
  Unit: GeV. Status: DERIVED.
\item
  \(k\): A wave number: the momentum of a plane wave along \(x_3\),
  \(\mathrm{e}^{\mathrm{i}(kx_3 - \omega x_4)}\); in the cosmological
  perturbations the comoving wave number of a Fourier mode,
  \(k = \lvert\vec k\rvert\). Unit: GeV (1/length). Status: STANDARD.
\end{itemize}

\[
\tilde\sigma = \mathrm{sgn}(\sigma)\,\frac{\mathrm{e}^{H\lvert\sigma\rvert} - 1}{H}:\qquad
g_{MN}\,\mathrm{d}X^M\mathrm{d}X^N = \frac{\mathrm{d}\tilde\sigma^2 + \mathrm{d}x_1^2 + \mathrm{d}x_2^2 + \mathrm{d}x_3^2 - \mathrm{d}x_4^2
 - \mathrm{d}x_5^2 - \mathrm{d}x_6^2 - \mathrm{d}x_7^2}{\big(1 + H\lvert\tilde\sigma\rvert\big)^2},\qquad \Delta x_4 \ge L \tag{FE-83}
\] \textbf{Symbols in (FE-83).}

\begin{itemize}
\tightlist
\item
  \(\tilde\sigma\): The conformal coordinate of the hidden direction on
  the static block,
  \(\tilde\sigma = \mathrm{sgn}(\sigma)(\mathrm{e}^{H\lvert\sigma\rvert} - 1)/H\),
  in which the static block is conformally flat. Unit: length. Status:
  DERIVED.
\item
  \(\mathrm{sgn}\): The sign function: \(+1\) for a positive argument
  and \(-1\) for a negative one. Status: STANDARD.
\item
  \(\sigma\): The proper coordinate of the hidden direction:
  \(\mathrm{d}\sigma = \lvert\cot z\rvert\,\mathrm{d}x_0\),
  \(\sigma = \mathrm{sgn}(\pi/2 - z)\ln(\sin z)/(6H)\) on one period
  \(0 < z < \pi\); \(\sigma = 0\) is the fold \(z = \pi/2\),
  \(\sigma \to \mp\infty\) the tips \(z \to 0, \pi\); one period of
  \(x_0\) is the whole line \(-\infty < \sigma < \infty\). Unit: length.
  Status: DERIVED.
\item
  \(\mathrm{e}\): Euler's number, the base of the natural logarithm;
  \(\mathrm{e}^{X}\) is the exponential of \(X\). Unit: dimensionless.
  Value: 2.718281828. Status: STANDARD.
\item
  \(H\): The author's constant (an inverse length) that fixes the period
  \(\pi/(6H)\) of \(x_0\), the warp factor
  \(\mathrm{e}^{-2H\lvert\sigma\rvert}\) of the metric, the cosmological
  constant \(\Lambda = -21H^2\) and the tension
  \(\kappa\lambda_{\mathrm{F}} = 12H\) of the fold; in UM-4 a free
  parameter, bounded from below by the propagation of the gauge fields
  in the bulk. Unit: GeV (1/length). Value: \(H \gtrsim 2.5\) TeV (from
  the invisible width of the Z boson). Status: AUTHOR.
\item
  \(g_{MN}\): The metric tensor of the eight-dimensional spacetime
  (signature \(+,+,+,+,-,-,-,-\) in the order \(x_0, \dots, x_7\)):
  \(g_{MN}\,\mathrm{d}X^M\mathrm{d}X^N\) is the squared interval,
  \(g^{MN}\) its inverse, \(\lvert g\rvert\) the absolute value of its
  determinant (\(\sqrt{\lvert g\rvert}\,\mathrm{d}^8X\) is the proper
  volume element). Unit: dimensionless. Status: AUTHOR.
\item
  \(M\): A coordinate index of eight dimensions (also with a numbered
  copy \(M_1, M_2, \dots\)): it runs over \(x_0, x_1, \dots, x_7\); a
  repeated index is summed. Status: STANDARD.
\item
  \(N\): A coordinate index of eight dimensions (also
  \(N_1, N_2, \dots\)), running over \(x_0, \dots, x_7\); a repeated
  index is summed. Status: STANDARD.
\item
  \(\mathrm{d}\): The differential: \(\mathrm{d}\) in front of a
  quantity is an infinitesimal change of it (as \(\mathrm{d}\sigma\), a
  small step of \(\sigma\)); the quotient of two such changes is a
  derivative; \(\mathrm{d}\) with an exponent, in front of several
  variables, is their volume element, the product of their
  differentials. Status: STANDARD.
\item
  \(X = \{x_0, \dots, x_7\}\): The author's coordinates of the
  eight-dimensional spacetime, \(X = \{x_0, x_1, \dots, x_7\}\) with the
  hidden direction first (Dirac\_Einstein-Lovelock\_4+4\_Nash.nb,
  section DEFCHART); \(X^M\) is the coordinate with the index \(M\) and
  \(\mathrm{d}^8X = \mathrm{d}x_0\,\mathrm{d}x_1\cdots\mathrm{d}x_7\)
  the coordinate volume element. Unit: length (each coordinate). Status:
  AUTHOR.
\item
  \(x_1\): The first coordinate of ordinary 3-space. Unit: length.
  Status: AUTHOR.
\item
  \(x_2\): The second coordinate of ordinary 3-space. Unit: length.
  Status: AUTHOR.
\item
  \(x_3\): The third coordinate of ordinary 3-space. Unit: length.
  Status: AUTHOR.
\item
  \(x_4\): The time; a prime on a function of \(x_4\) alone is
  \(\mathrm{d}/\mathrm{d}x_4\). Unit: length (= time, \(c = 1\)).
  Status: AUTHOR.
\item
  \(x_5\): The first of the three extra time-like coordinates. Unit:
  length. Status: AUTHOR.
\item
  \(x_6\): The second of the three extra time-like coordinates. Unit:
  length. Status: AUTHOR.
\item
  \(x_7\): The third of the three extra time-like coordinates. Unit:
  length. Status: AUTHOR.
\item
  \(\Delta\): A change, difference or small variation of the quantity
  that follows (for example an interval of the time, a perturbation of a
  density, a difference of two squared masses, or an addition to an
  effective number of species). Status: STANDARD.
\item
  \(L\): A distance between two points of the fold along ordinary space:
  for a gravitational wave the radius of a ring of free test masses
  (\(\Delta L\) the change of the distance of one of them from the
  centre), for a signal the distance between its emitter and its
  receiver. Unit: length. Status: STANDARD.
\end{itemize}

\[
c_{\mathrm{loc}}(\sigma) = 1 - \frac{a_4''}{5H^2}\big(\mathrm{e}^{2H\lvert\sigma\rvert} - 1\big),\qquad
\sigma(x_3) = \frac{\lvert a_4''\rvert}{5H}\,x_3\,(L - x_3),\qquad
\Delta x_4 = L - \frac{a_4''^2L^3}{150H^2} \tag{FE-84}
\] \textbf{Symbols in (FE-84).}

\begin{itemize}
\tightlist
\item
  \(c_{\mathrm{loc}}\): The coordinate speed of light along \(x_3\) at
  the depth \(\sigma\) (the rate of change of \(x_3\) with \(x_4\) along
  a null line at fixed \(\sigma\)), in units of that on the fold:
  \(\mathrm{e}^{H\lvert\sigma\rvert}\) on the static block of UM-3,
  whose time was not warped, 1 on the static block of UM-4, and on an
  expanding fold of UM-4 different from 1 in proportion to the
  acceleration \(a_4\) of the expansion. Unit: dimensionless. Status:
  DERIVED.
\item
  \(\sigma\): The proper coordinate of the hidden direction:
  \(\mathrm{d}\sigma = \lvert\cot z\rvert\,\mathrm{d}x_0\),
  \(\sigma = \mathrm{sgn}(\pi/2 - z)\ln(\sin z)/(6H)\) on one period
  \(0 < z < \pi\); \(\sigma = 0\) is the fold \(z = \pi/2\),
  \(\sigma \to \mp\infty\) the tips \(z \to 0, \pi\); one period of
  \(x_0\) is the whole line \(-\infty < \sigma < \infty\). Unit: length.
  Status: DERIVED.
\item
  \(a_4\): The author's metric function of the time \(x_4\) that scales
  ordinary 3-space: on the fold \(\mathrm{e}^{a_4}\) is the scale factor
  of cosmology and \(a_4'\) the Hubble rate. Unit: dimensionless.
  Status: AUTHOR.
\item
  \({}'\): The derivative of a function of the time \(x_4\) alone with
  respect to \(x_4\) (\(a_4' = \mathrm{d}a_4/\mathrm{d}x_4\); two primes
  the second derivative); on a function of \(\sigma\) alone (a profile)
  the derivative with respect to \(\sigma\). Status: STANDARD.
\item
  \(H\): The author's constant (an inverse length) that fixes the period
  \(\pi/(6H)\) of \(x_0\), the warp factor
  \(\mathrm{e}^{-2H\lvert\sigma\rvert}\) of the metric, the cosmological
  constant \(\Lambda = -21H^2\) and the tension
  \(\kappa\lambda_{\mathrm{F}} = 12H\) of the fold; in UM-4 a free
  parameter, bounded from below by the propagation of the gauge fields
  in the bulk. Unit: GeV (1/length). Value: \(H \gtrsim 2.5\) TeV (from
  the invisible width of the Z boson). Status: AUTHOR.
\item
  \(\mathrm{e}\): Euler's number, the base of the natural logarithm;
  \(\mathrm{e}^{X}\) is the exponential of \(X\). Unit: dimensionless.
  Value: 2.718281828. Status: STANDARD.
\item
  \(x_3\): The third coordinate of ordinary 3-space. Unit: length.
  Status: AUTHOR.
\item
  \(L\): A distance between two points of the fold along ordinary space:
  for a gravitational wave the radius of a ring of free test masses
  (\(\Delta L\) the change of the distance of one of them from the
  centre), for a signal the distance between its emitter and its
  receiver. Unit: length. Status: STANDARD.
\item
  \(\Delta\): A change, difference or small variation of the quantity
  that follows (for example an interval of the time, a perturbation of a
  density, a difference of two squared masses, or an addition to an
  effective number of species). Status: STANDARD.
\item
  \(x_4\): The time; a prime on a function of \(x_4\) alone is
  \(\mathrm{d}/\mathrm{d}x_4\). Unit: length (= time, \(c = 1\)).
  Status: AUTHOR.
\end{itemize}

Three facts. (1) dirac16complex is invariant under the boosts of the
fold (FE-78): the generator \(\mathsf{S}\) of the boosts along \(x_3\),
built from the author's matrices, commutes with \(\gamma^{(0)}\) and
with the chirality, so with the portal terms of the fold, and the warp
multiplies the frequency and the momentum by the same factor. Every
state at the fold therefore has \(\omega^2 = \omega_0^2 + k^2\) exactly
and moves slower than light, at every momentum; the same holds for every
massive mode of every field of the bulk (A-3), and the massless modes
travel at the speed of light at every wavelength. On the static block
the metric is the metric of the fold times the warp factor, plus
\(\mathrm{d}\sigma^2\), so every Lorentz transformation of the fold
leaves it unchanged, and loops of \(\chi\) cannot pass a Lorentz
violation to the Standard Model: the bound of UM-3 on
\(q_{\mathrm{b}}/H\) does not exist in UM-4. (2) Nothing crosses the
static bulk faster than light (FE-83): in the conformal coordinate
\(\tilde\sigma\) the static block is flat space times a factor, which
does not change which curves are null, so a signal that leaves a point
of the fold and comes back at the distance \(L\) along ordinary space
needs the time \(\Delta x_4 \ge L\), equality only along the fold.
Around an expanding fold (FE-84) the light cone of the bulk is slightly
wider than that of the fold where the expansion decelerates; the path
that gains the most arrives earlier than light along the fold by
\(a_4''^2L^3/(150H^2)\), which over the Hubble distance today is
\(4.5\times10^{-94}\) of the travel time for \(H = 2.5\) TeV. (3)
Causality is not protected by the time \(x_4\) alone: the four time-like
directions of the eight-dimensional metric allow closed time-like curves
through every point (every circle in the plane of the time and an extra
time at fixed depth has a tangent of negative squared length everywhere,
for any \(a_4\) and \(b_4\)), and \(x_4\) is a time function only on the
sections of fixed extra times, on which the fields of UM-4 live. Status:
DERIVED (\texttt{um4\_bulk.py}, B1, B2, B7 and B10;
\texttt{um4\_spectra.py}, U5; \texttt{um4\_causality.py}, C1 to C4).

The answer. No: in UM-4 no particle and no signal moves faster than
light on the fold. The massless fields of the bulk, the photon and the
graviton among them, travel at the speed of light at every depth and
every wavelength; the massive modes and the particles \(\chi\) are
slower; the static bulk offers no shortcut, and the expanding bulk one
of \(4.5\times10^{-94}\) of the travel time over the Hubble distance.
The Standard Model inherits no Lorentz violation, so the scale \(H\) is
no longer bounded from above by the speed of the electron; it is bounded
from below, by 2.5 TeV, by the decays of the Z boson into the bulk
(FE-82). The eight-dimensional spacetime keeps the closed time-like
curves of its four time-like directions: causality on the fold rests on
fields that do not depend on the extra times.

\hypertarget{the-primordial-era}{%
\subsection{The primordial era}\label{the-primordial-era}}

\[
a_4 = Hx_4,\quad b_4 = \text{constant}\ (\mathsf{A} = 1,\ \mathsf{U} = 0):\qquad
\hat R_{ab} \ne 0,\qquad \kappa_7\rho_7 = -3H^2,\quad \kappa_7P_7 = \kappa_7q_7 = -3H^2 \tag{FE-31}
\] \textbf{Symbols in (FE-31).}

\begin{itemize}
\tightlist
\item
  \(a_4\): The author's metric function of the time \(x_4\) that scales
  ordinary 3-space: on the fold \(\mathrm{e}^{a_4}\) is the scale factor
  of cosmology and \(a_4'\) the Hubble rate. Unit: dimensionless.
  Status: AUTHOR.
\item
  \(H\): The author's constant (an inverse length) that fixes the period
  \(\pi/(6H)\) of \(x_0\), the warp factor
  \(\mathrm{e}^{-2H\lvert\sigma\rvert}\) of the metric, the cosmological
  constant \(\Lambda = -21H^2\) and the tension
  \(\kappa\lambda_{\mathrm{F}} = 12H\) of the fold; in UM-4 a free
  parameter, bounded from below by the propagation of the gauge fields
  in the bulk. Unit: GeV (1/length). Value: \(H \gtrsim 2.5\) TeV (from
  the invisible width of the Z boson). Status: AUTHOR.
\item
  \(x_4\): The time; a prime on a function of \(x_4\) alone is
  \(\mathrm{d}/\mathrm{d}x_4\). Unit: length (= time, \(c = 1\)).
  Status: AUTHOR.
\item
  \(b_4\): The author's metric function of the time \(x_4\) in the
  extra-time components, \(b_4 > 0\), never identically 1 (the author's
  order of 2026-10-24); the scale factor (the factor of the lengths) of
  the extra times is \(\mathrm{e}^{-a_4}b_4\) on the fold and
  \(\mathrm{e}^{-a_4}b_4\,\mathrm{e}^{-H\lvert\sigma\rvert}\) at the
  distance \(\sigma\) from it (the square root of the warp factor times
  \(\mathrm{e}^{-a_4}b_4\)). Unit: dimensionless. Status: AUTHOR.
\item
  \(\mathsf{A}\): The expansion rate of ordinary space in units of
  \(H\), \(\mathsf{A} = a_4'/H\). Unit: dimensionless. Value: 0 at the
  static block; \(H_4^{(0)}/H\) today (\(H_4^{(0)}\) the Hubble
  constant); the first-order equations of the fold hold for
  \(\mathsf{A}^2 \ll 1\); 1 at the author's vacuum, which is not a
  solution of UM-4. Status: DERIVED.
\item
  \(\mathsf{U}\): The rate of change of \(b_4\) in units of \(H\),
  \(\mathsf{U} = b_4'/(Hb_4)\). Unit: dimensionless. Value:
  \(\mathsf{U} = \mathsf{A}\) with frozen extra times. Status: DERIVED.
\item
  \(\hat R\): The curvature of the metric \(\hat g_{ab}\) of the fold
  (FD-2): with four frame indices its Riemann tensor, with two its Ricci
  tensor \(\hat R_{ab}\), without indices its scalar curvature
  \(\hat R\); for frozen extra times the scalar curvature equals that of
  the four-dimensional metric \(\hat h_{\mu\nu}\). Unit: 1/length\(^2\).
  Status: STANDARD.
\item
  \(a\): A frame (orthonormal) index running over \(0, 1, \dots, 7\)
  (frame direction \((0)\) along \(\sigma\), \((1), (2), (3)\) along
  ordinary space, \((4)\) along the time, \((5), (6), (7)\) along the
  extra times); raised and lowered with the frame metric
  \(\eta_{ab} = \mathrm{diag}(1, 1, 1, 1, -1, -1, -1, -1)\); on the fold
  (\(a, b \in \{1, \dots, 7\}\)) the frame of the fold; in the tensors
  of the fold written with its coordinates \(x^a\) (its metric
  \(\hat g_{ab}\) and the tensors built from it) a coordinate index of
  the fold, \(a \in \{1, \dots, 7\}\); a repeated index is summed.
  Status: STANDARD.
\item
  \(b\): A frame index (as \(a\)). Status: STANDARD.
\item
  \(\kappa_7\): The gravitational coupling of the fold: to the first
  order in its curvature over \(H^2\) the fold obeys the Einstein
  equations of seven dimensions
  \(\hat G^a{}_b = \kappa_7\mathsf{t}^a{}_b\) with
  \(\kappa_7 = 5\kappa H/2\), the coupling of the massless graviton of
  the bulk bound to the fold; four-dimensional gravity on the fold has
  \(8\pi G_{\mathrm{N}} = \kappa_7/V_*\) (plus the induced term, which
  is negligible). Unit: GeV\(^{-5}\). Value: \(5\kappa H/2\). Status:
  DERIVED.
\item
  \(\rho_7\): The energy density of the matter of the fold above its
  tension in the seven-dimensional normalisation (minus the component
  \(\mathsf{t}^{(4)}{}_{(4)}\): an energy per unit volume of the six
  other directions of the fold); \(V_*\rho_7\) in the four-dimensional
  normalisation. Unit: GeV\(^7\). Status: DERIVED.
\item
  \(P_7\): The pressure along ordinary space of the matter of the fold
  above its tension in the seven-dimensional normalisation (the
  component \(\mathsf{t}^{(i)}{}_{(i)}\)). Unit: GeV\(^7\). Status:
  DERIVED.
\item
  \(q_7\): The pressure along the extra times of the matter of the fold
  above its tension in the seven-dimensional normalisation (the
  component \(\mathsf{t}^{(r)}{}_{(r)}\)). Unit: GeV\(^7\). Status:
  DERIVED.
\end{itemize}

The author's vacuum, a de Sitter expansion of ordinary space at the rate
\(H\) with a constant \(b_4\), which the earlier versions of the model
took as the primordial era, is not a solution of UM-4: the metric of its
fold is not Ricci flat, so the bulk of the author's form is not a vacuum
for it, and as a fold with matter it would need a negative energy
density. Status: DERIVED (\texttt{um4\_primordial.py}, PR0).

\[
\begin{aligned}
&g_{MN}\,\mathrm{d}X^M\mathrm{d}X^N = \mathrm{d}\sigma^2 + f_{\mathrm{dS}}^2\Big[\mathrm{e}^{2H_{\mathrm{i}}x_4}\big(\mathrm{d}\vec x^2 - \mathrm{d}\vec x_5^2\big) - \mathrm{d}x_4^2\Big],\qquad
 f_{\mathrm{dS}} = \frac{\sinh\big(H(\sigma_0 - \lvert\sigma\rvert)\big)}{\sinh(H\sigma_0)},\qquad H_{\mathrm{i}} = \frac{H}{\sinh(H\sigma_0)},\\
&\kappa\big(\lambda_{\mathrm{F}} + \rho_{\mathrm{i}}\big) = 12H\coth(H\sigma_0),\qquad
 H_{\mathrm{i}}^2 = \Big(\frac{\kappa(\lambda_{\mathrm{F}} + \rho_{\mathrm{i}})}{12}\Big)^2 - H^2 = \frac{\kappa_7\rho_{\mathrm{i}}}{15}\Big(1 + \frac{\rho_{\mathrm{i}}}{2\lambda_{\mathrm{F}}}\Big)
\end{aligned} \tag{FE-87}
\] \textbf{Symbols in (FE-87).}

\begin{itemize}
\tightlist
\item
  \(g_{MN}\): The metric tensor of the eight-dimensional spacetime
  (signature \(+,+,+,+,-,-,-,-\) in the order \(x_0, \dots, x_7\)):
  \(g_{MN}\,\mathrm{d}X^M\mathrm{d}X^N\) is the squared interval,
  \(g^{MN}\) its inverse, \(\lvert g\rvert\) the absolute value of its
  determinant (\(\sqrt{\lvert g\rvert}\,\mathrm{d}^8X\) is the proper
  volume element). Unit: dimensionless. Status: AUTHOR.
\item
  \(M\): A coordinate index of eight dimensions (also with a numbered
  copy \(M_1, M_2, \dots\)): it runs over \(x_0, x_1, \dots, x_7\); a
  repeated index is summed. Status: STANDARD.
\item
  \(N\): A coordinate index of eight dimensions (also
  \(N_1, N_2, \dots\)), running over \(x_0, \dots, x_7\); a repeated
  index is summed. Status: STANDARD.
\item
  \(\mathrm{d}\): The differential: \(\mathrm{d}\) in front of a
  quantity is an infinitesimal change of it (as \(\mathrm{d}\sigma\), a
  small step of \(\sigma\)); the quotient of two such changes is a
  derivative; \(\mathrm{d}\) with an exponent, in front of several
  variables, is their volume element, the product of their
  differentials. Status: STANDARD.
\item
  \(X = \{x_0, \dots, x_7\}\): The author's coordinates of the
  eight-dimensional spacetime, \(X = \{x_0, x_1, \dots, x_7\}\) with the
  hidden direction first (Dirac\_Einstein-Lovelock\_4+4\_Nash.nb,
  section DEFCHART); \(X^M\) is the coordinate with the index \(M\) and
  \(\mathrm{d}^8X = \mathrm{d}x_0\,\mathrm{d}x_1\cdots\mathrm{d}x_7\)
  the coordinate volume element. Unit: length (each coordinate). Status:
  AUTHOR.
\item
  \(\sigma\): The proper coordinate of the hidden direction:
  \(\mathrm{d}\sigma = \lvert\cot z\rvert\,\mathrm{d}x_0\),
  \(\sigma = \mathrm{sgn}(\pi/2 - z)\ln(\sin z)/(6H)\) on one period
  \(0 < z < \pi\); \(\sigma = 0\) is the fold \(z = \pi/2\),
  \(\sigma \to \mp\infty\) the tips \(z \to 0, \pi\); one period of
  \(x_0\) is the whole line \(-\infty < \sigma < \infty\). Unit: length.
  Status: DERIVED.
\item
  \(f_{\mathrm{dS}}\): The warp of the de Sitter fold: the factor of the
  lengths of the fold at the depth \(\sigma\). Unit: dimensionless.
  Status: DERIVED.
\item
  \(\mathrm{e}\): Euler's number, the base of the natural logarithm;
  \(\mathrm{e}^{X}\) is the exponential of \(X\). Unit: dimensionless.
  Value: 2.718281828. Status: STANDARD.
\item
  \(H_{\mathrm{i}}\): A constant expansion rate of ordinary space: of
  the de Sitter fold (whose extra times expand at the same rate) or of
  an inflation of the fold with frozen extra times. Unit: GeV. Status:
  DERIVED (FE-87), ASSUMED (FE-88).
\item
  \(x_4\): The time; a prime on a function of \(x_4\) alone is
  \(\mathrm{d}/\mathrm{d}x_4\). Unit: length (= time, \(c = 1\)).
  Status: AUTHOR.
\item
  \(\vec x\): The position in ordinary space,
  \(\vec x = (x_1, x_2, x_3)\). Unit: length. Status: AUTHOR.
\item
  \(\vec x_5\): The position along the extra times,
  \(\vec x_5 = (x_5, x_6, x_7)\);
  \(\mathrm{d}\vec x_5^2 = \mathrm{d}x_5^2 + \mathrm{d}x_6^2 + \mathrm{d}x_7^2\).
  Unit: length. Status: AUTHOR.
\item
  \(\sigma_0\): The depth at which the warp of the de Sitter fold
  vanishes (a horizon of the bulk),
  \(\sinh(H\sigma_0) = H/H_{\mathrm{i}}\). Unit: length. Status:
  DERIVED.
\item
  \(H\): The author's constant (an inverse length) that fixes the period
  \(\pi/(6H)\) of \(x_0\), the warp factor
  \(\mathrm{e}^{-2H\lvert\sigma\rvert}\) of the metric, the cosmological
  constant \(\Lambda = -21H^2\) and the tension
  \(\kappa\lambda_{\mathrm{F}} = 12H\) of the fold; in UM-4 a free
  parameter, bounded from below by the propagation of the gauge fields
  in the bulk. Unit: GeV (1/length). Value: \(H \gtrsim 2.5\) TeV (from
  the invisible width of the Z boson). Status: AUTHOR.
\item
  \(\kappa\): The eight-dimensional gravitational coupling,
  \(\kappa = 8\pi G_8\). Unit: GeV\(^{-6}\). Status: AUTHOR.
\item
  \(\lambda_{\mathrm{F}}\): The tension that the static block requires
  of the fold, \(\kappa\lambda_{\mathrm{F}} = 12H\): the total vacuum
  energy of the seven-dimensional fold, made of the tension of the fold
  fluid and the vacuum energy of the fields of the fold (both have the
  pressure minus their energy density along ordinary space and along the
  extra times), per unit volume of the six other directions of the fold
  (eight-dimensional normalisation); in the four-dimensional
  normalisation it is \(\lambda_{\mathrm{F}}V_*\), and the excess of the
  total over \(\lambda_{\mathrm{F}}V_*\) is the dark energy
  \(\rho_\Lambda\). Unit: GeV\(^7\). Value: \(12H/\kappa\);
  \(\lambda_{\mathrm{F}}V_* = 30H^2M_{\mathrm{red}}^2 = 1.1\times10^{45}\)
  GeV\(^4\) for \(H = 2.5\) TeV. Status: DERIVED.
\item
  \(\rho_{\mathrm{i}}\): The vacuum energy of the fold above the
  critical tension \(\lambda_{\mathrm{F}}\) that expands the de Sitter
  fold (eight-dimensional normalisation). Unit: GeV\(^7\). Status:
  DERIVED.
\item
  \(\kappa_7\): The gravitational coupling of the fold: to the first
  order in its curvature over \(H^2\) the fold obeys the Einstein
  equations of seven dimensions
  \(\hat G^a{}_b = \kappa_7\mathsf{t}^a{}_b\) with
  \(\kappa_7 = 5\kappa H/2\), the coupling of the massless graviton of
  the bulk bound to the fold; four-dimensional gravity on the fold has
  \(8\pi G_{\mathrm{N}} = \kappa_7/V_*\) (plus the induced term, which
  is negligible). Unit: GeV\(^{-5}\). Value: \(5\kappa H/2\). Status:
  DERIVED.
\end{itemize}

\[
\mathsf{c}' = 0,\quad a_4' = H_{\mathrm{i}}:\qquad \kappa_7\rho_7 = 3H_{\mathrm{i}}^2,\qquad P_7 = -\rho_7,\qquad q_7 = -2\rho_7,\qquad
\frac{H_{\mathrm{i}}}{H} < \frac{1}{10} \tag{FE-88}
\] \textbf{Symbols in (FE-88).}

\begin{itemize}
\tightlist
\item
  \(\mathsf{c}\): The logarithm of the scale factor of the extra times
  on the fold, \(\mathsf{c} = \ln(\mathrm{e}^{-a_4}b_4)\);
  \(\mathsf{c}' = \mathsf{V}H\); frozen extra times have
  \(\mathsf{c}' = 0\). Unit: dimensionless. Status: DERIVED.
\item
  \({}'\): The derivative of a function of the time \(x_4\) alone with
  respect to \(x_4\) (\(a_4' = \mathrm{d}a_4/\mathrm{d}x_4\); two primes
  the second derivative); on a function of \(\sigma\) alone (a profile)
  the derivative with respect to \(\sigma\). Status: STANDARD.
\item
  \(a_4\): The author's metric function of the time \(x_4\) that scales
  ordinary 3-space: on the fold \(\mathrm{e}^{a_4}\) is the scale factor
  of cosmology and \(a_4'\) the Hubble rate. Unit: dimensionless.
  Status: AUTHOR.
\item
  \(H_{\mathrm{i}}\): A constant expansion rate of ordinary space: of
  the de Sitter fold (whose extra times expand at the same rate) or of
  an inflation of the fold with frozen extra times. Unit: GeV. Status:
  DERIVED (FE-87), ASSUMED (FE-88).
\item
  \(\kappa_7\): The gravitational coupling of the fold: to the first
  order in its curvature over \(H^2\) the fold obeys the Einstein
  equations of seven dimensions
  \(\hat G^a{}_b = \kappa_7\mathsf{t}^a{}_b\) with
  \(\kappa_7 = 5\kappa H/2\), the coupling of the massless graviton of
  the bulk bound to the fold; four-dimensional gravity on the fold has
  \(8\pi G_{\mathrm{N}} = \kappa_7/V_*\) (plus the induced term, which
  is negligible). Unit: GeV\(^{-5}\). Value: \(5\kappa H/2\). Status:
  DERIVED.
\item
  \(\rho_7\): The energy density of the matter of the fold above its
  tension in the seven-dimensional normalisation (minus the component
  \(\mathsf{t}^{(4)}{}_{(4)}\): an energy per unit volume of the six
  other directions of the fold); \(V_*\rho_7\) in the four-dimensional
  normalisation. Unit: GeV\(^7\). Status: DERIVED.
\item
  \(P_7\): The pressure along ordinary space of the matter of the fold
  above its tension in the seven-dimensional normalisation (the
  component \(\mathsf{t}^{(i)}{}_{(i)}\)). Unit: GeV\(^7\). Status:
  DERIVED.
\item
  \(q_7\): The pressure along the extra times of the matter of the fold
  above its tension in the seven-dimensional normalisation (the
  component \(\mathsf{t}^{(r)}{}_{(r)}\)). Unit: GeV\(^7\). Status:
  DERIVED.
\item
  \(H\): The author's constant (an inverse length) that fixes the period
  \(\pi/(6H)\) of \(x_0\), the warp factor
  \(\mathrm{e}^{-2H\lvert\sigma\rvert}\) of the metric, the cosmological
  constant \(\Lambda = -21H^2\) and the tension
  \(\kappa\lambda_{\mathrm{F}} = 12H\) of the fold; in UM-4 a free
  parameter, bounded from below by the propagation of the gauge fields
  in the bulk. Unit: GeV (1/length). Value: \(H \gtrsim 2.5\) TeV (from
  the invisible width of the Z boson). Status: AUTHOR.
\end{itemize}

The two candidates for an inflation of the fold. The exact de Sitter
fold (FE-87), a fold with a tension above the critical one in the bulk
of constant curvature, expands at the rate \(H_{\mathrm{i}}\) in all its
space-like and time-like directions: the extra times inflate together
with ordinary space, the factor \(\mathcal{Z}\) of the couplings of the
fold grows like \(\mathrm{e}^{3N}\) after \(N\) e-folds, and every
coupling of the fold falls like \(\mathrm{e}^{-3N}\), by
\(6.7\times10^{-79}\) after 60 e-folds: it is not the primordial era of
our universe. The inflation with frozen extra times (FE-88), which the
observed constancy of the couplings requires, needs the pressure
\(q_7 = -2\rho_7\) along the extra times, so the fold fluid adds
\(q_{\mathrm{hold}} = -\rho_7\), a pressure of the size of the vacuum
energy of the inflation (in UM-3, with its small \(H\), it had to be
about \(\mathsf{A}^4/8\) times the tension of the fold, with
\(\mathsf{A}\) the rate of the inflation in units of \(H\): at least
\(10^{103}\) times for an inflation at the largest allowed rate); the
exact equations of such a fold in a bulk of constant curvature have no
solution for \(H_{\mathrm{i}} \ne 0\), so it needs a Weyl tensor of the
bulk, whose bulk solution is not computed. Status: DERIVED
(\texttt{um4\_primordial.py}, PR1, PR2 and PR5), OPEN (the Weyl tensor
of the bulk, the origin of the holding pressure).

\[
\begin{aligned}
&r_{\mathrm{T}} < 0.032:\qquad H_{\mathrm{i}} < \pi M_{\mathrm{red}}\sqrt{\frac{r_{\mathrm{T}}A_{\mathrm{s}}}{2}} = 4.43\times10^{13}\ \text{GeV},\\
&\frac{H_{\mathrm{i}}}{H} < \frac{1}{10}:\qquad H > 4.4\times10^{14}\ \text{GeV}\ \ \text{or}\ \ H_{\mathrm{i}} < 250\ \text{GeV}\ \ (H = 2.5\ \text{TeV})
\end{aligned} \tag{S-24}
\] \textbf{Symbols in (S-24).}

\begin{itemize}
\tightlist
\item
  \(r_{\mathrm{T}}\): The ratio of the amplitudes of the primordial
  spectra of gravitational waves (tensor) and of curvature (scalar) at
  the pivot scale. Unit: dimensionless. Value: below 0.032 (Planck PR4
  with BICEP/Keck 2018 and BAO, M. Tristram et al.~2022, 95 per cent).
  Status: STANDARD.
\item
  \(H_{\mathrm{i}}\): A constant expansion rate of ordinary space: of
  the de Sitter fold (whose extra times expand at the same rate) or of
  an inflation of the fold with frozen extra times. Unit: GeV. Status:
  DERIVED (FE-87), ASSUMED (FE-88).
\item
  \(\pi\): The number pi, the ratio of the circumference of a circle to
  its diameter. Unit: dimensionless. Value: 3.141592654. Status:
  STANDARD.
\item
  \(M_{\mathrm{red}}\): The reduced Planck mass,
  \(M_{\mathrm{red}} = (8\pi G_{\mathrm{N}})^{-1/2}\). Unit: GeV. Value:
  \(2.435\times10^{18}\) GeV. Status: STANDARD.
\item
  \(A_{\mathrm{s}}\): The amplitude of the primordial curvature spectrum
  at the pivot \(k_*\). Unit: dimensionless. Value:
  \(2.10\times10^{-9}\) (Planck 2018). Status: ASSUMED.
\item
  \(H\): The author's constant (an inverse length) that fixes the period
  \(\pi/(6H)\) of \(x_0\), the warp factor
  \(\mathrm{e}^{-2H\lvert\sigma\rvert}\) of the metric, the cosmological
  constant \(\Lambda = -21H^2\) and the tension
  \(\kappa\lambda_{\mathrm{F}} = 12H\) of the fold; in UM-4 a free
  parameter, bounded from below by the propagation of the gauge fields
  in the bulk. Unit: GeV (1/length). Value: \(H \gtrsim 2.5\) TeV (from
  the invisible width of the Z boson). Status: AUTHOR.
\item
  \(k_*\): The pivot wave number of the primordial spectrum. Unit: GeV
  (1/length). Value: 0.05 Mpc\(^{-1}\). Status: STANDARD.
\end{itemize}

\[
\xi_{\mathrm{H}} = N_{\mathrm{e}}\sqrt{\frac{\lambda}{72\pi^2A_{\mathrm{s}}}},\qquad n_{\mathrm{S}} = 1 - \frac{2}{N_{\mathrm{e}}},\qquad r_{\mathrm{T}} = \frac{12}{N_{\mathrm{e}}^2},\qquad
H_{\mathrm{i}} = \sqrt{\frac{\lambda}{12}}\,\frac{M_{\mathrm{red}}}{\xi_{\mathrm{H}}} \tag{S-25}
\] \textbf{Symbols in (S-25).}

\begin{itemize}
\tightlist
\item
  \(\xi_{\mathrm{H}}\): The coupling of the Higgs doublet to the scalar
  curvature of the fold, the term
  \(-\xi_{\mathrm{H}}\hat R\,\Phi^\dagger\Phi\) in the Lagrangian of the
  fold: an extension of UM-4, whose Higgs Lagrangian has no such term
  (\(\xi_{\mathrm{H}} = 0\)). Unit: dimensionless. Value: 0 in UM-4;
  about \(1.6\times10^4\) if the Higgs field drove inflation (ESTIMATE,
  an extension of the action). Status: ESTIMATE (an extension of the
  action).
\item
  \(N_{\mathrm{e}}\): The number of e-foldings of inflation after the
  pivot scale left the horizon. Unit: dimensionless. Value: 55. Status:
  STANDARD.
\item
  \(\lambda\): The quartic coupling of the Higgs potential,
  \(\lambda > 0\). Unit: dimensionless. Value:
  \(\lambda = m_h^2/(2v^2) = 0.129\). Status: AUTHOR.
\item
  \(\pi\): The number pi, the ratio of the circumference of a circle to
  its diameter. Unit: dimensionless. Value: 3.141592654. Status:
  STANDARD.
\item
  \(A_{\mathrm{s}}\): The amplitude of the primordial curvature spectrum
  at the pivot \(k_*\). Unit: dimensionless. Value:
  \(2.10\times10^{-9}\) (Planck 2018). Status: ASSUMED.
\item
  \(n_{\mathrm{S}}\): The spectral index of the primordial curvature
  spectrum. Unit: dimensionless. Value: 0.965 (Planck 2018). Status:
  ASSUMED.
\item
  \(r_{\mathrm{T}}\): The ratio of the amplitudes of the primordial
  spectra of gravitational waves (tensor) and of curvature (scalar) at
  the pivot scale. Unit: dimensionless. Value: below 0.032 (Planck PR4
  with BICEP/Keck 2018 and BAO, M. Tristram et al.~2022, 95 per cent).
  Status: STANDARD.
\item
  \(H_{\mathrm{i}}\): A constant expansion rate of ordinary space: of
  the de Sitter fold (whose extra times expand at the same rate) or of
  an inflation of the fold with frozen extra times. Unit: GeV. Status:
  DERIVED (FE-87), ASSUMED (FE-88).
\item
  \(M_{\mathrm{red}}\): The reduced Planck mass,
  \(M_{\mathrm{red}} = (8\pi G_{\mathrm{N}})^{-1/2}\). Unit: GeV. Value:
  \(2.435\times10^{18}\) GeV. Status: STANDARD.
\end{itemize}

The scales of the inflation (S-24) and Higgs inflation (S-25). The
tensor bound \(r_{\mathrm{T}} < 0.032\) (Planck PR4 with BICEP/Keck 2018
and BAO, M. Tristram et al.~2022, 95 per cent) with the measured
amplitude \(A_{\mathrm{s}} = 2.1\times10^{-9}\) of the scalar
perturbations limits any inflation to
\(H_{\mathrm{i}} < 4.4\times10^{13}\) GeV. The equations of the fold are
those of general relativity only for
\((12/25)(H_{\mathrm{i}}/H)^2 \ll 1\), that is
\(H_{\mathrm{i}} < H/10\): an inflation at the tensor bound needs
\(H > 4.4\times10^{14}\) GeV, and with \(H\) at its lower bound of 2.5
TeV only an inflation with \(H_{\mathrm{i}} < 250\) GeV (an energy scale
below \(3.2\times10^{10}\) GeV) is in that regime. If the Higgs doublet
had a coupling \(\xi_{\mathrm{H}}\) to the curvature of the fold, which
the Higgs Lagrangian of UM-4 does not contain, the Higgs field could
drive the inflation of the fold: with \(N_{\mathrm{e}} = 55\) e-folds,
\(\xi_{\mathrm{H}} = 16195\) (\(\lambda = 0.1294\)) or 4502
(\(\lambda = 0.01\), the running value near the scale of inflation),
\(n_{\mathrm{S}} = 0.9636\) and \(r_{\mathrm{T}} = 0.0040\) (0.3
standard deviations below Planck 2018, \(0.9649 \pm 0.0042\), and 3.5
below ACT DR6 combined with Planck, CMB lensing and DESI,
\(0.974 \pm 0.003\)), and \(H_{\mathrm{i}} = 1.6\times10^{13}\) GeV, so
that UM-4 would need \(H > 1.6\times10^{14}\) GeV for it. Status:
DERIVED (\texttt{um4\_primordial.py}, PR3), ESTIMATE (PR4, an extension
of the action), STANDARD (the measurements).

\[
\frac{1}{16\pi G_{\mathrm{ind}}} = \frac{N_{\mathrm{ind}}\,\Lambda_{\mathrm{c}}^2}{192\pi^2},\qquad
N_{\mathrm{ind}} = \sum_{\text{real scalars}}\big(1 - 6\xi\big) + N_{\mathrm{Weyl}} - 4N_{\mathrm{gauge}} = 52 - 24\xi_{\mathrm{H}} \tag{FE-64}
\] \textbf{Symbols in (FE-64).}

\begin{itemize}
\tightlist
\item
  \(\pi\): The number pi, the ratio of the circumference of a circle to
  its diameter. Unit: dimensionless. Value: 3.141592654. Status:
  STANDARD.
\item
  \(G_{\mathrm{ind}}\): The Newton constant of the induced
  Einstein-Hilbert term of the fold,
  \(\kappa_{\mathrm{ind}} = 8\pi G_{\mathrm{ind}}\). Unit: GeV\(^{-2}\).
  Status: ESTIMATE.
\item
  \(N_{\mathrm{ind}}\): The weighted number of the fields of the fold
  that induce the Einstein-Hilbert term: each real scalar counts
  \(1 - 6\xi\), each Weyl fermion 1, each gauge field \(-4\) (with its
  ghost fields). Unit: dimensionless. Value: 52 for the fields of the
  fold of UM-4 (the four real scalars of the Higgs doublet and 48 Weyl
  fermions; the gauge fields live in the bulk) with
  \(\xi_{\mathrm{H}} = 0\); \(52 - 24\xi_{\mathrm{H}}\) with the
  curvature coupling \(\xi_{\mathrm{H}}\) of an extension. Status:
  DERIVED.
\item
  \(\Lambda_{\mathrm{c}}\): The cutoff of the proper-time integral of
  the one-loop effective action of the fields of the fold: the highest
  energy of the quantum fluctuations that induce gravity. Unit: GeV.
  Value: taken of the order of \(H\) (2.5 TeV); the induced term alone
  would give the measured Newton constant for
  \(\Lambda_{\mathrm{c}} = \pi M_{\mathrm{red}}\sqrt{96/52} = 1.0\times10^{19}\)
  GeV. Status: ESTIMATE.
\item
  \(\xi\): The coupling of a real scalar field to the scalar curvature
  \(\hat R\) of the fold: the term \(-\frac12\xi\hat R\) times the
  square of the field in its Lagrangian (\(\xi = 1/6\) is the conformal
  coupling). Unit: dimensionless. Status: STANDARD.
\item
  \(N_{\mathrm{Weyl}}\): The number of Weyl (two-component) fermion
  fields of the fold. Unit: dimensionless. Value: 48 (three generations
  of quarks and leptons with three right-handed neutrinos). Status:
  STANDARD.
\item
  \(N_{\mathrm{gauge}}\): The number of gauge fields of the fold. Unit:
  dimensionless. Value: 0 in UM-4 (the gauge fields live in the bulk);
  12 (eight gluons, three \(W\), one \(B\)) if they lived on the fold.
  Status: STANDARD.
\item
  \(\xi_{\mathrm{H}}\): The coupling of the Higgs doublet to the scalar
  curvature of the fold, the term
  \(-\xi_{\mathrm{H}}\hat R\,\Phi^\dagger\Phi\) in the Lagrangian of the
  fold: an extension of UM-4, whose Higgs Lagrangian has no such term
  (\(\xi_{\mathrm{H}} = 0\)). Unit: dimensionless. Value: 0 in UM-4;
  about \(1.6\times10^4\) if the Higgs field drove inflation (ESTIMATE,
  an extension of the action). Status: ESTIMATE (an extension of the
  action).
\end{itemize}

The induced gravity of the fold. The Einstein-Hilbert term that the
fields of the fold induce at one loop (the Higgs doublet and the 48 Weyl
fermions of the quarks and leptons with three right-handed neutrinos;
the gauge fields are in the bulk) is negligible for a cutoff of the
order of \(H\):
\(1/(8\pi G_{\mathrm{ind}}) = 5.8\times10^{-32}M_{\mathrm{red}}^2\) at
\(\Lambda_{\mathrm{c}} = H = 2.5\) TeV. The gravity of the fold is that
of the bulk graviton, and a coupling \(\xi_{\mathrm{H}} \ne 0\) would
not be needed for it. Status: ESTIMATE (\texttt{um4\_foldsources.py},
FS3).

The answer. The author's vacuum is not a solution of UM-4. The exact de
Sitter fold inflates the extra times as well and is excluded by the
constancy of the couplings. The primordial era of UM-4 can be an
inflation of the fold with frozen extra times, which is the inflation of
four-dimensional general relativity, if the fold fluid supplies a
holding pressure of the size of its vacuum energy, if a Weyl tensor of
the bulk completes the solution, and if its rate is below \(H/10\),
which for \(H\) near its lower bound means an inflation far below the
scales of the usual models. UM-4 does not specify its inflaton: Higgs
inflation is an extension that would need \(H > 1.6\times10^{14}\) GeV.
What precedes inflation and the reheating after it are not computed, and
the Planck values remain inputs of the cosmological calculation (OPEN).

\hypertarget{the-lightest-neutrino-mass-and-the-dirac-phase}{%
\subsection{The lightest neutrino mass and the Dirac
phase}\label{the-lightest-neutrino-mass-and-the-dirac-phase}}

\[
\big(\lvert U_{\mathrm{e}i}\rvert^2\big)_{i = 1, 2, 3} = \big(\cos^2\theta_{12}\cos^2\theta_{13},\ \sin^2\theta_{12}\cos^2\theta_{13},\ \sin^2\theta_{13}\big) \tag{S-17}
\] \textbf{Symbols in (S-17).}

\begin{itemize}
\tightlist
\item
  \(U_{\mathrm{e}i}\): The electron row of the mixing matrix,
  \(U_{\mathrm{e}i} = (U_{\mathrm{PMNS}})_{ei}\):
  \(\lvert U_{\mathrm{e}i}\rvert^2\) is the electron-flavour fraction of
  the mass state \(i\). Unit: dimensionless. Value:
  \(\lvert U_{\mathrm{e}i}\rvert^2 = 0.682\), 0.296 and 0.022 for
  \(i = 1, 2, 3\) (NuFIT 5.2, the fit without the Super-Kamiokande
  atmospheric data: the values of the record; NuFIT 6.0 with them gives
  0.677, 0.301 and 0.022). Status: STANDARD.
\item
  \(i\): A generation index (\(i = 1, 2, 3\)), which also labels the
  three neutrino mass states; in a frame component such as
  \(\mathcal{S}^{(i)}{}_{(i)}\) one of the three directions of ordinary
  space (the three components are equal), and in a tensor of ordinary
  space such as \(\mathsf{h}_{ij}\) one of the directions
  \(x_1, x_2, x_3\); a repeated index is not summed unless a sum sign is
  written. Status: STANDARD.
\item
  \(\theta_{12},\ \theta_{13},\ \theta_{23}\): The three mixing angles
  of \(U_{\mathrm{PMNS}}\) (its fourth parameter is the Dirac phase,
  which does not enter the electron-flavour results of UM-4). Unit:
  dimensionless. Value: \(33.40^\circ\), \(8.54^\circ\), \(49.14^\circ\)
  in the computations of UM-4 (\(\sin^2\theta_{12} = 0.303\),
  \(\sin^2\theta_{13} = 0.02203\), \(\sin^2\theta_{23} = 0.572\): NuFIT
  5.2, November 2022, normal ordering, the fit without the
  Super-Kamiokande atmospheric data); \(33.71^\circ\), \(8.56^\circ\),
  \(43.28^\circ\) in NuFIT 6.0 (2024, with the Super-Kamiokande
  atmospheric data; \(\sin^2\theta_{12} = 0.308\),
  \(\sin^2\theta_{13} = 0.02215\), \(\sin^2\theta_{23} = 0.470\)).
  Status: STANDARD.
\item
  \(U_{\mathrm{PMNS}}\): The lepton mixing matrix
  (Pontecorvo-Maki-Nakagawa-Sakata) of the three Dirac neutrinos; its
  entries \((U_{\mathrm{PMNS}})_{\alpha i}\) connect the flavour state
  \(\nu_\alpha\) with the mass state \(\nu_i\). Unit: dimensionless.
  Value: \(\theta_{12} = 33.40^\circ\), \(\theta_{23} = 49.14^\circ\),
  \(\theta_{13} = 8.54^\circ\) (NuFIT 5.2, November 2022, normal
  ordering, the fit without the Super-Kamiokande atmospheric data); the
  Dirac phase is set to zero (ASSUMED; the fit prefers \(197^\circ\)).
  Status: STANDARD, ASSUMED (the phase).
\item
  \(\mathsf{h}_{ij}\): The transverse traceless perturbation of the
  metric of ordinary space by a gravitational wave as the
  four-dimensional world sees it:
  \(\hat h_{ij} = \mathrm{e}^{2a_4}(\delta_{ij} + \mathsf{h}_{ij})\)
  (\(i, j = 1, 2, 3\)). Unit: dimensionless. Status: STANDARD.
\end{itemize}

The Dirac phase \(\delta_{\mathrm{CP}}\). The mixing matrix is three
rotations by the angles \(\theta_{23}\), \(\theta_{13}\) and
\(\theta_{12}\), the phase attached to the middle one, and its electron
row (S-17) contains no phase. Every flavour result of UM-4 is a result
for electron neutrinos: the solar survival probabilities, the conversion
in the envelope of a supernova (the matter potential acts on \(\nu_e\)
only, so the phase can be moved into \(\nu_\mu\) and \(\nu_\tau\)), and
the capture of the cosmic neutrinos. Computed with the global fit NuFIT
6.0 (I. Esteban, M. C. Gonzalez-Garcia, M. Maltoni, I. Martinez-Soler,
J. P. Pinheiro and T. Schwetz, JHEP 12 (2024) 216; with the atmospheric
data of Super-Kamiokande, \(\delta_{\mathrm{CP}} = 212^\circ\) in the
normal ordering), the electron row is the same for
\(\delta_{\mathrm{CP}} = 0\), \(177^\circ\), \(212^\circ\) and
\(274^\circ\) to \(7\times10^{-18}\), and the envelope gives the same
survival probability as with \(\delta_{\mathrm{CP}} = 0\) to
\(6\times10^{-14}\). The phase appears only in the appearance channels,
which UM-4 does not compute (in vacuum at 295 km and 0.6 GeV, the
baseline and the peak energy of T2K, \(\nu_\mu \to \nu_e\) has the
probability 0.0421, 0.0478 and 0.0541 for \(\delta_{\mathrm{CP}} = 0\),
\(212^\circ\) and \(270^\circ\), \(\bar\nu_\mu \to \bar\nu_e\) 0.0421,
0.0346 and 0.0291), and in the Jarlskog invariant of CP violation,
\(-0.0178\) at \(212^\circ\). Status: DERIVED (\texttt{um4\_numass.py},
NM3).

\[
m_2 = \sqrt{m_1^2 + \Delta m^2_{21}},\qquad m_3 = \sqrt{m_1^2 + \Delta m^2_{31}},\qquad
m_1 < 23\ \text{meV}\ \ (\text{95 per cent}),\qquad
m_{\mathrm{eff}}^2 = \sum_i\lvert U_{\mathrm{e}i}\rvert^2m_i^2 \tag{S-18}
\] \textbf{Symbols in (S-18).}

\begin{itemize}
\tightlist
\item
  \(m_u, \dots, m_3\): The masses of the quarks (\(m_u, m_c, m_t\);
  \(m_d, m_s, m_b\)), of the charged leptons (\(m_e, m_\mu, m_\tau\))
  and of the three neutrino mass states (\(m_1, m_2, m_3\), Dirac
  masses). Unit: GeV. Value: measured values;
  \(\Delta m_{21}^2 = 7.41\times10^{-5}\) eV\(^2\),
  \(\lvert\Delta m_{31}^2\rvert = 2.511\times10^{-3}\) eV\(^2\) in the
  computations of UM-4 (the global fit NuFIT 5.2 of the oscillation
  data, November 2022), \(7.49\times10^{-5}\) and \(2.513\times10^{-3}\)
  eV\(^2\) in NuFIT 6.0 (2024); the lightest mass \(m_1 = 0\) in the
  computations, below 23 meV by DESI DR2 with the CMB and NuFIT 6.0.
  Status: STANDARD.
\item
  \(\Delta\): A change, difference or small variation of the quantity
  that follows (for example an interval of the time, a perturbation of a
  density, a difference of two squared masses, or an addition to an
  effective number of species). Status: STANDARD.
\item
  \(m\): A mass: \(m_i\) the mass of the neutrino mass state \(i\);
  \(\Delta m_{ij}^2 = m_i^2 - m_j^2\) the squared-mass difference of the
  states \(i\) and \(j\) (\(\Delta m_{21}^2\), \(\Delta m_{31}^2\),
  \(\Delta m_{32}^2\)); \(\Delta m^2\) without labels the difference of
  the two states of a resonance, \(m_2^2 - m_1^2\), multiplied by
  \(v_{\mathrm{loc}}^2/v^2\) where this differs from 1; \(m\) without a
  label the mass of a mode of a field of the bulk as the fold sees it,
  \(m^2 = \omega^2 - k^2\); for the W and Z bosons their complex pole
  mass, \(m^2 = m_0^2 - \mathrm{i}m_0\Gamma_{\mathrm{bulk}}\) to the
  leading order; in the relic abundance of the particles \(\chi\) their
  mass. Unit: GeV. Status: STANDARD.
\item
  \(i\): A generation index (\(i = 1, 2, 3\)), which also labels the
  three neutrino mass states; in a frame component such as
  \(\mathcal{S}^{(i)}{}_{(i)}\) one of the three directions of ordinary
  space (the three components are equal), and in a tensor of ordinary
  space such as \(\mathsf{h}_{ij}\) one of the directions
  \(x_1, x_2, x_3\); a repeated index is not summed unless a sum sign is
  written. Status: STANDARD.
\item
  \(m_{\mathrm{eff}}\): The effective mass of the electron neutrino
  measured in the beta decay of tritium,
  \(m_{\mathrm{eff}}^2 = \sum_i\lvert U_{\mathrm{e}i}\rvert^2m_i^2\).
  Unit: eV. Value: below 0.45 eV (KATRIN 2025, 90 per cent). Status:
  STANDARD.
\item
  \(U_{\mathrm{e}i}\): The electron row of the mixing matrix,
  \(U_{\mathrm{e}i} = (U_{\mathrm{PMNS}})_{ei}\):
  \(\lvert U_{\mathrm{e}i}\rvert^2\) is the electron-flavour fraction of
  the mass state \(i\). Unit: dimensionless. Value:
  \(\lvert U_{\mathrm{e}i}\rvert^2 = 0.682\), 0.296 and 0.022 for
  \(i = 1, 2, 3\) (NuFIT 5.2, the fit without the Super-Kamiokande
  atmospheric data: the values of the record; NuFIT 6.0 with them gives
  0.677, 0.301 and 0.022). Status: STANDARD.
\end{itemize}

\[
\Gamma_{\mathrm{C\nu B}} = N_{\mathrm{T}}\,\bar\sigma_{\mathrm{T}}\,\frac{n_\nu^{(0)}}{2}\sum_{i=1}^{3}\lvert U_{\mathrm{e}i}\rvert^2\big(1 + \mathsf{s}_i\big)\qquad(\text{Dirac neutrinos}) \tag{S-28}
\] \textbf{Symbols in (S-28).}

\begin{itemize}
\tightlist
\item
  \(\Gamma_{\mathrm{C\nu B}}\): The rate at which a tritium target
  captures neutrinos of the cosmic neutrino background. Unit: 1/time.
  Status: STANDARD.
\item
  \(N_{\mathrm{T}}\): The number of tritium nuclei of the target. Unit:
  dimensionless. Value: \(2.0\times10^{25}\) for 100 g. Status:
  STANDARD.
\item
  \(\bar\sigma_{\mathrm{T}}\): The averaged cross-section of the capture
  of an electron neutrino on tritium: a tritium nucleus captures
  left-handed electron neutrinos of mean speed \(\mathsf{s}\) and number
  density \(n\) at the rate
  \(\bar\sigma_{\mathrm{T}}(1 + \mathsf{s})\,c\,n\) (\(c\) the speed of
  light, 1 in the formulas; Long, Lunardini and Sabancilar). Unit:
  cm\(^2\). Value: \(3.834\times10^{-45}\) cm\(^2\). Status: STANDARD.
\item
  \(n_\nu\): The number density of the cosmic neutrino background per
  flavour (neutrinos plus antineutrinos, one helicity each);
  \(n_\nu^{(0)}\) its value today. Unit: GeV\(^3\). Value: 112
  cm\(^{-3}\) today. Status: STANDARD.
\item
  \(i\): A generation index (\(i = 1, 2, 3\)), which also labels the
  three neutrino mass states; in a frame component such as
  \(\mathcal{S}^{(i)}{}_{(i)}\) one of the three directions of ordinary
  space (the three components are equal), and in a tensor of ordinary
  space such as \(\mathsf{h}_{ij}\) one of the directions
  \(x_1, x_2, x_3\); a repeated index is not summed unless a sum sign is
  written. Status: STANDARD.
\item
  \(U_{\mathrm{e}i}\): The electron row of the mixing matrix,
  \(U_{\mathrm{e}i} = (U_{\mathrm{PMNS}})_{ei}\):
  \(\lvert U_{\mathrm{e}i}\rvert^2\) is the electron-flavour fraction of
  the mass state \(i\). Unit: dimensionless. Value:
  \(\lvert U_{\mathrm{e}i}\rvert^2 = 0.682\), 0.296 and 0.022 for
  \(i = 1, 2, 3\) (NuFIT 5.2, the fit without the Super-Kamiokande
  atmospheric data: the values of the record; NuFIT 6.0 with them gives
  0.677, 0.301 and 0.022). Status: STANDARD.
\item
  \(\mathsf{s}_i\): The mean speed of the neutrino mass state \(i\)
  today, averaged over its momentum distribution (0 for a state at rest,
  1 for a massless one). Unit: dimensionless. Value: 1, 0.061 and 0.011
  for the mass states 1, 2 and 3 with \(m_1 = 0\); 0.023 for the state 1
  with \(m_1 = 23\) meV. Status: STANDARD.
\end{itemize}

The lightest mass \(m_1\). DESI DR2 (the BAO) with the CMB, combined
with NuFIT 6.0 (W. Elbers et al., arXiv:2503.14744; the
cosmological-constant model, 95 per cent), bounds the lightest mass to
\(m_1 < 23\) meV in the normal ordering (24 meV in the inverted one) and
the sum of the masses to 0.101 eV (normal) and 0.133 eV (inverted); the
inverted ordering, whose smallest sum is 0.0989 eV, is disfavoured (the
probability of the normal ordering is 0.91) but not excluded. With
\(\Delta m^2_{21} = 7.49\times10^{-5}\) eV\(^2\) and
\(\Delta m^2_{31} = 2.513\times10^{-3}\) eV\(^2\) (S-18) the masses at
\(m_1 = 23\) meV are 23.00, 24.57 and 55.15 meV, and the effective mass
of tritium beta decay is 8.84 to 24.64 meV in this range, far below the
limit 0.45 eV of KATRIN. Within the range \(m_1\) changes one result of
UM-4, the capture rate of the cosmic neutrinos (S-28): on 100 g of
tritium it falls from 6.88 per year at \(m_1 = 0\) to 4.15 per year at
\(m_1 = 23\) meV, as the lightest mass state \(\nu_1\) becomes slow
(mean speed 0.023) and loses the doubling of the capture of relativistic
neutrinos (4.06 per year if all states were at rest). Status: DERIVED
(\texttt{um4\_numass.py}, NM1, NM2 and NM4), STANDARD (the bounds).

The answer. Setting the Dirac phase to zero was without consequence: it
drops out of every result of UM-4. The lightest mass is not fixed by any
measurement; DESI DR2 with NuFIT 6.0 bounds it below 23 meV, and within
this range it lowers the capture rate of the cosmic neutrinos by up to
40 per cent. The results computed with NuFIT 5.2 differ from NuFIT 6.0
by at most 1.5 per cent (NM5). The neutrino sector does not depend on
the warp of the time: the neutrinos live on the fold.

\hypertarget{what-the-author-may-have-left-out}{%
\section{What the author may have left
out}\label{what-the-author-may-have-left-out}}

The author asked to be told if something important was missing from his
specification. UM-4 had to add the following, and the analysis found the
following problems.

\begin{longtable}[]{@{}
  >{\raggedright\arraybackslash}p{(\columnwidth - 4\tabcolsep) * \real{0.3333}}
  >{\raggedright\arraybackslash}p{(\columnwidth - 4\tabcolsep) * \real{0.3333}}
  >{\raggedright\arraybackslash}p{(\columnwidth - 4\tabcolsep) * \real{0.3333}}@{}}
\toprule\noalign{}
\begin{minipage}[b]{\linewidth}\raggedright
item
\end{minipage} & \begin{minipage}[b]{\linewidth}\raggedright
why it is needed or what it means
\end{minipage} & \begin{minipage}[b]{\linewidth}\raggedright
where it stands
\end{minipage} \\
\midrule\noalign{}
\endhead
\bottomrule\noalign{}
\endlastfoot
right-handed quark singlets and the Yukawa couplings & the quark
doublets alone have no masses, and the hypercharges cancel the gauge
anomalies only with the singlets & added (a default proposed in the
notification of 2026-10-24, 17:20, which the author did not answer:
ASSUMED) \\
three right-handed neutrinos & Dirac neutrino masses and the mixing that
flavour conversion needs; the anomalies still cancel & added (accepted
by the author, 2026-10-24, 17:32); they must never thermalise, or they
would add 0.14 to 0.24 to the effective number of neutrino species \\
the warped time & in UM-3 the time was not warped, so the light cone of
the bulk was wider than that of the fold and every field of the bulk
outran light on the fold & the author's order of 2026-10-25, 13:15 PDT:
the time carries the warp factor of the other six directions of the fold
(FE-27); among the static warps of the author's form only equal rates in
all seven directions give a vacuum of gravity with a cosmological
constant (DERIVED, \texttt{um4\_geometry.py}, G10) \\
the gravity & the author's Einstein-Lovelock action has free weights &
Einstein gravity with \(\Lambda = -21H^2\) (the weights of the
Gauss-Bonnet and cubic terms set to zero), the only choice for which
gravitons travel on the light cone on every background, not only around
the static block (DERIVED, \texttt{um4\_gravity.py}, L1 to L4) \\
the gauge fields in the bulk & the order puts the photon in the bulk;
with it the electroweak fields, of which it is a combination & ASSUMED
for all gauge fields; their massless modes are bound to the fold and
travel at the speed of light, the W and Z bosons become resonances of
the bulk, and the invisible width of the Z boson bounds \(H > 2.5\) TeV
(FE-82) \\
a thin source at the fold & the kink of the warp factor
\(\mathrm{e}^{-2H\lvert\sigma\rvert}\) requires one (the fold fluid); at
the static block it is a pure tension,
\(\kappa\lambda_{\mathrm{F}} = 12H\), like the brane of Randall and
Sundrum; its excess over \(\lambda_{\mathrm{F}}\) is the dark energy,
tuned to one part in \(10^{92}\) (S-4) & ASSUMED (its values DERIVED);
OPEN (the tuning) \\
the gravity of the fold & the fold obeys the Einstein equations of seven
dimensions with the coupling of the massless graviton of the bulk, with
frozen extra times the Friedmann equations of general relativity, with
corrections of the relative size \((12/25)(a_4'/H)^2\); its
gravitational waves travel at the speed of light at every wavelength,
with the plus and cross polarisations; dust requires a projected Weyl
tensor of the bulk at the order of the square of its density over the
tension & DERIVED; OPEN (the Weyl tensor of the bulk) \\
the volume of the extra times & it is not frozen by itself: dust and
vacuum energy drive it, and a free volume would change Newton's constant
8500 times faster than lunar laser ranging allows; the fold fluid must
add the pressure \(-0.842\rho_{\mathrm{c}}^{(0)}\) along the extra times
today (with no pressure of the baryons along them; \(-0.793\) or
\(-0.891\rho_{\mathrm{c}}^{(0)}\) with \(-\rho_{\mathrm{b}}\) or
\(+\rho_{\mathrm{b}}\)), and minus the vacuum energy of an inflation of
the fold & ASSUMED (the holding pressure and the pressure of the
baryons); OPEN (its origin, a potential for the volume) \\
frozen extra times in the late universe & otherwise the couplings of the
Standard Model change: Newton's constant, the squares of the gauge and
Yukawa couplings and the quartic Higgs coupling of the fold are divided
by the factor \(\mathcal{Z}\) of the volume of the extra times &
required, and held by the fold fluid \\
a gauge field for the ad hoc Higgs vacuum & the literal ad hoc vacuum
carries an isospin charge that the Gauss law forbids; with the
pure-gauge field \(W^2\) along the time (the value at the fold of the
field of the bulk) it is a rotated vacuum with one radial residual, and
it must decay before nucleosynthesis & DERIVED \\
the Krein norm of dirac16complex & the signature \((4,4)\) of the
author's matrices makes the charge indefinite; quantised, it has ghosts
that no vacuum removes; the vacuum is stable, the classical stability of
the gas is not decided; chiral couplings at the fold produce ghosts &
DERIVED; OPEN: the stability of the gas, a unitary coupling
(chirality-blind at the fold is necessary, not shown sufficient) \\
the extra times & fields that depend on them are tachyonic; the
restriction to fields without a momentum along them is consistent but
unstable; compact extra times do not help (the \(n\)-th harmonic grows
\(n\) times faster than the first, and a compact time-like direction is
a closed time-like curve) & ASSUMED (no dependence); OPEN: a gauge
principle that removes that momentum \\
causality & nothing moves faster than light in UM-4, but the four
time-like directions allow closed time-like curves through every point &
DERIVED; causality on the fold rests on fields that do not depend on the
extra times \\
the binding of \(\chi\) to the fold & the states of dirac16complex at
the fold are resonances, which live longer than the universe only if
\(m_\Psi/H > 1.1\times10^3\); bound by the author's Higgs portal they
would have to be bound so weakly (direct detection) that they live long
enough only for \(H < 2.1\) GeV, below the bound \(H > 2.5\) TeV & OPEN:
a binding that does not couple \(\chi\) to nucleons (A-6) \\
the origin of the dark matter & the long-lived particles \(\chi\)
(heavier than \(2.4\times10^6\) GeV) are heavier than any thermal relic
that unitarity allows (\(8.87\times10^4\) GeV) & OPEN: a non-thermal
origin of the abundance \\
the dark energy & a vacuum energy with \(w = -1\); DESI DR2 prefers an
evolving one & OPEN (tension at 2.3 to 3.1 standard deviations) \\
the primordial era & the author's vacuum is not a solution; the de
Sitter fold inflates the extra times; an inflation with frozen extra
times needs a Weyl tensor of the bulk and, in the regime of general
relativity, a rate below \(H/10\) (below 250 GeV for \(H = 2.5\) TeV);
the Higgs field could drive it only with a curvature coupling that the
action does not contain & DERIVED; ESTIMATE (Higgs inflation, an
extension); OPEN (the inflaton, the Weyl tensor; Planck inputs) \\
the lightest neutrino mass and the Dirac phase & not fixed by the
specification; the phase drops out of every result of UM-4; DESI DR2
with NuFIT 6.0 bounds \(m_1 < 23\) meV, which changes only the capture
rate of the cosmic neutrinos & DERIVED (the phase is irrelevant);
ASSUMED: \(m_1 = 0\) within the bound \\
\end{longtable}

\hypertarget{um-4-compared-with-um-3}{%
\section{UM-4 compared with UM-3}\label{um-4-compared-with-um-3}}

UM-3 is kept unchanged as the record of the model before the order of
2026-10-25, 13:15 PDT. The table lists what the warped time changed; the
values of UM-3 are records.

\begin{longtable}[]{@{}
  >{\raggedright\arraybackslash}p{(\columnwidth - 4\tabcolsep) * \real{0.3333}}
  >{\raggedright\arraybackslash}p{(\columnwidth - 4\tabcolsep) * \real{0.3333}}
  >{\raggedright\arraybackslash}p{(\columnwidth - 4\tabcolsep) * \real{0.3333}}@{}}
\toprule\noalign{}
\begin{minipage}[b]{\linewidth}\raggedright
\end{minipage} & \begin{minipage}[b]{\linewidth}\raggedright
UM-3
\end{minipage} & \begin{minipage}[b]{\linewidth}\raggedright
UM-4
\end{minipage} \\
\midrule\noalign{}
\endhead
\bottomrule\noalign{}
\endlastfoot
the time component of the metric & \(-1\): the time not warped &
\(-(\sin^2z)^{1/6}\): the time warped like the other six directions of
the fold \\
the light cone of the bulk at the depth \(\sigma\) &
\(\mathrm{e}^{H\lvert\sigma\rvert}\) times wider than that of the fold &
that of the fold at every depth \\
the static block & a space whose curvature is \(-H^2\) in the planes
without the time and zero in the seven planes with the time & a space of
constant curvature \(-H^2\) in all 28 planes (locally the anti-de Sitter
space of eight dimensions) \\
the gravity & the Lovelock weights \(1, 1/20, 1/360\) with
\(\Lambda = -7H^2\): a degenerate vacuum without linear gravitons in the
static bulk & Einstein gravity with \(\Lambda = -21H^2\): a vacuum with
a massless graviton bound to the fold \\
bulk scalars and photons & 1.22 and 1.41 times the speed of light at
long wavelengths, faster than light at every wavelength & the massless
modes at the speed of light, the massive modes slower \\
the gravitational wave & the four-dimensional graviton at the speed of
light at linear order only & the massless graviton of the bulk at the
speed of light at every depth and wavelength, with the plus and cross
polarisations (helicity \(\pm2\)) \\
the gauge fields of the Standard Model & on the fold & in the bulk; the
photon a transverse wave of spin 1 bound to the fold \\
the spectrum of the hidden direction & gaps \(3H\) (scalars) and \(2H\)
(gauge fields) above the massless modes & no gap: a continuum from zero
mass (A-3) \\
the states of dirac16complex at the fold & eight bound states & eight
long-lived resonances that tunnel into the bulk \\
their speed & the speed coefficient
\(c_{\mathrm{b}}^2 = q_{\mathrm{b}}/(q_{\mathrm{b}} - H)\), a group
velocity up to 1.0046 times the speed of light & exactly relativistic,
\(\omega^2 = \omega_0^2 + k^2\) \\
the author's constant \(H\) & below about \(5\times10^{-13}\) GeV (the
Lorentz violation passed to the electron, for
\(\omega_{\mathrm{b}} = 1.5\) TeV) & above 2.5 TeV (the invisible width
of the Z boson) \\
the dark matter & a thermal relic, weakly bound and leptophilic, heavier
than about 1.5 TeV & not a thermal relic; heavier than \(2.4\times10^6\)
GeV; a binding without a nucleon coupling needed (OPEN) \\
the inflation of the fold with frozen extra times & a holding pressure
of about \(\mathsf{A}^4/8\) times the tension of the fold & a holding
pressure of the size of the vacuum energy of the inflation, a Weyl
tensor of the bulk, and a rate below \(H/10\) \\
unchanged & the free volume excluded; the Krein norm with its ghosts;
the tachyons of the extra times; the closed time-like curves; the
neutrino sector & the same \\
\end{longtable}

\hypertarget{what-was-verified}{%
\section{What was verified}\label{what-was-verified}}

Every result labelled DERIVED rests on a computation in the folder
\texttt{universe\_model\_4/code/} of the repository, which checks itself
and writes a report to \texttt{universe\_model\_4/reports/} (all
equations and results are written out above). Each script is run with
\texttt{python\ universe\_model\_4/code/NAME.py} (Python 3 with numpy,
scipy, sympy, mpmath and matplotlib, and CAMB for the cosmology), and no
single calculation runs longer than 423 seconds. The programs that this
document cites, with their checks, all passed:

\begin{longtable}[]{@{}
  >{\raggedright\arraybackslash}p{(\columnwidth - 4\tabcolsep) * \real{0.3333}}
  >{\raggedright\arraybackslash}p{(\columnwidth - 4\tabcolsep) * \real{0.3333}}
  >{\raggedright\arraybackslash}p{(\columnwidth - 4\tabcolsep) * \real{0.3333}}@{}}
\toprule\noalign{}
\begin{minipage}[b]{\linewidth}\raggedright
program
\end{minipage} & \begin{minipage}[b]{\linewidth}\raggedright
checks
\end{minipage} & \begin{minipage}[b]{\linewidth}\raggedright
what they verify
\end{minipage} \\
\midrule\noalign{}
\endhead
\bottomrule\noalign{}
\endlastfoot
\texttt{um4\_geometry.py} & G1, G10 & the proper chart (A-1); the static
warps and their vacua \\
\texttt{um4\_gravity.py} & L1 to L4 & the choice of Einstein gravity \\
\texttt{um4\_bulk.py} & B1, B2, B7, B10 & the massless modes of the
scalar and the gauge field, the boost invariance of dirac16complex, the
continuum from zero mass and its spectral density (A-3) \\
\texttt{um4\_propagators.py} & P1 to P5 & the propagators, Newton's and
Coulomb's laws, the W and Z resonances, the bounds on \(H\) \\
\texttt{um4\_spectra.py} & U1 to U7 & the resonances of dirac16complex,
their widths, their dispersion, the lifetime condition (S-6), their
Krein signature \\
\texttt{um4\_fields.py} & F3 & the Krein operator and the signature of
the charge \\
\texttt{um4\_kohnsham.py} & K1 to K10 & the Kohn-Sham states of the gas
(S-7, S-9), their equation of state, their width, the weakly bound
point, how dilute the gas of the universe is \\
\texttt{um4\_krein.py} & KR1 to KR7 & FE-67, the symmetries and the
mirror of the states at the fold, the flat directions and the partner
level of the gas, the selection rules at the fold \\
\texttt{um4\_thermal.py} & T2 to T5, T7, T8 & the relic density, the
portals and direct detection (S-21), the right-handed neutrinos, the
unitarity bound (S-20), the binding by the Higgs portal (S-22) \\
\texttt{um4\_foldsources.py} & FS3, FS4 & the induced gravity, the
tuning of the tension (S-4) \\
\texttt{um4\_cosmology.py} & C4, C5 & the CMB and the BAO with
\(w = -1\) \\
\texttt{um4\_extratimes.py} & EX1 to EX5 & the fields with a momentum
along the extra times (A-8) \\
\texttt{um4\_causality.py} & C1 to C4 & no shortcut through the static
bulk, the light cone of an expanding fold, the largest advance, the
closed time-like curves \\
\texttt{um4\_primordial.py} & PR0 to PR5 & the author's vacuum, the de
Sitter fold, the inflation with frozen extra times, its scales (S-24),
Higgs inflation (S-25), the fold factor \\
\texttt{um4\_numass.py} & NM1 to NM5 & S-17, S-18, the capture rate
S-28, and the comparison with NuFIT 6.0 \\
\end{longtable}

All 185 checks of the 20 programs of UM-4 pass. The Kohn-Sham module
runs in parts (\texttt{um4\_kohnsham.py\ part\ NAME} for NAME = K1, I,
K6, K9, then \texttt{um4\_kohnsham.py\ assemble}), so that no single
calculation takes longer than 423 seconds.

\hypertarget{further-symbols}{%
\section{Further symbols}\label{further-symbols}}

\textbf{Further symbols.} These symbols appear in the definitions of the
symbols listed after the equations of this document.

\begin{itemize}
\tightlist
\item
  \(M_{\mathrm{red}}\): The reduced Planck mass,
  \(M_{\mathrm{red}} = (8\pi G_{\mathrm{N}})^{-1/2}\). Unit: GeV. Value:
  \(2.435\times10^{18}\) GeV. Status: STANDARD.
\item
  \(h\): The Higgs field: the fluctuation of the lower component of the
  Higgs doublet about \(v\). Unit: GeV. Value: \(m_h = 125.25\) GeV.
  Status: STANDARD.
\item
  \(\Omega_M\): The spin connection acting on dirac16complex,
  \(\Omega_M = \frac12\varpi_{Mab}\Sigma^{ab}\). Unit: 1/length. Status:
  RECORD.
\item
  \(\beta\): The matrix of the four-dimensional Dirac adjoint,
  \(\beta = \hat\gamma^{(4)}\) (Hermitian, \(\beta^2 = \mathbf{1}\)).
  Unit: dimensionless. Status: STANDARD.
\item
  \(\epsilon^{IJK}\): The Levi-Civita symbol of three indices (the
  structure constants of \(SU(2)\)), \(\epsilon^{123} = 1\). Unit:
  dimensionless. Status: STANDARD.
\item
  \(\mathsf{V}\): The rate of change of the scale factor
  \(\mathrm{e}^{-a_4}b_4\) of the extra times on the fold in units of
  \(H\), \(\mathsf{V} = \mathsf{U} - \mathsf{A}\); zero with frozen
  extra times. Unit: dimensionless. Status: DERIVED.
\item
  \(\mathsf{g}_1,\ \mathsf{g}_2,\ \mathsf{g}_3\): The gauge couplings of
  eight dimensions of \(U(1)_Y\), \(SU(2)_L\) and \(SU(3)_C\): the
  coefficients of the kinetic terms of the gauge fields of the bulk in
  their geometric normalisation; the four-dimensional couplings are
  \(g_k^2 = 3H\mathsf{g}_k^2/(2V_*)\). Unit: GeV\(^{-2}\)
  (\(\mathsf{g}^2\): length\(^4\)). Status: ASSUMED.
\item
  \(\nabla\): The covariant derivative of a tensor with the Levi-Civita
  connection of the metric \(g_{MN}\) (Christoffel symbols
  \(\Gamma^P{}_{MN}\)). Status: STANDARD.
\item
  \(g_1\): The \(U(1)_Y\) gauge coupling of the four-dimensional world
  (usually \(g'\); not normalised for grand unification): the coupling
  of the massless mode of the hypercharge field of the bulk,
  \(g_1^2 = 3H\mathsf{g}_1^2/(2V_*)\). Unit: dimensionless. Value:
  \(g_1(m_Z) = 0.3574\). Status: STANDARD.
\item
  \(g_2\): The \(SU(2)_L\) gauge coupling of the four-dimensional world
  (usually \(g\)): the coupling of the massless mode of the \(SU(2)_L\)
  field of the bulk, \(g_2^2 = 3H\mathsf{g}_2^2/(2V_*)\). Unit:
  dimensionless. Value: \(g_2(m_Z) = 0.6517\). Status: STANDARD.
\item
  \(g_3\): The \(SU(3)_C\) gauge coupling of the four-dimensional world
  (the strong coupling \(g_s\)): the coupling of the massless mode of
  the gluon field of the bulk, \(g_3^2 = 3H\mathsf{g}_3^2/(2V_*)\).
  Unit: dimensionless. Value: \(g_3(m_Z) = 1.2177\)
  (\(g_3^2/(4\pi) = 0.1180\)). Status: STANDARD.
\item
  \(g_\rho\): The effective number of relativistic degrees of freedom of
  the energy density. Unit: dimensionless. Value: 106.75 in the Standard
  Model above the electroweak scale. Status: STANDARD.
\item
  \(q_{\mathrm{hold}}\): The pressure along the extra times that the
  fold fluid adds to its tension to hold the volume of the extra times
  fixed (four-dimensional normalisation). Unit: GeV\(^4\). Value:
  \(-(\rho - 3\mathcal{P})/2 - q_{\mathrm{SM}} - q_\chi - q_{\mathrm{P}}\):
  0 for radiation, \(-\rho/2\) for dust, \(-\rho\) for a vacuum energy;
  today \(-0.842\,\rho_{\mathrm{c}}^{(0)}\) if the baryons have no
  pressure along the extra times (ASSUMED). Status: DERIVED (the value),
  ASSUMED (that the fold fluid supplies it).
\item
  \(v_{\mathcal{A}}\): The amplitude of almightyAdHoc at \(x_4 = 0\);
  the author's almightyAdHoc, \(\mathrm{e}^{-2a_{40}x_4}\) (written
  Exp{[}-2 a40 x4{]} in his notebook), has \(v_{\mathcal{A}} = 1\) in
  the units of his notebook. Unit: GeV. Value: 1 in the units of the
  author's notebook; UM-4 keeps it free (the examples take
  \(v_{\mathcal{A}} = v\) and \(10v\)). Status: AUTHOR (the form),
  ASSUMED (a free amplitude).
\item
  \(\Gamma^P{}_{MN}\): The Christoffel symbols of the metric \(g_{MN}\)
  (the Levi-Civita connection). Unit: 1/length. Status: STANDARD.
\item
  \(\varpi_{Mab}\): The canonical (Levi-Civita) spin connection of the
  vielbein, antisymmetric in \(a, b\). Unit: 1/length. Status: STANDARD.
\item
  \(\mathcal{T}\): The energy-momentum tensor of any one term of the
  fold (in \(\mathcal{T}^r{}_r = \mathcal{L}_{\mathrm{fold}}\)). Unit:
  GeV\(^4\). Status: STANDARD.
\item
  \(\mathcal{W}\): The warp factor: in a warped metric, the function of
  the position along one direction that multiplies the squared steps of
  other directions (its square root multiplies their lengths); in UM-4
  the function of \(\sigma\) that multiplies all seven directions of the
  fold, ordinary space, the time and the extra times,
  \(\mathcal{W} = \mathrm{e}^{-2H\lvert\sigma\rvert} = (\sin^2z)^{1/6}\),
  equal to 1 on the fold. Unit: dimensionless. Status: DERIVED.
\item
  \(c_{\mathrm{b}}\): The speed coefficient of the particles \(\chi\) of
  dirac16complex along ordinary space, in units of the speed of light on
  the fold: at small momentum
  \(\omega^2 = \omega_0^2 + c_{\mathrm{b}}^2\vec k^2\). In UM-3, whose
  time was not warped, \(c_{\mathrm{b}}^2 = 1.19\) for
  \(H = 0.05\,m_\Psi\) and \(q_{\mathrm{b}} = 0.3\,m_\Psi\) (faster than
  light); in UM-4 the boosts of the fold are a symmetry and
  \(c_{\mathrm{b}} = 1\) exactly. Unit: dimensionless. Value: 1 in UM-4;
  1.09 in UM-3. Status: DERIVED.
\item
  \(\delta_{\mathrm{CP}}\): The Dirac phase of the mixing matrix
  \(U_{\mathrm{PMNS}}\). Unit: degree. Value: \(212^\circ\) (NuFIT 6.0,
  normal ordering); 0 in the computations of UM-4. Status: STANDARD.
\item
  \(\nu_\alpha,\ \nu_i\): A neutrino field: \(\nu_\alpha\) the flavour
  states (the flavour index runs over the electron, muon and tau
  flavours), \(\nu_i\) the mass states (\(i = 1, 2, 3\)). Unit:
  GeV\(^{3/2}\). Status: STANDARD.
\item
  \(\rho_{\mathrm{b}}\): The (rest-mass) energy density of the baryons:
  in a supernova the local density of the matter (quoted in g/cm\(^3\)),
  in cosmology the mean density of the baryons
  (\(\rho_{\mathrm{b}}^{(0)}\) today). Unit: GeV\(^4\). Status:
  STANDARD.
\item
  \(w_\chi\): The equation of state of the gas of the particles
  \(\chi\): its pressure divided by its energy density,
  \(w_\chi = P_{\mathrm{KS}}/(n_\chi\epsilon_{\mathrm{p}})\) (0 for cold
  dark matter, \(-1\) for a vacuum energy). Unit: dimensionless. Status:
  DERIVED.
\item
  \(\mathcal{L}_{\mathrm{fold}}\): The Lagrangian of any one term of the
  fold (four-dimensional normalisation). Unit: GeV\(^4\). Status:
  STANDARD.
\end{itemize}

\end{document}
